% The Algebra of Chemical Space — combined LaTeX for reading and review
% All 31 active LaTeX sources are expanded in reading order, including Appendix F.
% The complete 70-entry bibliography and published index are also inlined.
% Source-boundary comments preserve the original file names.
% This is the current revised book; excluded historical drafts are not reintroduced.
% Figure references and captions remain; binary images are not embedded.
% The rendered figures can be seen in https://austinclyde.com/scaffold-book
% Public source: https://austinclyde.com/scaffold-book-tex
%
% BEGIN SOURCE: main.tex
% !TeX program = pdflatex
\documentclass[10pt,openany]{book}

%-------Packages---------
\usepackage{amsmath,amssymb,amsfonts,bm}
\usepackage{amsthm} % <-- FIX: provides \theoremstyle and theorem styles
\usepackage[all,arc]{xy}
\usepackage{enumerate}
\usepackage{mathrsfs}
\usepackage{tikz}
\usetikzlibrary{cd}
\usepackage{graphicx}
\usepackage{caption}
\usepackage{hyperref}
\hypersetup{hidelinks,pdftitle={The Algebra of Chemical Space: Abstraction, Transfer, and Molecular Decisions},pdfauthor={Austin Clyde}}
\usepackage{cleveref}
\usepackage{microtype}
\usepackage{booktabs}
\usepackage{algorithm}
\usepackage{algpseudocode}
\usepackage{framed}
\usepackage{makeidx}
\makeindex

%-------Page Layout---------
\usepackage[margin=1.25in]{geometry}
\usepackage{fancyhdr}
\pagestyle{fancy}
\fancyhf{}
\fancyhead[LE]{\small\nouppercase{\leftmark}}
\fancyhead[RO]{\small\nouppercase{\rightmark}}
\fancyfoot[C]{\thepage}
\renewcommand{\headrulewidth}{0.4pt}

%-------Part/Chapter Styling---------
\usepackage{titlesec}
\titleformat{\part}[display]
  {\centering\Huge\bfseries}
  {\partname\ \thepart}{20pt}{\Huge}
\titleformat{\chapter}[display]
  {\normalfont\huge\bfseries}
  {\chaptertitlename\ \thechapter}{20pt}{\Huge}

%--------Theorem Environments--------
\newtheorem{thm}{Theorem}[chapter]
\newtheorem{cor}[thm]{Corollary}
\newtheorem{prop}[thm]{Proposition}
\newtheorem{lem}[thm]{Lemma}
\newtheorem{conj}[thm]{Conjecture}
\newtheorem{quest}[thm]{Question}

\theoremstyle{definition}
\newtheorem{defn}[thm]{Definition}
\newtheorem{defns}[thm]{Definitions}
\newtheorem{con}[thm]{Construction}
\newtheorem{exmp}[thm]{Example}
\newtheorem{exmps}[thm]{Examples}
\newtheorem{notn}[thm]{Notation}
\newtheorem{algor}[thm]{Algorithm}
\newtheorem{exper}[thm]{Experiment}

\theoremstyle{remark}
\newtheorem{rem}[thm]{Remark}
\newtheorem{rems}[thm]{Remarks}
\newtheorem{warn}[thm]{Warning}
\newtheorem{exer}{Exercise}[chapter]
\newtheorem{prob}{Problem}[chapter]
\newtheorem*{soln}{Solution}

%--------Equation Numbering--------
\numberwithin{equation}{chapter}

%--------Custom Commands--------
\newcommand{\Mol}{\mathrm{Mol}}
\newcommand{\Scaf}{\mathrm{Scaf}}
\newcommand{\Fib}{\mathrm{Fib}}
\newcommand{\Err}{\mathrm{Err}}
\newcommand{\Var}{\mathrm{Var}}
\newcommand{\Cov}{\mathrm{Cov}}
\newcommand{\E}{\mathbb{E}}
\newcommand{\R}{\mathbb{R}}
\newcommand{\N}{\mathbb{N}}
\newcommand{\cS}{\mathcal{S}}
\newcommand{\cA}{\mathcal{A}}
\newcommand{\cK}{\mathcal{K}}
\newcommand{\cD}{\mathcal{D}}
\newcommand{\ScafOverhead}{\mathrm{ScafOverhead}}


% Shared response-model notation.
\newcommand{\Sset}{\mathcal{S}}
\newcommand{\Dset}{\mathcal{D}}
\newcommand{\Aspace}{\mathcal{A}}
\newcommand{\one}{\mathbf{1}}
\newcommand{\eps}{\varepsilon}
\newcommand{\rank}{\operatorname{rank}}
\newcommand{\im}{\operatorname{im}}
\newcommand{\tr}{\operatorname{tr}}
\newcommand{\diag}{\operatorname{diag}}
\crefname{thm}{theorem}{theorems}
\Crefname{thm}{Theorem}{Theorems}
\crefname{prop}{proposition}{propositions}
\Crefname{prop}{Proposition}{Propositions}
\crefname{cor}{corollary}{corollaries}
\Crefname{cor}{Corollary}{Corollaries}
\crefname{lem}{lemma}{lemmas}
\Crefname{lem}{Lemma}{Lemmas}

%--------Meta Data--------
\title{\Huge\textbf{The Algebra of Chemical Space}\\[1em]
       \Large Abstraction, Transfer, and Molecular Decisions}
\author{Austin Clyde}
\date{September 15, 2026}
\setcounter{tocdepth}{1}
\setlength{\headheight}{24pt}

%--------Document--------
\begin{document}

\frontmatter

% Title page
\maketitle

% Dedication
\chapter*{Dedication}
% BEGIN SOURCE: frontmatter/dedication.tex
For all my teachers
% END SOURCE: frontmatter/dedication.tex


% Epigraph
\chapter*{}
% BEGIN SOURCE: frontmatter/epigraph.tex
\begin{quote}
``Thoughts without content are empty; intuitions without concepts are blind.''
\begin{flushright}
Immanuel Kant
\end{flushright}
\end{quote}
% END SOURCE: frontmatter/epigraph.tex


% Abstract
\chapter*{Abstract}
% BEGIN SOURCE: frontmatter/abstract.tex
A chemical abstraction groups molecules that a model will treat as equivalent.
This book asks what that choice preserves, what it discards, and which molecular
decisions remain justified after compression. Deterministic maps induce fibers,
closure operators, and conditional expectations. These constructions quantify
retrieval overhead and prediction error for a chosen representation; they do
not establish that a particular scaffold definition is useful for a target.

The mathematical center is a finite response model whose coordinates are
chemical contexts and consistently defined decorations. Mixed differences test
additive transfer. Sparse response graphs extend those tests to incomplete
chemical combinations, while a parent-response map prevents multiple
representations of one measurement from being counted as independent evidence.
Under a fixed linear model, estimable contrasts and robust ranking certificates
separate the measurements needed for a molecular decision from those needed to
recover an entire response table. The assumptions supporting each certificate
are explicit.

The algorithmic chapters distinguish output size from evaluation cost, scaffold
means from molecular extremes, and representation levels from executable
fidelity levels. They develop conditional screening and allocation rules and
counterexamples to stronger claims. Retrospective chemical data identifies
structural opportunities and rejects simple transfer models; it does not yet
demonstrate prospective discovery gains. The resulting research program uses
RLMM as a proposed construction and evaluation system for testing whether
representation checks can reduce the cost of a correct molecular decision.

\bigskip
\noindent\textbf{Keywords:} molecular abstraction, scaffold, response curvature,
transfer, experimental design, decision certificate, sequential inference
% END SOURCE: frontmatter/abstract.tex


% Acknowledgments
\chapter*{Acknowledgments}
% BEGIN SOURCE: frontmatter/acknowledgments.tex
While there are many to thank, there a few that cannot go without naming. Rick Stevens, thank you for believing in my work, entrusting my research sprawl, and letting me fly with curiosity.
% END SOURCE: frontmatter/acknowledgments.tex


% Table of contents
\tableofcontents
\listoffigures
\listoftables

% Notation
\chapter*{Notation}
\addcontentsline{toc}{chapter}{Notation}
% BEGIN SOURCE: frontmatter/notation.tex
\begin{tabular}{@{}ll@{}}
\toprule
\textbf{Symbol} & \textbf{Meaning} \\
\midrule
$\Mol$ & Space of molecules (isomorphism classes) \\
$S$ & Space of scaffolds \\
$\Scaf: \Mol \to S$ & Scaffold map \\
$\Fib(s)$ & Fiber of scaffold $s$: $\{m : \Scaf(m) = s\}$ \\
$\cS = \sigma(\Scaf)$ & Scaffold $\sigma$-algebra \\
$\mu$ & Probability measure on $\Mol$ (library distribution) \\
$F: \Mol \to \R$ & Scoring function \\
$A(s) = \E[F \mid \Scaf = s]$ & Scaffold signal \\
$\varepsilon = F - A \circ \Scaf$ & Residual unresolved by scaffold identity \\
$\Err(\Scaf)$ & Irreducible scaffold error: $\E[\Var(F \mid \Scaf)]$ \\
\midrule
$\preceq$ & Partial order (subgraph embedding) \\
$\uparrow s$, $\downarrow s$ & Upper cone, lower cone \\
$\cD(S)$ & Lattice of downsets of $S$ \\
\midrule
$\alpha: \mathcal{P}(\Mol) \to \mathcal{P}(S)$ & Abstraction map \\
$\gamma: \mathcal{P}(S) \to \mathcal{P}(\Mol)$ & Concretization map \\
$c = \gamma \circ \alpha$ & Closure operator \\
$\ScafOverhead(A)$ & $|c(A)|/|A|$ \\
\midrule
$\cS_0 \subseteq \cdots \subseteq \cS_L$ & Nested information from refinement \\
$A_\ell = \E[F \mid \cS_\ell]$ & Signal at level $\ell$ \\
$D_\ell = A_\ell - A_{\ell-1}$ & Martingale difference \\
$\sigma(F)$ & Variance spectrum \\
$L^*(\epsilon)$ & Intrinsic scaffold depth \\
\bottomrule
\end{tabular}
% END SOURCE: frontmatter/notation.tex


\mainmatter

%--------Part I: Foundations--------
\part{Molecular Domains and Abstractions}
\label{part:foundations}

\chapter{The Problem of Chemical Space}
\label{ch:problem}
% BEGIN SOURCE: chapters/ch01_problem.tex
% Chapter content; main.tex supplies the chapter title and label.
\section{A molecular decision under a budget}

A molecular campaign must choose what to construct or evaluate next. Its
candidate set may be large, its measurements expensive, and its response model
wrong in ways that matter only near the final decision. The useful question is
therefore more specific than how to enumerate chemical space:
\begin{quote}
Which representation lets us reach a declared molecular decision with less
measurement and computation, while exposing the errors that could change it?
\end{quote}

Consider a finite family of molecules made by combining chemical contexts with
allowed decorations. We want to select a set of candidates with the best value
of one specified endpoint. An additive model predicts that a decoration has the
same effect in different contexts. If that prediction is reliable, measurements
can be shared across the family. If it fails, the campaign must detect the
failure before using the prediction to discard a good candidate.

This example already separates construction, prediction, and verification.
A valid molecular graph need not have a practical synthesis. A docking score
need not estimate experimental potency. A small average prediction error need
not preserve the best molecule. Each connection needs its own assumptions and
evidence.

\section{What a scaffold does and does not provide}

Chemical graph theory organizes molecular structure; Cayley's work on trees
and chemical isomers is an early example\cite{cayley1857trees,cayley1874isomers}.
Scaffold extraction provides another kind of organization. A deterministic
map $\Scaf:\Mol\to S$ partitions a declared molecular domain into fibers
$\Fib(s)=\Scaf^{-1}(s)$. Bemis--Murcko scaffolds are an important chemical
choice\cite{bemis1996scaffold}, but the algebra of a map does not select that
choice for us.

Every deterministic abstraction induces the same basic constructions: images
and preimages, a closure that saturates a set by its fibers, and a conditional
expectation that gives the best prediction measurable with respect to the
abstraction. These are tools for auditing a representation. A representation
earns its role in a campaign if its preserved information reduces the cost of
a correct decision on measurements not used to select it.

We distinguish a literal graph-theoretic core from a software extraction
protocol. Labels, aromaticity, protonation, stereochemistry, and treatment of
acyclic molecules affect the latter. A theorem about a graph operation does not
automatically describe every cheminformatics implementation of a scaffold.
Chapter~\ref{ch:scaffolds} makes this distinction explicit.

\section{The claims developed in this book}

\subsection{Closure describes retained sets}
Given $A\subseteq\Mol$, its fiber closure is
\[
 c(A)=\Scaf^{-1}(\Scaf(A)).
\]
It is the smallest union of complete scaffold fibers containing $A$. For a
finite nonempty $A$, the ratio $|c(A)|/|A|$ measures the enlargement caused by
that restriction. It is an output-cardinality statement. A method can return
the entire library without querying any response, so this ratio alone cannot
be a lower bound on evaluation cost. Complexity results also need a success
criterion that excludes such uninformative outputs and a specified oracle model.

\subsection{Conditional variance describes lost predictive information}
For a fixed sampling measure and square-integrable response $F$,
\[
 \Var(F)=\Var(\E[F\mid\Scaf])+\E[\Var(F\mid\Scaf)].
\]
The second term is the minimum mean-squared error of a predictor restricted to
scaffold identity. Refining the available information cannot increase this
population error. Estimating a finer predictor from finite data can still make
its future error worse. Neither a small population mean-squared error nor a
favorable scaffold mean supplies a uniform molecular ranking guarantee.

\subsection{Transfer is a hypothesis with observable consequences}
For a response table $Y_{sd}$, additive transfer means $Y_{sd}=a_s+b_d$.
The four-corner contrast
\[
 Y_{sd}-Y_{sd'}-Y_{s'd}+Y_{s'd'}
\]
vanishes for every rectangle exactly when the complete table is additive.
This is classical interaction analysis and chemical double-transformation-cycle
analysis\cite{free1964qsar,tukey1949nonadditivity,kramer2015nonadditivity}.
Its use requires decorations with a declared common meaning and compatible
response measurements. On sparse graphs, longer cycles may be needed. Repeated
decompositions of one measured molecule must be accounted for at the parent
level rather than treated as fresh observations.

\subsection{A certificate is conditional on its evidence model}
A top-$k$ certificate asks whether the same set wins for every response allowed
by a fixed model and the observation uncertainty. It can succeed without
recovering every model coordinate. It can also fail because a needed contrast
is unobserved, the observed data conflicts with the model, or the remaining
uncertainty changes the ranking.

Chapter~\ref{ch:response-certificates} develops these finite statements. Exact
algebra, numerical verification, statistical calibration, and chemical
validation have different roles. A verified numerical identity supports an
implementation of an algebraic result. A confidence statement needs a valid
sampling model. A claim about chemical transfer needs compatible chemical
measurements. A computational certificate concerns its computational endpoint
unless an additional calibration justifies a biological interpretation.

\section{From theory to an executable study}

The first research objective is a finite candidate family, one homogeneous
endpoint, and a declared decision. We compare policies by their total cost to
make that decision, their error rate, and their ability to reject an inadequate
response model. Costs include initial training observations, preparation,
failed evaluations, and confirmatory measurements.

RLMM supplies a possible connection to executable molecular design: propose a
molecule, prepare a physical representation, evaluate it, and update the next
choice. The existing RLMM2 software and its stored benchmarks are useful
starting evidence, but this book does not infer a new guarantee from them.
The proposed integration must retain protocol-specific observations, distinguish
cached results from fresh replicates, and audit what its chemical actions
actually do. These software changes and new experiments are future work.

The current response-graph audits are retrospective. Their structural counts
show which comparisons can be formed, while their prediction errors reject
simple additive transfer on the aggregated data. Missing assay comparability
prevents those results from establishing chemical mechanisms or prospective
assay savings. These negative results determine what the next experiment must
control.

\section{How to read the revised volume}

Part~\ref{part:foundations} defines molecular domains and abstraction protocols.
Part~\ref{part:algebra} develops order, closure, and information loss.
Part~\ref{part:analysis} introduces probability and response models, culminating
in Chapter~\ref{ch:response-certificates}. Part~\ref{part:algorithms} studies
screening, multilevel estimation, and refinement under explicit assumptions.
Part~\ref{part:extensions} examines transfer and generation and ends with a
bounded research program.

A short route through the main argument is Chapters~\ref{ch:scaffolds},
\ref{ch:galois}, \ref{ch:variance}, \ref{ch:response-certificates},
\ref{ch:screening}, and \ref{ch:open}. The order-theoretic and information-geometric
chapters provide background rather than prerequisites for every decision result.

The appendices give supporting definitions and selected verified examples.
Earlier exploratory material on renormalization, tropical and quantum
perspectives, other molecular domains, and the unreconciled solution collection
remains in the project sources. It is outside this revised volume and is not
used as a premise for its decision guarantees.
% END SOURCE: chapters/ch01_problem.tex


\chapter{Molecules as Mathematical Objects}
\label{ch:molecules}
% BEGIN SOURCE: chapters/ch02_molecules.tex
% ======================================================================
% Chapter 2 content file: chapters/ch02_molecules.tex
% (Do NOT include \chapter{...} here; main.tex already does that.)
% ======================================================================

\section{Why formalize molecules?}
Chemistry already has excellent informal abstractions for ``structure'': rings, functional groups,
scaffolds, fragments, pharmacophores, and synthetic transforms. The purpose of this chapter is not
to replace that language, but to \emph{pin down} a mathematical object that:

\begin{enumerate}
\item is faithful to how cheminformatics represents small molecules (graphs with typed vertices/edges),
\item supports a canonical notion of \emph{containment} (substructure embedding),
\item is compatible with \emph{probability} (libraries as distributions), and
\item is sufficiently structured to support later algebraic and analytic constructions
(Galois connections, closure operators, conditional expectation\index{conditional expectation}, and refinement towers).
\end{enumerate}

A recurring theme is that many popular representations (strings, fingerprints, continuous
embeddings) are not wrong, but they are \emph{coordinate systems} or \emph{quotients} that
discard structure we will later need. The book will treat them as derived objects from a more
canonical substrate.

\subsection*{Why distinguish common substructures from a greatest one?}

Two molecules can have several incomparable maximal common substructures.
The explicit example below shows that even simple carbon skeletons need not have a meet or a join under subgraph embedding.
A maximum common substructure algorithm therefore needs an objective and a rule for resolving ties; the order alone may not select one answer.

Chapter~\ref{ch:lattice} retains the whole set of lower bounds as a downset\index{downset}.
Downsets have canonical unions and intersections, although representing them and deciding graph embeddings can remain expensive.
The fiber closure in Chapter~\ref{ch:galois} addresses a separate issue: the loss of resolution caused by a many-to-one abstraction.
That closure exists whether or not the molecule order is a lattice.
Neither lattice failure nor fiber closure proves that a generative model must produce invalid molecules; such a claim requires assumptions about its decoder and a separate evaluation.

\section{Chemical graphs: typed graphs satisfying rules}
\label{sec:ch2-chemical-graphs}

\subsection{Atom and bond type alphabets}
We begin with the discrete alphabets that specify the ``typing'' of graphs.

\begin{defn}[Atom types]
\label{def:atom-types}
Let $A$ be a finite set of \emph{atom types}. In practice, an atom type packages the information
that must be preserved under isomorphism:
atomic number, formal charge, aromaticity flag, and (optionally) isotope and chirality tags.
We write $a\in A$ for an atom type.
\end{defn}

\begin{defn}[Bond types]
\label{def:bond-types}
Let $B$ be a finite set of \emph{bond types} (single, double, triple, aromatic, etc.).
We write $b\in B$ for a bond type.
\end{defn}

The sets $A$ and $B$ can be enriched if desired: stereochemical edge annotations, dative bonds,
coordinate bonds, and ``unknown'' placeholders. The only requirement for the mathematics below
is that these are discrete types.

\subsection{Typed graphs}
\begin{defn}[Typed graph]
\label{def:typed-graph}
A \emph{typed graph} (over $(A,B)$) is a tuple
\[
G = (V,E,\tau_V,\tau_E),
\]
where $V$ is a finite vertex set,
$E \subseteq \{\{u,v\}: u,v\in V,\, u\neq v\}$ is an undirected edge set,
$\tau_V:V\to A$ assigns an atom type to each vertex,
and $\tau_E:E\to B$ assigns a bond type to each edge.
\end{defn}

\begin{rem}[Connectedness]
Most small-molecule\index{molecule} discussions assume connected graphs.
Salts, mixtures, and multi-component ``molecules'' can be modeled as disconnected typed graphs or
as connected graphs with a special ``ionic'' edge type.
Unless explicitly stated, we restrict attention to connected graphs.
\end{rem}

\subsection{Chemical validity as rule satisfaction}
Graph typing alone does not define chemistry. We also need a notion of validity.

\begin{defn}[Chemical rule]
\label{def:chemical-rule}
A \emph{chemical rule} is a predicate
\[
r:\{\text{typed graphs over }(A,B)\}\to\{0,1\}.
\]
A set $C$ of chemical rules defines validity by conjunction:
a typed graph $G$ is \emph{$C$-valid} if $r(G)=1$ for all $r\in C$.
\end{defn}

\begin{defn}[Molecule (graph-theoretic)]
\label{def:molecule-graph-theoretic}
Let $\widetilde{\Mol}$ denote the set of connected typed graphs over $(A,B)$ that satisfy the rule set $C$:
\[
\widetilde{\Mol} \;:=\; \{\, G \text{ a connected typed graph over }(A,B) : \forall r\in C,\ r(G)=1 \,\}.
\]
Elements of $\widetilde{\Mol}$ are \emph{molecules as typed graphs}.
\end{defn}

\begin{rem}[What is in $C$?]
The rule set $C$ is a modeling choice. In typical cheminformatics pipelines $C$ includes (at minimum)
valence\index{valence} constraints and basic aromaticity rules; it may also include allowed charge states,
ring closure constraints, forbidden motifs, or dataset-specific filters.
Later chapters will treat $C$ as fixed and focus on the structure induced by $\widetilde{\Mol}$.
\end{rem}

\begin{exmp}[Methane as a typed graph]
Let $A=\{\mathrm{C},\mathrm{H}\}$ and $B=\{\mathrm{single}\}$.
Methane corresponds to a star on five vertices with center type $\mathrm{C}$ and four leaves type $\mathrm{H}$.
The validity rules enforce that carbon is tetravalent and hydrogen monovalent.
\end{exmp}

\section{Isomorphism classes and canonical representatives}
\label{sec:ch2-isomorphism}

The same molecule can be represented by many graphs if vertex identities are treated as names.
Chemically, the names do not matter. The invariant object is the isomorphism class.

\begin{defn}[Typed graph isomorphism]
\label{def:graph-iso}
Let $G=(V,E,\tau_V,\tau_E)$ and $G'=(V',E',\tau_V',\tau_E')$ be typed graphs.
A \emph{typed graph isomorphism} is a bijection $\phi:V\to V'$ such that
\begin{enumerate}
\item $\tau_V(v)=\tau_V'(\phi(v))$ for all $v\in V$,
\item $\{u,v\}\in E \iff \{\phi(u),\phi(v)\}\in E'$,
\item and whenever $\{u,v\}\in E$, we have $\tau_E(\{u,v\})=\tau_E'(\{\phi(u),\phi(v)\})$.
\end{enumerate}
We write $G\cong G'$ if such an isomorphism exists.
\end{defn}

\begin{defn}[Molecule space as isomorphism classes]
\label{def:Mol-quotient}
Define
\[
\Mol \;:=\; \widetilde{\Mol}/\cong,
\]
the set of isomorphism classes of $C$-valid typed graphs.
By convention, throughout this book ``molecule'' means an element of $\Mol$ (an isomorphism class),
unless stated otherwise.
\end{defn}

\begin{rem}[SMILES and canonicalization]
String encodings such as SMILES are best viewed as \emph{representatives} of isomorphism classes.
Different SMILES strings can represent the same element of $\Mol$.
Canonical SMILES can be treated as a choice of section of the quotient map
$\widetilde{\Mol}\twoheadrightarrow \Mol$, but the mathematics below never depends on that choice.
\end{rem}

\section{Substructure embeddings and the induced preorder}
\label{sec:ch2-embedding}

The central order-theoretic primitive for scaffold theory is substructure inclusion, modeled by
typed subgraph embeddings.

\subsection{Embeddings}
\begin{defn}[Typed subgraph embedding]
\label{def:subgraph-embed}
Let $G=(V,E,\tau_V,\tau_E)$ and $G'=(V',E',\tau_V',\tau_E')$ be typed graphs.
A \emph{typed subgraph embedding} of $G$ into $G'$ is an injective map $f:V\to V'$ such that:
\begin{enumerate}
\item Vertex types are preserved: $\tau_V(v)=\tau_V'(f(v))$ for all $v\in V$.
\item Edge incidence is preserved: if $\{u,v\}\in E$ then $\{f(u),f(v)\}\in E'$.
\item Edge types are preserved: if $\{u,v\}\in E$ then $\tau_E(\{u,v\})=\tau_E'(\{f(u),f(v)\})$.
\end{enumerate}
We write $G\hookrightarrow G'$ if such an embedding exists.
\end{defn}

\begin{rem}[Extra edges allowed]
An embedding requires the image of each edge to be an edge of $G'$, but does not forbid $G'$
from having extra edges between embedded vertices. In other words, this is a \emph{(not necessarily induced)}
subgraph embedding. This matches common chemical practice: a substructure may appear as part of a denser motif.
\end{rem}

\subsection{Preorder on graphs; poset\index{poset} on isomorphism classes}

Define a relation on $\widetilde{\Mol}$ by
\[
G \preceq G' \quad \Longleftrightarrow \quad G\hookrightarrow G'.
\]

\begin{prop}[Embedding induces a preorder]
\label{prop:preorder}
The relation $\preceq$ on $\widetilde{\Mol}$ is a preorder: it is reflexive and transitive.
\end{prop}

\begin{proof}
Reflexivity holds by taking the identity embedding $f=\mathrm{id}_V$.
Transitivity holds because composition of injective type-preserving maps is injective and type-preserving:
if $G\hookrightarrow G'$ via $f$ and $G'\hookrightarrow G''$ via $g$, then $g\circ f$ witnesses $G\hookrightarrow G''$.
\end{proof}

\begin{prop}[Antisymmetry fails on labeled graphs]
\label{prop:antisymmetry-fails}
On $\widetilde{\Mol}$, antisymmetry can fail: it is possible to have $G\preceq G'$ and $G'\preceq G$ with $G\neq G'$ as graphs.
\end{prop}

\begin{proof}
Take two isomorphic graphs $G\cong G'$ on different vertex sets. Then $G\hookrightarrow G'$ and $G'\hookrightarrow G$,
but $G\neq G'$ as labeled objects.
\end{proof}

This motivates working on isomorphism classes.

\begin{cor}[Poset on isomorphism classes]
\label{cor:poset-on-classes}
Define on $\Mol=\widetilde{\Mol}/\cong$ the relation
\[
[m]\preceq [m'] \quad \Longleftrightarrow \quad
\exists\, G\in[m],\ G'\in[m'] \text{ with } G\hookrightarrow G'.
\]
This is well-defined and makes $(\Mol,\preceq)$ a poset.
\end{cor}

\begin{proof}
Well-definedness follows because if $G\cong H$ and $G'\cong H'$, then $G\hookrightarrow G'$ implies $H\hookrightarrow H'$ by conjugating
the embedding with the isomorphisms. Reflexivity and transitivity inherit from Proposition~\ref{prop:preorder}.
For antisymmetry, embeddings in both directions give equality of both vertex counts and edge counts.
Either embedding is therefore bijective on vertices and maps the entire edge set onto the host edge set.
It preserves types and has an inverse that preserves adjacency, so it is an isomorphism.
Equal vertex counts alone would not suffice for a non-induced embedding.
\end{proof}

\begin{rem}[Computational note]
Deciding whether $m\preceq m'$ is the typed subgraph isomorphism problem.
This is NP-complete for unrestricted input graphs. Atom and bond types, small patterns, and indexing can reduce the work in particular chemical libraries;
these practical advantages do not give a general polynomial-time guarantee.
We will treat $\preceq$ as a mathematical primitive and discuss algorithmic surrogates later.
\end{rem}

\section{Why $(\Mol,\preceq)$ is not a lattice}
\label{sec:ch2-not-lattice}

Later chapters build lattices by lifting to \emph{regions} of scaffold space.
This section explains why we cannot expect lattice operations to exist on molecules themselves.

\subsection{Lower bounds, upper bounds, meet and join}

\begin{defn}[Bounds in a poset]
Let $(P,\preceq)$ be a poset and let $x,y\in P$.
An element $z$ is a \emph{lower bound} of $\{x,y\}$ if $z\preceq x$ and $z\preceq y$.
A \emph{greatest lower bound} (meet) $x\wedge y$ is a lower bound such that any lower bound $w$ satisfies $w\preceq x\wedge y$.
Dually, $u$ is an \emph{upper bound} of $\{x,y\}$ if $x\preceq u$ and $y\preceq u$.
A \emph{least upper bound} (join) $x\vee y$ is an upper bound such that any upper bound $v$ satisfies $x\vee y\preceq v$.
\end{defn}

In chemical language, meets correspond to a canonical ``maximum common substructure,'' and joins to a canonical ``minimal common superstructure.''
Both fail in general.

\subsection{Meet failure via non-unique maximal common substructures}

\begin{exmp}[Two incomparable maximal common substructures]
\label{exmp:meet-fails}
Use carbon-skeleton typing: all vertices have type carbon, all edges are single, and implicit hydrogen counts are not part of the preserved vertex type.
Assume the chemical domain contains the following graphs and the two four-vertex trees discussed below:
\begin{enumerate}
\item $G$: a ``T with a tail'' on five vertices (one vertex of degree 3, one vertex of degree 2, three leaves).
Its edge set is $\{12,13,14,45\}$; chemically it is the skeleton of 2-methylbutane.
\item $H$: a triangle with a pendant vertex, with edge set $\{12,23,13,14\}$.
This four-vertex ``paw'' graph is the skeleton of methylcyclopropane.
\end{enumerate}
Both contain the star $K_{1,3}$ and the path $P_4$.
In $H$, retain edges $12,13,14$ for the star, or $41,12,23$ for the path.
Every common subgraph has at most four vertices because it embeds in $H$, and is acyclic because it embeds in the tree $G$.
A common subgraph containing either four-vertex tree therefore has four vertices and exactly three edges, so it is isomorphic to that tree.
Since $K_{1,3}$ and $P_4$ have different degree sequences, they are incomparable and both are maximal common substructures.
No common subgraph can contain both, so there is no greatest lower bound.
\end{exmp}

\begin{figure}[htbp]
\centering
\begin{tikzpicture}[every node/.style={circle,fill,inner sep=1.5pt},scale=0.9]
\begin{scope}[xshift=0cm]
\node (a) at (0,0) {}; \node (b) at (-.8,.7) {}; \node (c) at (-.8,-.7) {};
\node (d) at (1,0) {}; \node (e) at (2,0) {};
\draw (a)--(b) (a)--(c) (a)--(d)--(e);
\node[fill=none,rectangle] at (.5,-1.15) {$G$: 2-methylbutane};
\end{scope}
\begin{scope}[xshift=5.2cm]
\node (a) at (0,0) {}; \node (b) at (1,.7) {}; \node (c) at (1,-.7) {}; \node (d) at (-1,0) {};
\draw (a)--(b)--(c)--(a)--(d);
\node[fill=none,rectangle] at (0,-1.15) {$H$: methylcyclopropane};
\end{scope}
\end{tikzpicture}
\caption{\textbf{Lattice failure.} In $(\Mol,\preceq)$, two non-isomorphic molecules $G$ and $H$ share incomparable maximal common substructures ($K_{1,3}$ and $P_4$), so no greatest lower bound exists.}
\label{fig:lattice-failure}
\end{figure}

\begin{prop}[Meet need not exist]
\label{prop:meet-need-not-exist}
In a domain containing the four carbon skeletons of Example~\ref{exmp:meet-fails}, there are pairs with no meet.
\end{prop}

\begin{proof}
Take $m_1=[G]$ and $m_2=[H]$ from Example~\ref{exmp:meet-fails}.
Any meet would have to be a greatest lower bound, but there is no greatest element among the lower bounds.
\end{proof}

\subsection{Join failure via non-unique minimal common superstructures}

The same four graphs also witness join failure.

\begin{exmp}[Two incomparable minimal common superstructures]
\label{exmp:join-fails}
Let $x$ be the star $K_{1,3}$ and $y$ the path $P_4$ (both realizable as hydrocarbon skeletons).
Both $G$ and $H$ in Example~\ref{exmp:meet-fails} are upper bounds of $x$ and $y$.
If a least upper bound $U$ existed, it would embed in both $G$ and $H$ and contain both $x$ and $y$.
It would therefore be an acyclic graph on four vertices containing two non-isomorphic four-vertex trees, which is impossible by the preceding argument.
Thus $x\vee y$ does not exist.
This argument uses non-induced embeddings throughout; a host graph can have smaller diameter than a path it contains.
\end{exmp}

\begin{prop}[Join need not exist]
\label{prop:join-need-not-exist}
In a domain containing those four carbon skeletons, there are pairs with no join.
\end{prop}

\begin{proof}
Example~\ref{exmp:join-fails} gives two upper bounds below which no common upper bound of $x$ and $y$ exists.
\end{proof}

\begin{cor}[Not a lattice]
\label{cor:not-a-lattice}
$(\Mol,\preceq)$ is a poset but not a lattice in general.
\end{cor}

\begin{proof}
A lattice requires existence of meets and joins for every pair.
Propositions~\ref{prop:meet-need-not-exist} and~\ref{prop:join-need-not-exist} provide counterexamples.
\end{proof}

\begin{rem}[Why this matters]
This failure is not a technical annoyance; it is the reason the book repeatedly lifts structure
from \emph{elements} (molecules/scaffolds) to \emph{regions} (downsets, upper cones, closures),
where union and intersection are well-defined. Their computational cost still depends on the representation of those regions.
\end{rem}

\section{Probability on molecules: libraries as measures}
\label{sec:ch2-probability}

The algebraic structures above become useful only once paired with a notion of sampling and scoring.

\subsection{Library distributions}
In the finite-library regime, $\Mol$ is finite and the natural $\sigma$-algebra\index{$\sigma$-algebra@$\sigma$-algebra} is the power set.
A probability measure $\mu$ is then a probability mass function over molecules.
Uniform $\mu$ corresponds to ``pick a random molecule from the library,'' while nonuniform $\mu$
can encode synthetic priors, vendor availability, or exploration policies.

\begin{defn}[Probability space of molecules]
\label{def:prob-space}
A \emph{library distribution} is a probability measure $\mu$ on $\Mol$.
We treat $(\Mol,\mu)$ as a discrete probability space unless stated otherwise.
\end{defn}

\subsection{Scoring functions as random variables}
\begin{defn}[Scoring function\index{scoring function}]
\label{def:scoring}
A \emph{scoring function}\index{scoring function} is a measurable map
\[
F:\Mol\to\R.
\]
Examples include docking scores, surrogate energies, binding free-energy estimates, and experimental readouts.
\end{defn}

\begin{defn}[$L^2(\Mol,\mu)$]
\label{def:L2}
Define the Hilbert space
\[
L^2(\Mol,\mu) := \Big\{\, F:\Mol\to\R \ \Big|\ \E[F^2] < \infty \,\Big\}
\]
where functions equal $\mu$-almost surely are identified, with inner product $\langle F,G\rangle = \E[FG]$.
\end{defn}

\begin{rem}[Why $L^2$?]
The central quantitative objects later in the book are variances, conditional variances,
and orthogonal decompositions under refinement. These live naturally in $L^2$.
\end{rem}

\section{Representations as quotients: what common encodings forget}
\label{sec:ch2-representations}

This book will repeatedly distinguish between the underlying mathematical object $\Mol$
and the various coordinate systems used in practice.

\subsection{Strings (SMILES) are coordinates}
SMILES is a convenient serialization, but it is not invariant: many strings represent the same molecule.
Canonicalization selects a representative but does not create new structure;
in particular, string edit operations are not intrinsically meaningful in $\Mol$ without additional interpretation.

\subsection{Fingerprints are lossy quotients}
Binary fingerprints $f:\Mol\to\{0,1\}^d$ are designed for similarity search and fast indexing.
They are typically \emph{many-to-one}, hence quotient $\Mol$ into collision classes.
They also do not preserve the poset structure $(\Mol,\preceq)$ unless carefully engineered.

\begin{rem}[Order information is fragile]
A representation $\rho:\Mol\to X$ preserves substructure information only if it is (at least) monotone:
$m\preceq m'$ should imply $\rho(m)$ is ``contained in'' $\rho(m')$ under some order on $X$.
Most commonly used fingerprints do not satisfy a clean monotonicity law in this sense.
\end{rem}

\subsection{Continuous embeddings hide compositional structure}
Learned embeddings $\phi:\Mol\to\R^d$ are powerful for interpolation and regression,
but the map $\phi$ typically does not preserve:
\begin{itemize}
\item the existence-of-embedding order $\preceq$,
\item the combinatorial asymmetry between ``going down'' (fragmentation) and ``going up'' (generation),
\item or canonical closure operations induced by scaffold quotients (developed in Chapters~\ref{ch:scaffolds}--\ref{ch:galois}).
\end{itemize}
This is not a criticism; it clarifies the target:
the book seeks algebraic structure that remains visible after quotienting and conditioning,
rather than structure that is smoothed away.

\section{Molecules as objects in a category}
\label{sec:ch2-category}

There is a category-theoretic formulation that clarifies later constructions.

\begin{defn}[Category of molecules and embeddings]
\label{def:MolEmb}
Define a category $\mathbf{MolEmb}$ as follows:
\begin{itemize}
\item Objects are elements of $\widetilde{\Mol}$ (valid typed graphs).
\item Morphisms $G\to G'$ are typed subgraph embeddings $G\hookrightarrow G'$.
\item Composition is composition of embeddings; identities are identity embeddings.
\end{itemize}
\end{defn}

\begin{rem}[From categories to orders]
The preorder on $\widetilde{\Mol}$ is the decategorified statement ``there exists a morphism.''
Passing to isomorphism classes yields the poset $(\Mol,\preceq)$.
When we later build scaffold maps, closure operators, and refinement towers,
the resulting set maps induce structures on powersets and function spaces.
A functor on $\mathbf{MolEmb}$ additionally requires a compatible action on embeddings; an arbitrary deterministic abstraction need not have one.
The order and powerset constructions used below do not require a categorical extension to other chemical domains.
\end{rem}

\section*{Summary}
We now have the foundational object for the rest of the book:

\begin{enumerate}
\item A molecule is an isomorphism class in $\Mol$ of a $C$-valid typed graph over $(A,B)$.
\item Substructure containment induces a poset $(\Mol,\preceq)$. The carbon-skeleton example shows that this poset need not be a lattice.
\item A library is a probability measure $\mu$ on $\Mol$, and scores $F$ live in $L^2(\Mol,\mu)$.
\end{enumerate}

Chapter~\ref{ch:scaffolds} introduces the scaffold map $\Scaf:\Mol\to S$ as a structured quotient of $\Mol$,
and begins the central theme of the manuscript: \emph{abstraction induces algebra}.

\section*{Exercises}
\begin{exer}[Chemistry: when ``shared core'' is not unique]
\label{exer:ch2-chem}
Work with \emph{heavy-atom carbon skeletons} (ignore hydrogens; treat all atoms as type \texttt{C} and all bonds as single).
Consider the two hydrocarbons:
\[
m_1=\text{2-methylbutane (isopentane)},\qquad
m_2=\text{methylcyclopropane}.
\]
\begin{enumerate}
\item Draw the typed graphs (carbon skeletons) for $m_1$ and $m_2$.
\item Show that the following two non-isomorphic graphs are both common substructures of $m_1$ and $m_2$:
\[
K_{1,3}\quad\text{(a 3-star)},\qquad P_4\quad\text{(a 4-vertex path)}.
\]
\item Show that every common substructure has at most four vertices and is acyclic. Prove that no common substructure can contain both displayed trees, and conclude that no meet exists.
\item Explain in chemical language why both $K_{1,3}$ and $P_4$ are plausible ``cores'' depending on the intended downstream task (e.g.\ branching vs.\ ring/chain similarity).
\end{enumerate}
\end{exer}

\begin{exer}[Mathematics: preorder $\to$ poset, and lattice failure]
\label{exer:ch2-math}
Let $\widetilde{\Mol}$ be valid typed graphs and $\Mol=\widetilde{\Mol}/\cong$ their isomorphism classes.
Define $m\preceq m'$ iff a typed subgraph embedding exists.
\begin{enumerate}
\item Prove that $\preceq$ is a preorder on $\widetilde{\Mol}$.
\item Prove that $\preceq$ induces a partial order\index{partial order} on $\Mol$ (i.e.\ antisymmetry holds on isomorphism classes).
\item Using Exercise~\ref{exer:ch2-chem} (or any explicit counterexample you prefer), prove that $(\Mol,\preceq)$ is not a lattice by showing:
  \begin{enumerate}
  \item there exist $x,y\in\Mol$ with no meet $x\wedge y$,
  \item there exist $u,v\in\Mol$ with no join $u\vee v$.
  \end{enumerate}
\item (Optional) Identify a nontrivial subclass of molecules (or graphs) for which meets \emph{do} exist, and justify your claim.
\end{enumerate}
\end{exer}

\begin{exer}[Machine learning: irreducible error\index{irreducible error} from lossy representations]
\label{exer:ch2-ml}
Let $(\Mol,\mu)$ be a finite library with a probability mass function $\mu$.
Let $F:\Mol\to\R$ be a scoring function with $\E[F^2]<\infty$.
Let $\rho:\Mol\to X$ be any many-to-one representation (e.g.\ a fingerprint, descriptor map, or learned embedding).
Let $\sigma(\rho)$ denote the $\sigma$-algebra\index{$\sigma$-algebra@$\sigma$-algebra} generated by $\rho$.
\begin{enumerate}
\item Prove that among all $\rho$-measurable predictors $g(\rho(\cdot))$, the unique minimizer of
\[
\min_{g} \; \E\!\left[(F - g\circ \rho)^2\right]
\]
is $g^*(x)=\E[F\mid \rho=x]$.
\item Show that the minimum achievable mean-squared error is
\[
\E\!\left[\Var(F\mid \rho)\right] \;=\; \E\!\left[(F-\E[F\mid\sigma(\rho)])^2\right].
\]
\item Give an explicit collision example: construct two distinct molecules $m\neq m'$ such that $\rho(m)=\rho(m')$ but $F(m)\neq F(m')$.
Compute the contribution of this collision to the irreducible error\index{irreducible error}.
\item Interpret the result as an \emph{information bound}: explain why increasing model capacity cannot beat $\E[\Var(F\mid \rho)]$.
\end{enumerate}
\end{exer}

\begin{exer}[Computation: building the embedding order in practice]
\label{exer:ch2-comp}
Using a cheminformatics toolkit (e.g.\ RDKit) on a small library (e.g.\ $N\approx 10^3$--$10^4$ molecules):
\begin{enumerate}
\item Implement a predicate \texttt{embeds(m, m')} that returns true if $m\preceq m'$ via a substructure match.
\item Construct the directed graph $G_\preceq$ on molecules with an edge $m\to m'$ when $m\preceq m'$ and $m\neq m'$.
\item Empirically verify:
  \begin{enumerate}
  \item transitivity holds (spot-check on random triples),
  \item antisymmetry fails on raw labeled graphs but holds on canonicalized/isomorphism representatives,
  \item meets and joins frequently fail to exist (estimate the fraction of pairs lacking a meet or join under your chosen common-substructure heuristic).
  \end{enumerate}
\item Report how the ``meet failure rate'' changes as you restrict to molecules below a size cutoff (e.g.\ heavy atoms $\leq 15, 25, 35$).
\end{enumerate}
\end{exer}
% END SOURCE: chapters/ch02_molecules.tex


\chapter{Scaffold Abstraction}
\label{ch:scaffolds}
% BEGIN SOURCE: chapters/ch03_scaffolds.tex
% ======================================================================
% Chapter 3 content file: chapters/ch03_scaffolds.tex
% (Do NOT include \chapter{...} here; main.tex already does that.)
% ======================================================================

\section{Why scaffolds? Abstraction as a design primitive}
The central observation motivating this book is empirical but conceptually sharp:

\begin{quote}
\emph{In large molecular libraries, high-scoring molecules for a target are not scattered uniformly across chemical space\index{chemical space}; they concentrate in a small number of structural families.}
\end{quote}

Medicinal chemistry has always acted on this implicitly: the unit of reasoning is often not a single molecule,
but a \emph{series}---a scaffold\index{scaffold} with a cloud of substituent variations and an evolving structure--activity relationship (SAR).
The mathematical move in this chapter is to turn ``series'' into an explicit quotient:
a deterministic map from molecules to \emph{scaffolds}. The map induces fibers and a pushforward distribution; additional graph structure supports the other constructions:

\begin{itemize}
\item \textbf{Algebra:} A scaffold\index{scaffold} map defines an equivalence relation and partitions chemical space\index{chemical space} into fibers.
\item \textbf{Order:} Scaffolds inherit a natural partial order\index{partial order} from substructure embedding.
\item \textbf{Probability:} A library distribution on molecules pushes forward to a distribution on scaffolds.
\item \textbf{Computation:} ``Going down'' (fragmentation) is algorithmic; ``going up'' (generation) is combinatorial.
\end{itemize}

This chapter fixes a typed graph 2-core as a mathematical baseline for ring-and-linker abstraction, while treating each operational scaffold map as a \emph{choice of abstraction}.
Later chapters will quantify what any abstraction buys you and what it irreversibly destroys.

\section{The graph 2-core and Bemis--Murcko\index{Bemis--Murcko scaffold} conventions}
\label{sec:ch3-murcko}

\subsection{Heavy-atom graphs and implicit hydrogens}
Throughout, we treat molecules as typed graphs as in Chapter~\ref{ch:molecules}. In practice, most scaffold pipelines operate
on \emph{heavy-atom graphs} with implicit hydrogens. Formally, we can either:
(i) include hydrogen vertices and let rules $C$ enforce valence\index{valence}, or
(ii) drop hydrogen vertices and encode implicit hydrogen counts in the atom typing.
Both are compatible with the formalism in Chapter~\ref{ch:molecules}.

For definiteness, the scaffold extraction described below is easiest to state on heavy-atom graphs:
remove vertices that are chemically peripheral (side chains) while preserving ring systems and linkers.

\subsection{A graph-theoretic characterization: the 2-core\index{2-core}}
The graph 2-core gives a precise model of the ring-and-linker backbone used in scaffold abstractions.
We distinguish this operator from a toolkit's Bemis--Murcko implementation: chemical conventions may retain exocyclic atoms or alter atom and bond types.
The graph proofs below concern the stated 2-core map; a software pipeline must be checked against that definition before those proofs are applied to its output.

\begin{defn}[$k$-core]
Let $G=(V,E)$ be an undirected graph (ignore types for the moment).
The \emph{$k$-core} of $G$ is the maximal induced subgraph $G_k$ in which every vertex has degree at least $k$.
Equivalently, $G_k$ is obtained by iteratively removing vertices of degree $<k$ until none remain.
\end{defn}

For $k=2$, this removes all pendant trees and retains cyclic structure and any ``bridges'' that connect cycles.
This is the ring-and-linker backbone modeled here.

\begin{defn}[Typed 2-core scaffold extraction\index{2-core}]
\label{def:murcko-2core}
Let $m\in\Mol$ be a molecule\index{molecule} (isomorphism class of a valid typed graph).
Choose a representative typed graph $G=(V,E,\tau_V,\tau_E)$.
Let $H(G)$ denote the underlying untyped graph $(V,E)$.
Compute the 2-core $H(G)_2$, and let $V_2\subseteq V$ be its vertex set.
Define the \emph{scaffold graph} $\mathrm{ScafGraph}(G)$ to be the typed induced subgraph on $V_2$:
\[
\mathrm{ScafGraph}(G) \;:=\; (V_2,\,E|_{V_2},\,\tau_V|_{V_2},\,\tau_E|_{E|_{V_2}}),
\]
where $E|_{V_2}:=\{e\in E: e\subseteq V_2\}$ denotes the restriction of the edge set to vertex pairs within $V_2$.
If $V_2=\emptyset$, define $\mathrm{ScafGraph}(G)$ to be the empty graph.
\end{defn}

\begin{rem}[Acyclic molecules]
If a molecule\index{molecule} has no cycles, its 2-core is empty. Under this baseline, such molecules map to a trivial scaffold.
This is not a bug: it reflects a ring-centric notion of scaffold. If you want acyclic ``backbone'' scaffolds,
you must choose a different abstraction (see \S\ref{sec:ch3-variants}).

\emph{Convention for later chapters:} We treat the empty scaffold as a distinguished element $\bot\in S$.
When computing statistics over $S$, one may either include or exclude $\bot$ depending on context;
we will note explicitly when it matters.
The fibers, including the acyclic fiber, must partition the probability space used for conditional expectations.
If acyclic molecules are excluded from an analysis, restrict the molecular domain and renormalize its measure explicitly; do not simply omit a positive-mass fiber.
\end{rem}

\subsection{Well-definedness on isomorphism classes}
Definition~\ref{def:murcko-2core} used a chosen representative $G$ of the isomorphism class $m$.
We must verify that scaffold extraction is invariant under isomorphism.

\begin{prop}[Scaffold extraction is well-defined on $\Mol$]
\label{prop:scaf-well-defined}
If $G\cong G'$ as typed graphs, then $\mathrm{ScafGraph}(G)\cong \mathrm{ScafGraph}(G')$.
Hence the typed 2-core scaffold depends only on the isomorphism class $m\in\Mol$.
\end{prop}

\begin{proof}
A typed graph isomorphism $\phi:V\to V'$ preserves adjacency and vertex/edge types.
In particular, it preserves degrees, hence preserves the iterative pruning procedure defining the 2-core.
Therefore $\phi$ maps the 2-core vertex set $V_2$ of $G$ bijectively to the 2-core vertex set $V_2'$ of $G'$.
Restricting $\phi$ to $V_2$ gives an isomorphism of the induced typed subgraphs.
\end{proof}

\section{The scaffold map and scaffold space}
\label{sec:ch3-scaf-map}

\subsection{Defining the scaffold space $S$}
We now define the scaffold space as the set of all scaffolds realized by molecules under the chosen extraction.

\begin{defn}[Scaffold space and scaffold map]
\label{def:scaf-map}
Let $S$ be the set of isomorphism classes of scaffold graphs produced by Definition~\ref{def:murcko-2core}:
\[
S \;:=\; \{\, [\mathrm{ScafGraph}(G)] : [G]\in\Mol \,\}.
\]
Define the \emph{scaffold map} $\Scaf:\Mol\to S$ by
\[
\Scaf([G]) \;:=\; [\mathrm{ScafGraph}(G)].
\]
\end{defn}

By definition, $\Scaf$ is surjective onto its image (which is $S$). The important point is what it discards:
everything outside the 2-core, i.e.\ all pendant subgraphs (trees) attached to the cyclic backbone.

\subsection{Fibers and equivalence classes}
A scaffold map induces a canonical partition of chemical space.

\begin{defn}[Fiber\index{fiber}]
\label{def:fiber}
For $s\in S$, the \emph{fiber}\index{fiber} over $s$ is
\[
\Fib(s)\;:=\;\{\, m\in\Mol:\Scaf(m)=s \,\}.
\]
\end{defn}

\begin{defn}[Scaffold equivalence relation]
\label{def:scaf-equivalence}
Define $m\sim m'$ iff $\Scaf(m)=\Scaf(m')$.
\end{defn}

\begin{prop}[Fibers partition $\Mol$]
\label{prop:fibers-partition}
The relation $\sim$ is an equivalence relation on $\Mol$, and the equivalence classes are exactly the fibers $\Fib(s)$.
In particular,
\[
\Mol \;=\;\bigsqcup_{s\in S}\Fib(s).
\]
\end{prop}

\begin{proof}
Reflexivity and symmetry are immediate from equality; transitivity follows from transitivity of equality.
Each equivalence class is the preimage of a scaffold value, hence a fiber.
Disjointness and the union decomposition follow because $\Scaf$ is a function.
\end{proof}

\begin{prop}[Surjective, not injective]
\label{prop:surj-not-inj}
$\Scaf:\Mol\to S$ is surjective and typically not injective.
\end{prop}

\begin{proof}
Surjectivity holds because $S$ is defined as the image of $\Scaf$.
Non-injectivity holds whenever there exist two distinct molecules with the same scaffold (e.g.\ any scaffold with at least two distinct substituent patterns).
In any realistic library, this occurs ubiquitously.
\end{proof}

\begin{rem}[A hypergraph view]
You can view the partition $\{\Fib(s)\}_{s\in S}$ as a hypergraph on vertex set $\Mol$ whose hyperedges are the fibers.
This hypergraph is \emph{1-regular}: every molecule belongs to exactly one hyperedge.
This is mathematically trivial but conceptually useful: later, refinement towers correspond to hypergraph coarsenings and refinements.
\end{rem}

\subsection{Idempotence on the image}
A scaffold is already ``core-only'' by definition: re-extracting its scaffold does nothing.

\begin{prop}[Idempotence on scaffold space]
\label{prop:idempotent}
Extend the graph extraction rule to all finite typed graphs, including the empty graph.
For every scaffold graph $s\in S$, this extension satisfies $\Scaf(s)=s$.
The expression is an identity on graph representatives; it does not assume $S\subseteq\Mol$, since a chemical rule set may reject a core-only graph and the empty graph is not a connected molecule.
\end{prop}

\begin{proof}
A scaffold graph is a 2-core: it has no vertices of degree $<2$ unless it is empty.
Thus applying the 2-core pruning procedure leaves it unchanged (up to isomorphism).
\end{proof}

\section{Decomposing a molecule into core and decoration\index{decoration}}
\label{sec:ch3-core-decoration}

Figure~\ref{fig:core-decoration} illustrates the scaffold (2-core) and decoration split for representative molecules.

\begin{figure}[htbp]
\centering
\includegraphics[width=\linewidth]{experiments/ch03_scaffolds/fig_core_decoration.png}
\caption{\textbf{Core-decoration decomposition.} Each row shows the original structure (left) and its toolkit-extracted scaffold (blue), with decorations highlighted (orange) and attachment points marked by asterisks (red). Equivalence to the literal graph 2-core must be checked for the chosen chemical conventions.}
\label{fig:core-decoration}
\end{figure}

\subsection{The decoration\index{decoration} forest}
The 2-core decomposition induces a canonical structural split.

Let $G$ be a representative graph of $m$, and let $G_2$ be its 2-core.
Consider the subgraph $T$ induced by $V\setminus V_2$ (the removed vertices).
For $k=2$, this remainder has a strong property.

\begin{prop}[The remainder is a forest attached to the core]
\label{prop:remainder-forest}
Let $G$ be a finite connected graph and let $G_2$ be its 2-core (possibly empty).
Then the induced subgraph on $V\setminus V_2$ is a forest (acyclic).
If $G_2$ is nonempty, each connected component of this forest is joined to $G_2$ by exactly one edge.
\end{prop}

\begin{proof}
Vertices removed in the iterative $k$-core procedure have degree $0$ or $1$ at the moment of removal.
If the removed subgraph contained a cycle, vertices on that cycle would have degree at least $2$ throughout and would not be removed.
Hence the removed subgraph is acyclic, i.e.\ a forest.
Connectedness gives at least one edge from each removed component to the core.
Two such edges, together with the unique tree path between their endpoints in the component, would form either a path between core vertices or a cycle returning to the same core vertex.
The union with the core would have minimum degree at least two, contradicting removal of the path vertices.
Thus exactly one attachment edge exists.
\end{proof}

\begin{rem}[Chemistry interpretation]
Under the stated 2-core map, ``decorations'' are precisely those substituents that are tree-like with respect to the cyclic backbone.
Pendant rings are \emph{not} decorations: they remain in the 2-core and become part of the scaffold. This is deliberate and matches the ring-system emphasis
of the original definition.
\end{rem}

\subsection{Attachment points and scaffold boundary\index{boundary}}
The attachment points where decoration trees connect to the scaffold will matter later, especially for refinement and energy factorization\index{energy factorization}.

\begin{defn}[Boundary (attachment) set]
\label{def:boundary-set}
Fix a molecule $m$ with scaffold $s=\Scaf(m)$ and choose representatives $G$ and $G_2$ as above.
Define the \emph{boundary}\index{boundary} (or attachment) set $\partial_s(m)\subseteq V_2$ to be the set of core vertices in $G_2$
that have at least one neighbor in $V\setminus V_2$ in the original graph $G$.
\end{defn}

\begin{rem}[Boundary depends on the molecule, not just the scaffold]
The core graph $G_2$ is determined by $s$, but the set of which core vertices actually carry decorations depends on $m$.
Later refinements often promote boundary information into the scaffold abstraction itself by marking attachment sites.
\end{rem}

\subsection{A constructive view: molecules as scaffold-plus-decorations}
The previous propositions justify a convenient mental model:

\begin{quote}
\emph{For a nonempty scaffold $s$, a molecule in $\Fib(s)$ is obtained by attaching typed trees to vertices of $s$, subject to the chemical validity rules. The empty fiber consists of the acyclic molecules.}
\end{quote}

This is not a free product decomposition in any algebraic sense yet, but it is a canonical decomposition of the underlying graph.
In Chapter~\ref{ch:energy} we will tie this to additive scoring functionals:
core terms are reused across the fiber, and fiber variance arises from decoration and boundary terms.

\section{Order on scaffolds and inherited structure}
\label{sec:ch3-scaf-order}

The scaffold map produces a quotient $S$, but we also want $S$ to carry structure.
The most important is an order relation by embedding, analogous to the molecule order in Chapter~\ref{ch:molecules}.

\begin{defn}[Substructure order on scaffolds]
\label{def:ch3-scaffold-order}
For scaffolds $s,t\in S$, define
\[
s \preceq t
\quad\Longleftrightarrow\quad
\exists \text{ representatives } G_s\in s,\ G_t\in t \text{ with } G_s \hookrightarrow G_t
\]
where $\hookrightarrow$ denotes typed subgraph embedding.
\end{defn}

\begin{prop}[$(S,\preceq)$ is a poset\index{poset}]
\label{prop:S-poset}
The relation $\preceq$ on $S$ is a partial order\index{partial order}.
\end{prop}

\begin{proof}
This is the same argument as for molecules in Chapter~\ref{ch:molecules}.
Reflexivity and transitivity follow from identity embeddings and composition.
Antisymmetry holds on isomorphism classes: if $s\preceq t$ and $t\preceq s$, there are embeddings both ways between finite graphs,
forcing isomorphism, hence $s=t$.
\end{proof}

\begin{prop}[Monotonicity of the typed 2-core]
\label{prop:two-core-monotone}
For the extraction in Definition~\ref{def:murcko-2core}, $G\hookrightarrow H$ implies $\Scaf(G)\hookrightarrow\Scaf(H)$.
\end{prop}

\begin{proof}
The image of the 2-core of $G$ under an injective embedding has minimum degree at least two within its image.
No vertex of that image can be deleted during degree-zero-or-one pruning of $H$, so the image lies in the 2-core of $H$.
Restrict the original embedding to this image; its atom and bond types remain preserved.
The empty-core case is immediate.
\end{proof}

\begin{rem}[Scope of monotonicity]
This proof assumes unchanged graph types and the literal 2-core operator.
A normalization or scaffold-extraction pipeline with different semantics requires its own check.
The powerset Galois connection in Chapter~\ref{ch:galois} exists for every deterministic map and does not require graph-order monotonicity.
\end{rem}

\section{Scaffold variants, refinements, and alternative abstractions}
\label{sec:ch3-variants}

Bemis--Murcko is a baseline, not a law of nature. The point of the framework is that \emph{any} deterministic abstraction induces
fibers, closures, and irreducible errors. Different abstractions trade compression against fidelity.

\subsection{Coarsenings: more compression, more collisions}
Common coarsenings include:
\begin{itemize}
\item \textbf{Generic scaffolds:} collapse atom types (e.g.\ all heteroatoms to a single type) or collapse bond orders to a coarser alphabet.
This merges scaffolds that differ only in typing.
\item \textbf{Ring systems only:} discard linkers and retain only ring components (useful when linkers are highly variable).
\item \textbf{Topological scaffolds:} discard chemical typing entirely and retain only graph topology of the 2-core.
\end{itemize}

Mathematically, a coarsening is another map $\Scaf':\Mol\to S'$ such that $\Scaf'$ factors through $\Scaf$:
there exists $\pi:S\to S'$ with $\Scaf'=\pi\circ\Scaf$. This is exactly the refinement relation used later in Chapter~\ref{ch:variance}.

\subsection{Refinements: more fidelity, smaller fibers}
Refinements include:
\begin{itemize}
\item \textbf{Attachment-site marked scaffolds:} promote boundary information to the scaffold by marking which vertices are substituted and with what bond types.
\item \textbf{Substitution pattern scaffolds:} keep ring and linker core but add ring-substitution patterns (e.g.\ which ring positions are substituted).
\item \textbf{Stereochemical scaffolds:} include stereocenters that lie on the core.
\end{itemize}

These refinements shrink fibers and often improve predictive consistency within fibers, at the cost of increasing the number of scaffold classes.
Normalization before extraction is a separate change of representation: it may merge some classes and split others, so it is not automatically a refinement.
The factorization condition through the finer map must be checked.

\begin{rem}[Refinement towers]
Later chapters assume a tower of abstractions
\[
\Scaf_0 \subset \Scaf_1 \subset \cdots \subset \Scaf_L
\]
in the sense that each $\Scaf_{\ell+1}$ refines $\Scaf_\ell$ (equivalently, the induced $\sigma$-algebras satisfy $\cS_\ell \subset \cS_{\ell+1}$).
Chapter~\ref{ch:variance} shows that such towers induce an orthogonal variance spectrum\index{variance spectrum}.
This chapter only needs the qualitative idea: you can tune the abstraction knob from coarse to fine.
\end{rem}

\begin{framed}
\noindent\textbf{Example: The theory is not tied to Bemis--Murcko.}

\medskip
\noindent
To illustrate that the algebraic framework survives alternative scaffold definitions, consider two non-Bemis--Murcko abstractions:

\medskip
\noindent\emph{1. Pharmacophore-aware scaffolds.}
Define $\Scaf_{\mathrm{pharm}}(m)$ as the Bemis--Murcko core \emph{plus} a labeling of key pharmacophoric features (H-bond donors/acceptors, charged centers, hydrophobic regions) at attachment sites.
Two molecules with identical ring systems but different pharmacophoric profiles (e.g., an amine vs.\ an ether at the same position) now map to \emph{different} scaffolds.

This is a refinement of Bemis--Murcko: there exists a projection $\pi:S_{\mathrm{pharm}}\to S_{\mathrm{BM}}$ such that $\Scaf_{\mathrm{BM}}=\pi\circ\Scaf_{\mathrm{pharm}}$.
All theorems about fibers, partitions, and variance decomposition (Chapters~\ref{ch:variance}--\ref{ch:galois}) apply unchanged---only the fiber sizes shrink.

\medskip
\noindent\emph{2. Acyclic backbone scaffolds.}
For peptides, nucleotides, or linear natural products, Bemis--Murcko is nearly useless (most structures are acyclic, mapping to the trivial scaffold $\bot$).
Instead, use a chemically declared backbone extraction rule, ignoring side chains.
If the rule chooses a longest path, specify isomorphism-invariant tie breaking or retain all tied paths so that the result is deterministic.

This is \emph{not} a refinement or coarsening of Bemis--Murcko---it is a different abstraction entirely.
Yet the algebraic structure persists:
\begin{itemize}
\item Fibers $\Fib(s)$ partition $\Mol$ (Proposition~\ref{prop:fibers-partition}).
\item The scaffold space $S_{\mathrm{backbone}}$ inherits a partial order from backbone embedding.
\item Variance decomposes into between-backbone and within-backbone terms (Chapter~\ref{ch:variance}).
\item Closure operators and Galois connections exist (Chapter~\ref{ch:galois}).
\end{itemize}

\medskip
\noindent\textbf{Takeaway:}
Bemis--Murcko is the running example because it is widely implemented and well-understood.
The fiber partition, powerset closure, and conditional-variance identities hold for any deterministic map on the declared domain.
Graph order, energy factorization, and algorithmic guarantees require their separately stated assumptions.
The choice of scaffold is a modeling decision; the universal identities do not choose it.
\end{framed}

\section{Empirical compression: how much does $\Scaf$ reduce chemical space?}
\label{sec:ch3-compression}

The entire framework becomes compelling when one observes that $\Scaf$ compresses large libraries dramatically:
millions to billions of molecules reduce to orders-of-magnitude fewer scaffolds, with a heavy-tailed fiber size distribution.

\begin{exper}[Compression ratio and fiber statistics]
\label{exp:compression}
Given a library $\mathcal{L}\subseteq\Mol$ (finite), compute:
\begin{enumerate}
\item number of molecules: $N = |\mathcal{L}|$,
\item number of distinct scaffolds: $|S_\mathcal{L}| := |\Scaf(\mathcal{L})|$,
\item \emph{scaffold coverage} fraction: $|S_\mathcal{L}|/N$,
\item fiber size multiset: $\{\, |\Fib_\mathcal{L}(s)| : s\in S_\mathcal{L}\,\}$,
\item summary statistics: median, tail quantiles, and a log--log plot of frequency vs.\ size.
\end{enumerate}
Here $\Fib_\mathcal{L}(s)=\{m\in\mathcal{L}:\Scaf(m)=s\}$.
Run this across multiple libraries (vendor catalogs, enumerated synthesis spaces, screening\index{screening} collections).
\end{exper}

\begin{rem}[What to expect]
In most realistic collections, the fiber size distribution is heavy-tailed: many scaffolds appear once or a few times,
and a small number of scaffolds dominate the library mass. This is the structural origin of the ``series'' phenomenon:
a few scaffold families account for a disproportionate fraction of candidates.

Typical numerical ranges (order-of-magnitude estimates from the literature):
\begin{itemize}
\item \emph{Vendor catalogs} ($10^6$--$10^7$ molecules): scaffold coverage $|S_\mathcal{L}|/N \approx 5$--$15\%$, i.e.\ compression ratios of $7$--$20\times$.
\item \emph{Enumerated synthesis spaces} ($10^9$--$10^{12}$ molecules): scaffold coverage often drops to $0.1$--$1\%$, reflecting combinatorial decoration of a fixed set of cores.
\item \emph{Screening\index{screening} hits}: the top-scoring molecules for a given target frequently concentrate in $\lesssim 10$ scaffold families, even when drawn from libraries of millions.
\end{itemize}
The exact numbers depend on the scaffold definition and normalization conventions; treat these as ballpark anchors for Experiment~\ref{exp:compression}.
\end{rem}

\begin{table}[htbp]
\centering
\small
\begin{tabular}{@{}lrrrrrrrr@{}}
\toprule
Library & $N_{\mathrm{valid}}$ & $|S|$ & Coverage & Comp. & Median & $p_{99}$ & Gini & $\hat{\alpha}$ \\
\midrule
BDB & 1.80M & 267.6K & 14.9\% & $6.7\times$ & 2 & 74 & 0.711 & 2.39 \\
G13 & 919.1M & 52.47M & 5.7\% & $17.5\times$ & 1 & 137 & 0.926 & 2.07 \\
G17 & 48.83M & 11.42M & 23.4\% & $4.3\times$ & 1 & 29 & 0.744 & 2.02 \\
MOSES & 1.94M & 449.0K & 23.2\% & $4.3\times$ & 1 & 42 & 0.715 & 2.07 \\
PubChem & 94.05M & 10.74M & 11.4\% & $8.8\times$ & 1 & 71 & 0.834 & 2.11 \\
\bottomrule
\end{tabular}
\caption{Observed scaffold compression across processed libraries using the implemented RDKit extraction pipeline. Coverage is $|S|/N_{\mathrm{valid}}$; Comp. is $N_{\mathrm{valid}}/|S|$; $p_{99}$ is the 99th percentile of fiber size; $\hat{\alpha}$ is the fitted power-law tail exponent for fiber sizes. These are measurements of that operational map.}
\label{tab:compression-observed}
\end{table}

\begin{rem}[What the current data show]
The processed libraries in Table~\ref{tab:compression-observed} sharpen the qualitative picture. Across five collections, fitted tail exponents lie in the narrow range $2.02$--$2.39$, and all libraries have high fiber-size inequality (Gini $0.71$--$0.93$). This is strong evidence for a robust heavy-tail phenomenon, though not necessarily for a pure power law: in several runs the lognormal alternative is not decisively rejected, while the exponential tail is rejected.

The important point for the theory is distributional, not terminological. Most scaffolds are small fibers---the median fiber is one molecule in four of five libraries---but a few fibers carry a macroscopic fraction of the library. The largest single scaffold fiber accounts for about $6.8\%$ of valid PubChem molecules, $4.6\%$ of MOSES, and $2.1\%$ of G13. Scaffold abstraction is therefore a two-regime object: sparse at the level of typical scaffolds, but massively reusable in the tail. Algorithms that exploit only the typical fiber miss the amortization opportunity; algorithms that exploit only the tail miss rare scaffold discovery.

One operational caveat is also visible. Idempotence checks passed in all runs except PubChem, where $3$ of $1000$ scaffold-image checks failed under the implemented normalization pipeline. This should be treated as a protocol quality-control signal rather than as a mathematical contradiction: the theorem applies to the fixed map $\Scaf$, while software pipelines may change aromaticity, tautomer, or representation conventions between molecule and scaffold images.
\end{rem}

\subsection{A practical pipeline}
A reproducible pipeline needs deterministic conventions: aromaticity model, kekulization, tautomer standardization, salt stripping, and stereochemistry handling.
These choices affect $\Scaf$ and thus affect $S$ and the fiber distribution. The safest stance is:

\begin{quote}
\emph{Treat the scaffold map as part of the experimental protocol. Fix it, document it, and measure its induced compression and irreducible error\index{irreducible error}.}
\end{quote}

\begin{algorithm}
\caption{Compute scaffold classes and fiber sizes on a finite library}
\label{alg:scaffold-stats}
\begin{algorithmic}[1]
\Require Library $\mathcal{L}=\{m_i\}_{i=1}^N$ (molecules)
\Ensure Scaffold set $S_\mathcal{L}$, fiber sizes $\{|\Fib_\mathcal{L}(s)|\}$, and summary statistics
\State Initialize empty map \texttt{count} from scaffold to integer
\For{$i=1$ to $N$}
  \State $m \gets \mathrm{Normalize}(m_i)$ \Comment{tautomer/protonation/salt conventions}
  \State $s \gets \Scaf(m)$
  \State \texttt{count[$s$]} $\gets$ \texttt{count[$s$]} $+ 1$
\EndFor
\State $S_\mathcal{L} \gets \{ s : \texttt{count[$s$]} > 0 \}$
\State Return $S_\mathcal{L}$ and $\{\texttt{count[$s$]}: s\in S_\mathcal{L}\}$
\end{algorithmic}
\end{algorithm}

\subsection{Why heavy tails matter}
If the fiber size distribution were uniform, scaffolds would be a weak abstraction: every scaffold would represent only a few molecules.
Heavy tails produce the opposite: a small set of scaffolds carries enormous combinatorial mass. This implies:

\begin{enumerate}
\item \textbf{Optimization concentrates.} If high-scoring molecules concentrate in a small number of scaffold families,
search can be organized at the scaffold level.
\item \textbf{Reuse is real.} Computation on a scaffold can be amortized across many molecules in its fiber (Chapter~\ref{ch:energy}).
\item \textbf{Lower bounds appear.} Any method that is ``scaffold-sound'' inherits unavoidable overhead measured by closure (Chapter~\ref{ch:galois} and Chapter~\ref{ch:complexity}).
\end{enumerate}

\section{Pitfalls: what can destabilize a scaffold map?}
\label{sec:ch3-pitfalls}

A scaffold map is only as stable as its normalization conventions. Common destabilizers include:

\subsection{Aromaticity and kekulization}
Different aromaticity models can change atom and bond typing. If connectivity is unchanged, the untyped 2-core vertex set is unchanged, but the typed scaffold may differ. Fixing an aromaticity model is essential.

\subsection{Tautomerism and protonation}
Tautomers can change perceived ring membership and bond orders. Protonation states can change typing.
If the goal is to partition by ``core topology,'' a tautomer-canonicalization step is often necessary.

\subsection{Stereochemistry}
Bemis--Murcko scaffolds are often defined without stereochemistry, but for many targets, core stereochemistry matters.
Decide early whether stereochemistry is part of $A$ (atom types) and whether it is retained in $S$.

\subsection{Macrocycles and fused systems}
Macrocycles and fused ring systems interact with core extraction in subtle ways:
side chains that form part of the 2-core by degree constraints might be included even if chemically ``peripheral.''
This is not wrong, but it must be understood.

\begin{rem}[A practical rule]
The book will treat $\Scaf$ as fixed. In practice, you should treat $\Scaf$ as a tunable abstraction:
evaluate it by measuring (i) compression, (ii) class cohesion (variance within fibers), and (iii) downstream screening performance.
Chapters~\ref{ch:variance} and~\ref{ch:complexity} provide formal tools to do exactly that.
\end{rem}

\section*{Preview: why this abstraction is the right input to algebra}
The key output of this chapter is the object
\[
\Scaf:\Mol\to S
\]
and its induced partition $\{\Fib(s)\}_{s\in S}$.
From this point onward, the book develops structure \emph{on} $S$ and \emph{between} $\Mol$ and $S$:

\begin{itemize}
\item Chapter~\ref{ch:poset}: the scaffold poset $(S,\preceq)$, upper/lower cones, and the asymmetry of going down vs.\ going up.
\item Chapter~\ref{ch:lattice}\index{lattice}: lifting from points to regions via the downset lattice $\cD(S)$.
\item Chapter~\ref{ch:galois}: abstraction and concretization between $\mathcal{P}(\Mol)$ and $\mathcal{P}(S)$, and the closure operator\index{closure operator} $c$.
\item Chapter~\ref{ch:variance}: conditional expectation\index{conditional expectation} on the scaffold $\sigma$-algebra and the variance spectrum\index{variance spectrum} induced by refinement towers.
\end{itemize}

\section*{Summary}
A scaffold map is a deterministic abstraction that collapses chemical space into scaffold classes:

\begin{enumerate}
\item We fixed a typed 2-core operator on heavy-atom graphs and distinguished it from toolkit-specific scaffold extraction.
\item The scaffold space $S$ is the image of $\Scaf$; fibers $\Fib(s)$ partition $\Mol$ into scaffold classes.
\item Every graph with nonempty core decomposes into that core plus pendant trees attached by one edge each; acyclic graphs form the empty-core fiber.
\item The scaffold space inherits a natural partial order from substructure embedding, enabling order-theoretic structure developed next.
\end{enumerate}

In the next part of the book we stop treating scaffolds as a heuristic and treat them as a mathematical domain:
a graph poset, a downset completion, and an abstraction in the sense of Galois connections, with separate assumptions for each.

\section*{Exercises}

\begin{exer}[Chemistry: cores, decorations, and attachment points]
\label{exer:ch3-chem}
Work in the heavy-atom graph model with implicit hydrogens. Consider the molecules
\[
\begin{aligned}
m_1&=\text{toluene}, &m_2&=\text{ethylbenzene},\\
m_3&=\text{biphenyl}, &m_4&=\text{2-phenylethanol}.
\end{aligned}
\]
Use the literal graph 2-core abstraction from Chapter~\ref{ch:scaffolds}; keep toolkit normalization separate.
\begin{enumerate}
\item For each $m_i$, determine (up to isomorphism) its scaffold $\Scaf(m_i)$.
\item For each $m_i$, identify which heavy atoms belong to the 2-core (scaffold) and which belong to the decoration forest.
\item For each $m_i$, compute the attachment (boundary) set $\partial_{\Scaf(m_i)}(m_i)$ (the scaffold vertices incident to at least one decoration atom).
\item Explain, in chemical terms, why (i) $m_1$ and $m_2$ fall in the same scaffold class, (ii) $m_3$ has a different scaffold, and (iii) $m_4$ shares scaffold with $m_1$ despite having a polar terminal group.
\end{enumerate}
\end{exer}

\begin{exer}[Mathematics: the scaffold map as a quotient and a canonical decomposition]
\label{exer:ch3-math}
Let $\Mol$ be the set of molecule isomorphism classes from Chapter~2, and let $\Scaf:\Mol\to S$ be the scaffold map from Chapter~3.
\begin{enumerate}
\item Prove that the relation $m\sim m'$ defined by $\Scaf(m)=\Scaf(m')$ is an equivalence relation on $\Mol$ and that its equivalence classes are exactly the fibers $\Fib(s)$.
\item Prove that $\Scaf$ is surjective onto $S$. Give a library on which it is not injective and a library on which it is injective.
\item Let $m\in\Mol$ and choose a representative graph $G$. Let $G_2$ be its 2-core and let $T$ be the induced subgraph on the removed vertices.
Prove that $T$ is a forest and that each connected component of $T$ attaches to $G_2$ at exactly one vertex (when $G_2\neq \emptyset$).
\item Prove idempotence on the image: for any scaffold $s\in S$ (viewed as a core-only graph), $\Scaf(s)=s$.
\end{enumerate}
\end{exer}

\begin{exer}[Machine learning: pushforward measures and scaffold-level statistics]
\label{exer:ch3-ml}
Let $(\Mol,\mu)$ be a probability space (finite library with a probability mass function) and let $\Scaf:\Mol\to S$.
Define the pushforward (scaffold) measure $\nu$ on $S$ by $\nu(B)=\mu(\Scaf^{-1}(B))$ for $B\subseteq S$.
Let $F:\Mol\to\R$ be a square-integrable score.
\begin{enumerate}
\item Show that $\nu(s)=\mu(\Fib(s))$ for singleton scaffolds $s\in S$, and that $\sum_{s\in S}\nu(s)=1$ in the finite case.
\item Prove the law of total expectation in scaffold form:
\[
\E_\mu[F] = \sum_{s\in S} \nu(s)\, \E_\mu[F\mid \Scaf=s].
\]
\item Prove the law of total variance in scaffold form:
\[
\Var_\mu(F) = \Var_\nu\!\left(\E[F\mid \Scaf]\right) + \E_\nu\!\left[\Var(F\mid \Scaf)\right].
\]
\item Interpret $\E_\nu[\Var(F\mid \Scaf)]$ as ``within-fiber noise'' and explain why it is the natural quantity later called the \emph{irreducible scaffold error} (Chapter~7).
\end{enumerate}
\end{exer}

\begin{exer}[Computation: implement $\Scaf$ and measure compression]
\label{exer:ch3-comp}
Using a cheminformatics toolkit (e.g.\ RDKit) and a finite library $\mathcal{L}\subseteq\Mol$ (say $10^4$--$10^6$ molecules if feasible):
\begin{enumerate}
\item Implement Bemis--Murcko scaffolds and compute the set of scaffolds $S_{\mathcal{L}}=\Scaf(\mathcal{L})$.
\item Compute the compression ratio $|S_{\mathcal{L}}|/|\mathcal{L}|$ and the empirical fiber size distribution
\[
\bigl\{\,|\Fib_{\mathcal{L}}(s)| : s\in S_{\mathcal{L}}\,\bigr\},
\qquad
\Fib_{\mathcal{L}}(s)=\{m\in\mathcal{L}:\Scaf(m)=s\}.
\]
Plot the complementary CDF (log--log) of fiber sizes.
\item Compute the attachment count distribution: for each $m\in\mathcal{L}$, compute $|\partial_{\Scaf(m)}(m)|$ and plot its histogram.
\item Verify empirically that scaffold extraction is idempotent on the image: for random $m\in\mathcal{L}$, check that $\Scaf(\Scaf(m))=\Scaf(m)$.
\end{enumerate}
\end{exer}
% END SOURCE: chapters/ch03_scaffolds.tex


%--------Part II: The Algebraic Structure--------
\part{Order and Information Loss}
\label{part:algebra}

\chapter{The Scaffold Poset}
\label{ch:poset}
% BEGIN SOURCE: chapters/ch04_poset.tex
% Chapter 4: graph order for a declared scaffold domain.
\section{Scaffold order and its scope}

A scaffold graph supports a substructure order once atom types, bond types, and the embedding convention are fixed.
This order records structural containment. It does not rank binding affinity, synthetic accessibility, or the cost of discovering a useful molecule.

\section{The scaffold poset}
\begin{defn}[Scaffold order]
\label{def:ch4-scaffold-order}
For $s,t\in S$, write $s\preceq t$ when a typed, not necessarily induced, subgraph embedding exists between representative scaffold graphs.
This definition agrees with Definition~\ref{def:ch3-scaffold-order}.
\end{defn}

\begin{prop}[$(S,\preceq)$ is a poset]
\label{prop:scaffold-poset}
Typed subgraph embedding induces a partial order on the isomorphism classes in $S$.
\end{prop}
\begin{proof}
Identity and composition give reflexivity and transitivity.
Embeddings in both directions force equality of vertex counts and edge counts, so either embedding is an isomorphism.
\end{proof}

\begin{rem}[When molecule witnesses give the same order]
For the literal 2-core map, molecule witnesses
$m_s\in\Fib(s)$, $m_t\in\Fib(t)$ with $m_s\hookrightarrow m_t$
imply $s\preceq t$ by Proposition~\ref{prop:two-core-monotone}.
The converse holds on a graph domain closed under taking 2-cores, including the empty graph: take the core graphs themselves as witnesses.
A restricted chemical library need not contain these witnesses, so we define the scaffold order directly rather than assert that equivalence for every library.
Both $\Mol$ and $S$ already consist of isomorphism classes and are posets; the preorder is on their labeled graph representatives.
\end{rem}

\section{Upper and lower cones}
\begin{defn}[Lower cone]
\label{def:lower-cone}
The lower cone is $\downarrow s=\{t\in S:t\preceq s\}$.
It contains exactly the scaffold types in the declared domain that embed in $s$.
\end{defn}
\begin{defn}[Upper cone]
\label{def:upper-cone}
The upper cone is $\uparrow s=\{t\in S:s\preceq t\}$.
It contains the admitted scaffold extensions of $s$.
\end{defn}

\begin{exmp}[Cones depend on typing and domain]
\label{ex:benzimidazole-cones}
A sub-scaffold of benzimidazole must preserve the chosen atom and bond types.
The presence of a five-membered ring does not permit replacing its nitrogen atoms by carbon or moving them to adjacent positions.
A pictured chemical family is therefore not automatically a cone or a Hasse diagram.
In a finite library, upper and lower cones both depend on which scaffold types the library contains.
\end{exmp}

\begin{prop}[Cone properties]
\label{prop:cone-properties}
Every lower cone is finite: its members embed in one finite typed graph.
Also $s\in\downarrow s\cap\uparrow s$ and
$t\in\downarrow s$ if and only if $s\in\uparrow t$.
If $S$ is finite, all upper cones are finite.
An upper cone is infinite only if the chosen domain admits extensions of unbounded size.
\end{prop}

\section{Enumeration and the size of the search domain}

The finite host graph bounds downward enumeration.
Upward enumeration requires a domain or a size bound before it becomes a finite problem.
This distinction does not make downward enumeration polynomial in the number of atoms.

\subsection{Going down: sub-scaffold enumeration}
\begin{prop}[A finite exhaustive lower-cone procedure]
\label{prop:lower-cone-enum}
Let a representative of $s$ have $n$ vertices and $e$ edges.
Enumerating every vertex subset and every subset of its inherited edges examines at most $2^{n+e}$ labeled subgraphs.
Keep the empty scaffold when admitted, and the connected subgraphs of minimum degree at least two whose types belong to $S$.
Deduplicating these graphs by typed isomorphism gives $\downarrow s$ for the 2-core scaffold domain.
\end{prop}
\begin{proof}
Every embedded scaffold is represented by some vertex subset and some edge subset of the host.
Conversely, every retained subgraph embeds in the host and is a fixed point of the 2-core operator.
Domain filtering and isomorphism deduplication then give exactly the declared lower cone.
\end{proof}

\begin{algorithm}
\caption{Exhaustive lower-cone enumeration}
\label{alg:lower-cone}
\begin{algorithmic}[1]
\Require Scaffold graph $G=(V,E)$ and a declared scaffold domain $S$
\Ensure $\downarrow s$ in that domain
\State Initialize \texttt{result} with the empty scaffold if it belongs to $S$
\For{each nonempty vertex subset $W\subseteq V$}
  \For{each edge subset $D\subseteq E[W]$}
    \State $H\gets (W,D)$ with inherited types
    \If{$H$ is connected, has minimum degree at least two, and $[H]\in S$}
      \State Add $[H]$ to \texttt{result}, deduplicating by typed isomorphism
    \EndIf
  \EndFor
\EndFor
\State \Return \texttt{result}
\end{algorithmic}
\end{algorithm}

\begin{rem}[Why vertex-induced subgraphs are insufficient]
A four-cycle embeds in a four-cycle with a diagonal, but it is not an induced subgraph of that host.
Both are 2-cores. An induced-subgraph-only enumeration misses the four-cycle.
The bound above counts labeled candidates; it is neither a polynomial atom-count bound nor an output-sensitive bound, and it excludes the cost of isomorphism and chemical-validity checks.
For a finite observed $S$, an alternative is to test each admitted scaffold for embedding in $s$.
\end{rem}

\subsection{Going up: admissible extensions}
\begin{prop}[An explicit condition for an infinite upper cone]
\label{prop:upper-cone-explosion}
If the domain contains a sequence $s\preceq t_1\preceq t_2\preceq\cdots$
with strictly increasing vertex counts, then $\uparrow s$ is infinite.
A finite type alphabet and a vertex bound make the set of admitted extensions finite.
\end{prop}
\begin{proof}
The sequence contains pairwise non-isomorphic graphs. With at most $n$ vertices and finite type alphabets there are only finitely many typed graphs.
\end{proof}

\begin{exmp}[A bounded extension problem]
\label{ex:benzene-explosion}
Starting from benzene, an extension search must specify the atom palette, allowed bonds, valence rules, maximum size, and permitted construction operations.
Fused rings and linked rings give different candidates.
Their number is a property of that declared search domain; no universal numerical count or growth rate follows from the order alone.
\end{exmp}

\begin{rem}[Implications for screening]
Lower-cone enumeration organizes simplifications, and upper-cone construction organizes elaborations.
Neither containment nor direction of movement guarantees retained activity.
A useful search policy must also model the response and pay for its evaluations.
\end{rem}

\section{Size functions and level sets}
\begin{defn}[Strict order size]
\label{def:grading}
A strict order size is a map $\rho:S\to\mathbb{Z}_{\geq0}$ such that
$s\prec t$ implies $\rho(s)<\rho(t)$.
This is weaker than the usual definition of a graded poset, which requires a rank increase of exactly one along every cover relation.
\end{defn}

\begin{prop}[Cycle rank for connected 2-cores]
\label{prop:ring-grading}
On the nonempty connected 2-core scaffold domain, the cycle rank
$\beta(s)=|E_s|-|V_s|+1$ strictly increases under a proper embedding.
Set $\beta(\bot)=0$ when the empty scaffold is admitted.
\end{prop}
\begin{proof}
Identify $s$ with its image in $t$.
Contract the connected subgraph $s$ to a root, retaining loops and parallel edges.
The resulting connected multigraph has cycle rank $\beta(t)-\beta(s)\geq0$.
If this difference were zero, it would be a tree.
If any vertices outside $s$ remained, that tree would have a non-root leaf, contradicting minimum degree at least two in $t$.
If no outside vertices remained, zero cycle rank would permit no added edges.
Either case rules out a proper inclusion with equal cycle rank.
\end{proof}

\begin{rem}[Atom count and other descriptors]
Atom count is nondecreasing but not strictly increasing: adding a diagonal to a four-cycle preserves all vertices.
The size $|V_s|+|E_s|$ is strictly increasing for all finite typed subgraph embeddings between distinct isomorphism classes.
Ring-system counts and physicochemical descriptors require separate checks; they are not automatically monotone.
Cycle rank here is the graph invariant $|E|-|V|+1$, not an unspecified ring-perception convention.
\end{rem}

\begin{defn}[Level sets]
\label{def:level-sets}
For a declared descriptor $\rho$, define $S_k=\{s\in S:\rho(s)=k\}$.
These are a partition into descriptor classes; arbitrary descriptors need not respect the order.
\end{defn}
\begin{exmp}[Cycle-rank levels]
\label{ex:ring-levels}
Nonempty connected 2-cores of cycle rank one are single cycles.
Their sizes can be arbitrarily large, so a finite atom alphabet alone does not make a level finite.
Adding a vertex bound does.
\end{exmp}

\section{A non-lattice example within scaffold graphs}

\begin{exmp}[Cycles with no least common scaffold]
\label{ex:scaffold-lattice-failure}
Take uniform carbon/single-bond typing with implicit hydrogens.
Let $x=C_3$ and $y=C_4$.
One common upper bound is the diamond $D$, a four-cycle with one diagonal.
Another is the six-vertex figure-eight graph $Q$ obtained by making a triangle and a square share exactly one vertex.
All four graphs are connected 2-cores with maximum degree at most four.

Any common upper bound of $x,y$ that embeds in $D$ must use four vertices and at least five edges:
four edges give only the square, without a triangle.
It must therefore be isomorphic to $D$.
But $D$ contains two distinct triangles, whereas $Q$ contains only one, so $D$ does not embed in $Q$.
No least upper bound exists in any scaffold domain containing these four graphs.
The same argument shows that $D,Q$ have no meet, because a meet would have to contain both $C_3$ and $C_4$.
\end{exmp}

\section{Fragment-based drug design as cone intersection}

Fragment-based drug discovery (FBDD) has a natural interpretation in the scaffold poset.

\begin{defn}[Fragment hits as lower-cone anchors]
\label{def:fragment-anchors}
A fragment hit $f$ with scaffold $s_f$ defines the upper cone $\uparrow s_f$ of all scaffolds that could elaborate $f$.
Multiple fragment hits $f_1, \ldots, f_k$ with scaffolds $s_1, \ldots, s_k$ constrain the search to:
\[
\bigcap_{i=1}^k \uparrow s_i = \{t \in S : s_i \preceq t \text{ for all } i\}.
\]
This is the set of scaffolds containing \emph{all} fragment cores.
\end{defn}

\begin{rem}[Structural and physical compatibility]
\label{prop:empty-intersection}
The intersection is empty exactly when no scaffold in the declared domain contains both patterns.
A size limit or a chemical-validity rule may cause this.
Conflicting binding poses are an additional physical constraint; they do not by themselves make the graph-theoretic intersection empty.
\end{rem}

\begin{rem}[Fragment linking as join]
When fragments $f_1$ and $f_2$ can be linked (their scaffolds $s_1, s_2$ have a common superstructure), the resulting linked scaffold is an upper bound.
The \emph{join} $s_1 \vee s_2$, if it exists, is the smallest such upper bound.
Example~\ref{ex:scaffold-lattice-failure} shows that joins need not exist even among connected 2-core scaffolds.
Fragment linking often requires explicit construction when no canonical join exists.
\end{rem}

\begin{exper}[Fragment elaboration paths]
\label{exp:fragment-paths}
Given a fragment hit with scaffold $s$:
\begin{enumerate}
\item Enumerate $\uparrow s \cap S_{\leq n}$ for $n = |s| + k$ (scaffolds with up to $k$ additional atoms).
\item For each super-scaffold $t$, generate molecules in $\Fib(t)$ consistent with the fragment binding mode.
\item Score/dock the elaborated molecules.
\item Analyze which \emph{paths} through the poset (sequences $s = t_0 \prec t_1 \prec \cdots \prec t_m$) yield the best evaluated scores.
\end{enumerate}
The poset structure constrains and organizes the combinatorial search.
\end{exper}

\begin{figure}[htbp]
\centering
\includegraphics[width=0.95\textwidth]{experiments/ch04_poset/output/fig_4_3_elaboration_paths.pdf}
\caption{Stored FKBP12 docking example. Nodes represent scaffold classes and their best evaluated molecular docking scores, not the scores of bare cores. The displayed best terminal candidate scores $-5.99$ kcal/mol along the benzene--phenyl-imidazole--phenyl-imidazole-pyrrole branch. This selected finite search does not establish activity monotonicity or a search-policy advantage.}
\label{fig:elaboration-paths}
\end{figure}

\begin{rem}[What the stored docking results establish]
The stored summary in \path{experiments/ch04_poset/output/experiment_4_1_summary.json}
records 20 scaffold classes and 382 evaluated molecules, with \texttt{dry\_run=false}.
Its best terminal score is $-5.991$ kcal/mol.
The best score among the eight evaluated benzene-scaffold molecules is $-3.821$ kcal/mol; the associated molecule is isopropylbenzene.
The CSV records bare benzene separately at $-2.979$ kcal/mol.
Thus a scaffold-class maximum must not be reported as the score of its undecorated core.

Among the ten evaluated level-one classes, biphenyl has the best score, $-5.661$ kcal/mol.
Phenyl-imidazole ranks ninth at $-5.028$ kcal/mol, but its sampled descendants include the best terminal candidate.
This illustrates a distinction between immediate and observed downstream score.
Because the level-two panel contains selected extensions of that branch rather than matched exploration of every branch, it does not estimate the benefit of a lookahead policy over a greedy policy.

The snapshot has unequal numbers of candidates per class and contains no reported affinity measurements, runtime, or size-matched control comparison.
Its CSV also lacks the ligand-efficiency columns now produced by the experiment code.
We therefore retain it as a docking example, without inferring an affinity multiplier, size-adjusted improvement, runtime advantage, or prospective discovery gain.
A reproducible rerun must freeze the preparation and docking protocol and compare policies at equal evaluation cost.
\end{rem}

\section{Computational aspects}

We briefly address implementation considerations for working with the scaffold poset.

\subsection{Representing the poset}

\begin{itemize}
\item \textbf{Explicit Hasse diagram:} Store edges $(s, t)$ where $s \prec t$ with no intermediate scaffold.
Suitable for small scaffold sets (up to $\sim 10^4$).
\item \textbf{Implicit representation:} Store scaffolds as canonical SMILES; compute $s \preceq t$ on demand via substructure matching.
Suitable for large or infinite scaffold spaces.
\item \textbf{Hybrid:} Precompute the Hasse diagram for common scaffolds; use on-demand queries for rare ones.
\end{itemize}

\subsection{Algorithms}

\begin{itemize}
\item \textbf{Substructure query} ($s \preceq t$?): Use VF2 or similar subgraph isomorphism algorithm. Report its measured cost on the declared library; small chemical patterns can help without removing the general worst-case difficulty.
\item \textbf{Lower cone enumeration:} Algorithm~\ref{alg:lower-cone} gives a complete finite procedure; benchmark its cost before using it on large cores.
\item \textbf{Upper cone sampling:} When exhaustive enumeration exceeds the budget, consider:
\begin{itemize}
\item Generative models conditioned on containing $s$
\item Explicit enumeration libraries (e.g., REAL Space) filtered by substructure
\item Retrosynthetic expansion from $s$
\end{itemize}
\end{itemize}

\subsection{Database considerations}

Chemical databases (ChEMBL, PubChem, ZINC) can be indexed by scaffold to support:
\begin{itemize}
\item Fast $\uparrow s$ queries: ``Find all molecules whose scaffold contains $s$''
\item Fast $\downarrow s$ queries: ``Find all molecules whose scaffold is contained in $s$''
\end{itemize}
An exact scaffold-key index answers equality queries. Cone queries additionally need containment matching or a precomputed order index; their latency depends on that implementation and the library.

\section*{Summary}

Scaffold containment gives a partial order and finite lower cones.
Upper cones are finite in a finite library and can be infinite in an unrestricted construction domain.
Non-induced embeddings require checking both vertex and edge subsets; the order does not supply a polynomial-time enumeration guarantee.
Cycle rank strictly increases for proper embeddings between connected 2-cores, whereas atom count need not.
These facts organize candidate construction but leave biological response and evaluation cost to separate models.

\section*{Exercises}

\begin{exer}[Cone cardinality]
\label{exer:cone-cardinality}
Use the domain of all connected, uniformly typed 2-core graphs, together with the empty scaffold.
For a four-cycle with one diagonal:
\begin{enumerate}
\item Enumerate the lower cone up to isomorphism.
\item Identify a lower-cone element missed by vertex-induced-subgraph enumeration.
\item Explain why a vertex bound and finite type alphabet make its upper-cone search finite.
\end{enumerate}
\end{exer}

\begin{exer}[Size-function comparison]
\label{exer:grading-comparison}
\begin{enumerate}
\item Give $s\prec t$ with strictly increasing cycle rank and unchanged atom count.
\item Explain why equal cycle rank is impossible for a proper embedding of nonempty connected 2-cores.
\item Show that atom-count level sets and cycle-rank level sets need not refine one another, using pairs that need not be comparable.
\end{enumerate}
\end{exer}

\begin{exer}[Fragment merging]
\label{exer:fragment-merging}
Use the four scaffold graphs in Example~\ref{ex:scaffold-lattice-failure}.
\begin{enumerate}
\item Exhibit embeddings of $C_3$ and $C_4$ into both $D$ and $Q$.
\item Prove that $D$ and $Q$ are minimal common upper bounds and are incomparable.
\item Explain why a common upper bound alone says nothing about synthesis or binding compatibility.
\end{enumerate}
\end{exer}

\begin{exer}[Computational complexity]
\label{exer:complexity}
\begin{enumerate}
\item Explain why equal vertex counts do not make a one-way non-induced embedding an isomorphism. What additional count suffices?
\item For a pattern with $k$ vertices and a host with $n$ vertices, give an exhaustive algorithm over injective vertex maps and its running-time bound.
\item Explain why polynomial time for fixed $k$ does not imply polynomial time when $k$ is part of the input. Why does a small pattern treewidth alone not remove the difficulty of injective path matching?
\end{enumerate}
\end{exer}
% END SOURCE: chapters/ch04_poset.tex


\chapter{The Downset Completion of Scaffold Space}
\label{ch:lattice}
% BEGIN SOURCE: chapters/ch05_lattice.tex
% ======================================================================
% Chapter 5 content file: chapters/ch05_lattice.tex
% (Do NOT include \chapter{...} here; main.tex already does that.)
% ======================================================================

\section{Why lift from points to regions?}

Chapter~\ref{ch:poset} defined scaffold containment and gave an explicit pair with no meet or join.
A set of common sub-scaffolds is still well-defined even when none is greatest.
We therefore study downward-closed sets, whose unions and intersections provide exact operations on structural search regions.

This is the \emph{downset completion} of a poset.
Some authors call downsets order ideals; the term \emph{ideal completion} also has a narrower meaning that requires directedness.
We keep that distinction explicit below.

\begin{framed}
\noindent\textbf{When a design region is a downset.}

The sub-scaffolds of a lead form a downset, as does the union of such regions for several leads.
A monotone upper bound on graph size also defines a downset.
A requirement to contain a named core instead defines an upset, and a restriction to exactly two cycles need be neither.
Downward closure is a chosen structural constraint, not a property of every chemical campaign or of its active molecules.
\end{framed}

\section{Downsets: definition and first examples}
\label{sec:downsets-def}

We begin with the central definition of this chapter.

\begin{defn}[Downset (order ideal)]
\label{def:downset}
Let $(P, \preceq)$ be a poset\index{poset}.
A subset $I \subseteq P$ is a \emph{downset}\index{downset} (also called an \emph{order ideal} or \emph{down-closed set}) if it satisfies the following closure property:
\[
t \preceq s \text{ and } s \in I \implies t \in I.
\]
Equivalently, $I$ is a downset if $I = {\downarrow} I$, where ${\downarrow} I := \{t \in P : \exists s \in I,\, t \preceq s\}$ denotes the \emph{downward closure} of $I$.
\end{defn}

\begin{notn}
We write $\cD(P)$ for the set of all downsets of the poset $P$.
When the poset is the scaffold space $S$, we write $\cD(S)$.
\end{notn}

\begin{exmp}[Elementary downsets]
\label{ex:elementary-downsets}
Let $S$ be a scaffold poset.
The following are downsets:
\begin{enumerate}
\item The \emph{empty set} $\emptyset$ is vacuously a downset: there is no $s \in \emptyset$ for which the closure condition could fail.
\item The \emph{entire space} $S$ is a downset: if $s \in S$, then certainly $t \in S$ for any $t \preceq s$.
\item For any scaffold $s$, the \emph{lower cone}\index{lower cone} ${\downarrow} s = \{t \in S : t \preceq s\}$ is a downset.
This is immediate from transitivity: if $u \preceq t$ and $t \preceq s$, then $u \preceq s$, so $u \in {\downarrow} s$.
\item If $I$ and $J$ are downsets, so is $I \cap J$: if $s \in I \cap J$ and $t \preceq s$, then $t \in I$ (since $I$ is a downset) and $t \in J$ (since $J$ is a downset), so $t \in I \cap J$.
\item If $I$ and $J$ are downsets, so is $I \cup J$: if $s \in I \cup J$ and $t \preceq s$, then $s$ lies in at least one of $I, J$; say $s \in I$.
Since $I$ is a downset, $t \in I \subseteq I \cup J$.
\end{enumerate}
\end{exmp}

\begin{exmp}[Downsets in scaffold space: pharmaceutical interpretation]
\label{ex:pharma-downsets}
Consider a drug discovery campaign with the following design constraints:
\begin{enumerate}
\item Let $I = \{s \in S : \beta(s) \leq 2\}$, where $\beta$ is the cycle rank on the connected 2-core domain.
This is a downset by Proposition~\ref{prop:ring-grading}, including the empty scaffold with cycle rank zero when admitted.
This represents the hypothesis ``focus on simple scaffolds.''

\item Let $J = \{s \in S : s \preceq t \text{ for some hit scaffold } t \in H\}$ where $H$ is a finite set of HTS hit scaffolds.
This is $J = \bigcup_{t \in H} {\downarrow} t$, the union of the lower cones of hit scaffolds---a downset representing ``all sub-scaffolds appearing in our hits.''

\item The intersection $I \cap J$ consists of scaffolds with at most two rings that also appear as sub-scaffolds of our hits.
This combines both constraints into a refined design hypothesis.
\end{enumerate}
\end{exmp}

\subsection{Principal ideals}

The most fundamental downsets are generated by single elements.

\begin{defn}[Principal ideal]
\label{def:principal-ideal}
For $s \in S$, the \emph{principal ideal} (or \emph{principal downset}) generated by $s$ is the lower cone\index{lower cone}:
\[
{\downarrow} s := \{t \in S : t \preceq s\}.
\]
\end{defn}

Principal ideals generate all downsets by union; they need not be atoms of the lattice.
They represent the simplest nontrivial design hypothesis: ``consider all scaffolds up to and including $s$.''

\begin{prop}[Principal ideals are the smallest downsets containing their generator]
\label{prop:principal-minimal}
For any $s \in S$, the principal ideal ${\downarrow} s$ is the smallest downset containing $s$:
\[
{\downarrow} s = \bigcap \{I \in \cD(S) : s \in I\}.
\]
\end{prop}

\begin{proof}
Let $I$ be any downset containing $s$.
If $t \in {\downarrow} s$, then $t \preceq s$ and $s \in I$, so $t \in I$ by downward closure.
Thus ${\downarrow} s \subseteq I$ for all such $I$, hence ${\downarrow} s \subseteq \bigcap \{I : s \in I\}$.
The reverse inclusion holds because ${\downarrow} s$ itself is a downset containing $s$.
\end{proof}

\begin{rem}[Terminology: ideals and downsets]
Here an \emph{order ideal} means a downset.
In domain theory an \emph{ideal} usually means a nonempty directed downset: every finite subset has an upper bound inside it.
Arbitrary downsets need not be directed, and a union of directed ideals need not be directed.
The complete lattice constructed here is the lattice of all downsets, not the directed-ideal completion.

The dual notion is an \emph{upset}: if $s\in U$ and $s\preceq t$, then $t\in U$.
\end{rem}

\section{The downset lattice $\cD(S)$}
\label{sec:downset-lattice}

We now establish the central algebraic result of this chapter: the collection of all downsets forms a complete distributive lattice.

\subsection{Order structure}

\begin{defn}[Downset ordering]
\label{def:downset-ordering}
Order the set of downsets $\cD(S)$ by set inclusion:
\[
I \leq J \iff I \subseteq J.
\]
\end{defn}

\begin{prop}[$\cD(S)$ is a poset]
\label{prop:downset-poset}
$(\cD(S), \subseteq)$ is a partially ordered set.
\end{prop}

\begin{proof}
This is immediate from the properties of set inclusion: reflexivity ($I \subseteq I$), antisymmetry ($I \subseteq J$ and $J \subseteq I$ implies $I = J$), and transitivity ($I \subseteq J$ and $J \subseteq K$ implies $I \subseteq K$).
\end{proof}

\subsection{Lattice operations}

The key observation is that downsets are closed under arbitrary intersections and arbitrary unions.

\begin{lem}[Closure under intersection]
\label{lem:intersection-closed}
Let $\{I_\alpha\}_{\alpha \in A}$ be any family of downsets (possibly infinite).
Then $\bigcap_{\alpha \in A} I_\alpha$ is a downset.
\end{lem}

\begin{proof}
Let $I = \bigcap_\alpha I_\alpha$.
Suppose $s \in I$ and $t \preceq s$.
Then $s \in I_\alpha$ for all $\alpha$.
Since each $I_\alpha$ is a downset, $t \in I_\alpha$ for all $\alpha$.
Therefore $t \in \bigcap_\alpha I_\alpha = I$.
\end{proof}

\begin{lem}[Closure under union]
\label{lem:union-closed}
Let $\{I_\alpha\}_{\alpha \in A}$ be any family of downsets.
Then $\bigcup_{\alpha \in A} I_\alpha$ is a downset.
\end{lem}

\begin{proof}
Let $I = \bigcup_\alpha I_\alpha$.
Suppose $s \in I$ and $t \preceq s$.
Then $s \in I_\beta$ for some $\beta \in A$.
Since $I_\beta$ is a downset, $t \in I_\beta \subseteq I$.
\end{proof}

These closure properties immediately yield the lattice structure.

\begin{thm}[The downset lattice is complete]
\label{thm:downset-complete-lattice}
$(\cD(S), \subseteq)$ is a \emph{complete lattice}:
\begin{enumerate}
\item \textbf{Meet (greatest lower bound):} For any family $\{I_\alpha\}_{\alpha \in A}$ of downsets,
\[
\bigwedge_{\alpha \in A} I_\alpha = \bigcap_{\alpha \in A} I_\alpha.
\]
\item \textbf{Join (least upper bound):} For any family $\{I_\alpha\}_{\alpha \in A}$ of downsets,
\[
\bigvee_{\alpha \in A} I_\alpha = \bigcup_{\alpha \in A} I_\alpha.
\]
\item \textbf{Bottom element:} $\bot = \emptyset$.
\item \textbf{Top element:} $\top = S$.
\end{enumerate}
\end{thm}

\begin{proof}
By Lemmas~\ref{lem:intersection-closed} and~\ref{lem:union-closed}, arbitrary intersections and unions of downsets are downsets, so meets and joins exist in $\cD(S)$.

For meets: The intersection $\bigcap_\alpha I_\alpha$ is a lower bound for the family (it is contained in each $I_\alpha$).
Any other lower bound $J$ must satisfy $J \subseteq I_\alpha$ for all $\alpha$, hence $J \subseteq \bigcap_\alpha I_\alpha$.
Thus the intersection is the \emph{greatest} lower bound.

For joins: The union $\bigcup_\alpha I_\alpha$ is an upper bound (it contains each $I_\alpha$).
Any other upper bound $J$ must satisfy $I_\alpha \subseteq J$ for all $\alpha$, hence $\bigcup_\alpha I_\alpha \subseteq J$.
Thus the union is the \emph{least} upper bound.

The empty set is a downset (vacuously) and is contained in every downset, so $\bot = \emptyset$.
The entire space $S$ is a downset and contains every downset, so $\top = S$.
\end{proof}

\begin{rem}[No downward closure needed for joins]
In many lattice constructions, the join requires a closure operation: one takes the union and then ``closes'' it.
For downsets, this is unnecessary---the union of downsets is already a downset.
This is a special feature of downsets in any poset, stemming from the ``hereditary'' nature of the downward-closure property.
\end{rem}

\subsection{Distributivity}

Because its operations are union and intersection, the downset lattice satisfies the following distributive identities.

\begin{thm}[Distributivity of $\cD(S)$]
\label{thm:downset-distributive}
The lattice $(\cD(S), \subseteq)$ is \emph{distributive}: for all downsets $I, J, K$,
\begin{align}
I \cap (J \cup K) &= (I \cap J) \cup (I \cap K), \label{eq:dist1} \\
I \cup (J \cap K) &= (I \cup J) \cap (I \cup K). \label{eq:dist2}
\end{align}
\end{thm}

\begin{proof}
Equations~\eqref{eq:dist1} and~\eqref{eq:dist2} are set-theoretic identities that hold for any subsets of any set, regardless of the downset property.
\end{proof}

\begin{cor}[Frame structure]
\label{cor:frame}
$\cD(S)$ is a \emph{frame} (also called a \emph{complete Heyting algebra}): finite meets distribute over arbitrary joins.
That is, for any downset $I$ and any family $\{J_\alpha\}$,
\[
I \cap \bigcup_\alpha J_\alpha = \bigcup_\alpha (I \cap J_\alpha).
\]
\end{cor}

\begin{proof}
Again, this is a set-theoretic identity holding for arbitrary subsets.
\end{proof}

\begin{rem}[Why distributivity matters]
Distributivity is not automatic for lattices.
For instance, the lattice of all subgroups of a group is typically \emph{not} distributive.
The distributivity of $\cD(S)$ has important consequences:
\begin{itemize}
\item Union and intersection have explicit set-theoretic definitions; their implementation cost depends on how downsets and the order are represented.
\item There is a Birkhoff representation theorem: every finite distributive lattice arises as the downset lattice of some poset.
\item The frame structure supports logical interpretation: downsets can be viewed as ``propositions'' about scaffold space, with $\cap$ as conjunction and $\cup$ as disjunction.
\end{itemize}
\end{rem}

\section{Atoms and the structure of $\cD(S)$}
\label{sec:atoms}

We now examine the fine structure of the downset lattice, beginning with its atoms.

\begin{defn}[Atom]
\label{def:atom}
In a lattice $L$ with bottom element $\bot$, an element $a \in L$ is an \emph{atom} if:
\begin{enumerate}
\item $a \neq \bot$, and
\item there is no element strictly between $\bot$ and $a$: if $\bot \leq x \leq a$, then $x = \bot$ or $x = a$.
\end{enumerate}
Equivalently, $a$ covers $\bot$ in the lattice order.
\end{defn}

For the downset lattice, atoms have a clean characterization.

\begin{prop}[Atoms of $\cD(S)$]
\label{prop:atoms}
The atoms of $\cD(S)$ are precisely the principal ideals of minimal scaffolds:
\[
\mathrm{Atoms}(\cD(S)) = \{{\downarrow} s : s \text{ is minimal in } S\}.
\]
Here ``minimal'' means $s$ has no proper predecessors in $S$: there is no $t \in S$ with $t \prec s$.
\end{prop}

\begin{proof}
Let $s$ be minimal in $S$.
The principal ideal ${\downarrow} s = \{t : t \preceq s\}$ contains only $s$ itself (since $s$ is minimal, any $t \preceq s$ must satisfy $t = s$).
Thus ${\downarrow} s = \{s\}$, a singleton downset.

There is no nonempty proper subset of that singleton, so no downset lies strictly between $\emptyset$ and $\{s\}$.
Hence $\{s\} = {\downarrow} s$ is an atom.

Conversely, let $I$ be an atom and choose $s\in I$.
The nonempty downset $\downarrow s$ is contained in $I$, hence equals $I$.
If $t\prec s$, then $\emptyset\neq\downarrow t\subsetneq\downarrow s$, contradicting that $I$ is an atom.
Thus $s$ is minimal and $I=\{s\}$.
\end{proof}

\begin{exmp}[Minimal scaffolds]
\label{ex:minimal-scaffolds}
In the full typed 2-core graph domain, with all typed cycles admitted:
\begin{itemize}
\item If we include an ``empty scaffold'' $\bot$ for acyclic molecules, then $\bot$ is the unique minimum of $S$, and $\{\bot\}$ is the unique atom of $\cD(S)$.
\item If we exclude $\bot$ and consider only cyclic scaffolds, the minimal scaffolds are the \emph{monocyclic} scaffolds with no sub-scaffolds: cyclopropane, cyclobutane, cyclopentane, cyclohexane, benzene, pyridine, pyrrole, furan, thiophene, etc.
Each gives rise to an atom $\{s\}$ in $\cD(S)$.
\item A finite atom alphabet alone does not bound the number of atoms of $\cD(S)$: uniformly typed cycles of every length are incomparable minimal nonempty 2-cores. A size bound or a finite library makes their number finite.
\end{itemize}
In a restricted observed scaffold set, a larger graph can be minimal simply because none of its proper sub-scaffolds is present.
\end{exmp}

\begin{rem}[Atomic vs. atomistic]
A lattice is \emph{atomic} if every nonzero element lies above an atom; it is \emph{atomistic} if every element is a join of atoms.
For finite $S$, $\cD(S)$ is atomic.
It is atomistic if and only if $S$ is an antichain: the atoms are singletons of minimal elements, and their unions cannot contain a nonminimal element.
For the two-element chain $a<b$, the lattice
$\{\emptyset,\{a\},\{a,b\}\}$ has one atom and is not atomistic.
Generation by principal ideals is the general statement used in this chapter.
\end{rem}

\section{Generation by principal ideals}
\label{sec:generation}

Principal ideals are the natural ``basis'' for constructing arbitrary downsets.

\begin{prop}[Every downset is a union of principal ideals]
\label{prop:union-of-principals}
For any downset $I \in \cD(S)$,
\[
I = \bigcup_{s \in I} {\downarrow} s.
\]
\end{prop}

\begin{proof}
($\supseteq$) If $s \in I$ and $t \in {\downarrow} s$ (i.e., $t \preceq s$), then $t \in I$ by downward closure of $I$.
Thus ${\downarrow} s \subseteq I$ for each $s \in I$, and therefore $\bigcup_{s \in I} {\downarrow} s \subseteq I$.

($\subseteq$) If $t \in I$, then $t \in {\downarrow} t$ (since $t \preceq t$), and $t \in I$ means ${\downarrow} t$ appears in the union.
Hence $t \in \bigcup_{s \in I} {\downarrow} s$.
\end{proof}

\begin{cor}[Join representation]
\label{cor:join-representation}
In the lattice $\cD(S)$, every downset is a join of principal ideals:
\[
I = \bigvee_{s \in I} {\downarrow} s.
\]
\end{cor}

\begin{proof}
Since join in $\cD(S)$ is union, this is a restatement of Proposition~\ref{prop:union-of-principals}.
\end{proof}

This representation theorem says that principal ideals are the ``building blocks'' of $\cD(S)$: any abstract search region can be expressed as a union of ``everything up to scaffold $s$'' regions.

\subsection{Maximal elements and efficient representation}

While Proposition~\ref{prop:union-of-principals} expresses $I$ as a union over all elements of $I$, this is redundant.
A more efficient representation uses only the maximal elements.

\begin{defn}[Maximal elements of a downset]
\label{def:maximal-elements}
For a downset $I$, define
\[
\max(I) := \{s \in I : \text{there is no } t \in I \text{ with } s \prec t\}.
\]
These are the ``top'' elements of $I$---scaffolds in $I$ that have no larger scaffolds in $I$.
\end{defn}

\begin{prop}[Downsets are determined by their maximal elements]
\label{prop:max-determines}
For a finite downset $I$, or more generally a downset in which every element lies below a maximal element,
\[
I = \bigcup_{s \in \max(I)} {\downarrow} s = {\downarrow}(\max(I)).
\]
\end{prop}

\begin{proof}
By Proposition~\ref{prop:union-of-principals}, $I = \bigcup_{s \in I} {\downarrow} s$.
For any $s \in I$, either $s \in \max(I)$, or there exists $t \in I$ with $s \prec t$.
In the latter case, $s \in {\downarrow} t$, so ${\downarrow} s \subseteq {\downarrow} t$.
In a finite downset, an ascending chain must terminate at a maximal element. The more general stated assumption gives the same conclusion directly.
Thus the union over all of $I$ equals the union over $\max(I)$.
\end{proof}

\begin{rem}[Scope and cost of antichain representation]
The maximal elements of a finite downset form an antichain and determine it by lower closure.
Membership reduces to checking $t\preceq s$ for some stored maximal $s$.
This can save storage, but the maximal antichain can be as large as the full downset and each comparison may be expensive.
The representation need not exist for an arbitrary infinite downset: the downset $\mathbb{N}$ in its usual order has no maximal element.
Finiteness or the cofinal-maximal-element condition must therefore be stated.
\end{rem}

\begin{exmp}[Antichain representation of a hit set]
\label{ex:antichain-hits}
Suppose an HTS campaign yields 50 hit molecules with 12 distinct scaffolds.
Let $I = \bigcup_{i=1}^{12} {\downarrow} s_i$ be the downset of all sub-scaffolds of hits.
Depending on the overlap structure:
\begin{itemize}
\item If the 12 scaffolds form an antichain (none is a sub-scaffold of another), then $\max(I) = \{s_1, \ldots, s_{12}\}$ has size 12.
\item If some scaffolds are sub-scaffolds of others (e.g., $s_3 \preceq s_7$), then $\max(I)$ is smaller.
\end{itemize}
In either case, storing 12 (or fewer) scaffolds suffices to represent all sub-scaffolds of all hits.
\end{exmp}

\section{Algebraic lattices and the finite-library regime}
\label{sec:algebraic}

For computational purposes, we need to understand the size and structure of $\cD(S)$.

\subsection{Size bounds}

\begin{prop}[Bounds on $|\cD(S)|$]
\label{prop:size-bounds}
Let $S$ be a finite poset with $n$ elements.
\begin{enumerate}
\item $|\cD(S)| \leq 2^n$ (at most one downset per subset).
\item $|\cD(S)| \geq n + 1$ (at least the principal ideals plus $\emptyset$).
\item Equality $|\cD(S)| = 2^n$ holds iff $S$ is an antichain (no two elements are comparable).
\item Equality $|\cD(S)| = n + 1$ holds iff $S$ is a total order (chain).
\end{enumerate}
\end{prop}

\begin{proof}
(1) Every downset is a subset of $S$.

(2) The $n$ principal downsets are nonempty and pairwise distinct: equality
${\downarrow}s={\downarrow}t$ implies $s\preceq t$ and $t\preceq s$, hence
$s=t$. Together with $\emptyset$ they give at least $n+1$ downsets.

(3) If $S$ is an antichain, every subset is a downset (downward closure is vacuous), so $|\cD(S)| = 2^n$.
Conversely, if $s \prec t$ for some $s, t \in S$, then $\{t\}$ is not a downset (missing $s$), so $|\cD(S)| < 2^n$.

(4) In the chain $s_1\prec\cdots\prec s_n$, every nonempty downset is an
initial segment $\{s_1,\ldots,s_j\}$ for one $1\leq j\leq n$.
Including the empty set gives exactly $n+1$ downsets.
\end{proof}

\begin{rem}[Scaffold spaces are between extremes]
Real scaffold spaces are neither antichains nor chains.
A scaffold space graded by ring count has $|S_k|$ scaffolds at level $k$, with comparability only between levels.
For such spaces, $|\cD(S)|$ is intermediate---typically much larger than $n + 1$ but much smaller than $2^n$.

A rough estimate: if $S$ has $L$ levels with $n_k$ scaffolds at level $k$, and scaffolds at level $k$ are roughly incomparable, then $|\cD(S)| \approx \prod_{k=0}^{L} (n_k + 1)$ (choosing independently which scaffolds to include at each level).
For drug-like scaffolds with $L \approx 5$ levels and $n_k \approx 10^2$ to $10^4$ per level, this gives $|\cD(S)| \sim 10^{10}$ to $10^{20}$---huge, but far from $2^{|S|}$.
\end{rem}

\subsection{Compact and algebraic structure}

\begin{defn}[Compact element]
\label{def:compact}
In a complete lattice $L$, an element $c$ is \emph{compact} (or \emph{finite}) if: whenever $c \leq \bigvee D$ for a \emph{directed} set $D$ (i.e., every finite subset of $D$ has an upper bound in $D$), there exists $d \in D$ with $c \leq d$.

Informally: $c$ is ``finitely determined''---if $c$ is below a directed join, it is already below one of the joinands.
\end{defn}

\begin{defn}[Algebraic lattice]
\label{def:algebraic-lattice}
A complete lattice $L$ is \emph{algebraic} if every element is a join of compact elements below it:
\[
\forall x \in L, \quad x = \bigvee \{c \in L : c \text{ compact}, c \leq x\}.
\]
\end{defn}

\begin{prop}[Principal ideals are compact]
\label{prop:principals-compact}
In $\cD(S)$, every principal ideal ${\downarrow} s$ is a compact element.
\end{prop}

\begin{proof}
Suppose ${\downarrow} s \subseteq \bigcup_\alpha I_\alpha$ where $\{I_\alpha\}$ is a directed family of downsets.
In particular, $s \in \bigcup_\alpha I_\alpha$, so $s \in I_\beta$ for some $\beta$.
Since $I_\beta$ is a downset, ${\downarrow} s \subseteq I_\beta$.
Thus ${\downarrow} s \leq I_\beta$, and we have found a single element of the directed family above ${\downarrow} s$.
\end{proof}

\begin{thm}[$\cD(S)$ is algebraic]
\label{thm:algebraic}
For any poset $S$, the downset lattice $\cD(S)$ is an algebraic lattice.
The compact elements are precisely the \emph{finitely generated} downsets---those of the form ${\downarrow} F = \bigcup_{s \in F} {\downarrow} s$ for some finite $F \subseteq S$.
\end{thm}

\begin{proof}
By Corollary~\ref{cor:join-representation}, every downset $I$ is a join of principal ideals: $I = \bigvee_{s \in I} {\downarrow} s$.
By Proposition~\ref{prop:principals-compact}, principal ideals are compact.
Thus every element is a join of compact elements below it.

For the characterization: ${\downarrow} F$ is a finite union of principal ideals, hence a finite join of compact elements.
Finite joins of compact elements are compact (exercise).
Conversely, write a compact downset $I$ as the directed union
\[
I=\bigcup_{F\subseteq I,\;F\ \mathrm{finite}}\downarrow F.
\]
The family is directed because $\downarrow F\cup\downarrow F'=\downarrow(F\cup F')$.
Compactness therefore gives a finite $F\subseteq I$ with $I\subseteq\downarrow F$.
The reverse inclusion holds because $I$ is a downset, so $I=\downarrow F$.
\end{proof}

\begin{cor}[Finite scaffold spaces yield finite downset lattices]
\label{cor:finite-finite}
If $S$ is finite, then $\cD(S)$ is finite (and hence trivially algebraic, since every element is compact).
\end{cor}

\section{Operations and their interpretation}
\label{sec:operations}

We now summarize the lattice operations and their meaning for drug discovery.

\begin{table}[htbp]
\centering
\caption{Lattice operations in $\cD(S)$ and their pharmaceutical interpretation.}
\label{tab:operations}
\begin{tabular}{@{}lll@{}}
\toprule
\textbf{Operation} & \textbf{Formula} & \textbf{Interpretation} \\
\midrule
Meet & $I \wedge J = I \cap J$ & Scaffolds in \emph{both} regions \\
Join & $I \vee J = I \cup J$ & Scaffolds in \emph{either} region \\
Bottom & $\bot = \emptyset$ & Empty design space \\
Top & $\top = S$ & All scaffolds under consideration \\
Principal ideal & ${\downarrow} s$ & All scaffolds up to $s$ \\
Order & $I \leq J$ iff $I \subseteq J$ & Region $I$ is contained in region $J$ \\
\bottomrule
\end{tabular}
\end{table}

\begin{exmp}[Combining design hypotheses]
\label{ex:combining}
A medicinal chemistry team has two competing hypotheses:
\begin{itemize}
\item Team A proposes focusing on bicyclic scaffolds related to hit $s_A$.
Let $I_A = {\downarrow} s_A$.
\item Team B proposes a different bicyclic series around hit $s_B$.
Let $I_B = {\downarrow} s_B$.
\end{itemize}

The combined design space is:
\begin{itemize}
\item $I_A \cup I_B$: scaffolds in \emph{either} series (union = join).
\item $I_A \cap I_B$: scaffolds in \emph{both} series (intersection = meet)---the common sub-scaffolds shared by both leads.
\end{itemize}

If $s_A$ and $s_B$ share a common sub-scaffold $s_C$, then ${\downarrow} s_C \subseteq I_A \cap I_B$.
The meet $I_A \cap I_B$ reveals the structural overlap between competing hypotheses.
\end{exmp}

\begin{framed}
\noindent\textbf{Why the lattice structure matters.}

\medskip
\noindent
Without the downset lattice, combining design hypotheses is ad hoc: one must manually enumerate scaffolds in each region and compute overlaps.
With the lattice:
\begin{itemize}
\item \textbf{Combining} two hypotheses is a join (union).
\item \textbf{Refining} a hypothesis by adding a constraint is a meet (intersection).
\item \textbf{Comparing} hypotheses uses the order ($I \subseteq J$ means $I$ is more restrictive).
\item \textbf{Finite representation} via maximal antichains can reduce storage; the saving depends on the width of the represented region.
\end{itemize}
The lattice provides a formal algebra for navigating the space of design hypotheses.
\end{framed}

\section{Connection to upsets and the duality}
\label{sec:upsets}

We briefly introduce the dual notion, which will be important in later chapters.

\begin{defn}[Upset (order filter)]
\label{def:upset}
A subset $U \subseteq S$ is an \emph{upset} (or \emph{order filter}, or \emph{up-closed set}) if:
\[
s \in U \text{ and } s \preceq t \implies t \in U.
\]
The set of all upsets of $S$ is denoted $\mathcal{U}(S)$.
\end{defn}

\begin{prop}[Duality]
\label{prop:duality}
The complement map $I \mapsto S \setminus I$ is an order-reversing bijection between $\cD(S)$ and $\mathcal{U}(S)$.
Hence $\mathcal{U}(S)$ is also a complete distributive lattice, with:
\begin{itemize}
\item Meet in $\mathcal{U}(S)$ = intersection (corresponding to join in $\cD(S)$).
\item Join in $\mathcal{U}(S)$ = union (corresponding to meet in $\cD(S)$).
\end{itemize}
\end{prop}

\begin{proof}
If $I$ is a downset, we show $U = S \setminus I$ is an upset.
Suppose $s \in U$ and $s \preceq t$.
If $t \in I$, then $s \in I$ (downward closure), contradicting $s \in U = S \setminus I$.
Hence $t \in U$.

The map is clearly a bijection.
Order-reversal: $I \subseteq J$ iff $S \setminus J \subseteq S \setminus I$.
Lattice duality follows.
\end{proof}

\begin{rem}[Upsets in fragment-based design]
Recall from Chapter~\ref{ch:poset} that fragment hits with scaffolds $s_1, \ldots, s_k$ constrain the elaboration space to $\bigcap_i {\uparrow} s_i$---an intersection of upper cones, which is an upset.
The duality between downsets and upsets reflects the asymmetry between:
\begin{itemize}
\item Downsets: ``what simpler scaffolds are included?'' (fragmentation direction).
\item Upsets: ``what more complex scaffolds are allowed?'' (elaboration direction).
\end{itemize}
Chapter~\ref{ch:galois} will connect both to the Galois connection between molecules and scaffolds.
\end{rem}

\section{Preview: from downsets to the Galois connection}
\label{sec:preview}

Chapter~\ref{ch:galois} starts with arbitrary scaffold subsets.
An algorithm that accepts molecules according only to scaffold identity selects a subset of $S$; it selects a downset only if downward closure is an additional requirement.

\begin{defn}[Preview: abstraction and concretization]
\label{def:preview-galois}
For a deterministic scaffold map, define
\[
\alpha(A)=\{\Scaf(m):m\in A\},\qquad
\gamma(B)=\{m:\Scaf(m)\in B\}.
\]
The fiber closure $c(A)=\gamma(\alpha(A))$ is the smallest scaffold-saturated superset of $A$.
If the admissible abstract regions must be downsets, the corresponding closure is instead
$c_\downarrow(A)=\gamma(\downarrow\alpha(A))$.
These two closures need not coincide.
\end{defn}

Inverse image preserves unions and intersections, so both set operations on scaffold regions have exact molecular counterparts.
For a nonempty target set in a finite library, $|c(A)|/|A|$ measures the output expansion forced by a whole-fiber selection rule.
It does not lower-bound the number of evaluations: a rule can return the entire library without evaluating its members.

\section*{Summary}

The downset lattice $\cD(S)$ is complete and distributive, with intersection as meet and union as join.
Every downset is a union of principal ideals.
A finite downset is determined by its maximal antichain; infinite downsets require an additional condition for that representation.
The compact elements are precisely the finitely generated downsets.

This algebra specifies operations on structural regions.
It does not establish activity monotonicity, cheap graph matching, or a sampling-complexity bound.
Those require separate response and observation models.

\section*{Exercises}
\begin{exer}[Chemistry: downsets from hit scaffolds]
\label{exer:ch5-chem}
Consider four scaffolds: benzene, pyridine, quinoline (benzene fused to pyridine), and isoquinoline (isomer\index{isomer} of quinoline with different nitrogen position).
\begin{enumerate}
\item Draw the Hasse diagram of these four scaffolds under the substructure order $\preceq$.
\item List all downsets of this 4-element poset.
How many are there?
\item For the downset $I = \{$benzene, pyridine$\}$, determine whether there exists a downset $J$ such that $J$ is a proper superset of $I$ but does not contain quinoline.
\item Interpret the downset $\{$benzene$\}$ in terms of a design constraint a medicinal chemist might impose.
\end{enumerate}
\end{exer}

\begin{exer}[Mathematics: verifying lattice properties]
\label{exer:ch5-math}
Let $(P, \preceq)$ be a finite poset and let $\cD(P)$ be its downset lattice.
\begin{enumerate}
\item Prove that the intersection of any family of downsets is a downset.
\item Prove that the union of any family of downsets is a downset.
\item Conclude that $\cD(P)$ is a complete lattice with meet = intersection and join = union.
\item Prove that $\cD(P)$ is distributive by verifying $I \cap (J \cup K) = (I \cap J) \cup (I \cap K)$ directly from set theory.
\item Prove that if $P$ has a unique minimal element $\bot$, then $\cD(P)$ has a unique atom $\{\bot\}$.
\end{enumerate}
\end{exer}

\begin{exer}[Machine learning: downsets as hypothesis classes]
\label{exer:ch5-ml}
Consider a supervised learning setting where the input is a molecule $m$ and the label is binary (active/inactive).
A scaffold-based classifier predicts ``active'' for all molecules whose scaffold lies in some downset $I \subseteq S$.
\begin{enumerate}
\item For a fixed scaffold poset $S$ with $n$ elements, how many distinct scaffold-based classifiers are there?
(Express in terms of $|\cD(S)|$.)
\item Prove that the set of scaffold-based classifiers is closed under conjunction and disjunction: if $h_I$ and $h_J$ are classifiers corresponding to downsets $I$ and $J$, then $h_I \wedge h_J$ corresponds to $I \cap J$ and $h_I \vee h_J$ corresponds to $I \cup J$.
\item Suppose the true active set $A \subseteq \Mol$ has scaffold image $\alpha(A) = \{s_1, \ldots, s_k\}$.
What is the smallest downset $I$ such that the classifier $h_I$ has perfect recall (no false negatives)?
What is its false positive rate in terms of the scaffold-to-molecule distribution?
\item Relate this to the ``scaffold overhead'' concept previewed for Chapter~\ref{ch:galois}.
\end{enumerate}
\end{exer}

\begin{exer}[Computation: implementing downset operations]
\label{exer:ch5-comp}
Using a cheminformatics toolkit (e.g., RDKit) and a small library $\mathcal{L}$ (say $10^3$--$10^4$ molecules):
\begin{enumerate}
\item Compute the scaffold set $S = \Scaf(\mathcal{L})$.
\item Implement the substructure order: a function \texttt{leq(s, t)} returning \texttt{True} iff $s \preceq t$.
\item Implement downset representation: store a downset as the list of its maximal elements.
\item Implement \texttt{meet(I, J)} and \texttt{join(I, J)} for downsets in antichain representation.
\item Test: pick two scaffolds $s_1, s_2$ from your library, compute $I = {\downarrow} s_1$, $J = {\downarrow} s_2$, and verify that $I \cap J = {\downarrow}(s_1 \wedge s_2)$ when the meet exists, or is the union of maximal common sub-scaffolds otherwise.
\item Measure the compression: compare $|I|$ (number of scaffolds in the downset) to $|\max(I)|$ (number of maximal elements).
\end{enumerate}
\end{exer}

\begin{exer}[Advanced: Birkhoff's theorem]
\label{exer:birkhoff}
Birkhoff's representation theorem states that every finite distributive lattice is isomorphic to the downset lattice of some poset.
\begin{enumerate}
\item State the theorem precisely, including how to recover the poset from the lattice.
\item Verify the theorem for the lattice $\{0 < a < b < 1\}$ (a 4-element chain): what is the poset?
\item Verify for the lattice of divisors of 6 under divisibility.
\item Explain why this theorem implies that distributive lattices are ``essentially'' downset lattices.
\end{enumerate}
\end{exer}
% END SOURCE: chapters/ch05_lattice.tex


\chapter{The Galois Connection}
\label{ch:galois}
% BEGIN SOURCE: chapters/ch06_galois.tex
% ======================================================================
% Chapter 6 content file: chapters/ch06_galois.tex
% (Do NOT include \chapter{...} here; main.tex already does that.)
% ======================================================================

\section{Abstraction as a mathematical object}

The previous chapters built two complementary views of chemical space\index{chemical space}:
\begin{itemize}
\item A \emph{concrete} space of molecules $\Mol$, equipped with a scaffold\index{scaffold} map $\Scaf:\Mol\to S$ (Chapter~\ref{ch:scaffolds}).
\item An \emph{algebra} of scaffold\index{scaffold} \emph{regions} (downsets) $\cD(S)$, where meets and joins exist even when $S$ itself is not a lattice\index{lattice} (Chapter~\ref{ch:lattice}).
\end{itemize}

This chapter makes precise the sense in which a scaffold map is an \emph{abstraction}:
it induces an adjunction between sets of molecules and sets of scaffolds.
The key technical tool is a \emph{Galois connection}\index{Galois connection}, familiar from order theory and abstract interpretation.
Once this structure is explicit, one obtains an unavoidable closure operator\index{closure operator}, and from it a quantitative notion of \emph{screening\index{screening} overhead}.

\begin{framed}
\noindent\textbf{The operational question.}

\medskip
\noindent
Suppose there is a (possibly unknown) set $A\subseteq \Mol$ of molecules that we would like to find:
the true actives, the top-$k$ dockers, the molecules above a potency threshold, etc.
A scaffold-based strategy often works at the level of \emph{regions of scaffolds}:
choose a set of scaffolds $B\subseteq S$, and then enumerate or screen molecules in $\gamma(B)$.

The Galois connection\index{Galois connection} tells you the smallest \emph{scaffold-saturated} superset of $A$ that any such strategy can ever hope to isolate deterministically.
That superset is the closure $c(A)=\gamma(\alpha(A))$.
\end{framed}

\section{The concrete and abstract domains}

We work with powersets ordered by inclusion.
These are complete lattices, with meets and joins given by intersections and unions.

\begin{defn}[Concrete and abstract domains]
\label{def:concrete-abstract-domains}
Let $\Mol$ be the set of molecules (isomorphism classes; Chapter~\ref{ch:molecules}) and let $S$ be the set of scaffolds (Chapter~\ref{ch:scaffolds}).
Define
\[
(\mathcal{P}(\Mol),\subseteq),\qquad
(\mathcal{P}(S),\subseteq),
\]
the posets of subsets ordered by inclusion.
\end{defn}

\begin{rem}
In the \emph{finite-library regime} (Chapter~\ref{ch:poset}\index{poset}), one replaces $\Mol$ by a finite library $\mathcal{L}\subseteq \Mol$ and $S$ by $S_{\mathcal{L}}:=\Scaf(\mathcal{L})$.
All statements below remain true verbatim after this restriction, and cardinalities such as $|c(A)|$ are then finite and computable.
We will make this restriction explicit when defining numerical overhead.
\end{rem}

\section{Abstraction and concretization}

The scaffold map $\Scaf:\Mol\to S$ induces two canonical maps between powersets: direct image (abstraction) and inverse image (concretization).

\begin{defn}[Abstraction and concretization]
\label{def:alpha-gamma}
Define maps
\[
\alpha:\mathcal{P}(\Mol)\to\mathcal{P}(S),\qquad
\gamma:\mathcal{P}(S)\to\mathcal{P}(\Mol)
\]
by
\begin{align*}
\alpha(A) &:= \{\Scaf(m): m\in A\}, \\
\gamma(B) &:= \{m\in \Mol : \Scaf(m)\in B\} = \Scaf^{-1}(B).
\end{align*}
\end{defn}

\begin{rem}[What these maps mean]
\label{rem:alpha-gamma-meaning}
\begin{itemize}
\item $\alpha(A)$ forgets everything about $A$ except \emph{which scaffold classes appear}.
\item $\gamma(B)$ expands a scaffold hypothesis $B$ back to the set of all molecules consistent with it.
\end{itemize}
In abstract interpretation language, $\alpha$ computes the \emph{best abstraction} of a concrete set, while $\gamma$ gives the \emph{concretization} of an abstract set.
\end{rem}

\begin{prop}[Monotonicity]
\label{prop:alpha-gamma-monotone}
Both $\alpha$ and $\gamma$ are monotone maps:
\[
A\subseteq A' \implies \alpha(A)\subseteq \alpha(A'),
\qquad
B\subseteq B' \implies \gamma(B)\subseteq \gamma(B').
\]
\end{prop}

\begin{proof}
If $A\subseteq A'$ and $s\in \alpha(A)$, then $s=\Scaf(m)$ for some $m\in A\subseteq A'$, hence $s\in\alpha(A')$.
Similarly, if $B\subseteq B'$ and $m\in\gamma(B)$, then $\Scaf(m)\in B\subseteq B'$, so $m\in\gamma(B')$.
\end{proof}

\section{The Galois connection}

We now state the central structural theorem.

\begin{thm}[Galois connection induced by a scaffold map]
\label{thm:galois-connection}
For all $A\subseteq \Mol$ and $B\subseteq S$,
\[
\alpha(A)\subseteq B \quad\Longleftrightarrow\quad A\subseteq \gamma(B).
\]
Equivalently, $\alpha$ is left adjoint to $\gamma$ (written $\alpha \dashv \gamma$) as maps of posets.
\end{thm}

\begin{proof}
($\Rightarrow$) Assume $\alpha(A)\subseteq B$ and let $m\in A$.
Then $\Scaf(m)\in \alpha(A)\subseteq B$, so by definition $m\in\gamma(B)$.
Hence $A\subseteq \gamma(B)$.

($\Leftarrow$) Assume $A\subseteq\gamma(B)$ and let $s\in \alpha(A)$.
Then $s=\Scaf(m)$ for some $m\in A$.
Since $m\in\gamma(B)$, we have $\Scaf(m)\in B$, hence $s\in B$.
Therefore $\alpha(A)\subseteq B$.
\end{proof}

\begin{figure}[htbp]
\centering
\begin{tikzcd}[column sep=huge,row sep=large]
\mathcal{P}(\Mol) \arrow[r, shift left=1.2, "\alpha"] &
\mathcal{P}(S) \arrow[l, shift left=1.2, "\gamma"]
\end{tikzcd}
\caption{The scaffold map induces a Galois connection between sets of molecules and sets of scaffolds.}
\label{fig:galois-connection}
\end{figure}

\subsection{Adjoints preserve lattice\index{lattice} structure}

Because $\mathcal{P}(\Mol)$ and $\mathcal{P}(S)$ are complete lattices, the adjunction immediately implies preservation properties.
We record the two that are used repeatedly later.

\begin{prop}[Unions and intersections]
\label{prop:alpha-gamma-union-intersection}
For any family $\{A_i\}_{i\in I}\subseteq \mathcal{P}(\Mol)$ and any family $\{B_j\}_{j\in J}\subseteq \mathcal{P}(S)$,
\[
\alpha\Big(\bigcup_{i\in I} A_i\Big)=\bigcup_{i\in I}\alpha(A_i),
\qquad
\gamma\Big(\bigcap_{j\in J} B_j\Big)=\bigcap_{j\in J}\gamma(B_j).
\]
\end{prop}

\begin{proof}
For $\alpha$, a scaffold $s$ lies in $\alpha(\bigcup_i A_i)$ iff $s=\Scaf(m)$ for some $m\in \bigcup_i A_i$, i.e.\ for some $m\in A_i$ for some $i$, which is equivalent to $s\in \bigcup_i \alpha(A_i)$.
For $\gamma$, a molecule\index{molecule} $m$ lies in $\gamma(\bigcap_j B_j)$ iff $\Scaf(m)\in \bigcap_j B_j$, i.e.\ iff $\Scaf(m)\in B_j$ for all $j$, which is equivalent to $m\in \bigcap_j \gamma(B_j)$.
\end{proof}

\begin{rem}
In order-theoretic terms: left adjoints preserve arbitrary joins, and right adjoints preserve arbitrary meets.
Here this reduces to elementary facts about direct and inverse image.
\end{rem}

\subsection{A Galois insertion}

Because $S$ was defined as the image of $\Scaf$ (Definition~\ref{def:scaf-map}), the adjunction is in fact an insertion.

\begin{prop}[Insertion property]
\label{prop:galois-insertion}
For every $B\subseteq S$,
\[
\alpha(\gamma(B)) = B.
\]
Consequently, $\gamma$ is injective and order-reflecting: if $\gamma(B)\subseteq \gamma(B')$ then $B\subseteq B'$.
\end{prop}

\begin{proof}
Let $s\in \alpha(\gamma(B))$.
Then $s=\Scaf(m)$ for some $m$ with $\Scaf(m)\in B$, hence $s\in B$.
So $\alpha(\gamma(B))\subseteq B$.

Conversely, let $s\in B$.
Since $s\in S$ there exists $m_s\in \Mol$ with $\Scaf(m_s)=s$.
Then $m_s\in\gamma(B)$, so $s=\Scaf(m_s)\in \alpha(\gamma(B))$.
Thus $B\subseteq \alpha(\gamma(B))$.
\end{proof}

\begin{rem}[Interpretation]
The abstract domain $\mathcal{P}(S)$ embeds into the concrete domain $\mathcal{P}(\Mol)$ as the family of \emph{scaffold-saturated} subsets of molecules (defined below).
The embedding is exactly $\gamma$.
\end{rem}

\section{Closure and scaffold saturation}

The compositional operator $c=\gamma\circ \alpha$ is the mathematical form of ``losing information by going to scaffolds and then coming back.''

\begin{defn}[Scaffold closure]
\label{def:scaffold-closure}
Define the \emph{scaffold closure} operator
\[
c:\mathcal{P}(\Mol)\to\mathcal{P}(\Mol),
\qquad
c(A):=(\gamma\circ\alpha)(A)=\gamma(\alpha(A)).
\]
\end{defn}

\begin{prop}[Explicit form]
\label{prop:closure-explicit}
For any $A\subseteq \Mol$,
\[
c(A)=\bigcup_{s\in \alpha(A)} \Fib(s)
=\{m\in\Mol : \Scaf(m)\in \alpha(A)\}.
\]
\end{prop}

\begin{proof}
By definition, $m\in c(A)$ iff $\Scaf(m)\in \alpha(A)$.
This is equivalent to saying: there exists $m'\in A$ with $\Scaf(m)=\Scaf(m')$.
Equivalently, $m$ lies in the fiber\index{fiber} of some scaffold appearing in $A$, hence in $\bigcup_{s\in \alpha(A)}\Fib(s)$.
\end{proof}

\begin{defn}[Scaffold-saturated subset]
\label{def:scaffold-saturated}
A subset $C\subseteq \Mol$ is \emph{scaffold-saturated} if it is a union of whole scaffold fibers:
\[
m\in C \ \text{and}\ \Scaf(m')=\Scaf(m)\ \Longrightarrow\ m'\in C.
\]
Equivalently, $C=\gamma(B)$ for some $B\subseteq S$.
\end{defn}

\begin{prop}[Closed sets are exactly scaffold-saturated sets]
\label{prop:closed-iff-saturated}
A set $C\subseteq \Mol$ is scaffold-saturated if and only if it is a fixed point of $c$:
\[
C \text{ is scaffold-saturated } \quad\Longleftrightarrow\quad c(C)=C.
\]
\end{prop}

\begin{proof}
If $C$ is scaffold-saturated, then $C=\gamma(B)$ for some $B\subseteq S$.
Using Proposition~\ref{prop:galois-insertion},
\[
c(C)=\gamma(\alpha(\gamma(B)))=\gamma(B)=C.
\]
Conversely, if $c(C)=C$, then $C=\gamma(\alpha(C))$, hence $C$ is of the form $\gamma(B)$ and is scaffold-saturated.
\end{proof}

\subsection{Closure operator\index{closure operator} axioms}

\begin{prop}[Closure axioms]
\label{prop:closure-axioms}
The scaffold closure $c$ is a closure operator on $\mathcal{P}(\Mol)$:
for all $A,A'\subseteq \Mol$,
\begin{enumerate}
\item \textbf{Extensive:} $A\subseteq c(A)$.
\item \textbf{Monotone:} $A\subseteq A' \implies c(A)\subseteq c(A')$.
\item \textbf{Idempotent:} $c(c(A))=c(A)$.
\end{enumerate}
\end{prop}

\begin{proof}
Extensiveness: if $m\in A$, then $\Scaf(m)\in \alpha(A)$, hence $m\in\gamma(\alpha(A))=c(A)$.

Monotonicity follows from monotonicity of $\alpha$ and $\gamma$ (Proposition~\ref{prop:alpha-gamma-monotone}).

Idempotence: using Proposition~\ref{prop:galois-insertion} with $B=\alpha(A)$,
\[
c(c(A))=\gamma(\alpha(\gamma(\alpha(A))))=\gamma(\alpha(A))=c(A).
\]
\end{proof}

\subsection{The minimal scaffold-saturated superset}

The closure $c(A)$ is the smallest scaffold-saturated set containing $A$.

\begin{thm}[Minimal closed superset]
\label{thm:minimal-closed-superset}
For any $A\subseteq \Mol$:
\begin{enumerate}
\item $c(A)$ is scaffold-saturated and contains $A$.
\item If $C$ is scaffold-saturated and $A\subseteq C$, then $c(A)\subseteq C$.
\end{enumerate}
Equivalently, $c(A)$ is the unique smallest scaffold-saturated superset of $A$.
\end{thm}

\begin{proof}
(1) We already have $A\subseteq c(A)$ from Proposition~\ref{prop:closure-axioms}.
Also $c(A)$ is scaffold-saturated because $c(A)=\gamma(\alpha(A))$ is of the form $\gamma(B)$.

(2) If $C$ is scaffold-saturated, then $C=\gamma(B)$ for some $B\subseteq S$.
Assume $A\subseteq C=\gamma(B)$.
By the Galois connection (Theorem~\ref{thm:galois-connection}), this is equivalent to $\alpha(A)\subseteq B$.
Applying $\gamma$ (monotone) gives $\gamma(\alpha(A))\subseteq \gamma(B)$, i.e.\ $c(A)\subseteq C$.
\end{proof}

\begin{rem}[Closure as ``loss of resolution'']
If $A$ contains a single molecule\index{molecule} $m$, then $c(\{m\})=\Fib(\Scaf(m))$:
the closure forgets which decorations produced $m$ and returns the entire scaffold class.
At the other extreme, if $A$ already contains whole fibers, then $c(A)=A$ and there is no loss.
\end{rem}

\section{Scaffold overhead and output-cardinality bounds}

\subsection{Overhead in the finite-library regime}

To quantify the practical impact of closure, we now restrict to a finite library.

\begin{defn}[Library-restricted closure]
\label{def:library-closure}
Fix a finite library $\mathcal{L}\subseteq \Mol$ and define $S_{\mathcal{L}}:=\Scaf(\mathcal{L})$.
For $s\in S_{\mathcal{L}}$, write $\Fib_{\mathcal{L}}(s):=\{m\in \mathcal{L}:\Scaf(m)=s\}$ for the fiber\index{fiber} within the library.
For $A\subseteq \mathcal{L}$ and $B\subseteq S_{\mathcal{L}}$, define
\begin{align*}
\alpha_{\mathcal{L}}(A) &:= \{\Scaf(m): m\in A\},\\
\gamma_{\mathcal{L}}(B) &:= \{m\in \mathcal{L}: \Scaf(m)\in B\},\\
c_{\mathcal{L}}(A) &:= \gamma_{\mathcal{L}}(\alpha_{\mathcal{L}}(A)).
\end{align*}
\end{defn}

\begin{defn}[Scaffold overhead]
\label{def:scaffold-overhead}
For a nonempty set $A\subseteq \mathcal{L}$, define the \emph{scaffold overhead}
\[
\ScafOverhead(A) := \frac{|c_{\mathcal{L}}(A)|}{|A|}.
\]
\end{defn}

\begin{rem}
$\ScafOverhead(A)\ge 1$ always, with equality if and only if $A$ is scaffold-saturated within the library.
\end{rem}

\begin{prop}[Overhead decomposes by scaffold fibers]
\label{prop:overhead-fiber-decomp}
Let $A\subseteq \mathcal{L}$ be nonempty and write $a_s:=|A\cap \Fib_{\mathcal{L}}(s)|$ for the number of selected molecules in scaffold class $s$ (within the library).
Then
\[
|c_{\mathcal{L}}(A)|=\sum_{s:\, a_s>0} |\Fib_{\mathcal{L}}(s)|,
\qquad
|A|=\sum_{s} a_s.
\]
In particular,
\[
\ScafOverhead(A)=\frac{\sum_{s:\,a_s>0} |\Fib_{\mathcal{L}}(s)|}{\sum_s a_s}.
\]
\end{prop}

\begin{proof}
By Proposition~\ref{prop:closure-explicit}, $c_{\mathcal{L}}(A)$ is the union of entire fibers for those scaffolds that appear in $A$.
These fibers are disjoint, so cardinalities add.
The expression for $|A|$ is the disjoint union decomposition of $A$ by scaffold classes.
\end{proof}

\begin{exmp}[A two-fiber toy calculation]
\label{ex:toy-overhead}
Suppose the library decomposes into two scaffold fibers of sizes $|\Fib_{\mathcal{L}}(s_1)|=1000$ and $|\Fib_{\mathcal{L}}(s_2)|=5$.
If $A$ contains one molecule from each fiber, then $|A|=2$ while $|c_{\mathcal{L}}(A)|=1005$, so
\[
\ScafOverhead(A)=\frac{1005}{2}=502.5.
\]
If instead $A$ contains \emph{all} molecules from both fibers, then $c_{\mathcal{L}}(A)=A$ and $\ScafOverhead(A)=1$.
This illustrates the two drivers of overhead:
\begin{enumerate}
\item how many scaffold classes $A$ touches, and
\item how much of each touched class $A$ already includes.
\end{enumerate}
\end{exmp}

\subsection{A deterministic lower bound for region-selection methods}

We now formalize the sense in which closure is a lower bound.

\begin{thm}[Minimal output for scaffold region selection]
\label{thm:min-output-region}
Let $\mathcal{L}$ be a finite library and let $A\subseteq \mathcal{L}$.
Consider any method that selects molecules only by choosing a scaffold set $B\subseteq S_{\mathcal{L}}$ and then returning \emph{all} molecules in those scaffold classes, i.e.\ it outputs $\gamma_{\mathcal{L}}(B)$.
If the method guarantees $A\subseteq \gamma_{\mathcal{L}}(B)$, then necessarily
\[
c_{\mathcal{L}}(A)\subseteq \gamma_{\mathcal{L}}(B),
\]
and hence the smallest possible output size among such methods is $|c_{\mathcal{L}}(A)|$.
\end{thm}

\begin{proof}
Assume $A\subseteq \gamma_{\mathcal{L}}(B)$.
By the Galois connection (Theorem~\ref{thm:galois-connection} applied to the restricted maps), this is equivalent to $\alpha_{\mathcal{L}}(A)\subseteq B$.
Applying $\gamma_{\mathcal{L}}$ gives
\[
c_{\mathcal{L}}(A)=\gamma_{\mathcal{L}}(\alpha_{\mathcal{L}}(A))\subseteq \gamma_{\mathcal{L}}(B).
\]
Since $c_{\mathcal{L}}(A)$ itself is attainable by choosing $B=\alpha_{\mathcal{L}}(A)$, it is the smallest possible scaffold-region output containing $A$.
\end{proof}

\subsection{Scaffold-sound algorithms}

The region-selection model above is slightly stylized.
In practice, one might imagine an algorithm that streams molecules and filters them, but only using scaffold information.
The natural formalization is that the algorithm's accept/reject decision is constant on fibers.

\begin{defn}[Scaffold-sound algorithm]
\label{def:scaffold-sound-algorithm}
Fix a scaffold map $\Scaf$ and a library $\mathcal{L}$.
An algorithm $\mathcal{A}$ that outputs a subset $\mathcal{A}(\mathcal{L})\subseteq \mathcal{L}$ is \emph{scaffold-sound} if membership depends only on scaffold class:
for all $m,m'\in \mathcal{L}$,
\[
\Scaf(m)=\Scaf(m')
\quad\Longrightarrow\quad
\big(m\in \mathcal{A}(\mathcal{L}) \iff m'\in \mathcal{A}(\mathcal{L})\big).
\]
Equivalently, $\mathcal{A}(\mathcal{L})$ is scaffold-saturated (Definition~\ref{def:scaffold-saturated}).
\end{defn}

\begin{prop}[Scaffold-sound outputs are closed]
\label{prop:sound-closed}
An algorithm is scaffold-sound if and only if its output $O\subseteq \mathcal{L}$ satisfies $c_{\mathcal{L}}(O)=O$.
\end{prop}

\begin{proof}
By Proposition~\ref{prop:closed-iff-saturated}, a set is scaffold-saturated iff it is a fixed point of closure.
\end{proof}

\begin{cor}[Deterministic overhead lower bound for scaffold-sound methods]
\label{cor:scaffold-sound-lower-bound}
Let $A\subseteq \mathcal{L}$.
Any scaffold-sound algorithm that guarantees $A\subseteq O$ for its output $O$ must output at least $|c_{\mathcal{L}}(A)|$ molecules:
\[
A\subseteq O,\;\; c_{\mathcal{L}}(O)=O
\quad\Longrightarrow\quad
|O|\ge |c_{\mathcal{L}}(A)|.
\]
Equivalently, the best possible deterministic overhead for covering $A$ under scaffold-soundness is $\ScafOverhead(A)$.
\end{cor}

\begin{proof}
If $A\subseteq O$, then by monotonicity of closure, $c_{\mathcal{L}}(A)\subseteq c_{\mathcal{L}}(O)=O$.
Taking cardinalities yields $|O|\ge |c_{\mathcal{L}}(A)|$.
\end{proof}

\begin{rem}[Output size is not oracle complexity]
\label{rem:not-universal-lower-bound}
The bound concerns the number of molecules in a saturated output. It gives no
lower bound on response evaluations: returning the entire library costs zero
response queries and contains every target set. Exact identification,
precision, or a bound on output size changes that problem and requires a
separate oracle argument. Randomization does not change the pathwise closure
bound on a successful output that remains fiber-saturated. A method using
additional molecular information may return an unsaturated subset, whether
its guarantee is deterministic or statistical.
Chapter~\ref{ch:complexity} specifies these different decision problems.
\end{rem}

\section{From sets of scaffolds to downset\index{downset} regions}

Chapter~\ref{ch:lattice} argued that practical scaffold ``hypotheses'' are often \emph{downward closed} in the scaffold order.
This corresponds to restricting the abstract domain from $\mathcal{P}(S)$ to the downset\index{downset} lattice $\cD(S)$.

\begin{defn}[Downward-closed abstraction]
\label{def:alpha-down}
Let ${\downarrow}(\cdot)$ denote downward closure in the scaffold poset\index{poset} $(S,\preceq)$ (Chapter~\ref{ch:poset}).
Define
\[
\alpha_{\downarrow}:\mathcal{P}(\Mol)\to \cD(S),\qquad
\alpha_{\downarrow}(A):={\downarrow}\alpha(A).
\]
\end{defn}

\begin{prop}[Galois connection with downsets]
\label{prop:galois-downsets}
Let $\gamma$ be as in Definition~\ref{def:alpha-gamma}, but viewed as a map $\gamma:\cD(S)\to \mathcal{P}(\Mol)$ by restriction of its domain.
Then for all $A\subseteq \Mol$ and all downsets $I\in \cD(S)$,
\[
\alpha_{\downarrow}(A)\subseteq I
\quad\Longleftrightarrow\quad
A\subseteq \gamma(I).
\]
\end{prop}

\begin{proof}
($\Rightarrow$) If $\alpha_{\downarrow}(A)\subseteq I$ then in particular $\alpha(A)\subseteq I$.
By Theorem~\ref{thm:galois-connection}, this implies $A\subseteq \gamma(I)$.

($\Leftarrow$) If $A\subseteq \gamma(I)$ then $\alpha(A)\subseteq I$ by Theorem~\ref{thm:galois-connection}.
Taking downward closures gives ${\downarrow}\alpha(A)\subseteq {\downarrow}I=I$ since $I$ is a downset.
Thus $\alpha_{\downarrow}(A)\subseteq I$.
\end{proof}

\begin{rem}[Two closures, two notions of ``conservatism'']
There are now two different ``closures'' that can appear in scaffold reasoning:
\begin{itemize}
\item \textbf{Fiber saturation} (this chapter): $c(A)=\gamma(\alpha(A))$, which closes under $\Scaf$-equivalence.
\item \textbf{Order saturation} (Chapter~\ref{ch:lattice}): replacing a scaffold set $B$ by its downward closure ${\downarrow}B$ to represent a region stable under simplification.
\end{itemize}
Imposing both simultaneously yields the composite $\gamma({\downarrow}\alpha(A))$, which is larger than $c(A)$ and corresponds to a stronger notion of abstraction.
Which closure is appropriate is a modeling choice, not a theorem.
\end{rem}

\section{Experiment: the overhead curve for top-$k$ sets}

The closure operator becomes meaningful only when we measure it on real screening outcomes.
The signature plot for this chapter is the scaffold overhead of top-$k$ molecules as a function of $k$.

\begin{defn}[Top-$k$ set]
\label{def:topk-set}
Fix a finite library $\mathcal{L}$ and a scoring function\index{scoring function} $F:\mathcal{L}\to \R$ where smaller is better (e.g.\ docking energies).
For $k\in\{1,\dots,|\mathcal{L}|\}$, define $T_k\subseteq \mathcal{L}$ to be the set of the $k$ best-scoring molecules, using one fixed tie-breaking order for all $k$.
\end{defn}

\begin{exper}[Overhead curve]
\label{exp:overhead-curve}
Compute $\ScafOverhead(T_k)$ for a range of $k$ (e.g.\ $k=10,100,1000,\dots$) and plot the curve $k\mapsto \ScafOverhead(T_k)$.
\end{exper}

\begin{rem}[Expected qualitative behavior]
If high-scoring molecules concentrate in a small number of scaffold fibers, then as $k$ increases the top-$k$ set tends to fill out those same fibers before it starts touching many new ones.
In that regime, $|c_{\mathcal{L}}(T_k)|$ stays roughly constant while $|T_k|=k$ increases, so overhead decreases.
When the top-$k$ list begins to include molecules from new scaffold classes, the closure size jumps and the curve can exhibit upward discontinuities.
Thus $\ScafOverhead(T_k)$ is typically \emph{piecewise decreasing with jumps}.
\end{rem}

\begin{algorithm}[htbp]
\caption{Compute scaffold overhead for top-$k$ sets}
\label{alg:overhead-topk}
\begin{algorithmic}[1]
\Require Library $\mathcal{L}$ with scores $F(m)$ and scaffold map $\Scaf(m)$
\Require A list of $k$ values: $k_1<k_2<\cdots<k_r$
\State Sort molecules $m_1,\dots,m_{|\mathcal{L}|}$ by increasing score $F$
\State Precompute fiber sizes $|\Fib_{\mathcal{L}}(s)|$ for each scaffold $s\in S_{\mathcal{L}}$
\For{$\ell=1$ to $r$}
    \State $T \gets \{m_1,\dots,m_{k_\ell}\}$
    \State $B \gets \{\Scaf(m): m\in T\}$ \Comment{$B=\alpha_{\mathcal{L}}(T)$}
    \State $C \gets \sum_{s\in B} |\Fib_{\mathcal{L}}(s)|$ \Comment{$C=|c_{\mathcal{L}}(T)|$}
    \State $\mathrm{overhead}(k_\ell)\gets C/k_\ell$
\EndFor
\State \Return $\{(k_\ell,\mathrm{overhead}(k_\ell))\}_{\ell=1}^r$
\end{algorithmic}
\end{algorithm}

\begin{figure}[htbp]
\centering
\begin{tikzpicture}[scale=0.95]
% axes
\draw[->] (0,0) -- (7.0,0) node[below] {$k$};
\draw[->] (0,0) -- (0,4.2) node[left] {$\ScafOverhead(T_k)$};
% curve (schematic)
\draw (0.6,3.8) .. controls (2.2,3.2) and (3.0,2.6) .. (3.6,2.2);
\draw (3.6,2.2) -- (3.6,2.9); % jump
\draw (3.6,2.9) .. controls (4.6,2.6) and (5.4,2.0) .. (6.6,1.4);
\node at (4.9,3.6) {\small piecewise decrease};
\node at (3.95,2.65) {\small jump};
\end{tikzpicture}
\caption{Schematic shape of the overhead curve $k\mapsto \ScafOverhead(T_k)$ in a regime where top hits concentrate within a few scaffold families.}
\label{fig:overhead-schematic}
\end{figure}

\section*{Exercises}
\begin{exer}[Fixed points of closure]
\label{exer:fixed-points-closure}
Let $c=\gamma\circ\alpha$ be the scaffold closure on $\mathcal{P}(\Mol)$.
\begin{enumerate}
\item Prove directly (without using Proposition~\ref{prop:closed-iff-saturated}) that $c(A)=A$ iff $A$ is a union of scaffold fibers.
\item Show that the collection of closed sets $\{A: c(A)=A\}$ is closed under arbitrary unions and intersections.
\end{enumerate}
\end{exer}

\begin{exer}[Galois connections from functions]
\label{exer:galois-from-functions}
Let $f:X\to Y$ be any function between sets, and define $\alpha_f(A)=f(A)$ and $\gamma_f(B)=f^{-1}(B)$ on powersets.
\begin{enumerate}
\item Prove that $\alpha_f \dashv \gamma_f$ always holds.
\item Characterize when the adjunction is an insertion, i.e.\ $\alpha_f\circ \gamma_f=\mathrm{id}$.
\end{enumerate}
\end{exer}

\begin{exer}[Overhead bounds]
\label{exer:overhead-bounds}
Let $\mathcal{L}$ be finite and $A\subseteq \mathcal{L}$.
Write $a_s=|A\cap \Fib_{\mathcal{L}}(s)|$.
\begin{enumerate}
\item Show that $\ScafOverhead(A)\le \max_{s:\,a_s>0}\frac{|\Fib_{\mathcal{L}}(s)|}{a_s}$.
\item Give a condition on the $a_s$ under which $\ScafOverhead(A)$ is within a factor of two of 1.
\item Construct an example where $\ScafOverhead(A)$ grows linearly with $|S_{\mathcal{L}}|$.
\end{enumerate}
\end{exer}

\begin{exer}[Downsets vs arbitrary scaffold sets]
\label{exer:downsets-vs-arbitrary}
Let $\alpha_{\downarrow}(A)={\downarrow}\alpha(A)$ be the downset abstraction.
\begin{enumerate}
\item Show that for any $A$, $c(A)\subseteq \gamma(\alpha_{\downarrow}(A))$.
\item Give an example where the inclusion is strict.
\item Interpret the strictness chemically (what extra molecules are being forced in, and why?).
\end{enumerate}
\end{exer}
% END SOURCE: chapters/ch06_galois.tex


\chapter{The Variance Spectrum and Intrinsic Depth}
\label{ch:variance}
% BEGIN SOURCE: chapters/ch07_variance.tex
% Chapter 7: The Variance Spectrum and Intrinsic Depth

\begin{quote}
\textit{This chapter contains the mathematical heart of the book: the orthogonal decomposition of scoring functions along refinement towers, and the intrinsic limits this places on scaffold\index{scaffold}-based methods.}
\end{quote}

\section{From Structure to Information}

The preceding chapters established the algebraic scaffolding---literally and figuratively---for reasoning about chemical space\index{chemical space}. Chapter~\ref{ch:lattice} showed that the scaffold poset completes to a lattice\index{lattice} $\mathcal{D}(S)$ of search regions; Chapter~\ref{ch:galois} connected molecule\index{molecule} sets to scaffold regions via a Galois connection\index{Galois connection}, characterizing the enlargement of a target set when outputs must contain whole scaffold fibers. But these results are \emph{structural}: they tell us about containment and closure, not about the \emph{information} carried by scaffold\index{scaffold} abstractions.

This chapter pivots from algebra to analysis. We ask: given a scoring function $F$ (a docking score, a binding affinity\index{binding affinity}, an ADMET property), how much of $F$ can be predicted from scaffold identity alone? The answer involves conditional expectation\index{conditional expectation}, variance decomposition, and the theory of martingales. These tools yield the \textbf{variance spectrum}\index{variance spectrum}---the signature diagnostic that reveals where information lives in the scaffold hierarchy.

The decomposition quantifies how much variance a chosen abstraction captures under a fixed molecular distribution. That amount can be zero. Its usefulness for a decision must be assessed separately.

\begin{framed}
\noindent\textbf{For the medicinal chemist:} This chapter answers a practical question. Suppose your docking campaign scores a million molecules. If you group them by scaffold, how much of the variation in scores is explained by scaffold identity versus the decorations? The variance spectrum measures these two contributions under the chosen sampling law. Choosing a search policy also requires a decision objective, evaluation costs, and evidence about the relevant tails or molecular contrasts.
\end{framed}

%------------------------------------------------------------------
\section{Measure-Theoretic Setup}
\label{sec:ch7-setup}

We work in the probability space $(\Mol, \mathcal{F}, \mu)$ where $\Mol$ is the set of molecules, $\mathcal{F}$ is an appropriate $\sigma$-algebra\index{$\sigma$-algebra@$\sigma$-algebra} (typically generated by molecular descriptors), and $\mu$ is a probability measure representing the distribution of molecules we care about---perhaps uniform over a screening\index{screening} library, or weighted by synthetic accessibility, or the empirical distribution of a particular compound collection.

Let $F: \Mol \to \R$ be a square-integrable \textbf{scoring function}\index{scoring function}: $F \in L^2(\mu)$. Think of $F$ as:
\begin{itemize}
    \item A docking score\index{docking score} (lower is better for binding)
    \item A predicted binding affinity\index{binding affinity} ($\Delta G$ or $pK_i$)
    \item An ADMET property (solubility, permeability, metabolic stability)
    \item Any real-valued molecular property we wish to predict or optimize
\end{itemize}

The $L^2$ assumption is mild: it requires $\E[F^2] < \infty$, which holds for finite-valued scores on a finite library. On a larger population, finite second moment must be assumed or established for the specified endpoint; a typical observed docking range is not a population bound.

\begin{defn}[Scaffold $\sigma$-algebra\index{$\sigma$-algebra@$\sigma$-algebra}]
\label{defn:scaffold-sigma}
Let $\Scaf: \Mol \to S$ be a scaffold map with countable image. The \emph{scaffold $\sigma$-algebra} is
\[
\cS = \sigma(\Scaf) = \{\Scaf^{-1}(B) : B \subseteq S\},
\]
the coarsest $\sigma$-algebra making $\Scaf$ measurable.
\end{defn}

Intuitively, $\cS$ represents the information available from knowing only the scaffold of a molecule\index{molecule}. A function $g: \Mol \to \R$ is $\cS$-measurable if and only if it depends on molecules only through their scaffolds: $g(m) = h(\Scaf(m))$ for some $h: S \to \R$.

\begin{exmp}[Scaffold classes as partition]
\label{ex:ch7-partition}
The scaffold map partitions $\Mol$ into fibers $\Fib(s) = \Scaf^{-1}(s)$. The $\sigma$-algebra $\cS$ consists of all unions of fibers. Knowing which event in $\cS$ occurred tells you the scaffold but nothing about the decorations.
\end{exmp}

\subsection{Refinement Towers}

A single scaffold abstraction may be too coarse. We often have access to a \emph{hierarchy} of scaffold representations, from very coarse (e.g., ``contains aromatic ring'') to very fine (e.g., complete Bemis-Murcko\index{Bemis--Murcko scaffold} scaffold with stereochemistry).

\begin{defn}[Refinement tower]
\label{defn:refinement-tower}
A \emph{refinement tower} is a sequence of scaffold maps
\[
\Scaf_0, \Scaf_1, \ldots, \Scaf_L
\]
such that $\Scaf_{\ell+1}$ \emph{refines} $\Scaf_\ell$: knowing $\Scaf_{\ell+1}(m)$ determines $\Scaf_\ell(m)$. Equivalently, the associated $\sigma$-algebras satisfy
\[
\cS_0 \subseteq \cS_1 \subseteq \cdots \subseteq \cS_L.
\]
\end{defn}

\begin{exmp}[Bemis-Murcko\index{Bemis--Murcko scaffold} hierarchy]
\label{ex:bm-hierarchy}
A candidate refinement tower can use the following descriptors. Define every level as the cumulative tuple of descriptors through that level, with a fixed ring-perception and canonicalization convention:
\begin{itemize}
    \item $\Scaf_0$: Ring system\index{ring system} count (``monocyclic,'' ``bicyclic,'' etc.)
    \item $\Scaf_1$: Ring connectivity graph (which rings share atoms/bonds)
    \item $\Scaf_2$: Ring sizes and aromaticity
    \item $\Scaf_3$: Heteroatom positions (e.g., pyridine vs.\ benzene)
    \item $\Scaf_4$: Full Bemis-Murcko scaffold with linkers
\end{itemize}
The cumulative construction makes each earlier tuple a projection of the next, so the sigma-algebras are nested. A list of independently computed descriptor names alone does not establish this nesting. If a proposed terminal representation cannot recover an earlier descriptor, retain that descriptor explicitly in the terminal tuple.
\end{exmp}

\begin{framed}
\noindent\textbf{Pharmaceutical interpretation:} A refinement tower mirrors how medicinal chemists think about scaffold families. At the coarsest level: ``we're looking at bicyclic compounds.'' More refined: ``specifically, fused 6-6 systems.'' More refined still: ``quinoline-type, with nitrogen at position 1.'' The finest level specifies the exact scaffold. Different levels of the hierarchy capture different amounts of SAR information.
\end{framed}

%------------------------------------------------------------------
\section{The Canonical Signal-Noise Decomposition}
\label{sec:canonical-decomp}

The central construction of this chapter is the decomposition of any scoring function\index{scoring function} into its scaffold-predictable part (the ``signal'') and the residual (the ``noise'').

\begin{defn}[Scaffold signal and decoration\index{decoration} noise]
\label{defn:signal-noise}
For a scoring function $F \in L^2(\mu)$ and scaffold map $\Scaf$, define:
\begin{enumerate}
    \item The \emph{scaffold signal} (or \emph{scaffold activity}):
    \[
    A(s) := \E[F \mid \Scaf = s]
    \]
    This is the expected score among all molecules sharing scaffold $s$.

    \item The \emph{decoration\index{decoration} noise} (or \emph{residual}):
    \[
    \varepsilon(m) := F(m) - A(\Scaf(m))
    \]
    This is the deviation of a molecule's score from its scaffold's average.
\end{enumerate}
\end{defn}

The scaffold signal $A: S \to \R$ is a function on scaffolds; when composed with the scaffold map, $A \circ \Scaf: \Mol \to \R$ gives the best scaffold-based prediction of $F$. The decoration noise $\varepsilon$ captures everything that scaffold identity cannot explain.

\begin{thm}[Canonical decomposition]
\label{thm:canonical-decomp}
Every $F \in L^2(\mu)$ admits a decomposition unique up to $\mu$-null sets
\[
F = A \circ \Scaf + \varepsilon
\]
where:
\begin{enumerate}
    \item $A \circ \Scaf$ is $\cS$-measurable (depends only on scaffold)
    \item $\E[\varepsilon \mid \Scaf] = 0$ almost surely (noise averages to zero within each scaffold class)
    \item $A \circ \Scaf$ and $\varepsilon$ are orthogonal in $L^2$: $\langle A \circ \Scaf, \varepsilon \rangle = 0$
\end{enumerate}
\end{thm}

\begin{proof}
This is the fundamental property of conditional expectation in $L^2$. The conditional expectation $\E[F \mid \cS]$ is the orthogonal projection of $F$ onto the closed subspace $L^2(\cS) \subset L^2(\mu)$ of $\cS$-measurable functions. Since $\Scaf$ generates $\cS$, any $\cS$-measurable function has the form $h \circ \Scaf$, and the projection is $A \circ \Scaf$ where $A(s) = \E[F \mid \Scaf = s]$.

The residual $\varepsilon = F - \E[F \mid \cS]$ lies in the orthogonal complement $(L^2(\cS))^\perp$, which consists precisely of functions $g$ with $\E[g \mid \cS] = 0$. Orthogonality follows from the projection theorem.
\end{proof}

\begin{exmp}[Docking scores for the quinoline family]
\label{ex:quinoline-decomp}
Consider docking scores $F$ for SARS-CoV-2 Mpro across molecules containing the quinoline scaffold. Let $s_Q$ denote quinoline. The scaffold signal $A(s_Q) = \E[F \mid \Scaf = s_Q]$ might be $-8.2$ kcal/mol---the average docking score\index{docking score} for quinoline-containing compounds.

Individual molecules deviate from this average:
\begin{itemize}
    \item A 4-aminoquinoline derivative might score $-9.5$ kcal/mol, giving $\varepsilon = -1.3$ kcal/mol (better than scaffold average)
    \item A 2-methylquinoline might score $-7.1$ kcal/mol, giving $\varepsilon = +1.1$ kcal/mol (worse than average)
\end{itemize}

The key insight: $\varepsilon$ averages to zero over all quinoline derivatives. The scaffold tells you the baseline; decorations move you up or down from there.
\end{exmp}

\subsection{Variance Decomposition}

The canonical decomposition leads directly to a partition of variance.

\begin{prop}[Variance decomposition]
\label{prop:var-decomp}
For $F = A \circ \Scaf + \varepsilon$:
\[
\Var(F) = \Var(A \circ \Scaf) + \Var(\varepsilon)
\]
or equivalently:
\[
\Var(F) = \underbrace{\Var(\E[F \mid \Scaf])}_{\text{between-scaffold variance}} + \underbrace{\E[\Var(F \mid \Scaf)]}_{\text{within-scaffold variance}}
\]
This is the law of total variance.
\end{prop}

\begin{proof}
By orthogonality of $A \circ \Scaf$ and $\varepsilon$:
\[
\Var(F) = \|F - \E[F]\|^2 = \|(A \circ \Scaf - \E[A \circ \Scaf]) + \varepsilon\|^2 = \|A \circ \Scaf - \E[F]\|^2 + \|\varepsilon\|^2
\]
since $\E[A \circ \Scaf] = \E[F]$. The two terms are precisely $\Var(A \circ \Scaf)$ and $\Var(\varepsilon) = \E[\Var(F \mid \Scaf)]$.
\end{proof}

\begin{framed}
\noindent\textbf{The medicinal chemist's interpretation:} Total variability in scores comes from two sources:
\begin{enumerate}
    \item \emph{Between-scaffold variance}: Different scaffolds have different average activities. Quinolines average $-8.2$, benzimidazoles average $-7.5$, etc.
    \item \emph{Within-scaffold variance}: Within each scaffold family, different decorations cause spread around the scaffold average.
\end{enumerate}
These are contributions to population squared-error variation. Neither contribution alone establishes which search operation has the greatest expected value for a particular decision.
\end{framed}

%------------------------------------------------------------------
\section{The Irreducibility Theorem}
\label{sec:irreducibility}

The within-scaffold variance represents an irreducible floor: no scaffold-based predictor can do better.

\begin{defn}[Irreducible scaffold error]
\label{defn:irred-error}
The \emph{irreducible scaffold error} for scoring function $F$ under scaffold map $\Scaf$ is:
\[
\Err(\Scaf) := \E[\Var(F \mid \Scaf)] = \|\varepsilon\|_{L^2}^2 = \Var(F) - \Var(A \circ \Scaf)
\]
\end{defn}

The term ``irreducible'' is justified by the following fundamental result.

\begin{thm}[Irreducibility]
\label{thm:irreducibility}
For any function $g: S \to \R$ (any scaffold-level predictor):
\[
\E\left[(F - g \circ \Scaf)^2\right] \geq \Err(\Scaf)
\]
with equality if and only if $g = A$ almost everywhere (with respect to the pushforward measure\index{pushforward measure} on $S$).
\end{thm}

\begin{proof}
Using the tower property of conditional expectation:
\[
\E[(F - g \circ \Scaf)^2] = \E\left[\E[(F - g(s))^2 \mid \Scaf = s]\right]
\]
For fixed scaffold $s$, the inner expectation is
\[
\E[(F - g(s))^2 \mid \Scaf = s] = \Var(F \mid \Scaf = s) + (A(s) - g(s))^2
\]
by the bias-variance decomposition (the optimal constant predictor for a random variable is its mean).

The term $(A(s) - g(s))^2 \geq 0$ with equality iff $g(s) = A(s)$. Taking expectations over $s$:
\[
\E[(F - g \circ \Scaf)^2] = \underbrace{\E[\Var(F \mid \Scaf)]}_{= \Err(\Scaf)} + \underbrace{\E[(A(s) - g(s))^2]}_{\geq 0}
\]
The minimum is achieved when $g=A$ almost everywhere under the scaffold pushforward measure.
\end{proof}

\begin{rem}[The theorem has teeth]
\label{rem:teeth}
Theorem~\ref{thm:irreducibility} establishes a \emph{fundamental limit}: no matter how sophisticated your machine learning model---random forests, gradient boosting, graph neural networks, transformers---if it predicts scores using only scaffold information, it cannot achieve mean squared error below $\Err(\Scaf)$.

This is not a statement about current algorithms; it's a mathematical impossibility. To break the floor, you must use information beyond scaffold identity.
\end{rem}

\begin{exmp}[Quantifying the limit]
\label{ex:quantifying-limit}
Suppose docking scores have total variance $\Var(F) = 4.0$ (kcal/mol)$^2$, and the irreducible error\index{irreducible error} is $\Err(\Scaf) = 1.5$ (kcal/mol)$^2$. Then:
\begin{itemize}
    \item A perfect scaffold-level predictor achieves RMSE $= \sqrt{1.5} \approx 1.22$ kcal/mol
    \item No scaffold-based method can do better
    \item The scaffold signal explains $\frac{4.0 - 1.5}{4.0} = 62.5\%$ of total variance
\end{itemize}
\end{exmp}

%------------------------------------------------------------------
\section{Monotonicity Under Refinement}
\label{sec:monotonicity}

Intuitively, finer scaffold representations should capture more information and reduce the irreducible error\index{irreducible error}. This intuition is correct and is captured by a monotonicity theorem.

\begin{thm}[Monotonicity]
\label{thm:monotonicity}
If $\Scaf_2$ refines $\Scaf_1$ (i.e., $\cS_1 \subseteq \cS_2$), then:
\[
\Err(\Scaf_2) \leq \Err(\Scaf_1)
\]
Finer scaffold maps have smaller (or equal) irreducible error.
\end{thm}

\begin{proof}
By the conditional variance decomposition (a consequence of the tower property):
\[
\Var(F \mid \cS_1) = \E[\Var(F \mid \cS_2) \mid \cS_1] + \Var(\E[F \mid \cS_2] \mid \cS_1)
\]
Taking expectations over $\cS_1$:
\[
\Err(\Scaf_1) = \E[\Var(F \mid \cS_1)] = \E[\Var(F \mid \cS_2)] + \E[\Var(A_2 \mid \cS_1)]
\]
where $A_2 = \E[F \mid \cS_2]$. Since variance is non-negative:
\[
\Err(\Scaf_1) = \Err(\Scaf_2) + \underbrace{\E[\Var(A_2 \mid \cS_1)]}_{\geq 0}
\]
Therefore $\Err(\Scaf_2) \leq \Err(\Scaf_1)$.
\end{proof}

The difference between errors has a name and an interpretation.

\begin{defn}[Refinement gain]
\label{defn:refinement-gain}
The \emph{refinement gain} from $\Scaf_1$ to $\Scaf_2$ (where $\Scaf_2$ refines $\Scaf_1$) is:
\[
\Delta(\Scaf_1 \to \Scaf_2) := \Err(\Scaf_1) - \Err(\Scaf_2) = \E[\Var(A_2 \mid \Scaf_1)] \geq 0
\]
\end{defn}

\begin{framed}
\noindent\textbf{Pharmaceutical interpretation:} The refinement gain measures how much additional predictive power you get from distinguishing scaffolds at the finer level.

Example: Suppose $\Scaf_1$ groups all quinoline derivatives together, while $\Scaf_2$ distinguishes 4-substituted from 8-substituted quinolines. The refinement gain $\Delta(\Scaf_1 \to \Scaf_2)$ measures how much the average activities of 4- and 8-substituted quinolines differ---if they're similar, little is gained by the distinction.
\end{framed}

\begin{exmp}[Diminishing returns]
\label{ex:diminishing}
Consider the hierarchy from Example~\ref{ex:bm-hierarchy}:
\begin{center}
\begin{tabular}{lcc}
\toprule
Level & Example distinction & Typical $\Delta$ \\
\midrule
$0 \to 1$ & Bicyclic fused vs.\ spiro & Large \\
$1 \to 2$ & 6-6 vs.\ 6-5 ring system\index{ring system} & Moderate \\
$2 \to 3$ & Quinoline vs.\ isoquinoline & Small to moderate \\
$3 \to 4$ & Different linker lengths & Small \\
\bottomrule
\end{tabular}
\end{center}
Typically, early refinements capture large differences; later refinements yield diminishing returns. This pattern is not universal---some targets may have large gains at fine levels---but it is common.
\end{exmp}

%------------------------------------------------------------------
\section{The Orthogonal Decomposition Along Towers}
\label{sec:orthogonal}

The real power of the refinement framework emerges when we consider an entire tower simultaneously. The signal at each level can be decomposed into orthogonal ``innovations'' that capture what each refinement step contributes.

\begin{defn}[Conditional expectations along a tower]
For a refinement tower $\cS_0 \subseteq \cS_1 \subseteq \cdots \subseteq \cS_L$, define:
\[
A_\ell := \E[F \mid \cS_\ell]
\]
the conditional expectation of $F$ given information at level $\ell$. By convention, $A_{-1} := \E[F]$ (the global mean, which is $\cS_{-1}$-measurable for the trivial $\sigma$-algebra).
\end{defn}

\begin{defn}[Martingale\index{martingale} differences]
\label{defn:martingale-diff}
The \emph{martingale\index{martingale} differences} (or \emph{innovations}) are:
\[
D_\ell := A_\ell - A_{\ell-1}, \quad \ell = 0, 1, \ldots, L
\]
These represent the ``new information'' gained at level $\ell$: the adjustment to the predicted score when we learn the finer scaffold.
\end{defn}

The term ``martingale difference'' comes from probability theory: the sequence $(A_\ell)$ forms a martingale with respect to the filtration $(\cS_\ell)$, and the differences are the increments.

\begin{thm}[Orthogonal decomposition]
\label{thm:orthogonal-decomp}
The martingale differences $\{D_\ell\}_{\ell=0}^L$ are pairwise orthogonal in $L^2$:
\[
\langle D_j, D_k \rangle_{L^2} = 0 \quad \text{for } j \neq k
\]
Moreover, the scoring function decomposes as:
\[
F = \E[F] + \sum_{\ell=0}^{L} D_\ell + \varepsilon_L
\]
where $\varepsilon_L = F - A_L$ is the residual at the finest level. The variance decomposes as:
\[
\Var(F) = \sum_{\ell=0}^{L} \|D_\ell\|^2 + \Err(\Scaf_L)
\]
\end{thm}

\begin{proof}
For $j < k$, note that $D_j = A_j - A_{j-1}$ is $\cS_j$-measurable (both $A_j$ and $A_{j-1}$ are). Also, $D_k = A_k - A_{k-1}$ satisfies:
\[
\E[D_k \mid \cS_j] = \E[A_k \mid \cS_j] - \E[A_{k-1} \mid \cS_j] = A_j - A_j = 0
\]
using the tower property for $\cS_j\subseteq\cS_{k-1}\subseteq\cS_k$. Measurability with respect to the finer sigma-algebra does not imply measurability with respect to the coarser one. Therefore:
\[
\langle D_j, D_k \rangle = \E[D_j D_k] = \E[\E[D_j D_k \mid \cS_j]] = \E[D_j \E[D_k \mid \cS_j]] = \E[D_j \cdot 0] = 0
\]

The decomposition $F = A_{-1} + \sum_{\ell=0}^L D_\ell + \varepsilon_L$ is telescoping:
\[
A_{-1} + \sum_{\ell=0}^L (A_\ell - A_{\ell-1}) + (F - A_L) = A_{-1} + (A_L - A_{-1}) + F - A_L = F
\]

Since all terms are orthogonal, the variance decomposes additively.
\end{proof}

\begin{exmp}[Decomposing quinoline scores]
\label{ex:quinoline-tower}
Consider a 3-level tower for quinoline-type compounds:
\begin{itemize}
    \item Level 0: ``Bicyclic aromatic'' ($A_0 = -7.8$ kcal/mol)
    \item Level 1: ``Fused 6-6 with one nitrogen'' ($A_1 = -8.0$ kcal/mol)
    \item Level 2: ``Quinoline (N at position 1)'' vs.\ ``Isoquinoline (N at position 2)''
\end{itemize}

For quinoline ($A_2(Q) = -8.2$) and isoquinoline ($A_2(I) = -7.7$):
\begin{itemize}
    \item $D_0 = A_0 - \E[F]$: deviation of bicyclic aromatics from global mean
    \item $D_1 = A_1 - A_0 = -0.2$: adjustment for being a 6-6 nitrogen heterocycle
    \item $D_2(Q) = A_2(Q) - A_1 = -0.2$: further adjustment for quinoline specifically
    \item $D_2(I) = A_2(I) - A_1 = +0.3$: adjustment for isoquinoline
\end{itemize}

For these illustrative values to be conditional means of a two-child parent, the within-parent frequencies must be $0.6$ for quinoline and $0.4$ for isoquinoline. Those weights give $0.6(-0.2)+0.4(0.3)=0$. Frequencies cannot be chosen independently of the stated parent mean.
\end{exmp}

%------------------------------------------------------------------
\section{The Variance Spectrum}
\label{sec:variance-spectrum}

The orthogonal decomposition naturally leads to a ``spectrum'' that characterizes where variance lives in the hierarchy.

\begin{defn}[Variance spectrum]
\label{defn:variance-spectrum}
The \emph{variance spectrum}\index{variance spectrum} of scoring function $F$ with respect to a refinement tower $(\cS_\ell)_{\ell=0}^L$ is the vector:
\[
\sigma(F) = \left(\|D_0\|^2, \|D_1\|^2, \ldots, \|D_L\|^2, \Err(\Scaf_L)\right) \in \R_{\geq 0}^{L+2}
\]
The components sum to $\Var(F)$.
\end{defn}

\begin{defn}[Normalized variance spectrum]
\label{defn:normalized-spectrum}
When $\Var(F)>0$, the \emph{normalized variance spectrum} is:
\[
\hat{\sigma}_\ell = \frac{\|D_\ell\|^2}{\Var(F)}, \quad \hat{\sigma}_{\text{res}} = \frac{\Err(\Scaf_L)}{\Var(F)}
\]
These are nonnegative and sum to one. If $\Var(F)=0$, all unnormalized components are zero and these normalized ratios are undefined; report the constant-response case separately.
\end{defn}

The normalized spectrum enables comparison across different targets and scoring functions. A spectrum concentrated at low levels indicates that coarse scaffold information captures most predictable variance; a flat or increasing spectrum indicates that fine details matter.

\begin{framed}
\noindent\textbf{The signature plot:} The bar chart of $\hat{\sigma}_\ell$ versus level $\ell$ is the \textbf{signature plot} of this book. It answers, at a glance:
\begin{itemize}
    \item Where does predictive information live?
    \item At what level do refinements stop paying off?
    \item How much is fundamentally unpredictable from scaffold structure?
\end{itemize}
Computing a population spectrum can itself require many response evaluations. A sample spectrum is a diagnostic whose estimation uncertainty and acquisition cost must be reported.
\end{framed}

\begin{defn}[Intrinsic scaffold depth]
\label{defn:intrinsic-depth}
Fix a terminal tower level $L_{\max}$ and assume $\Var(F)>0$. For $\epsilon\in(0,1)$, define
\[
L^*(\epsilon)=\min\{k\in\{0,\ldots,L_{\max}\}:
\Err(\Scaf_k)<\epsilon\Var(F)\},
\]
with $L^*(\epsilon)=\infty$ if the set is empty. Equivalently, replace $\Err(\Scaf_k)$ by
\[
\sum_{\ell=k+1}^{L_{\max}}\|D_\ell\|^2+
\Err(\Scaf_{L_{\max}}).
\]
The terminal residual is counted once. For a constant response, the unnormalized error is zero at every level; no variance-normalized depth is defined.
\end{defn}

\begin{exmp}[Interpreting depth]
\label{ex:depth}
A value $L^*(0.1)=2$ means levels through two explain more than $90\%$ of the total response variance under the chosen measure. It does not establish top-hit recovery or an efficient sampling algorithm. A value of infinity means the stated tower does not reduce the residual below $10\%$; it can still explain a substantial fraction of variance. The depth depends on the response, molecular measure, tower, and tolerance.
\end{exmp}

%------------------------------------------------------------------
\section{Properties of the Variance Spectrum}
\label{sec:spectrum-properties}

The variance spectrum has several useful properties that follow from its construction.

\begin{prop}[Non-negativity and normalization]
\label{prop:spectrum-basic}
For $F\in L^2(\mu)$ and a refinement tower (assuming $\Var(F)>0$ for the normalized statement):
\begin{enumerate}
    \item Each component $\|D_\ell\|^2 \geq 0$ and $\Err(\Scaf_L) \geq 0$
    \item $\sum_{\ell=0}^L \|D_\ell\|^2 + \Err(\Scaf_L) = \Var(F)$
    \item $\sum_{\ell=0}^L \hat{\sigma}_\ell + \hat{\sigma}_{\mathrm{res}} = 1$
\end{enumerate}
\end{prop}

\begin{prop}[Explained variance is cumulative]
\label{prop:cumulative}
For $\Var(F)>0$, define the \emph{cumulative explained variance} at level $k$:
\[
R^2_k := \frac{\sum_{\ell=0}^k \|D_\ell\|^2}{\Var(F)} = 1 - \frac{\Err(\Scaf_k)}{\Var(F)}
\]
Then $R^2_k$ is non-decreasing in $k$, with $R^2_k \in [0, 1]$.
\end{prop}

\begin{proof}
By monotonicity (Theorem~\ref{thm:monotonicity}), $\Err(\Scaf_k)$ is non-increasing in $k$, so $R^2_k$ is non-decreasing.
\end{proof}

The quantity $R^2_k$ is the analog of the coefficient of determination in regression: it measures what fraction of total variance is explained by scaffold information up to level $k$.

\begin{prop}[Spectrum determines irreducible errors]
\label{prop:spectrum-errors}
The irreducible error at any level $k$ is recoverable from the spectrum:
\[
\Err(\Scaf_k) = \sum_{\ell > k} \|D_\ell\|^2 + \Err(\Scaf_L) = \Var(F) \cdot \left( \sum_{\ell > k} \hat{\sigma}_\ell + \hat{\sigma}_{\mathrm{res}} \right)
\]
\end{prop}

\begin{thm}[Spectrum invariance under tower extension]
\label{thm:invariance}
If $(\cS_\ell)_{\ell=0}^L$ is a refinement tower and we extend it to $(\cS_\ell)_{\ell=0}^{L'}$ with $L' > L$, then:
\begin{enumerate}
    \item The first $L+1$ components $(\|D_0\|^2, \ldots, \|D_L\|^2)$ are unchanged
    \item The residual $\Err(\Scaf_L)$ is split into $\sum_{\ell=L+1}^{L'} \|D_\ell\|^2 + \Err(\Scaf_{L'})$
\end{enumerate}
The spectrum ``extends'' rather than changes; computed components are stable.
\end{thm}

\begin{proof}
The martingale differences $D_0, \ldots, D_L$ depend only on $A_0, \ldots, A_L$, which are determined by $\cS_0, \ldots, \cS_L$. The new levels $L+1, \ldots, L'$ refine $\cS_L$, so they split the residual but don't change the earlier terms.
\end{proof}

%------------------------------------------------------------------
\section{Computational Considerations}
\label{sec:computation}

Computing the variance spectrum requires estimating conditional expectations---a statistical task with both theoretical and practical aspects.

\subsection{Exact Computation (Finite Populations)}

For a finite molecule library $\Mol = \{m_1, \ldots, m_N\}$ with uniform measure $\mu$, the conditional expectations are simply averages over fibers:
\[
A(s) = \frac{1}{|\Fib(s)|} \sum_{m \in \Fib(s)} F(m)
\]

The spectrum components are then:
\begin{align}
\|D_\ell\|^2 &= \frac{1}{N} \sum_{m} (A_\ell(\Scaf_\ell(m)) - A_{\ell-1}(\Scaf_{\ell-1}(m)))^2 \\
\Err(\Scaf_L) &= \frac{1}{N} \sum_m (F(m) - A_L(\Scaf_L(m)))^2
\end{align}

For $L+1$ levels, the sums require $O(N(L+1))$ work given the scaffold assignments and scores; this is linear in $N$ for a fixed-depth tower. Computing the assignments and acquiring the scores have separate costs.

\subsection{Estimation from Samples}

When the full library is too large or the scoring function is expensive (e.g., docking), we estimate from a sample $\{(m_i, F(m_i))\}_{i=1}^n$.

\begin{prop}[Unbiased estimation under fixed within-class sampling]
\label{prop:estimation}
For $n_s\geq1$ independent samples from the conditional law
$\mu(\cdot\mid\Scaf=s)$, the mean
$\hat A(s)=n_s^{-1}\sum_{i:s_i=s}F(m_i)$ is unbiased for $A(s)$.
The same statement holds for a fixed-size uniform sample without replacement
from a finite fiber. It does not assert unbiasedness after response-dependent
stopping or selection. Plug-in squared means and variance components generally
require separate bias corrections.
\end{prop}

For reliable estimates, each scaffold class should have multiple representatives. A sample fiber with one observation has zero empirical residual variance, which need not equal its population variance. The uncorrected algorithm below computes the exact spectrum of the empirical distribution; it is not an unbiased estimator of every population component.

\begin{algor}[Variance spectrum estimation]
\label{alg:spectrum}
~
\begin{enumerate}
    \item \textbf{Input:} Scored molecules $\{(m_i, F_i)\}_{i=1}^n$, refinement tower extractors $(\Scaf_\ell)_{\ell=0}^L$
    \item \textbf{Compute scaffolds:} For each molecule and level, compute $s_{i,\ell} = \Scaf_\ell(m_i)$
    \item \textbf{Estimate conditional expectations:} For each level $\ell$ and scaffold $s$ at that level:
    \[
    \hat{A}_\ell(s) = \frac{\sum_{i: s_{i,\ell} = s} F_i}{\sum_{i: s_{i,\ell} = s} 1}
    \]
    \item \textbf{Compute differences:} Set $\hat A_{-1}$ to the sample mean, then $\hat{D}_{i,\ell} = \hat{A}_\ell(s_{i,\ell}) - \hat{A}_{\ell-1}(s_{i,\ell-1})$
    \item \textbf{Estimate spectrum:} $\widehat{\|D_\ell\|^2} = \frac{1}{n} \sum_i \hat{D}_{i,\ell}^2$
    \item \textbf{Output:} Spectrum estimate $(\widehat{\|D_0\|^2}, \ldots, \widehat{\|D_L\|^2}, \widehat{\Err})$
\end{enumerate}
\end{algor}

\begin{rem}[Bootstrap confidence intervals]
A bootstrap can assess sampling variability under a specified sampling design. Resampling independent molecules is inappropriate when rows repeat a parent, share assay effects, or were selected adaptively. Resample at the actual independent unit and validate coverage; a bootstrap interval is not a distribution-free guarantee.
\end{rem}

%------------------------------------------------------------------
\section{Experimental Validation}
\label{sec:experiments}

The variance spectrum is not just a theoretical construct; it is an empirical diagnostic. We describe two signature experiments.

\begin{exper}[Variance spectrum for docking]
\label{exper:variance-spectrum}
\textbf{Objective:} Compute the normalized variance spectrum for a docking campaign and visualize where predictive information resides.

\textbf{Input:}
\begin{itemize}
    \item Docking scores for SARS-CoV-2 Mpro (main protease) against a library of $10^6$ molecules
    \item Bemis-Murcko scaffold extraction plus a 4-level refinement tower:
    \begin{enumerate}
        \item[0.] Ring count
        \item[1.] Ring connectivity graph
        \item[2.] Ring sizes and aromaticity pattern
        \item[3.] Full Bemis-Murcko scaffold
    \end{enumerate}
\end{itemize}

\textbf{Computation:}
\begin{enumerate}
    \item Extract scaffolds at each level for all molecules
    \item Compute sample means $\hat{A}_\ell(s)$ for each scaffold at each level
    \item Compute martingale differences and their squared norms
    \item Normalize by total variance
\end{enumerate}

\textbf{Output:} Bar chart of $\hat{\sigma}_\ell$ versus level $\ell$, with error bars from bootstrap.

\textbf{Analysis:} Report the observed spectrum with the sampling design and uncertainty. No amount of variance capture or residual range is assumed in advance. Validate the nesting of the implemented descriptors before calling their increments a martingale decomposition.

\textbf{Interpretation:} A coarse representation with low population residual is informative for average squared-error prediction. Whether it supports molecular prioritization must be tested on held-out decisions under a matched budget; there is no universal variance-ratio cutoff for that claim.
\end{exper}

\begin{exper}[Intrinsic depth across targets]
\label{exper:intrinsic-depth}
\textbf{Objective:} Compare intrinsic scaffold depth across diverse docking targets under a shared tower, sampling measure, and tolerance.

\textbf{Input:}
\begin{itemize}
    \item Ten diverse protein targets with docking scores (e.g., kinases, proteases, GPCRs, nuclear receptors)
    \item Same library and refinement tower for each
\end{itemize}

\textbf{Computation:} For each target $t$:
\begin{enumerate}
    \item Compute the variance spectrum $\hat{\sigma}^{(t)}$
    \item Compute intrinsic depth $L^*(0.1)$: smallest $L$ such that $\sum_{\ell > L} \hat{\sigma}_\ell^{(t)} + \hat{\sigma}_{\text{res}}^{(t)} < 0.1$
\end{enumerate}

\textbf{Output:}
\begin{itemize}
    \item Table of $L^*(0.1)$ for each target
    \item Heatmap of normalized spectra across targets (targets as rows, levels as columns)
\end{itemize}

\textbf{Analysis:} Compare the measured depths under the same library measure, tower, and tolerance. Report nonqualifying targets as infinity and constant responses separately. Use independent evaluation data to test whether a depth diagnostic predicts actual decision error or cost. The experiment does not presume a typical depth or a preferred binding-site category.
\end{exper}

%------------------------------------------------------------------
\section{Connections and Implications}
\label{sec:connections}

The variance spectrum connects to several themes developed elsewhere in this book and beyond.

\subsection{Relation to Scaffold Overhead (Chapter~\ref{ch:galois})}

Chapter~\ref{ch:galois} introduced the scaffold overhead $\ScafOverhead(A) = |c(A)|/|A|$, measuring the expansion required to capture a target set via scaffold selection. The variance spectrum provides a complementary view: overhead measures \emph{set-theoretic} cost; the spectrum measures \emph{information-theoretic} content.

High residual variance ($\hat{\sigma}_{\text{res}}$ large) suggests that even capturing the right scaffolds leaves much variation unexplained---scaffold-based algorithms will be noisy predictors even with no overhead. Low residual variance controls average squared prediction error under the stated measure; it does not guarantee a correct ranking of the most extreme molecules.

\subsection{Relation to multilevel estimation}
Chapter~\ref{ch:mlmc} distinguishes the exact projections studied here from
executable random estimators. The variance spectrum gives
$\Var(\E[F\mid\cS_\ell]-\E[F\mid\cS_{\ell-1}])$. An implementation based on
estimated conditional means or different physical evaluators has its own
correction variance and bias. Multilevel savings require measured costs and a
coupling that actually reduces that correction variance; spectral decay alone
does not establish a saving.

\subsection{Relation to molecular selection}
A scaffold mean is an arm value for a mean-selection problem. It need not
identify the class containing an extreme molecule. Likewise, population
mean-squared error does not control the largest error in a library. A molecular
decision needs an additional bridge, such as uniform prediction intervals or
explicit within-fiber search. Chapters~\ref{ch:response-certificates} and
\ref{ch:screening} state those conditions.

\subsection{Universality Across Property Types}

While we've emphasized docking scores, the framework applies to any square-integrable property:
\begin{itemize}
    \item \textbf{Binding affinity (experimental):} Spectrum reveals how much affinity variation is scaffold-driven
    \item \textbf{ADMET properties:} Solubility, permeability, metabolic stability each have their own spectra
    \item \textbf{Selectivity:} Differences in affinity across targets can be decomposed
\end{itemize}

Comparing spectra across property types reveals which properties are more ``scaffold-determined'' and which require finer modeling.

%------------------------------------------------------------------
\section*{Summary}
\label{sec:ch7-summary}

This chapter establishes population squared-error identities for a declared abstraction and sampling law.

\begin{enumerate}
    \item \textbf{Canonical decomposition.} Any scoring function splits uniquely into a scaffold signal (predictable from scaffold identity) and decoration noise (unpredictable from scaffold alone).

    \item \textbf{Irreducibility theorem.} The within-scaffold variance $\Err(\Scaf)$ is a hard floor: no scaffold-based predictor---no matter how sophisticated---can achieve lower mean squared error.

    \item \textbf{Monotonicity.} Finer scaffold representations never increase irreducible error; they can only help (or stay flat).

    \item \textbf{Orthogonal decomposition.} Along a refinement tower, the variance decomposes into orthogonal contributions from each level plus the residual.

    \item \textbf{Variance spectrum.} The vector $(\|D_0\|^2, \ldots, \|D_L\|^2, \Err(\Scaf_L))$ characterizes where predictive information lives. Its normalized version enables cross-target comparison.

    \item \textbf{Intrinsic depth.} Intrinsic scaffold depth records the first qualifying level of a specified tower and tolerance; its value is also measure-dependent.
\end{enumerate}

\medskip

\noindent\textbf{For medicinal chemists:} The variance spectrum answers the question ``how much does scaffold identity matter for this target?'' Its components show the response variation captured by each declared level. Search priorities additionally depend on hit tails, decision margins, and evaluation costs, which the spectrum does not specify.

\medskip

\noindent\textbf{For modelers:} The variance spectrum is your fundamental diagnostic. Estimate it when that diagnostic justifies its response-acquisition cost; the population quantity gives a squared-error floor for the specified information set. Use it as one diagnostic when proposing refinements. Validate actual estimator costs for multilevel schemes and decision error for molecular selection. The irreducibility theorem is not a suggestion---it is a mathematical limit that no algorithm can escape.

\bigskip

\noindent This chapter provides the theoretical foundation; Part~\ref{part:algorithms} translates it into actionable algorithms.

%------------------------------------------------------------------
\section*{Exercises}

\begin{exer}[Chemistry: Interpreting the spectrum]
\label{exer:ch7-chem}
A docking campaign against Target X yields the following normalized variance spectrum:
\[
\hat{\sigma} = (0.35, 0.25, 0.15, 0.10, 0.05, 0.10)
\]
where the first five components correspond to refinement levels 0--4 and the last is the residual.

\begin{enumerate}
    \item[(a)] What is the intrinsic depth $L^*(0.1)$?
    \item[(b)] What fraction of variance is explained by the coarsest two levels?
    \item[(c)] If you could only distinguish scaffolds at one level, which would you choose and why?
    \item[(d)] What does this spectrum establish about average squared-error prediction, and what additional evidence would be needed for a top-hit recovery claim?
\end{enumerate}
\end{exer}

\begin{exer}[Mathematics: Orthogonality proof]
\label{exer:ch7-math}
Let $(\cS_\ell)_{\ell=0}^L$ be a refinement tower (filtration) and $F \in L^2$.
\begin{enumerate}
    \item[(a)] Prove that $\E[D_k \mid \cS_j] = 0$ for all $j < k$, where $D_k = A_k - A_{k-1}$.
    \item[(b)] Use (a) to prove $\langle D_j, D_k \rangle = 0$ for $j \neq k$.
    \item[(c)] Show that the decomposition $F = \E[F] + \sum_{\ell=0}^L D_\ell + \varepsilon_L$ is the unique orthogonal decomposition of $F$ with respect to the subspaces $\{L^2(\cS_\ell) \ominus L^2(\cS_{\ell-1})\}_\ell$ and $(L^2(\cS_L))^\perp$.
\end{enumerate}
\end{exer}

\begin{exer}[ML/Statistics: Estimation bias]
\label{exer:ch7-ml}
Consider estimating $\Err(\Scaf)$ from a sample where some scaffold classes have only one molecule.
\begin{enumerate}
    \item[(a)] If scaffold $s$ has only one molecule $m_s$ in the sample, what is the contribution of this scaffold to the naive plug-in estimate $\widehat{\Err} = \frac{1}{n}\sum_m (F(m) - \hat{A}(\Scaf(m)))^2$?
    \item[(b)] Explain why this leads to systematic underestimation of $\Err(\Scaf)$.
    \item[(c)] Propose a correction using Bessel's correction (using $n_s - 1$ instead of $n_s$ in within-class variance estimates). Under what conditions is this correction valid?
    \item[(d)] Suggest a bootstrap-based approach to estimate and correct for bias.
\end{enumerate}
\end{exer}

\begin{exer}[Computation: Implementing spectrum calculation]
\label{exer:ch7-comp}
Implement a function that computes the variance spectrum from scored molecules.
\begin{enumerate}
    \item[(a)] Write pseudocode (or actual code) for the algorithm in Section~\ref{sec:computation}.
    \item[(b)] Analyze the time complexity in terms of $n$ (number of molecules) and $L$ (number of levels).
    \item[(c)] How would you modify the algorithm to handle weighted molecules (non-uniform $\mu$)?
    \item[(d)] Implement bootstrap confidence intervals for the spectrum components.
\end{enumerate}
\end{exer}

\begin{exer}[Theory: Spectrum bounds]
\label{exer:ch7-theory}
\begin{enumerate}
    \item[(a)] Prove that $\Err(\Scaf) \geq 0$ with equality if and only if $F$ is $\cS$-measurable (i.e., $F = g \circ \Scaf$ for some $g$).
    \item[(b)] Prove that $\Err(\Scaf) \leq \Var(F)$ with equality if and only if $\E[F \mid \Scaf]$ is constant (scaffold provides no information).
    \item[(c)] If the tower has $L$ levels and the scaffold at level $L$ has $K$ classes, give an upper bound on $\sum_{\ell=0}^L \|D_\ell\|^2$ in terms of $K$ and $\Var(F)$.
    \item[(d)] When can equality hold in (c)?
\end{enumerate}
\end{exer}

%------------------------------------------------------------------
\section{Bibliographic Notes}

The decomposition of random variables by conditional expectation is classical in probability theory; see Williams (1991) for an accessible treatment. The martingale difference construction appears throughout the theory of stochastic processes.

The conditional-expectation identities are classical; their value here is to audit explicitly chosen chemical abstractions. Related ideas appear in the QSAR literature under ``matched molecular pair analysis'' (Leach et al., 2006), which implicitly estimates within-scaffold variance by comparing molecules differing only in decorations. Chapter~\ref{ch:galois} places the same representation alongside its set-theoretic closure; no novelty claim for the underlying algebra is needed.

The variance spectrum as a diagnostic for molecular property prediction has analogs in spatial statistics (variograms) and signal processing (power spectra). The intrinsic depth concept parallels effective dimension in quasi-Monte Carlo theory (Caflisch, 1998).

Multilevel Monte Carlo\index{multilevel Monte Carlo} methods, which exploit similar hierarchical variance decompositions, were pioneered by Heinrich (2001) and Giles (2008); their application to molecular screening is developed in Chapter~\ref{ch:mlmc}.
% END SOURCE: chapters/ch07_variance.tex


%--------Part III: The Analytic Structure--------
\part{Response Models and Decision Certificates}
\label{part:analysis}

\chapter{Probability on Scaffold Classes}
\label{ch:probability}
% BEGIN SOURCE: chapters/ch08_probability.tex
% ======================================================================
% Chapter 8 content file: chapters/ch08_probability.tex
% (Do NOT include \chapter{...} here; main.tex already does that.)
% ======================================================================

\begin{quote}
\textit{Chapter~\ref{ch:variance} established that scaffold\index{scaffold} identity captures a quantifiable fraction of score variance---the variance spectrum\index{variance spectrum} tells us how much. But variance is only the second moment. This chapter descends from moments to distributions: the full probability law of scores within each scaffold\index{scaffold} class. The distinction matters profoundly: two scaffolds with identical mean and variance can have wildly different tail behaviors, and it is the tails that determine hit-finding success.}
\end{quote}

\section{From Moments to Distributions}
\label{sec:ch8-moments-to-distributions}

The variance spectrum\index{variance spectrum} of Chapter~\ref{ch:variance} provided a powerful summary of where predictive information lives in the scaffold hierarchy. For each scaffold $s$, we characterized the conditional mean $A(s) = \E[F \mid \Scaf = s]$ and the within-scaffold variance $\sigma_s^2 = \Var(F \mid \Scaf = s)$. The irreducible error\index{irreducible error} $\Err(\Scaf) = \E_s[\sigma_s^2]$ quantified the fundamental limit of scaffold-based prediction.

But mean and variance are only the first two moments of a distribution. Two random variables can share the same mean and variance while differing dramatically in their full distributional behavior. Consider:

\begin{itemize}
    \item A Gaussian distribution $\mathcal{N}(\mu, \sigma^2)$
    \item A uniform distribution on $[\mu - \sqrt{3}\sigma, \mu + \sqrt{3}\sigma]$
    \item A mixture of two point masses at $\mu \pm \sigma$
\end{itemize}

All three have mean $\mu$ and variance $\sigma^2$, yet their tail behaviors---the probabilities of extreme values---differ qualitatively. For hit-finding in drug discovery, it is precisely the tails that matter.

\begin{framed}
\noindent\textbf{Why tails matter for screening\index{screening}.}

\medskip
\noindent
In virtual screening\index{screening}, we seek molecules with scores below a threshold $\tau$ (for docking, lower is better). The key quantity is the \emph{hit probability}:
\[
\Pr(F \leq \tau \mid \Scaf = s)
\]
This is a tail probability, and it depends on the full distribution $p_s = \mathrm{Law}(F \mid \Scaf = s)$, not just on $(\mu_s, \sigma_s^2)$.

A scaffold with moderate mean $\mu_s$ but heavy tails may yield more extreme hits than one with excellent mean but tight concentration. Conversely, a scaffold with good mean and light tails offers reliable (if less spectacular) hit-finding. The variance spectrum tells us the spread; this chapter tells us the shape.
\end{framed}

\begin{exmp}[Two scaffolds, same moments, different tails]
\label{ex:ch8-two-scaffolds}
Consider two scaffold classes with identical summary statistics:
\begin{center}
\begin{tabular}{lcc}
\toprule
& Scaffold A & Scaffold B \\
\midrule
Mean $\mu_s$ (kcal/mol) & $-7.5$ & $-7.5$ \\
Std.\ dev.\ $\sigma_s$ (kcal/mol) & $1.2$ & $1.2$ \\
\bottomrule
\end{tabular}
\end{center}

Suppose Scaffold A has a unimodal, approximately Gaussian distribution, while Scaffold B is bimodal: half of its molecules cluster around $-8.5$ kcal/mol (a subpopulation with favorable binding geometry) and half around $-6.5$ kcal/mol (a subpopulation with steric clashes).

For a threshold $\tau = -10.0$ kcal/mol:
\begin{itemize}
    \item \textbf{Scaffold A:} The Gaussian tail gives $\Pr(F \leq -10.0) \approx \Phi\left(\frac{-10.0 - (-7.5)}{1.2}\right) = \Phi(-2.08) \approx 0.019$.
    \item \textbf{Scaffold B:} The favorable subpopulation (mean $-8.5$, smaller local variance) has heavier probability mass near $-10.0$; the hit probability may be substantially higher.
\end{itemize}

Despite identical $(\mu, \sigma)$, the scaffolds have different hit-finding potential. The bimodal structure of Scaffold B represents exploitable SAR information invisible to moment-based analysis.
\end{exmp}

This chapter develops the distributional framework needed to move beyond moments. We introduce the class distribution $p_s$, establish conditions (sub-Gaussian, sub-exponential) under which tail bounds hold, address the important case of multimodal distributions, and provide diagnostic tools to test whether assumptions hold in practice.

%------------------------------------------------------------------
\section{Class Distributions and Their Parameters}
\label{sec:ch8-class-distributions}

We formalize the distributional view of scaffold classes.

\begin{defn}[Class distribution]
\label{defn:class-distribution}
For a scaffold $s \in S$ with $\mu(\Fib(s)) > 0$, the \emph{class distribution} is the conditional law
\[
p_s := \mathrm{Law}(F \mid \Scaf = s),
\]
a probability measure on $\R$ describing the distribution of scores among molecules sharing scaffold $s$.
\end{defn}

The class distribution $p_s$ is determined by the scoring function\index{scoring function} $F$, the molecule distribution $\mu$, and the scaffold map $\Scaf$. Different scoring functions induce different class distributions on the same scaffold partition.

\begin{defn}[Class mean and variance]
\label{defn:class-mean-variance}
For scaffold $s$ with class distribution $p_s$:
\begin{enumerate}
    \item The \emph{class mean} is
    \[
    \mu_s := \E[F \mid \Scaf = s] = \int_{\R} x \, dp_s(x).
    \]
    This equals the scaffold signal $A(s)$ from Chapter~\ref{ch:variance}.

    \item The \emph{class variance} is
    \[
    \sigma_s^2 := \Var(F \mid \Scaf = s) = \int_{\R} (x - \mu_s)^2 \, dp_s(x).
    \]

    \item The \emph{class standard deviation} is $\sigma_s = \sqrt{\sigma_s^2}$.
\end{enumerate}
\end{defn}

\begin{defn}[Higher moments]
\label{defn:higher-moments}
The $k$-th centered moment of class $s$ is
\[
m_k(s) := \E[(F - \mu_s)^k \mid \Scaf = s].
\]
Of particular diagnostic interest:
\begin{itemize}
    \item The \emph{skewness}: $\gamma_s := m_3(s) / \sigma_s^3$
    \item The \emph{excess kurtosis}: $\kappa_s := m_4(s) / \sigma_s^4 - 3$
\end{itemize}
For a Gaussian distribution, $\gamma = 0$ (symmetric) and $\kappa = 0$ (mesokurtic). These quantities require finite moments and positive variance. Kurtosis summarizes a fourth moment; its sign does not classify asymptotic tail decay.
\end{defn}

The irreducible error\index{irreducible error} from Chapter~\ref{ch:variance} has a simple expression in terms of class variances.

\begin{prop}[Irreducible error as expected class variance]
\label{prop:irred-error-class-var}
\[
\Err(\Scaf) = \E_s[\sigma_s^2] = \sum_{s \in S} \mu(\Fib(s)) \cdot \sigma_s^2,
\]
where the expectation is over scaffolds weighted by their fiber\index{fiber} probabilities.
\end{prop}

\begin{proof}
By the law of total variance:
\[
\Var(F) = \E[\Var(F \mid \Scaf)] + \Var(\E[F \mid \Scaf]).
\]
The first term is $\E_s[\sigma_s^2]$, the expected within-class variance. By definition (Chapter~\ref{ch:variance}), this equals $\Err(\Scaf)$.
\end{proof}

\begin{rem}[Heterogeneity of class variances]
\label{rem:heterogeneity}
The irreducible error $\Err(\Scaf)$ is an \emph{average} over scaffold classes. Individual classes may have variances far from this average. A single high-variance scaffold with many molecules can dominate the error; conversely, a high-variance scaffold with few molecules contributes little to the overall error but may still be important for hit-finding (it could contain rare extreme values).
\end{rem}

\begin{framed}
\noindent\textbf{For the medicinal chemist:} Some scaffold families have ``tight SAR''---small $\sigma_s$, predictable behavior, limited room for optimization. Others are ``wild cards''---large $\sigma_s$, with both duds and stars among their decorated variants. Knowing which is which guides resource allocation: tight-SAR scaffolds are good for reliable leads; wild-card scaffolds are good for finding breakthrough compounds (at higher risk).
\end{framed}

%------------------------------------------------------------------
\section{Sub-Gaussian and Sub-Exponential Distributions}
\label{sec:ch8-subgaussian}

To derive useful tail bounds, we need assumptions on the class distributions. The natural starting point is the Gaussian distribution, but real molecular scores are rarely Gaussian: they may be bounded, asymmetric, multimodal, or heavy-tailed. We therefore work with the weaker notion of \emph{sub-Gaussian} distributions, which retain the essential tail decay of Gaussians without requiring exact Gaussianity.

\subsection{Sub-Gaussian Random Variables}

\begin{defn}[Sub-Gaussian: moment generating function characterization]
\label{defn:subgaussian-mgf}
A centered random variable $X$ (with $\E[X] = 0$) is \emph{sub-Gaussian with parameter $\sigma > 0$} if for all $\lambda \in \R$:
\[
\E[\exp(\lambda X)] \leq \exp\left(\frac{\lambda^2 \sigma^2}{2}\right).
\]
The parameter $\sigma$ is a sub-Gaussian scale; $\sigma^2$ is a variance proxy and need not equal $\Var(X)$.
\end{defn}

The condition states that the moment generating function of $X$ is dominated by that of a centered Gaussian with variance $\sigma^2$. This is the defining property that enables Gaussian-like concentration.

\begin{defn}[Sub-Gaussian: tail characterization]
\label{defn:subgaussian-tail}
The MGF condition with parameter $\sigma$ implies, for all $t>0$,
\[
\Pr(|X| \geq t) \leq 2 \exp\left(-\frac{t^2}{2\sigma^2}\right).
\]
The converse holds with universal changes in constants, not necessarily the same $\sigma$.
\end{defn}

\begin{defn}[Sub-Gaussian norm (Orlicz $\psi_2$ norm)]
\label{defn:subgaussian-norm}
The \emph{sub-Gaussian norm} of a random variable $X$ is
\[
\|X\|_{\psi_2} := \inf\left\{ t > 0 : \E\left[\exp\left(\frac{X^2}{t^2}\right)\right] \leq 2 \right\}.
\]
A random variable is sub-Gaussian if and only if $\|X\|_{\psi_2} < \infty$.
\end{defn}

The displayed Luxemburg--Orlicz convention is a norm on its domain and satisfies the triangle inequality. Equivalent sub-Gaussian scale conventions can differ by universal constants.

\begin{prop}[Equivalence of sub-Gaussian characterizations]
\label{prop:subgaussian-equivalence}
For centered $X$, the following characterizations are equivalent up to universal changes in scale, with constants $c_1, c_2, c_3, c_4 > 0$:
\begin{enumerate}
    \item $\|X\|_{\psi_2} \leq K$
    \item $\Pr(|X| \geq t) \leq 2 \exp(-c_1 t^2 / K^2)$ for all $t > 0$
    \item $\E[\exp(\lambda X)] \leq \exp(c_2 \lambda^2 K^2)$ for all $\lambda \in \R$
    \item $(\E[|X|^p])^{1/p} \leq c_3 K \sqrt{p}$ for all $p \geq 1$
\end{enumerate}
\end{prop}

\begin{proof}[Proof sketch]
The equivalences follow from standard arguments in high-dimensional probability. The key steps are:
\begin{itemize}
    \item $(1) \Rightarrow (2)$: Markov's inequality applied to $\exp(X^2/t^2)$.
    \item $(2) \Rightarrow (3)$: Integration by parts for the MGF.
    \item $(3) \Rightarrow (4)$: Use a Chernoff tail bound and integrate $p t^{p-1}\Pr(|X|>t)$; differentiating an inequality does not preserve it.
    \item $(4) \Rightarrow (1)$: Taylor expand the exponential and bound moments.
\end{itemize}
See Vershynin\cite{vershynin2018highdim}, Proposition 2.5.2, for complete details with the centering and constant changes made explicit.
\end{proof}

\begin{exmp}[Sub-Gaussian distributions]
\label{ex:subgaussian-examples}
~
\begin{enumerate}
    \item \textbf{Gaussian:} If $X \sim \mathcal{N}(0, \sigma^2)$, then $X$ is sub-Gaussian with parameter $\sigma$ and $\|X\|_{\psi_2} = c\sigma$ for a universal constant $c$.

    \item \textbf{Bounded:} If $|X| \leq B$ almost surely, then $X$ is sub-Gaussian with $\|X\|_{\psi_2} \leq cB$ for a universal constant $c$. This is Hoeffding's lemma.

    \item \textbf{Rademacher:} If $X \in \{-1, +1\}$ with equal probability, then $\|X\|_{\psi_2} \leq c$ for a universal constant $c$.

    \item \textbf{Sum of independent sub-Gaussians:} If $X_1, \ldots, X_n$ are independent and centered with $\|X_i\|_{\psi_2} \leq K_i$, then $\|\sum_i X_i\|_{\psi_2} \leq c\sqrt{\sum_i K_i^2}$.
\end{enumerate}
\end{exmp}

\subsection{Sub-Exponential Random Variables}

Some distributions have tails heavier than Gaussian but still light enough for useful concentration. These are captured by the sub-exponential class.

\begin{defn}[Sub-exponential (Orlicz $\psi_1$ norm)]
\label{defn:subexponential}
A centered random variable $X$ is \emph{sub-exponential} if
\[
\|X\|_{\psi_1} := \inf\left\{ t > 0 : \E\left[\exp\left(\frac{|X|}{t}\right)\right] \leq 2 \right\} < \infty.
\]
\end{defn}

\begin{prop}[Sub-exponential tail bound]
\label{prop:subexponential-tail}
If $\|X\|_{\psi_1} \leq K$, then for all $t > 0$:
\[
\Pr(|X| \geq t) \leq 2 \exp\left(-c \min\left(\frac{t^2}{K^2}, \frac{t}{K}\right)\right).
\]
The bound transitions from Gaussian-like ($t^2$) behavior for small $t$ to exponential ($t$) behavior for large $t$.
\end{prop}

\begin{rem}[Relationship between sub-Gaussian and sub-exponential]
\label{rem:subgaussian-subexponential}
Every sub-Gaussian random variable is sub-exponential: $\|X\|_{\psi_1} \leq c\|X\|_{\psi_2}$. The converse fails. Importantly, the square of a sub-Gaussian random variable is sub-exponential: if $\|X\|_{\psi_2} \leq K$, then $\|X^2\|_{\psi_1} \leq cK^2$.
\end{rem}

\subsection{When Do These Assumptions Hold for Molecular Scores?}

The key question for applications: when can we assume that $F \mid \Scaf = s$ is sub-Gaussian?

\begin{prop}[Bounded scores are sub-Gaussian]
\label{prop:bounded-subgaussian}
If the scoring function\index{scoring function} satisfies $F(m) \in [a, b]$ for all molecules $m$ (almost surely), then for any scaffold $s$ of positive probability, the centered conditional response $F-\mu_s$ is sub-Gaussian with parameter at most $(b-a)/2$.
\end{prop}

\begin{proof}
This is Hoeffding's lemma. Any bounded random variable $X \in [a,b]$ satisfies
\[
\E[\exp(\lambda(X - \E[X]))] \leq \exp\left(\frac{\lambda^2(b-a)^2}{8}\right),
\]
which is the sub-Gaussian MGF condition with $\sigma = (b-a)/2$.
\end{proof}

\begin{rem}[Docking scores are bounded]
\label{rem:docking-bounded}
A finite score range must be a property of the specified endpoint, not an
inference from the observed scores. Docking outputs are not universally bounded
by a typical reported range. If the protocol clips scores to a fixed interval,
concentration applies to that clipped endpoint, whose interpretation must be
stated. A sample minimum and maximum do not bound future observations.
\end{rem}

\begin{warn}[When sub-Gaussian may fail]
\label{warn:subgaussian-fail}
Sub-Gaussian assumptions may be inappropriate for:
\begin{enumerate}
    \item \textbf{Unbounded ML predictions:} Unboundedness alone does not disprove sub-Gaussianity: a Gaussian is sub-Gaussian. For an unbounded predictor, justify a scale under the intended molecular sampling law; finite observed outputs cannot establish a population tail bound.

    \item \textbf{Heavy-tailed properties:} Some ADMET properties (e.g., certain metabolic rates) may have heavy tails even after transformation.

    \item \textbf{Multimodal distributions:} A finite mixture with bounded component means and sub-Gaussian scales is sub-Gaussian with a possibly inflated scale; for such a mixture, the simple sub-Gaussian bound may be far from tight. See Section~\ref{sec:ch8-multimodality}.
\end{enumerate}
\end{warn}

\begin{rem}[Conditional vs.\ marginal distributions]
\label{rem:conditional-marginal}
The sub-Gaussian assumption is on the \emph{conditional} distribution $F \mid \Scaf = s$, not the marginal distribution of $F$. Even if the marginal is highly non-Gaussian (because scaffold means $\mu_s$ vary widely), each class distribution $p_s$ may still be sub-Gaussian. The variance spectrum (Chapter~\ref{ch:variance}) separates between-scaffold variation from within-scaffold variation; the sub-Gaussian assumption addresses only the latter.
\end{rem}

%------------------------------------------------------------------
\section{Tail Bounds for Scaffold Classes}
\label{sec:ch8-tail-bounds}

We now derive the tail bounds that underpin scaffold-stratified screening algorithms.

\begin{thm}[Conditional sub-Gaussian tail bound]
\label{thm:conditional-tail}
Assume the centered conditional response $F-\mu_s$ satisfies the MGF condition with a known or independently justified scale upper bound $\nu_s>0$. This need not equal the class standard deviation $\sigma_s$. Then for any $t > 0$:
\[
\Pr(F \leq \mu_s - t \mid \Scaf = s) \leq \exp\left(-\frac{t^2}{2\nu_s^2}\right).
\]
Similarly, $\Pr(F \geq \mu_s + t \mid \Scaf = s) \leq \exp(-t^2/(2\nu_s^2))$.
\end{thm}

\begin{proof}
Let $X = F - \mu_s$, so $X$ is centered. The sub-Gaussian assumption gives
\[
\Pr(X \leq -t) \leq \exp\left(-\frac{t^2}{2\nu_s^2}\right)
\]
by the one-sided version of Definition~\ref{defn:subgaussian-tail}. Substituting $X = F - \mu_s$ yields the result.
\end{proof}

\begin{cor}[Hit probability bound]
\label{cor:hit-probability}
Under the MGF assumption of Theorem~\ref{thm:conditional-tail}, for $\tau < \mu_s$, the probability of finding a molecule\index{molecule} scoring below $\tau$ in scaffold class $s$ is bounded by:
\[
\Pr(F \leq \tau \mid \Scaf = s) \leq \exp\left(-\frac{(\mu_s - \tau)^2}{2\nu_s^2}\right).
\]
\end{cor}

\begin{proof}
Apply Markov's inequality to $\exp(-\lambda(F-\mu_s))$ and optimize over $\lambda>0$; this gives the one-sided bound without the two-sided factor two.
\end{proof}

\begin{rem}[Interpretation of the bound]
\label{rem:bound-interpretation}
The ratio $(\mu_s-\tau)/\nu_s$ measures the gap in units of the justified sub-Gaussian scale. It is a standard-deviation $z$-score only when that scale equals $\sigma_s$, as for a Gaussian with its sharp proxy. Gaps of $\nu_s$, $2\nu_s$, and $3\nu_s$ give upper bounds $e^{-1/2}$, $e^{-2}$, and $e^{-9/2}$, respectively. For $\tau\geq\mu_s$, use the trivial bound one. An upper bound alone gives no positive guarantee of finding a hit.
\end{rem}

\begin{exmp}[Comparing scaffolds for hit-finding]
\label{ex:comparing-scaffolds}
Consider three scaffold classes with the following illustrative means and justified sub-Gaussian scales:
\begin{center}
\begin{tabular}{lccc}
\toprule
Scaffold & $\mu_s$ (kcal/mol) & $\nu_s$ (kcal/mol) & $(\mu_s - \tau)/\nu_s$ \\
\midrule
A & $-8.0$ & $0.8$ & $2.5$ \\
B & $-7.0$ & $1.5$ & $2.0$ \\
C & $-9.5$ & $0.4$ & $1.25$ \\
\bottomrule
\end{tabular}
\end{center}
for threshold $\tau = -10.0$ kcal/mol.

The sub-Gaussian bounds on hit probability are:
\begin{itemize}
    \item Scaffold A: $\exp(-2.5^2/2) = \exp(-3.125) \approx 0.044$
    \item Scaffold B: $\exp(-2.0^2/2) = \exp(-2.0) \approx 0.135$
    \item Scaffold C: $\exp(-1.25^2/2) = \exp(-0.78) \approx 0.458$
\end{itemize}

The displayed upper bounds are ordered C, B, A. They do not establish the ordering of the actual hit probabilities: an upper bound can be loose, and a large variance proxy need not imply a productive favorable tail. Direct hit estimates or additional distributional assumptions are needed for that conclusion.
\end{exmp}

\subsection{Empirical Bernstein Bounds}

When sub-Gaussian assumptions are uncertain or the MGF scale $\nu_s$ is unknown, we can use bounds that depend only on empirical estimates and boundedness.

\begin{thm}[Empirical Bernstein inequality]
\label{thm:empirical-bernstein}
The following two-sided form follows by applying the one-sided bound of Maurer and Pontil\cite{maurer2009bernstein} at failure probability $\delta/2$ to each sign.
Let $X_1, \ldots, X_n$ be i.i.d.\ samples from a distribution on $[a, b]$ with mean $\mu$ and variance $\sigma^2$. Let $\hat{\mu} = \frac{1}{n}\sum_i X_i$ and $\hat{\sigma}^2 = \frac{1}{n-1}\sum_i (X_i - \hat{\mu})^2$ be the sample mean and sample variance. For $n\geq2$ and any $\delta \in (0,1)$, with probability at least $1-\delta$:
\[
|\mu - \hat{\mu}| \leq \sqrt{\frac{2\hat{\sigma}^2 \log(4/\delta)}{n}} + \frac{7(b-a) \log(4/\delta)}{3(n-1)}.
\]
\end{thm}

\begin{rem}[Advantage of empirical Bernstein]
\label{rem:bernstein-advantage}
The empirical Bernstein bound:
\begin{enumerate}
    \item Uses the \emph{observed} variance $\hat{\sigma}^2$, which may be much smaller than the worst-case $(b-a)^2/4$.
    \item Requires the stated independent identically distributed sample, a fixed sample count and known bounds; no parametric distributional form is required.
    \item Adapts to the actual concentration of the data.
\end{enumerate}
The price is slightly more complex form and the need for $n \geq 2$ samples per class.
\end{rem}

\begin{framed}
\noindent\textbf{Practical guidance.}

\medskip
\noindent
When exploring a new scaffold class with limited data:
\begin{enumerate}
    \item Use empirical Bernstein bounds for confidence intervals on $\mu_s$.
    \item As more samples accumulate, test the sub-Gaussian assumption (Section~\ref{sec:ch8-diagnostics}).
    \item Use tighter sub-Gaussian bounds only with a justified scale bound and inference that accounts for this choice. Failure to reject a tail diagnostic does not validate the assumption.
\end{enumerate}
Use intervals valid for the intended sequence of looks; the displayed fixed-count theorem alone does not justify adaptive stopping.
\end{framed}

%------------------------------------------------------------------
\section{Multimodality and Mixture Distributions}
\label{sec:ch8-multimodality}

Not all class distributions are unimodal. Multimodality within a scaffold class represents hidden structure---and an opportunity for refined analysis.

\subsection{Sources of Multimodality}

Why might the score distribution $p_s$ be multimodal for a single scaffold?

\begin{enumerate}
    \item \textbf{Binding mode heterogeneity.} Molecules with the same scaffold may bind the target in different poses. Each binding mode induces a distinct score distribution; the mixture across modes is multimodal.

    \item \textbf{Substituent clusters.} Different decoration\index{decoration} families---e.g., electron-donating vs.\ electron-withdrawing R-groups at a key position---may cluster in score space, creating modes.

    \item \textbf{Stereochemistry.} If the scaffold map ignores stereochemistry, enantiomers or diastereomers with different activities are lumped together.

    \item \textbf{Conformational effects.} Flexible molecules may adopt different bioactive conformations, leading to score heterogeneity.

    \item \textbf{Scoring function artifacts.} Docking scoring functions may have discontinuities (e.g., hydrogen bond on/off) that create artificial modes.
\end{enumerate}

\subsection{Detecting Multimodality}

Several diagnostics can detect multimodality in class distributions.

\begin{defn}[Bimodality coefficient]
\label{defn:bimodality-coefficient}
The \emph{bimodality coefficient} of a distribution is
\[
BC = \frac{\gamma^2 + 1}{\kappa + 3}
\]
where $\gamma$ is skewness and $\kappa$ is excess kurtosis. For a sample of size $n$, the finite-sample corrected version is:
\[
BC_n = \frac{\gamma_n^2 + 1}{\kappa_n + 3 \cdot \frac{(n-1)^2}{(n-2)(n-3)}}.
\]
A uniform distribution has $BC = 5/9 \approx 0.555$. Values $BC > 5/9$ suggest bimodality; values $BC < 5/9$ suggest unimodality.
\end{defn}

\begin{rem}[Limitations of bimodality coefficient]
\label{rem:bc-limitations}
The bimodality coefficient is a simple heuristic, not a formal test. It can miss multimodality when modes are asymmetric or when there are more than two modes. For rigorous analysis, supplement with:
\begin{itemize}
    \item \textbf{Hartigan's dip test:} A formal statistical test against unimodality.
    \item \textbf{Kernel density estimation:} Visual inspection of smoothed density with mode detection.
    \item \textbf{Gaussian mixture modeling:} Fit mixtures and use model selection (BIC/AIC) to determine the number of components.
\end{itemize}
\end{rem}

\subsection{Handling Multimodal Classes}

Once multimodality is detected, several strategies are available.

\begin{enumerate}
    \item \textbf{Refine the scaffold abstraction.}

    If multimodality arises from decoration\index{decoration} patterns, a finer scaffold representation may separate the modes. For example, distinguishing substitution positions (para vs.\ meta vs.\ ortho) or including stereochemistry in the scaffold definition.

    This approach leverages the refinement framework of Chapter~\ref{ch:variance}: move to a finer level of the tower where the modes become distinct classes.

    \item \textbf{Mixture modeling.}

    Model the class distribution as a finite mixture:
    \[
    p_s = \sum_{j=1}^{K_s} w_{s,j} \cdot p_{s,j}
    \]
    where $p_{s,j}$ are component distributions (often Gaussian) and $w_{s,j}$ are mixing weights summing to 1.

    Mixture models can be fit by expectation-maximization (EM) or Bayesian methods. The number of components $K_s$ can be selected by information criteria.

    \item \textbf{Flag and investigate.}

    Sometimes multimodality reflects genuine, irreducible heterogeneity within a scaffold class. Flag such classes for manual investigation---the multimodality may reveal important SAR insights (e.g., a subpopulation with a distinct binding mode worth pursuing).
\end{enumerate}

\begin{prop}[Tail bound for mixtures]
\label{prop:mixture-tail}
Let $p_s = \sum_{j=1}^{K_s} w_{s,j} \cdot p_{s,j}$ where each component $p_{s,j}$ has mean $\mu_{s,j}$ and a centered MGF scale bound $\nu_{s,j}>0$. Then for any threshold $\tau$:
\[
\Pr(F \leq \tau \mid \Scaf = s) = \sum_{j=1}^{K_s} w_{s,j} \cdot \Pr(F \leq \tau \mid \Scaf = s, \text{component } j).
\]
If $\tau < \mu_{s,j}$ for all $j$, each component tail probability is bounded by $\exp(-(\mu_{s,j} - \tau)^2 / (2\nu_{s,j}^2))$.
\end{prop}

\begin{proof}
The first statement is the law of total probability. The second follows from Theorem~\ref{thm:conditional-tail} applied to each component.
\end{proof}

\begin{rem}[Mixtures can have heavier tails than components]
\label{rem:mixture-tails}
Even if each component is sub-Gaussian with small parameter, the mixture can have effectively heavier tails due to the spread of component means. A mixture of two Gaussians $0.5 \cdot \mathcal{N}(-1, 0.1) + 0.5 \cdot \mathcal{N}(1, 0.1)$ has variance $\approx 1.1$ (dominated by the between-component variance), even though each component has variance only $0.1$.

This is another manifestation of the variance decomposition: total variance = within-component variance + between-component variance.
\end{rem}

\begin{framed}
\noindent\textbf{For the medicinal chemist:} Multimodal score distributions within a scaffold often signal unexplored SAR. The ``modes'' may correspond to:
\begin{itemize}
    \item Different binding poses---investigate with docking pose clustering
    \item Functional group effects---stratify by substituent type
    \item Stereochemical preferences---resolve enantiomers
\end{itemize}
Rather than a nuisance, multimodality is diagnostic information pointing toward refinement opportunities. The chemist who understands the source of multimodality can exploit it; the one who ignores it averages over hidden structure.
\end{framed}

%------------------------------------------------------------------
\section{Aggregating Across Classes}
\label{sec:ch8-aggregation}

So far we have focused on individual scaffold classes. In practice, we search across multiple classes simultaneously. This section connects class-level distributions to library-wide hit-finding.

\begin{prop}[Law of total probability for hits]
\label{prop:total-probability-hits}
The library-wide hit probability at threshold $\tau$ is:
\[
\Pr(F \leq \tau) = \sum_{s \in S} \mu(\Fib(s)) \cdot \Pr(F \leq \tau \mid \Scaf = s).
\]
\end{prop}

\begin{proof}
This is the law of total probability applied to the partition $\{\Fib(s)\}_{s \in S}$ of $\Mol$.
\end{proof}

\begin{cor}[Library-wide hit rate bound]
\label{cor:library-hit-bound}
For justified conditional MGF scales $\nu_s>0$, define
\[
u_s(\tau)=\begin{cases}
\exp(-(\mu_s-\tau)^2/(2\nu_s^2)),&\tau<\mu_s,\\
1,&\tau\geq\mu_s.
\end{cases}
\]
Then $\Pr(F\leq\tau)\leq\sum_s\mu(\Fib(s))u_s(\tau)$. The class weight multiplies the bound one for classes whose means lie below the threshold.
\end{cor}

\begin{rem}[Library prevalence and sampling decisions]
\label{rem:strategic}
The weighted sum gives the hit probability of a draw from the declared library measure. It does not prescribe an optimal adaptive allocation. If one can sample a chosen class at fixed cost, the immediate expected hit yield depends on that class's hit probability, not its library frequency; finite availability and unequal costs change the decision. An upper probability bound is useful for ruling out sufficiently unpromising classes only under a valid decision contract, and cannot certify a positive hit rate.
\end{rem}

\begin{exmp}[Resource allocation]
\label{ex:resource-allocation}
Consider a screening budget that allows testing 1000 molecules from a library partitioned into 50 scaffold classes. Some allocation strategies:

\begin{enumerate}
    \item \textbf{Uniform:} 20 molecules per class. Ignores class-level information.

    \item \textbf{Proportional to size:} Allocate proportional to $\mu(\Fib(s))$. Matches the library composition.

    \item \textbf{Hit-rate weighted:} Allocate proportional to $\mu(\Fib(s)) \cdot \Pr(F \leq \tau \mid s)$ (estimated). Focuses on promising classes.

    \item \textbf{Adaptive (bandit\index{bandit algorithm}):} Start with exploration, then shift allocation toward classes with observed good performance. Uses the tail bounds to balance exploration and exploitation.
\end{enumerate}

Chapter~\ref{ch:screening} separates mean identification from molecular hit finding and states the additional assumptions needed for each.
\end{exmp}

%------------------------------------------------------------------
\section{Empirical Diagnostics}
\label{sec:ch8-diagnostics}

Theory is only as good as its assumptions. This section gives diagnostics that can reveal incompatibility with a proposed model. A finite sample cannot establish an unrestricted population tail assumption.

\begin{exper}[Distribution diagnostics within fibers]
\label{exper:distribution-diagnostics}
\textbf{Objective:} For a docking campaign, assess whether class distributions satisfy sub-Gaussian assumptions and detect multimodality.

\textbf{Input:}
\begin{itemize}
    \item Docking scores $\{F(m)\}$ for a library of $N$ molecules
    \item Scaffold assignments $\{\Scaf(m)\}$ for each molecule\index{molecule}
    \item A declared minimum sample size $n_{\min}$; no fixed count guarantees reliable extreme-tail estimation
\end{itemize}

\textbf{Protocol:}

For each scaffold $s$ with $|\Fib(s)| \geq n_{\min}$:

\begin{enumerate}
    \item \textbf{Compute summary statistics:}
    \begin{itemize}
        \item Sample mean: $\hat{\mu}_s = \frac{1}{n_s} \sum_{m \in \Fib(s)} F(m)$
        \item Sample variance: $\hat{\sigma}_s^2 = \frac{1}{n_s - 1} \sum_{m \in \Fib(s)} (F(m) - \hat{\mu}_s)^2$
        \item Sample skewness $\hat{\gamma}_s$ and kurtosis $\hat{\kappa}_s$
    \end{itemize}

    \item \textbf{Histogram:} Plot histogram of $\{F(m) : m \in \Fib(s)\}$ with kernel density overlay.

    \item \textbf{QQ plot:} Compare empirical quantiles against $\mathcal{N}(\hat{\mu}_s, \hat{\sigma}_s^2)$. Deviations in tails indicate non-Gaussianity.

    \item \textbf{Sub-Gaussian coefficient estimation:} Estimate $\|F - \mu_s\|_{\psi_2}$ via Algorithm~\ref{alg:subgaussian-estimation}. Compare to $\hat{\sigma}_s$; for Gaussian distributions, $\hat{K}_s \approx c \cdot \hat{\sigma}_s$ with population reference $c=\sqrt{8/3}$. The finite-sample ratio is a diagnostic, not a valid population scale bound.

    \item \textbf{Multimodality test:} Compute bimodality coefficient $BC_s$. Flag if $BC_s > 0.555$.

    \item \textbf{Tail diagnostic:} For a range of thresholds $\tau$, compare empirical tail probabilities
    \[
    \widehat{\Pr}(F \leq \tau \mid s) = \frac{|\{m \in \Fib(s) : F(m) \leq \tau\}|}{n_s}
    \]
    against the Gaussian reference curve or a bound based on a separately justified MGF scale. A curve obtained by replacing that scale with $\hat\sigma_s$ is uncalibrated and must not be used as a certified tail bound.
\end{enumerate}

\textbf{Output:}
\begin{itemize}
    \item Diagnostic flags: compatibility with a fitted reference model, observed tail deviations, and evidence of multimodality. These labels do not certify a population tail class
    \item Summary statistics: fraction of scaffolds in each category, median ratio $\hat{K}_s / \hat{\sigma}_s$
    \item Flagged scaffolds for investigation
\end{itemize}

\textbf{Expected results:}
\begin{itemize}
    \item Assess compatibility with the declared tail model, without interpreting a finite sample as proof of boundedness.
    \item Report the observed prevalence of multimodality; no population percentage is assumed in advance.
    \item Record outliers and failures under the same prespecified protocol for all methods.
\end{itemize}
\end{exper}

\begin{algor}[Sub-Gaussian parameter estimation]
\label{alg:subgaussian-estimation}
~
\begin{algorithmic}[1]
\Require Samples $\{F_1, \ldots, F_n\}$ from class $s$
\Ensure Estimated sub-Gaussian parameter $\hat{K}_s$
\State Compute sample mean: $\hat{\mu} \gets \frac{1}{n} \sum_{i=1}^{n} F_i$
\State Center samples: $X_i \gets F_i - \hat{\mu}$ for $i = 1, \ldots, n$
\State Initialize: $K_{\text{low}} \gets 0$, $K_{\text{high}} \gets \max_i |X_i|/\sqrt{\log 2}$
\If{$K_{\text{high}}=0$}\State \Return $0$\EndIf
\While{$K_{\text{high}} - K_{\text{low}} > \epsilon$} \Comment{Binary search}
    \State $K_{\text{mid}} \gets (K_{\text{low}} + K_{\text{high}}) / 2$
    \State $\psi \gets \frac{1}{n} \sum_{i=1}^{n} \exp(X_i^2 / K_{\text{mid}}^2)$
    \If{$\psi \leq 2$}
        \State $K_{\text{high}} \gets K_{\text{mid}}$
    \Else
        \State $K_{\text{low}} \gets K_{\text{mid}}$
    \EndIf
\EndWhile
\State \Return $K_{\text{high}}$
\end{algorithmic}
\end{algor}

\begin{rem}[Interpreting the sub-Gaussian parameter]
\label{rem:interpret-K}
The algorithm computes the Orlicz norm of the centered empirical
distribution. It is not an upper confidence bound for a population tail scale.
For a centered Gaussian with variance $\sigma^2$, the population norm under
this definition is $\sqrt{8/3}\,\sigma$. Ratios to sample standard deviation
therefore cannot be interpreted using a Gaussian reference value of one.
Histograms, tail plots, and this empirical norm are diagnostics; they cannot
exclude an unobserved rare component with arbitrary tail behavior.
\end{rem}

%------------------------------------------------------------------
\section{Connections to Subsequent Chapters}
\label{sec:ch8-connections}

The distributional framework of this chapter provides foundations for several developments later in the book.

\subsection{Chapter~\ref{ch:energy}: Energy Factorization\index{energy factorization}}

Chapter~\ref{ch:energy} decomposes the scoring function as $F = F_{\text{core}} + F_{\text{decor}} + F_{\text{boundary\index{boundary}}}$, separating contributions from the scaffold (core), decorations, and their interactions (boundary\index{boundary}). This decomposition induces structure on the class distribution $p_s$:
\begin{itemize}
    \item Under the fixed-core reuse assumptions of Chapter~\ref{ch:energy}, $F_{\text{core}}(s)$ is constant across the class. Shared scaffold identity alone does not fix coordinates, parameters, or receptor state.
    \item $F_{\text{decor}}$ varies within the class (depends on decorations).
    \item $F_{\text{boundary}}$ captures interactions.
\end{itemize}
If the core is fixed and the residual is controlled in a specified norm, the same norm can bound the discrepancy from a shifted decoration response. A small residual variance alone does not control its tail scale. Bounds on the centered component MGF or Orlicz scales need to be supplied separately.

\subsection{Chapter~\ref{ch:mlmc}: Multilevel Monte Carlo\index{multilevel Monte Carlo}}

The exact projection increments $D_\ell$ in Chapter~\ref{ch:variance} describe a response under a molecular measure. They are not automatically cheap evaluators. Executable multilevel estimation requires specified coupled corrections, their bias and variance, and their evaluation costs. Concentration bounds apply only when the implemented samples satisfy the corresponding sampling and tail assumptions. Chapter~\ref{ch:mlmc} states that separate estimator contract.

\subsection{Chapter~\ref{ch:screening}: Bandit\index{bandit algorithm} Algorithms}

A scaffold can be treated as an arm only after its sampling rule and reward are fixed. Mean-score identification uses reward $F$; hit-probability identification uses the indicator $\mathbf1\{F\leq\tau\}$. These objectives need not rank scaffolds alike and neither automatically recovers a molecular top set.

For an independent sample stream with justified sub-Gaussian scale $\nu_s$, a fixed-count mean radius is $\nu_s\sqrt{2\log(2/\delta)/n}$. This is not obtained by substituting the sample standard deviation without further assumptions. Adaptive arm selection and repeated stopping decisions need intervals valid simultaneously over the relevant counts and arms, or an appropriate confidence-sequence construction. Empirical Bernstein is available with known bounds and its stated sampling assumptions. Chapter~\ref{ch:screening} defines the selection problem and inference conditions.

%------------------------------------------------------------------
\section*{Summary}
\label{sec:ch8-summary}

This chapter established the distributional framework for scaffold-stratified molecular scoring.

\begin{enumerate}
    \item \textbf{Class distributions.} Each scaffold $s$ induces a conditional distribution $p_s = \mathrm{Law}(F \mid \Scaf = s)$ with mean $\mu_s$, variance $\sigma_s^2$, and higher moments. The irreducible error $\Err(\Scaf)$ equals the expected class variance: $\Err(\Scaf) = \E_s[\sigma_s^2]$.

    \item \textbf{Sub-Gaussian framework.} For tail bounds and concentration inequalities, we assume class distributions are sub-Gaussian (Definition~\ref{defn:subgaussian-mgf}). This is justified for bounded scores (Proposition~\ref{prop:bounded-subgaussian}) and otherwise needs a justified tail assumption. Diagnostics can test compatibility but cannot establish that assumption. The Orlicz norm $\|\cdot\|_{\psi_2}$ provides a quantitative measure of sub-Gaussianity.

    \item \textbf{Tail bounds.} For a justified centered conditional MGF scale $\nu_s>0$ and $\tau<\mu_s$, the probability of scores below threshold $\tau$ satisfies (Theorem~\ref{thm:conditional-tail}):
    \[
    \Pr(F \leq \tau \mid \Scaf = s) \leq \exp\left(-\frac{(\mu_s - \tau)^2}{2\nu_s^2}\right).
    \]
    For $\tau\geq\mu_s$, use the bound one. Empirical Bernstein instead gives a mean interval under known bounds and the theorem's sampling assumptions; it is not a substitute for a tail-probability model.

    \item \textbf{Multimodality.} Class distributions may be multimodal due to binding mode heterogeneity, substituent clustering, stereochemistry, or scoring artifacts. Detecting multimodality (via bimodality coefficient, dip tests, or mixture modeling) reveals hidden SAR structure. Handling multimodality requires refinement, mixture modeling, or flagging for investigation.

    \item \textbf{Aggregation.} Library-wide hit rates aggregate class-level hit probabilities weighted by scaffold frequencies (Proposition~\ref{prop:total-probability-hits}). This reveals the structure of efficient scaffold-stratified screening.

    \item \textbf{Diagnostics.} Always test distributional assumptions via histograms, QQ plots, sub-Gaussian parameter estimation, and multimodality tests (Experiment~\ref{exper:distribution-diagnostics}) before relying on tail bounds.
\end{enumerate}

\medskip

\noindent\textbf{For medicinal chemists:} Some scaffolds are predictable (tight distributions, reliable SAR), others are wild cards (broad distributions, surprising outliers). The class standard deviation $\sigma_s$ quantifies spread under the specified sampling law; it does not by itself quantify hit yield. Multimodal distributions signal unexplored SAR---the modes may correspond to different binding poses, functional group effects, or stereochemical preferences. Understanding distributional structure guides resource allocation: tight-SAR scaffolds for reliable leads, wild-card scaffolds for breakthrough discoveries.

\medskip

\noindent\textbf{For modelers and statisticians:} Keep actual variance, an MGF variance proxy, and the empirical Orlicz norm distinct. A Gaussian has population Orlicz-to-standard-deviation ratio $\sqrt{8/3}$ under this chapter's convention. No empirical ratio validates a population tail assumption. Inference must account for dependence, adaptive collection, parameter selection, and repeated looks.

%------------------------------------------------------------------
\section*{Exercises}

\begin{exer}[Chemistry: Interpreting class distributions]
\label{exer:ch8-chem}
A docking campaign against Target Y yields the following class statistics for three scaffolds:

\begin{center}
\begin{tabular}{lcccc}
\toprule
Scaffold & $n_s$ & $\hat{\mu}_s$ (kcal/mol) & $\hat{\sigma}_s$ (kcal/mol) & $BC_s$ \\
\midrule
A & 150 & $-8.5$ & $0.8$ & $0.35$ \\
B & 80 & $-7.8$ & $1.5$ & $0.62$ \\
C & 200 & $-9.0$ & $0.6$ & $0.41$ \\
\bottomrule
\end{tabular}
\end{center}

\begin{enumerate}
    \item[(a)] Which scaffold is most likely to yield a hit scoring below $\tau = -10.5$ kcal/mol? Compute the sub-Gaussian bound for each scaffold.

    \item[(b)] Scaffold B has $BC_s > 0.555$. What does this suggest about its score distribution? Propose an explanation in terms of SAR.

    \item[(c)] If you could screen 100 additional molecules from exactly one scaffold, which would you choose and why? Consider both hit probability and the value of information.

    \item[(d)] How would your answer to (c) change if the threshold were $\tau = -9.5$ kcal/mol instead?
\end{enumerate}
\end{exer}

\begin{exer}[Mathematics: Deriving Proposition~\ref{prop:irred-error-class-var}]
\label{exer:ch8-math-1}
Prove that $\Err(\Scaf) = \E_s[\sigma_s^2]$ directly from the law of total variance:
\[
\Var(F) = \E[\Var(F \mid \Scaf)] + \Var(\E[F \mid \Scaf]).
\]
Identify which term corresponds to the irreducible error and which to the between-scaffold variance.
\end{exer}

\begin{exer}[Mathematics: Sub-Gaussian properties]
\label{exer:ch8-math-2}
~
\begin{enumerate}
    \item[(a)] For centered $X$ with $|X|\leq B$ almost surely, prove that $B$ is an admissible sub-Gaussian scale. Show with a Rademacher variable taking values $\pm B$ that $B/\sqrt2$ is not generally admissible. Explain why an already admissible scale can always be enlarged.

    \item[(b)] Show that if $X_1, \ldots, X_n$ are independent centered sub-Gaussian random variables with parameters $\sigma_1, \ldots, \sigma_n$, then $\sum_i X_i$ is sub-Gaussian with parameter $\sqrt{\sum_i \sigma_i^2}$.

    \item[(c)] Let the centered conditional response $F-\mu_s$ have MGF scale bound $\nu_s$. For a sample of $n$ i.i.d.\ molecules from class $s$, show that the sample mean $\bar{F} = \frac{1}{n}\sum_{i=1}^n F(m_i)$ satisfies:
    \[
    \Pr(|\bar{F} - \mu_s| \geq t) \leq 2 \exp\left(-\frac{nt^2}{2\nu_s^2}\right).
    \]

    \item[(d)] How many samples $n$ are needed to ensure $|\bar{F} - \mu_s| \leq 0.5$ kcal/mol with probability at least $0.95$, if the justified MGF scale bound is $\nu_s = 1.0$ kcal/mol?
\end{enumerate}
\end{exer}

\begin{exer}[ML/Statistics: Testing the sub-Gaussian assumption]
\label{exer:ch8-ml}
~
\begin{enumerate}
    \item[(a)] Implement Algorithm~\ref{alg:subgaussian-estimation} for estimating the sub-Gaussian parameter $\hat{K}_s$.

    \item[(b)] For a Gaussian distribution $\mathcal{N}(0, \sigma^2)$, what is the theoretical value of $\|X\|_{\psi_2}$ in terms of $\sigma$? Verify empirically with simulated data ($n = 1000$ samples, various $\sigma$).

    \item[(c)] Generate $n = 500$ samples from a mixture $0.5 \cdot \mathcal{N}(-2, 0.5^2) + 0.5 \cdot \mathcal{N}(2, 0.5^2)$. Estimate $\hat{K}$ using your implementation. Compare to (i) the standard deviation of the mixture and (ii) the sub-Gaussian parameter of a single Gaussian with the same variance.

    \item[(d)] Propose a diagnostic criterion based on $\hat{K}_s / \hat{\sigma}_s$ for flagging scaffolds where the sub-Gaussian assumption may be inappropriate. Test your criterion on simulated data from Gaussian, uniform, and bimodal distributions.
\end{enumerate}
\end{exer}

\begin{exer}[Computation: Full diagnostic pipeline]
\label{exer:ch8-comp}
Implement the full Experiment~\ref{exper:distribution-diagnostics} protocol:

\begin{enumerate}
    \item[(a)] Given a dataset of (molecule, scaffold, score) triples, compute class statistics $(\hat{\mu}_s, \hat{\sigma}_s, BC_s, \hat{K}_s)$ for all scaffolds with $n_s \geq 30$.

    \item[(b)] Generate QQ plots for a random sample of 10 scaffolds.

    \item[(c)] Classify scaffolds as ``well-behaved,'' ``heavy-tailed,'' or ``multimodal'' based on your diagnostics. Define explicit criteria for each category.

    \item[(d)] Report the fraction of scaffolds in each category and discuss implications for screening strategy.

    \item[(e)] For scaffolds classified as multimodal, propose a refinement strategy (e.g., subdivide by substitution pattern) and estimate whether refinement would reduce within-class variance.
\end{enumerate}
\end{exer}

\begin{exer}[Theory: Mixture models and tail bounds]
\label{exer:ch8-theory}
~
\begin{enumerate}
    \item[(a)] Let $p = \sum_{j=1}^{K} w_j \cdot p_j$ be a finite mixture where each $p_j$ is $\mathcal{N}(\mu_j, \sigma_j^2)$. Derive the mean, variance, skewness, and kurtosis of $p$ in terms of $\{w_j, \mu_j, \sigma_j\}$.

    \item[(b)] Show that a finite mixture of centered-at-their-own-means sub-Gaussian components has a centered sub-Gaussian scale bound when both component means and scales are bounded. Account for the spread of component means; explain why an arbitrary infinite mixture need not have this property.

    \item[(c)] For a bimodal mixture $p = 0.5 \cdot \mathcal{N}(\mu_1, \sigma^2) + 0.5 \cdot \mathcal{N}(\mu_2, \sigma^2)$ with $\mu_1 < \mu_2$, compute $\Pr(X \leq \tau)$ for threshold $\tau < \mu_1$. Compare to the sub-Gaussian bound using the mixture's overall variance.

    \item[(d)] Under what conditions (on $\mu_1, \mu_2, \sigma$) is the exact mixture tail probability significantly larger than the sub-Gaussian bound? Interpret in terms of scaffold screening.
\end{enumerate}
\end{exer}

%------------------------------------------------------------------
\section{Bibliographic Notes}

The theory of sub-Gaussian and sub-exponential random variables is developed thoroughly in Vershynin (2018, \emph{High-Dimensional Probability}) and Wainwright (2019, \emph{High-Dimensional Statistics}). The Orlicz norm characterization provides a unified framework for concentration inequalities; we follow the conventions of Vershynin, Chapter 2.

Concentration inequalities for bounded random variables originate with Hoeffding (1963). The empirical Bernstein inequality, which uses observed variance, appears in Maurer and Pontil (2009) and is widely used in bandit algorithms and adaptive experimental design.

The bimodality coefficient and other multimodality diagnostics are surveyed in Pfister et al.\ (2013). Hartigan's dip test for unimodality (Hartigan and Hartigan, 1985) provides a formal statistical test. Gaussian mixture modeling via EM is classical; see McLachlan and Peel (2000).

The application of sub-Gaussian bounds to molecular scoring functions appears to be novel in this systematic form. Related ideas in computational chemistry include the analysis of scoring function distributions by Chaput et al.\ (2016), discussions of docking score\index{docking score} variability by Trott and Olson (2010), and the use of confidence intervals in virtual screening by Riniker and Landrum (2013).

The connection between distributional assumptions and bandit algorithm validity is central to the theoretical guarantees in Chapter~\ref{ch:screening}; see Lattimore and Szepesv\'{a}ri (2020, \emph{Bandit Algorithms}) for the general theory. The application of bandit methods to drug discovery is surveyed in Korovina et al.\ (2020) and Graff et al.\ (2021).
% END SOURCE: chapters/ch08_probability.tex


\chapter{Energy Decomposition and Limits of Reuse}
\label{ch:energy}
% BEGIN SOURCE: chapters/ch09_energy.tex
% Energy decomposition and limits of reuse. Chapter heading is in main.tex.

\section{Decomposition is not yet reuse}
\label{sec:ch9-thesis}

A score can be partitioned into core, decoration, and residual terms. The computational question is whether a term evaluated for one molecule can be reused for another with a controlled error. An exact decomposition with an unrestricted residual does not answer that question: the residual may contain almost all of the variation.

This chapter separates three claims. Algebra partitions a specified score exactly. Assumptions about coordinates, parameters, and response bound the error of reusing its parts. A computational comparison then measures the cost saved against a suitable baseline. Chapter~\ref{ch:response-certificates} develops the corresponding response-transfer certificates.

\section{Scoring protocols}
\label{sec:ch9-hierarchy}

\begin{defn}[Scoring protocol]
\label{defn:scoring-hierarchy}
A protocol specifies a molecular identity convention, receptor and chemical context, preparation procedure, score or response functional, and sampling rule. Empirical docking, MM-GBSA/PBSA, and alchemical FEP/TI are different protocols. Their outputs and computational costs are not interchangeable. A more detailed physical model can still have sampling, parameterization, and systematic errors.
\end{defn}

This chapter's deterministic formulas apply either to fixed configurations or to a declared scalar response after a specified averaging procedure. Physical free energies require equilibrium ensembles and a thermodynamic definition. Finite simulations provide estimates with uncertainty. Computation budgets should record measured wall time, hardware, preparation failures, and effective sampling rather than assign one universal cost per molecule.

\section{A fixed-configuration energy partition}
\label{sec:ch9-anatomy}
\label{sec:ch9-partition}

\begin{defn}[Pairwise scoring model]
\label{defn:generic-docking}
For a ligand $m$, fixed receptor $P$, and coordinates $\xi$, consider
\begin{equation}
\label{eq:generic-docking}
F(m,P,\xi)=\sum_{i\in m,j\in P}\phi_{ij}(r_{ij})
+\sum_{i<k\in m}\psi_{ik}(r_{ik})+\Omega(m,P,\xi).
\end{equation}
The final term includes contributions not represented by the two pairwise sums. The atom types, charges, and other parameters are part of the protocol.
\end{defn}

This is a model class, not a claim that every docking score or force field is pairwise additive. Many-body terms, solvent models, and learned scoring functions require their own factors or remain in $\Omega$.

\begin{defn}[Molecular partition]
\label{defn:molecular-partition}
Choose an atom mapping identifying the scaffold core $C$ and decoration atoms $D=m\setminus C$. Let $A\subseteq C$ be the attachment atoms. Symmetry, stereochemistry, protonation, and hydrogen handling must be resolved by the mapping convention.
\end{defn}

\begin{defn}[Energy partition]
\label{defn:energy-partition}
\begin{align}
F_{\mathrm{core}}&=\sum_{i\in C,j\in P}\phi_{ij}
+\sum_{i<k\in C}\psi_{ik},\label{eq:F-core}\\
F_{\mathrm{decor}}&=\sum_{i\in D,j\in P}\phi_{ij}
+\sum_{i<k\in D}\psi_{ik},\label{eq:F-decor}\\
F_{\mathrm{boundary}}&=\sum_{i\in C,k\in D}\psi_{ik}
+\Omega(m,P,\xi).\label{eq:F-boundary}
\end{align}
Thus
\begin{equation}
\label{eq:energy-partition}
F=F_{\mathrm{core}}+F_{\mathrm{decor}}+F_{\mathrm{boundary}}.
\end{equation}
This convention assigns all remaining terms to the boundary term. Other conventions are possible if every contribution is assigned exactly once.
\end{defn}

\begin{thm}[Exact decomposition at fixed geometry]
\label{sec:ch9-exact-decomp}
\label{thm:exact-factorization}
\label{eq:exact-factorization}
The partition in Definition~\ref{defn:energy-partition} is exact for the specified score. The core contribution is reusable across molecules only when its mapped core coordinates, receptor coordinates, and every parameter used in its terms are identical.
\end{thm}
\begin{proof}
Each ligand--receptor pair has its ligand atom in exactly one of $C,D$. Each intraligand pair is either wholly in $C$, wholly in $D$, or crosses the partition. The remainder is included explicitly. Equality of the inputs to each core term makes those terms equal across molecules.
\end{proof}

Shared scaffold identity alone does not establish equality of core scores. A decoration can change the core pose, its charge assignment, protonation, or receptor environment. Even at fixed geometry, reparameterizing the full molecule can change the core terms.

\begin{cor}[Scaffold signal under a reusable core]
\label{cor:scaffold-signal}
\label{prop:projection-factorized}
Suppose the protocol permits $F=a(S)+Y+B$, where $S=\Scaf(M)$ and all terms are integrable. Then
\[
\mathbb E[F\mid S=s]=a(s)+\mathbb E[Y+B\mid S=s].
\]
The scaffold signal equals $a(s)$ only when the conditional mean of the remaining terms is zero.
\end{cor}

\section{M\"obius coefficients and response interactions}
\label{sec:ch9-incidence}

M\"obius inversion gives exact coordinates for a response table on a finite ordered family\cite{rota1964mobius,stanley2011enumerative}. The coordinates depend on the poset, chosen reference structures, and scoring protocol; they do not identify a physical mechanism without further evidence.

\begin{defn}[Incidence algebra and scaffold coefficients]
\label{defn:incidence-algebra}
\label{defn:scaffold-cumulants}
For a finite poset $P$, convolution is
\[
(h*k)(x,z)=\sum_{x\preceq y\preceq z}h(x,y)k(y,z).
\]
The zeta function is $\zeta(x,y)=1$ on comparable pairs, and its inverse is $\mu_P$. For a scalar response $E:P\to\mathbb R$, define
\[
\kappa_E(x)=\sum_{y\preceq x}\mu_P(y,x)E(y).
\]
We call these scaffold coefficients, or scaffold cumulants in the finite-difference sense. They are not probabilistic cumulants.
\end{defn}

\begin{prop}[M\"obius decomposition]
\label{prop:mobius-energy}
The coefficients above are the unique coefficients satisfying
\[
E(x)=\sum_{y\preceq x}\kappa_E(y).
\]
\end{prop}
\begin{proof}
The zeta matrix in any linear extension of the finite poset is triangular with unit diagonal. Its inverse is the M\"obius matrix. Applying that inverse gives existence and uniqueness.
\end{proof}

\begin{exmp}[Two-site response interaction]
\label{ex:boolean-cumulants}
Fix a scaffold and a designated substitution at each of $r$ sites. Assume every required combination is defined and valid. For $I\subseteq[r]$, let $E_s(I)$ be the response with the substitutions in $I$ and reference groups elsewhere. On the Boolean lattice,
\[
\kappa_s(I)=\sum_{J\subseteq I}(-1)^{|I|-|J|}E_s(J).
\]
For two sites,
\[
\kappa_s(\{i,j\})=E_s(\{i,j\})-E_s(\{i\})-E_s(\{j\})+E_s(\varnothing).
\]
This compares one substitution's effect in two chemical contexts. It can be nonzero even if every $E_s(I)$ is an exact equilibrium free energy. It is not the sum of free-energy differences around a closed thermodynamic cycle.
\end{exmp}

\begin{thm}[Additivity as finite-difference truncation]
\label{thm:sar-cumulant-truncation}
On this complete Boolean table, $E_s(I)$ has an additive baseline-plus-single-site form if and only if $\kappa_s(I)=0$ for all $|I|\geq2$. More generally, a multilinear expansion of degree at most $q$ exists if and only if $\kappa_s(I)=0$ for $|I|>q$.
\end{thm}
\begin{proof}
M\"obius inversion expresses the response as the sum of coefficients on subsets. Removing all terms above degree $q$ gives the claimed expansion. Uniqueness proves the converse.
\end{proof}

\begin{rem}[Physical interpretation requires a comparison]
\label{rem:cumulant-stronger}
An interaction coefficient records failure of the specified additive response model. It does not distinguish direct atom interactions, solvation, conformational redistribution, assay context, or measurement error. Controlled comparisons across protocols can test these alternatives. Missing combinations require an estimation model and cannot be filled by the algebra alone.
\end{rem}

\begin{conj}[Locality of response interactions]
\label{conj:cumulant-locality}
For a specified family with stable poses and controlled sampling, low-order response interactions may dominate higher-order ones. Spatial separation may help predict their size. This is a hypothesis to test under declared protocols; finite interaction range alone does not establish it after minimization or ensemble averaging.
\end{conj}

\section{What locality can bound}
\label{sec:ch9-boundary-decay}

\begin{prop}[A bound for direct interactions]
\label{thm:boundary-decay}
Let two fixed atom groups have sizes $n_1,n_2$, all cross distances at least $R$, and cross interaction terms bounded by $|\phi_{ij}(r)|\leq C e^{-r/\lambda}$. Their direct interaction obeys
\begin{equation}
\label{eq:boundary-decay-screened}
|E_{12}|\leq n_1n_2 C e^{-R/\lambda}.
\end{equation}
If instead $|\phi_{ij}(r)|\leq Cr^{-\alpha}$, then
\begin{equation}
\label{eq:boundary-decay-powerlaw}
|E_{12}|\leq n_1n_2 C R^{-\alpha}.
\end{equation}
For a finite cutoff, the direct interaction is zero when every cross distance exceeds that cutoff.
\end{prop}
\begin{proof}
Apply the bound to each cross pair and sum absolute values.
\end{proof}

\begin{rem}[Limits of the distance bound]
\label{rem:practical-decay}
\label{cor:distant-independence}
The bound concerns direct interactions between the two fixed groups. It does not bound induced fit, long-range collective response, changes in solvation, or the interaction coefficient of a minimized or ensemble-averaged response. An electrostatic summation algorithm does not itself supply an exponential screening assumption. Attachment-point separation also need not lower-bound all atom-pair distances for flexible decorations.
\end{rem}

\begin{prop}[A conditional truncation bound]
\label{rem:cluster-expansion}
If a response expansion has total order-$n$ contribution $R_n$ with $|R_n|\leq C\rho^n$ for some $0<\rho<1$, then
\[
\left|\sum_{n>N}R_n\right|\leq\frac{C\rho^{N+1}}{1-\rho}.
\]
The geometric bound is an assumption on total contributions, including the number of terms at each order. It is not a consequence of pairwise locality alone.
\end{prop}

\section{Stability under pose and parameter changes}
\label{sec:ch9-pose-stability}

\begin{defn}[Pose deviation]
\label{defn:pose-space}
\label{defn:pose-core}
\label{defn:rigidity-index}
Compare mapped core coordinates $x_i,x_i'$ in the same receptor frame and let
\[
\delta=\left(\frac1{n_C}\sum_{i\in C}\|x_i-x_i'\|^2\right)^{1/2}.
\]
A free alignment that moves the ligand relative to the receptor would remove a physically relevant part of this displacement. For a distribution of decorations, define $\rho(s)=\mathbb E[\delta\mid S=s]$ relative to a declared reference pose.
\end{defn}

\begin{thm}[Lipschitz core stability]
\label{thm:lipschitz-stability}
Hold receptor coordinates and interaction parameters fixed. Suppose every core--receptor term is $L$-Lipschitz in its core coordinate on the domain of interest and every core--core term is $L'$-Lipschitz with respect to the sum of its two coordinate displacements. Then
\[
|F_{\mathrm{core}}(x)-F_{\mathrm{core}}(x')|
\leq [|P|L+(n_C-1)L']n_C\delta.
\]
\end{thm}
\begin{proof}
Put $d_i=\|x_i-x_i'\|$. Summation bounds the change by
$[|P|L+(n_C-1)L']\sum_i d_i$.
Cauchy--Schwarz gives $\sum_i d_i\leq n_C\delta$. An RMSD bound does not assert that each $d_i\leq\delta$.
\end{proof}

\begin{cor}[Expected core reuse error]
\label{cor:tight-binding}
Under the same uniform constants, expected absolute core reuse error is at most $[|P|L+(n_C-1)L']n_C\rho(s)$. Parameter and receptor changes need separate bounds. No universal RMSD threshold certifies energy accuracy.
\end{cor}

\section{MM-GBSA/PBSA and residual solvation}
\label{sec:ch9-mmgbsa}

\begin{defn}[End-point score convention]
\label{defn:mmgbsa}
For a specified complex, receptor, and ligand protocol, an MM-GBSA estimate is schematically
\[
\widehat{\Delta G}_{\mathrm{bind}}=
\langle\Delta E_{\mathrm{MM}}+\Delta G_{\mathrm{solv}}\rangle
-T\widehat{\Delta S},
\]
where each $\Delta$ is complex minus receptor minus ligand. Trajectory selection, protonation, dielectric treatment, and entropy estimation are part of the definition.
\end{defn}

\begin{prop}[Mechanical term bookkeeping]
\label{thm:mm-decomp}
For an explicitly listed molecular-mechanics energy, assign each term to the core, decoration, or a residual according to all the atoms and parameters it depends on. This gives an exact sum at each configuration. Averaging preserves this equality when all terms are averaged under the same measure. It does not make the averages reusable across different ensembles.
\end{prop}

\begin{defn}[Generalized Born dependence]
\label{defn:gb-solvation}
A schematic GB contribution has the form
\[
G_{\mathrm{GB}}=-\frac12(1-\varepsilon^{-1})
\sum_{i,j}\frac{q_iq_j}{f_{\mathrm{GB}}(r_{ij},R_i,R_j)},
\]
with units and prefactors fixed by the implementation. Effective radii $R_i$ depend on the molecular geometry. Adding decoration atoms can change radii associated with core atoms.
\end{defn}

\begin{prop}[Solvation decomposition with an explicit residual]
\label{thm:solvation-nonadd}
For any chosen reference functions $a(s),b(d)$,
\[
G_{\mathrm{solv}}(s,d)=a(s)+b(d)+r(s,d)
\]
holds by defining $r$ as the difference. Dependence of core Born radii on decoration prevents an automatic zero-residual assertion. No bound on $r$ follows from the displayed identity.
\end{prop}

\begin{rem}[Solvent exposure as a testable diagnostic]
\label{prop:solvation-linear}
\label{cor:solvent-exposed}
Solvent exposure, charge changes, and conserved geometry are plausible predictors of reuse error. A universal bound by overlapping solvent-accessible area is not established here. Fit and test such a model against held-out complete calculations; retain its residual and domain of validity.
\end{rem}

\begin{prop}[A conditional error budget]
\label{thm:mmgbsa-factorization}
If replacing the mechanical, solvation, and entropy contributions incurs absolute errors at most $\epsilon_{\mathrm{MM}},\epsilon_{\mathrm{solv}},\epsilon_{\mathrm{ent}}$, respectively, the resulting end-point score error is at most their sum. Each bound must be supplied for the protocol and domain being used.
\end{prop}
\begin{proof}
Apply the triangle inequality to the three differences.
\end{proof}

\section{Alchemical free energies and graph baselines}
\label{sec:ch9-alchemical}

\begin{defn}[Relative binding free energy]
\label{defn:relative-binding}
\label{defn:rgroup-perturbation}
For fixed thermodynamic conditions and state definitions, write $G(m)=\Delta G_{\mathrm{bind}}(m)$ and $\Delta\Delta G_{A\to B}=G(m_B)-G(m_A)$. An alchemical estimate combines the corresponding transformations in bound and unbound environments. Its uncertainty depends on phase-space overlap, sampling, and model assumptions.
\end{defn}

\begin{prop}[Cancellation of a chosen scaffold baseline]
\label{thm:scaffold-cancellation}
If $G(s,d)=a(s)+b(d)+r(s,d)$ under a specified reference convention, then
\[
G(s,d_B)-G(s,d_A)=b(d_B)-b(d_A)+r(s,d_B)-r(s,d_A).
\]
The chosen baseline $a(s)$ cancels. The residual difference need not be small.
\end{prop}
\begin{proof}
Subtract the two equalities.
\end{proof}

This identity does not derive an additive free-energy model from a pairwise Hamiltonian. In general, $-\beta^{-1}\log\int e^{-\beta(H_c+H_d+H_b)}$ does not split into separate core and decoration free energies. Product configuration spaces, product reference measures, and absence of coupling are sufficient for partition-function factorization; shared scaffold labels are not.

\begin{rem}[Within-scaffold transformations]
\label{cor:fep-within-scaffold}
Closely related molecules may admit useful atom mappings and alchemical schedules. Their shared scaffold is a candidate-selection feature. It does not guarantee efficient equilibration or transferable substituent effects.
\end{rem}

\begin{prop}[Path independence of exact free energies]
\label{thm:cross-scaffold-mapping}
\label{defn:scaffold-hopping-fep}
For any sequence of well-defined states $m_0,\ldots,m_k$,
\[
\sum_{i=0}^{k-1}[G(m_{i+1})-G(m_i)]=G(m_k)-G(m_0).
\]
In particular, every closed state-function loop sums to zero. Approximate edge estimates can have nonzero closure residuals because of estimation error or inconsistent state/protocol definitions.
\end{prop}
\begin{proof}
All intermediate values cancel in the sum.
\end{proof}

\begin{rem}[Alchemical paths and costs]
\label{prop:optimal-path}
\label{rem:topological-obstructions}
\label{warn:scaffold-hopping-difficulty}
An alchemical path is a path of Hamiltonians or thermodynamic states, not a continuous path in a discrete set of scaffold graphs. Absence of a common substructure does not prove nonexistence of such a path. Structural similarity can help propose manageable transformations, but path costs and overlap must be evaluated. Statistical similarity of marginal docking-score distributions does not certify alchemical feasibility.
\end{rem}

\begin{thm}[Graph counts and genuine additive-model savings]
\label{thm:combinatorial-reduction}
Let $K$ scaffold classes contain $L$ molecules each, so $N=KL$.
\begin{enumerate}
\item Without a response model relating the $N$ free energies, exact relative differences on any connected comparison graph determine them up to an additive constant. Such a graph needs at least $N-1$ edges, and a spanning tree attains that count.
\item Within-class spanning trees plus a tree connecting one representative from each class use $K(L-1)+(K-1)=N-1$ edges. This reorganizes the baseline; it does not reduce its edge count.
\item Complete within-class graphs plus a complete graph of representatives use $K\binom L2+\binom K2$ edges. This is fewer than the all-pairs count $\binom{KL}{2}$, but generally far more than $N-1$.
\item If all $K\times L$ combinations have a common decoration indexing and the exact model $G(s,d)=a_s+b_d$ holds, a reference scaffold row and decoration column determine all relative values using $(K-1)+(L-1)$ edges. This additional saving relies on the additive response assumption.
\end{enumerate}
\end{thm}
\begin{proof}
A connected graph has an incidence matrix of rank $N-1$, so its differences determine one potential per node up to a constant. A tree has exactly that many edges. The second and third counts follow by adding the listed edges. Under the additive model,
\[
G(s,d)=G(s,d_0)+G(s_0,d)-G(s_0,d_0),
\]
so the reference row and column suffice.
\end{proof}

\begin{exmp}[A fair count comparison]
\label{ex:concrete-savings}
For $K=100$ and $L=1000$, the all-pairs graph has $4{,}999{,}950{,}000$ edges; the two-level complete graph has $49{,}954{,}950$; a spanning tree has $99{,}999$. Under the extra exact additive response model, $1098$ reference edges suffice. These are edge counts, not measured runtimes or equal-accuracy comparisons. A noisy spanning tree can propagate errors, redundant cycles can help diagnose them, and some edges may be infeasible or much more expensive than others.
\end{exmp}

For the reference-row reconstruction, the exact error is the response rectangle
\[
G(s,d)-G(s,d_0)-G(s_0,d)+G(s_0,d_0).
\]
It need not vanish when thermodynamic cycles close. Chapter~\ref{ch:response-certificates} specifies how a bound on this quantity can support approximate reuse, including uncertainty from measured reference values.

\section{Induced fit and protocol state}
\label{sec:ch9-induced-fit}

\begin{defn}[Relaxation response]
\label{defn:response-operator}
A deterministic preparation-and-relaxation protocol maps an initial environment $z_0$ to a selected output $\mathcal R_m(z_0)$. It need not be a unique global energy minimizer or an equilibrium sample. A stochastic simulation protocol is instead a kernel $K_m(z_0,dz)$, whose definition includes duration, temperature, restraints, and initialization.
\end{defn}

\begin{defn}[Induced-fit deviation]
\label{defn:induced-fit-deviation}
For a declared common frame and initial state, define
$\Delta\mathcal R(d,d_0;s)=\|\mathcal R_{(s,d)}(z_0)-\mathcal R_{(s,d_0)}(z_0)\|$.
This compares selected protocol outputs, not whole equilibrium distributions.
\end{defn}

\begin{prop}[A coordinate-response bound]
\label{thm:factorization-induced-fit}
If $F(m,z)$ is $L_P$-Lipschitz in environment coordinates on the compared states, then
\[
|F(m,z_m)-F(m,z_{\mathrm{ref}})|\leq L_P\|z_m-z_{\mathrm{ref}}\|.
\]
For ensemble responses, a coupling of the two state distributions gives the corresponding expectation bound whenever the same Lipschitz condition holds.
\end{prop}
\begin{proof}
Apply the Lipschitz inequality pointwise and integrate under the declared coupling for the ensemble version.
\end{proof}

\begin{rem}[Protocol dependence]
\label{cor:induced-fit-diagnostic}
\label{rem:operator-theoretic}
Changing ligand identity generally requires new parameterization and an explicit initialization of solvent, velocities, and other state. An RLMM transition should record those operations as part of its stochastic protocol. A proposed noncommutativity diagnostic requires defined maps on compatible state spaces; the words ``bind'' and ``relax'' alone do not define an operator commutator.
\end{rem}

\begin{algor}[Reuse diagnostic]
\label{alg:factorization-diagnostic}
\label{warn:factorization-fails}
Declare the reference protocol and intended error tolerance. Check structural mapping, pose and parameter changes, ensemble overlap, and held-out response rectangles. Compare observed errors to the proposed bound, including uncertainty. Record unsupported contexts and preparation failures. Geometric diagnostics can prioritize these checks; fixed RMSD or boundary-variance thresholds are not universal certificates.
\end{algor}

\section{Statistical independence and energy additivity}
\label{sec:ch9-info-theory}

\begin{defn}[Conditional mutual information]
\label{defn:mi-coupling}
For random variables $X,Y$ under a declared molecular and simulation sampling law,
\[
I(X;Y\mid S=s)=D_{\mathrm{KL}}(P_{XY\mid s}\|P_{X\mid s}\otimes P_{Y\mid s}).
\]
This is a property of a probability law, not of the notation used to partition an energy.
\end{defn}

\begin{prop}[A scaffold-constant core has zero conditional information]
\label{prop:mi-factorization}
If $X=a(S)$, then $I(X;Y\mid S=s)=0$ for every conditioning value with positive probability, regardless of any boundary term $B$. Thus this quantity cannot detect nonadditive response within an exact scaffold fiber.
\end{prop}
\begin{proof}
Conditional on $S=s$, $X$ is constant, so $P_{XY\mid s}=\delta_{a(s)}\otimes P_{Y\mid s}$.
\end{proof}

\begin{rem}[When information measures are informative]
\label{rem:mi-canonical}
If the core score varies with sampled poses or environments, its dependence with other observables can be measured. The resulting MI depends on the ensemble and the chosen variables. Zero MI means independence of those variables; it does not imply zero energy residual. Conversely, an additive response evaluated on correlated molecular proposals can have dependent components. Chapter~\ref{ch:geometry} develops these distinctions.
\end{rem}

\section{Factor graphs and computation}
\label{sec:ch9-factor-graphs}
\label{sec:ch9-complexity}

\begin{defn}[Factor graph]
\label{defn:factor-graph}
A factor graph records variables and the terms of an energy $E(x)=\sum_a E_a(x_a)$. Its Gibbs law, when normalizable, has the product form $p(x)=Z^{-1}\prod_a e^{-\beta E_a(x_a)}$. This product representation is valid with or without cycles.
\end{defn}

\begin{prop}[A sufficient independence condition]
\label{prop:tree-factorization}
If core and decoration variables have a product configuration space and product reference measure, and $E(c,d)=E_c(c)+E_d(d)$ with finite partition functions, then their Gibbs law factors as $p_C(c)p_D(d)$.
\end{prop}
\begin{proof}
The integrand and measure split, so the partition function is $Z_CZ_D$ and the normalized density factors.
\end{proof}

\begin{rem}[Trees can contain strong coupling]
\label{rem:loop-corrections}
For $c,d\in\{-1,1\}$, the energy $E(c,d)=-Jcd$ has a single pair factor and a tree factor graph, yet it couples the variables whenever $J\ne0$. Tree structure can make marginalization easier under suitable state-space assumptions; it does not eliminate a boundary term. Cycles alone do not quantify its magnitude.
\end{rem}

\begin{prop}[Cost of assigning explicit terms]
\label{thm:polynomial-factorization}
Given the core atom mapping and an explicit list of $T$ evaluable energy terms, assigning and summing those terms takes $O(T)$ operations when each term is evaluated in constant time. For all ligand--receptor and intraligand pairs, $T=O(|m||P|+|m|^2)$. This does not include finding the mapping, optimizing the pose, or equilibrating an ensemble.
\end{prop}

\begin{rem}[Scope of a complexity claim]
\label{rem:hardness}
\label{cor:structure-tractability}
Evaluating a supplied decomposition, choosing a decomposition, minimizing an energy, and estimating a partition function are different computational problems. No hardness claim for one follows merely from the others. Pairwise additivity, locality, and tree structure are useful assumptions in particular algorithms, not universal necessary conditions for tractable score evaluation.
\end{rem}

\section{Variance accounting}
\label{sec:ch9-unified}

\begin{defn}[Residual variance ratio]
\label{defn:factorization-quality}
For a chosen decomposition and a fiber with positive score variance, define
\[
\eta(s)=\frac{\operatorname{Var}(B\mid S=s)}{\operatorname{Var}(F\mid S=s)}.
\]
This ratio is nonnegative and can exceed one. It is not a fraction of explained variance unless an orthogonal decomposition supplies that interpretation.
\end{defn}

\begin{rem}[No universal ordering by scoring method]
\label{thm:factorization-hierarchy}
Docking, end-point, and alchemical responses need not have ordered residual variance ratios. They use different score functions, sampling laws, and residual conventions. Compare the ratios empirically under explicit matched definitions.
\end{rem}

\begin{prop}[Exact conditional variance accounting]
\label{thm:physical-variance-spectrum}
If $F=a(S)+Y+B$ is square-integrable, then
\begin{align*}
\operatorname{Var}(F\mid S)
&=\operatorname{Var}(Y\mid S)+\operatorname{Var}(B\mid S)
+2\operatorname{Cov}(Y,B\mid S),\\
\operatorname{Var}(F)
&=\operatorname{Var}(\mathbb E[F\mid S])+
\mathbb E[\operatorname{Var}(F\mid S)].
\end{align*}
Between-scaffold variance need not equal $\operatorname{Var}(a(S))$, because $\mathbb E[Y+B\mid S]$ can vary with $S$.
\end{prop}

\begin{cor}[What a small projection residual establishes]
\label{cor:spectral-diagnostic}
If $\Err(\Scaf)=\mathbb E[\operatorname{Var}(F\mid S)]<\epsilon\operatorname{Var}(F)$, scaffold identity explains more than $1-\epsilon$ of squared-error variance under this sampling law. This gives no standalone guarantee about tail retrieval, response transfer, or a proposed decomposition's residual.
\end{cor}

\section{Experiments that distinguish the claims}
\label{sec:ch9-experiments}

\begin{exper}[Partition identity and reuse error]
\label{exper:docking-decomp}
\label{exper:mmgbsa-decomp}
Verify numerically that an implemented term partition reconstructs the full score at identical inputs. Separately hold out analogs and measure the error of reusing reference core or solvation terms. Record coordinates, atom parameters, ensemble provenance, and all terms retained in the residual. Passing the identity check does not establish reuse accuracy.
\end{exper}

\begin{exper}[Thermodynamic closure versus response transfer]
\label{exper:fep-closure}
On a set of well-defined thermodynamic states, compare independent edge estimates and their cycle sums with uncertainty. Separately evaluate a complete scaffold--decoration rectangle and estimate its mixed response difference. An exact cycle sum is zero for any state function; a mixed response difference can be nonzero. Test the additive-model reduction against the latter, not the former.
\end{exper}

\begin{exper}[Locality, cumulants, and failure detection]
\label{exper:boundary-decay}
\label{exper:scaffold-cumulants}
\label{exper:failure-detection}
Choose valid matched substitution combinations before looking at scores. Compare direct cross-interaction decay at fixed geometry with the mixed differences after relaxation and after ensemble averaging. Measure whether pose or environment changes explain failures of the fixed-geometry approximation. Evaluate any proposed locality or reuse threshold on held-out combinations, with no preset success rate.
\end{exper}

\section*{Exercises}

\begin{exer}[Partition a score]
\label{exer:ch9-chem}
\label{exer:ch9-comp}
Given an explicit pairwise score and mapped scaffold atoms, assign every term once and verify reconstruction. List which parameters or coordinates must be fixed before any term can be reused across analogs.
\end{exer}

\begin{exer}[RMSD and Lipschitz bounds]
\label{exer:ch9-math-1}
Prove the core stability bound using Cauchy--Schwarz. Give coordinates for which one atom moves by more than the RMSD, and explain why minimizing RMSD by independent ligand alignment changes the physical comparison.
\end{exer}

\begin{exer}[Truncation needs a total bound]
\label{exer:ch9-math-2}
Prove the geometric tail estimate. Explain why a bound on each term does not imply the same bound on the total order-$n$ contribution when the number of terms grows with $n$.
\end{exer}

\begin{exer}[A constant core and a variable residual]
\label{exer:ch9-ml}
Fix one scaffold and choose random decorations $D$. Construct $F=a+Y+B$ with $B$ nonzero and varying, and verify $I(a;Y)=0$. Give a second example with $B=0$ but dependent core and decoration contributions under an across-scaffold sampling law.
\end{exer}

\begin{exer}[Coupling on a tree]
\label{exer:ch9-theory}
For $E(c,d)=-Jcd$ on $\{-1,1\}^2$, compute the Gibbs distribution and show that it is independent exactly when $J=0$ at positive inverse temperature. Compare the factor graph with the number of energy terms.
\end{exer}

\begin{exer}[FEP counts and a response rectangle]
For a $K\times L$ table, derive all four counts in Theorem~\ref{thm:combinatorial-reduction}. Add a single nonzero interaction at one table entry. Show that thermodynamic loops still close while reference-row reconstruction fails there.
\end{exer}
% END SOURCE: chapters/ch09_energy.tex


\chapter{Information Geometry}
\label{ch:geometry}
% BEGIN SOURCE: chapters/ch10_geometry.tex
% ======================================================================
% Chapter 10 content file: chapters/ch10_geometry.tex
% Information Geometry
% (Do NOT include \chapter{...} here; main.tex already does that.)
% ======================================================================

\begin{quote}
\textit{A scaffold\index{scaffold} is not close to another because it looks similar, but because it induces a similar distribution of outcomes.}
\end{quote}

\section{Beyond Structural Similarity}
\label{sec:ch10-intro}

The preceding chapters developed two complementary views of scaffold\index{scaffold} abstraction. Chapter~\ref{ch:variance} gave the \emph{Hilbert-space} view: the conditional expectation\index{conditional expectation} $A(s) = \E[F \mid \Scaf = s]$ is an orthogonal projection, and the variance spectrum\index{variance spectrum} quantifies information content at each refinement level. Chapter~\ref{ch:energy} gave the \emph{physical} view: energy decomposition into core, decoration, and residual terms identifies assumptions that could support reuse across analogs.

This chapter supplies a third, geometric perspective: the \emph{information geometry}\index{information geometry} of scaffold classes. We treat the collection of class distributions $\{p_s : s \in S\}$ as a finite sample from a statistical manifold and ask: what is the natural geometry on this space? Statistical distances describe these laws under a declared molecular sampling measure. Their use for chemical or alchemical decisions requires a separate validation:

\begin{itemize}
    \item \textbf{Scaffold similarity:} When are two scaffolds ``close''? Structural metrics (Tanimoto, MCS) measure chemical resemblance; statistical distances measure \emph{functional} similarity---closeness in the space of binding outcomes.

    \item \textbf{Scaffold hopping\index{scaffold hopping}:} Can an empirically calibrated edge-cost model help select alchemical comparisons? A shortest path minimizes the supplied costs; score-distribution geometry alone does not establish those costs.

    \item \textbf{Factorization quality:} Mutual information measures statistical dependence between specified random variables. We distinguish that dependence from energy nonadditivity and show why a scaffold-constant core has zero conditional mutual information.
\end{itemize}

\begin{framed}
\noindent\textbf{Why geometry matters for drug discovery.}

\medskip
\noindent
Standard cheminformatics treats scaffold similarity structurally: two scaffolds are similar if they share substructures, have high Tanimoto similarity on fingerprints, or have small graph edit distance. But structural similarity is a poor proxy for \emph{functional} similarity.

Consider: a pyridine and a phenyl ring are structurally similar (both aromatic 6-membered rings), but if one accepts a hydrogen bond from a key protein residue while the other does not, their binding score distributions may be radically different. Conversely, structurally dissimilar scaffolds may induce nearly identical score distributions if they happen to make equivalent interactions.

Information geometry\index{information geometry} lets us define similarity by \emph{what matters}: the distribution of outcomes under the scoring function\index{scoring function}. This is not a replacement for structural metrics but a complement---and for many applications, the more relevant one.
\end{framed}

\begin{rem}[Chapter prerequisites and scope]
\label{rem:ch10-prerequisites}
This chapter is marked as optional/advanced. It assumes familiarity with:
\begin{itemize}
    \item Basic probability (Chapter~\ref{ch:probability})
    \item Energy decomposition (Chapter~\ref{ch:energy})
    \item Elementary differential geometry (tangent spaces, metrics) is helpful but not essential
\end{itemize}
Readers seeking only the main algorithmic content may skip to the key results: Definition~\ref{defn:statistical-distance} (statistical distances), Theorem~\ref{thm:mi-geometry} (MI as geometric deviation), and Proposition~\ref{prop:scaffold-hopping-geodesic} (scaffold hopping\index{scaffold hopping} paths). The Fisher-Rao geometry sections (\S\ref{sec:ch10-fisher-rao}) can be skipped on first reading.
\end{rem}

%------------------------------------------------------------------
\section{Structural versus Statistical Metrics}
\label{sec:ch10-structural-vs-statistical}

We begin by contrasting the two paradigms for measuring scaffold similarity.

\subsection{Structural Metrics}

\begin{defn}[Common structural metrics]
\label{defn:structural-metrics}
For scaffolds $s, t \in S$:
\begin{enumerate}
    \item \textbf{Tanimoto similarity:}
    \[
    \text{Tan}(s, t) = \frac{|F_s \cap F_t|}{|F_s \cup F_t|}
    \]
    where $F_s$ is the fingerprint bit-vector of scaffold $s$. The corresponding distance is $d_{\text{Tan}}(s,t) = 1 - \text{Tan}(s,t)$.

    \item \textbf{Maximum common substructure (MCS):}
    \[
    d_{\text{MCS}}(s, t) = 1 - \frac{|MCS(s,t)|}{\max(|s|, |t|)}
    \]
    where $|MCS(s,t)|$ is the size of the largest common subgraph.

    \item \textbf{Graph edit distance:}
    \[
    d_{\text{GED}}(s, t) = \min_{\text{edit sequence}} \sum_i c_i
    \]
    the minimum cost to transform $s$ into $t$ via insertions, deletions, and substitutions.
\end{enumerate}
\end{defn}

These metrics capture \emph{chemical resemblance}: scaffolds that look alike, share functional groups, or can be interconverted with few edits are considered similar.

\begin{exmp}[Structural similarity misses function]
\label{ex:structural-misses-function}
Consider three scaffolds binding to a kinase:
\begin{enumerate}
    \item $s_1$: Pyrimidine (makes two H-bonds to hinge)
    \item $s_2$: Pyridine (makes one H-bond to hinge)
    \item $s_3$: Thieno[2,3-d]pyrimidine (makes two H-bonds, different geometry)
\end{enumerate}

Structural distances:
\begin{itemize}
    \item $d_{\text{Tan}}(s_1, s_2)$ is small (both simple heterocycles)
    \item $d_{\text{Tan}}(s_1, s_3)$ is moderate (fused vs.\ single ring)
\end{itemize}

But functionally:
\begin{itemize}
    \item $s_1$ and $s_3$ may have similar score distributions (both make equivalent H-bonds)
    \item $s_2$ has a shifted distribution (weaker binding due to missing H-bond)
\end{itemize}

Structural distance ranks $(s_1, s_2)$ as more similar than $(s_1, s_3)$; functional similarity may reverse this ranking.
\end{exmp}

\subsection{Statistical Metrics}

\begin{defn}[Scaffold class law]
\label{defn:scaffold-class-law}
For scaffold $s \in S$, the \emph{class law} is:
\[
p_s := \mathrm{Law}(F \mid \Scaf = s)
\]
the probability distribution of scores among molecules sharing scaffold $s$. This is the class distribution from Chapter~\ref{ch:probability}.
\end{defn}

\begin{defn}[Scaffold statistical distance]
\label{defn:statistical-distance}
For scaffolds $s, t \in S$ with class distributions $p_s, p_t$, the statistical comparison induced by a divergence $\mathcal{D}$ is:
\[
d_{\mathcal{D}}(s, t) := \mathcal{D}(p_s, p_t)
\]
\end{defn}

Even when $\mathcal D$ is a metric on distributions, its pullback to scaffolds is generally only a pseudometric: distinct scaffolds can have identical class laws. The sampling distribution within each fiber is part of the definition. The choice of $\mathcal D$ depends on the application.

\begin{defn}[Statistical divergences and distances]
\label{defn:divergences}
For probability distributions $p, q$ on $\R$:
\begin{enumerate}
    \item \textbf{Kullback-Leibler\index{KL divergence} divergence:}
    \[
    D_{\text{KL}}(p \| q) = \int p(x) \log \frac{p(x)}{q(x)} \, dx
    \]
    Note: not symmetric, not a metric.

    \item \textbf{Jensen-Shannon divergence:}
    \[
    D_{\text{JS}}(p, q) = \frac{1}{2} D_{\text{KL}}(p \| m) + \frac{1}{2} D_{\text{KL}}(q \| m)
    \]
    where $m = (p + q)/2$. Symmetric, bounded by $\log 2$.

    \item \textbf{Hellinger distance:}
    \[
    H(p, q) = \sqrt{\frac{1}{2} \int \left(\sqrt{p(x)} - \sqrt{q(x)}\right)^2 dx}
    \]
    Symmetric, bounded by 1.

    \item \textbf{Wasserstein-$p$ distance:}
    \[
    W_p(p, q) = \left( \inf_{\gamma \in \Gamma(p,q)} \int |x - y|^p \, d\gamma(x,y) \right)^{1/p}
    \]
    where $\Gamma(p,q)$ is the set of couplings; take $p\geq1$ and finite $p$th moments. For 1D distributions, this simplifies.

    \item \textbf{Maximum Mean Discrepancy (MMD):}
    \[
    \text{MMD}(p, q) = \sup_{\|f\|_{\mathcal{H}} \leq 1} \left| \E_{p}[f(X)] - \E_{q}[f(X)] \right|
    \]
    for a reproducing kernel Hilbert space $\mathcal{H}$. A characteristic kernel is required for zero MMD to imply equality of distributions.
\end{enumerate}
\end{defn}

\begin{prop}[Wasserstein-2 for 1D distributions]
\label{prop:wasserstein-1d}
For 1D distributions with finite second moments, CDFs $F_p,F_q$, and generalized quantile functions $F_p^{-1},F_q^{-1}$:
\[
W_2(p, q) = \sqrt{\int_0^1 \left( F_p^{-1}(u) - F_q^{-1}(u) \right)^2 du}
\]
This is the $L^2$ distance between quantile functions.
\end{prop}

\begin{proof}
For $U\sim\mathrm{Unif}(0,1)$, the pair $(F_p^{-1}(U),F_q^{-1}(U))$ is a coupling with the required marginals. Monotone quantile coupling minimizes squared transport cost on the line, giving the displayed integral. Using a shared uniform variable also handles atoms, for which a deterministic transport map from $p$ to $q$ need not exist.

\end{proof}

%------------------------------------------------------------------
\section{Which Distance for Which Question}
\label{sec:ch10-which-distance}

Different screening\index{screening} objectives call for different distance measures.

\begin{center}
\begin{tabular}{p{5cm}p{6cm}}
\toprule
\textbf{Screening\index{screening} Objective} & \textbf{Candidate Statistic} \\
\midrule
Hit-finding (extreme tails) & Threshold probabilities, selected quantiles, restricted quantile integrals \\
Mean ranking & Mean difference with uncertainty \\
Robustness under multimodality & Jensen-Shannon, MMD \\
Computational simplicity & Hellinger (closed form for many families) \\
Metric properties required & $W_p$, Hellinger (true metrics) \\
\bottomrule
\end{tabular}
\end{center}

\begin{rem}[Tails versus bulk]
\label{rem:tails-vs-bulk}
For hit-finding, we care about the lower tail of the score distribution (assuming lower is better). Standard distances like KL divergence\index{KL divergence} weight the entire distribution; modifications are needed to emphasize tails.

\emph{Tail-weighted Wasserstein:} Define
\[
W_2^{(\tau)}(p, q) = \sqrt{\int_0^{\alpha_\tau} (F_p^{-1}(u) - F_q^{-1}(u))^2 \, du}
\]
where $0<\alpha_\tau<1$ is fixed in advance for all comparisons, for example from a declared reference distribution. This is a pseudometric on quantile restrictions; it does not distinguish differences outside that tail.

\emph{Quantile distance:} Compare specific quantiles:
\[
d_{q_\alpha}(p, q) = |F_p^{-1}(\alpha) - F_q^{-1}(\alpha)|
\]
for fixed $\alpha$ (e.g., $\alpha = 0.01$ for top 1\%).
\end{rem}

\begin{exmp}[Ranking by a declared hit threshold]
\label{ex:kinase-distance}
If a hit is defined by $F\leq\tau$, the relevant scaffold quantity is $p_s(( -\infty,\tau])$. Estimate and rank these probabilities under the same within-fiber sampling law, with uncertainty. A docking-score threshold is a computational criterion; it cannot be labeled a nanomolar affinity threshold without experimental calibration. Distributional distances may support a model for these probabilities but need not preserve their ordering.
\end{exmp}

%------------------------------------------------------------------
\section{Empirical Sanity Checks}
\label{sec:ch10-sanity-checks}

Before trusting a statistical distance, we should verify its behavior on real data.

\begin{defn}[Metric properties]
\label{defn:metric-properties}
A distance $d: S \times S \to [0, \infty)$ is a \emph{metric} if:
\begin{enumerate}
    \item $d(s, t) = 0 \iff s = t$ (identity of indiscernibles)
    \item $d(s, t) = d(t, s)$ (symmetry)
    \item $d(s, u) \leq d(s, t) + d(t, u)$ (triangle inequality)
\end{enumerate}
\end{defn}

\begin{rem}[Non-metric divergences]
\label{rem:non-metric}
KL divergence\index{KL divergence} fails symmetry and the triangle inequality. Using $d_{\text{KL}}(s,t) = D_{\text{KL}}(p_s \| p_t)$ directly as a ``distance'' can lead to inconsistent rankings: we might have $d(A, B) < d(A, C)$ but $d(B, A) > d(C, A)$.

For scaffold similarity, consider symmetric distances ($\sqrt{\mathrm{JS}}$, Hellinger, $W_p$) unless there's a specific asymmetric interpretation (e.g., $D_{\text{KL}}(p_s \| p_{\text{reference}})$ measures deviation from a reference).
\end{rem}

\begin{exper}[Triangle inequality verification]
\label{exper:triangle}
\textbf{Protocol:}
\begin{enumerate}
    \item Compute $d_{\mathcal{D}}(s, t)$ for all scaffold pairs in a dataset
    \item For each triple $(s, t, u)$, check: $d(s, u) \leq d(s, t) + d(t, u)$?
    \item Report violation frequency and maximum violation magnitude
\end{enumerate}

\textbf{Checks under the stated definitions:}
\begin{itemize}
    \item Wasserstein, Hellinger: zero violations (true metrics)
    \item JS: zero violations (metric after taking square root)
    \item KL: violations are possible; their frequency depends on the data
\end{itemize}
\end{exper}

\begin{exper}[Stability under subsampling]
\label{exper:stability}
\textbf{Protocol:}
\begin{enumerate}
    \item For scaffold $s$, split molecules into two halves: $\Fib(s) = A \cup B$
    \item Estimate $\hat{p}_s^{(A)}$ from $A$, $\hat{p}_s^{(B)}$ from $B$
    \item Compute $d_{\mathcal{D}}(\hat{p}_s^{(A)}, \hat{p}_s^{(B)})$ and quantify its variation across repeated splits
    \item Compare to $d_{\mathcal{D}}(\hat{p}_s, \hat{p}_t)$ for different scaffolds $t$
\end{enumerate}

\textbf{Interpretation:} If within-scaffold distance (due to sampling noise) is comparable to between-scaffold distance, the metric lacks discriminative power for scaffold comparison.
\end{exper}

%------------------------------------------------------------------
\section{The Statistical Manifold of Scaffold Classes}
\label{sec:ch10-manifold}

We now develop the geometric viewpoint more formally.

\begin{defn}[Statistical manifold]
\label{defn:statistical-manifold}
The collection $\{p_s : s \in S\}$ of class distributions can be viewed as a finite sample from a \emph{statistical manifold} $\mathcal{M} \subseteq \mathcal{P}(\R)$, the space of probability measures on $\R$.

If distributions are parameterized by a vector $\theta \in \Theta \subseteq \R^d$, so that $p_s \approx p(\cdot; \theta_s)$, then a regular, identifiable parameterization with nondegenerate Fisher information can define a $d$-dimensional statistical manifold. Arbitrary parameterizations need not: mixture models have singular and nonidentifiable cases.
\end{defn}

\begin{defn}[Parametric surrogate family]
\label{defn:parametric-surrogate}
To work with the geometry of $\{p_s\}$, we often fit a parametric family:
\begin{enumerate}
    \item \textbf{Gaussian:} $p_s \approx \mathcal{N}(\mu_s, \sigma_s^2)$, $\theta_s = (\mu_s, \sigma_s)$
    \item \textbf{Skew-normal:} $p_s \approx \text{SN}(\xi_s, \omega_s, \alpha_s)$, capturing asymmetry
    \item \textbf{Two-component mixture:} $p_s \approx \pi_s \mathcal{N}(\mu_{s,1}, \sigma_{s,1}^2) + (1-\pi_s)\mathcal{N}(\mu_{s,2}, \sigma_{s,2}^2)$
    \item \textbf{Tail-specific models:} Fit to declared extremes or exceedances, with an appropriate tail model and diagnostics
\end{enumerate}

The choice depends on the empirical distribution shapes observed in Chapter~\ref{ch:probability}.
\end{defn}

\begin{rem}[Dimension of the scaffold manifold]
\label{rem:manifold-dimension}
Even if each $p_s$ is parameterized by $d$ parameters, the collection $\{p_s : s \in S\}$ may lie on a lower-dimensional submanifold. For instance, if $\sigma_s \approx c \cdot \mu_s$ (constant coefficient of variation), the effective dimension is reduced.

Empirically determining the ``intrinsic dimension'' of the scaffold distribution space is an interesting open question connecting to manifold learning.
\end{rem}

%------------------------------------------------------------------
\section{Fisher-Rao Geometry}
\label{sec:ch10-fisher-rao}

\emph{This section is optional/advanced. Readers may skip to \S\ref{sec:ch10-dependence} on first reading.}

The Fisher-Rao metric is the natural Riemannian metric on a parametric statistical manifold.

\begin{defn}[Fisher information\index{Fisher information} matrix]
\label{defn:fisher-information}
For a parametric family $\{p(\cdot; \theta) : \theta \in \Theta\}$, the \emph{Fisher information\index{Fisher information} matrix} at $\theta$ is:
\[
g_{ij}(\theta) := \E_{X \sim p(\cdot; \theta)}\left[ \frac{\partial \log p(X; \theta)}{\partial \theta_i} \cdot \frac{\partial \log p(X; \theta)}{\partial \theta_j} \right]
\]
\end{defn}

The Fisher matrix defines a Riemannian metric on $\Theta$: infinitesimal distance between $\theta$ and $\theta + d\theta$ is
\[
ds^2 = \sum_{i,j} g_{ij}(\theta) \, d\theta_i \, d\theta_j
\]

\begin{prop}[Fisher metric for Gaussian family]
\label{prop:fisher-gaussian}
For the Gaussian family $p(x; \mu, \sigma) = \frac{1}{\sqrt{2\pi}\sigma} \exp\left(-\frac{(x-\mu)^2}{2\sigma^2}\right)$, the Fisher metric is:
\[
g = \begin{pmatrix} \sigma^{-2} & 0 \\ 0 & 2\sigma^{-2} \end{pmatrix}
\]
in coordinates $(\mu, \sigma)$.
\end{prop}

\begin{proof}
The log-likelihood is $\log p(x; \mu, \sigma) = -\log(\sqrt{2\pi}\sigma) - \frac{(x-\mu)^2}{2\sigma^2}$.

Partial derivatives:
\begin{align}
\frac{\partial}{\partial \mu} \log p &= \frac{x - \mu}{\sigma^2} \\
\frac{\partial}{\partial \sigma} \log p &= -\frac{1}{\sigma} + \frac{(x-\mu)^2}{\sigma^3}
\end{align}

Computing expectations:
\begin{align}
g_{\mu\mu} &= \E\left[\frac{(X-\mu)^2}{\sigma^4}\right] = \frac{1}{\sigma^4} \cdot \sigma^2 = \frac{1}{\sigma^2} \\
g_{\sigma\sigma} &= \E\left[\left(-\frac{1}{\sigma} + \frac{(X-\mu)^2}{\sigma^3}\right)^2\right] = \frac{2}{\sigma^2} \\
g_{\mu\sigma} &= \E\left[\frac{(X-\mu)}{\sigma^2} \cdot \left(-\frac{1}{\sigma} + \frac{(X-\mu)^2}{\sigma^3}\right)\right] = 0
\end{align}
The cross-term vanishes by symmetry.
\end{proof}

\begin{cor}[Interpretation of Gaussian Fisher metric]
\label{cor:fisher-interpretation}
~
\begin{enumerate}
    \item \textbf{Mean changes are ``cheaper'' when variance is large:} The metric $g_{\mu\mu} = \sigma^{-2}$ means that moving $\mu$ by $\delta$ has local squared line element $\delta^2/\sigma^2$. High variance makes the same mean shift less statistically detectable.

    \item \textbf{Variance changes have constant relative cost:} The metric $g_{\sigma\sigma} = 2\sigma^{-2}$ means that \emph{relative} changes in $\sigma$ (i.e., $d\sigma/\sigma$) have constant cost.
\end{enumerate}
\end{cor}

\begin{defn}[Geodesic distance]
\label{defn:geodesic-distance}
The \emph{geodesic distance} (or Rao distance) between $\theta_1$ and $\theta_2$ is:
\[
d_{\text{FR}}(\theta_1, \theta_2) = \inf_{\gamma} \int_0^1 \sqrt{\sum_{ij} g_{ij}(\gamma(t)) \dot{\gamma}_i(t) \dot{\gamma}_j(t)} \, dt
\]
where the infimum is over smooth paths $\gamma: [0,1] \to \Theta$ with $\gamma(0) = \theta_1$, $\gamma(1) = \theta_2$.
\end{defn}

\begin{prop}[Geodesic distance for 1D Gaussians]
\label{prop:gaussian-geodesic}
For $\sigma_1,\sigma_2>0$, the Fisher--Rao metric on the full Gaussian location-scale family gives
\[
d_{\mathrm{FR}}=\sqrt2\,\operatorname{arcosh}\!\left(
1+\frac{(\mu_1-\mu_2)^2+2(\sigma_1-\sigma_2)^2}{4\sigma_1\sigma_2}
\right).
\]
For equal standard deviations $\sigma$ this is
$2\sqrt2\,\operatorname{arsinh}(|\mu_1-\mu_2|/(2\sqrt2\sigma))$.
The quantity $|\mu_1-\mu_2|/\sigma$ is instead the length of the path constrained to constant $\sigma$, and agrees to first order for small mean shifts.
\end{prop}
\begin{proof}
Set $x=\mu/\sqrt2$ and $y=\sigma$. The line element becomes $2(dx^2+dy^2)/y^2$, twice the upper-half-plane hyperbolic metric. Substitution in its distance formula yields the result.
\end{proof}

\begin{rem}[Fisher vs.\ Wasserstein]
\label{rem:fisher-vs-wasserstein}
Both Fisher-Rao and Wasserstein-2 are Riemannian metrics on the space of distributions. They capture different notions of ``closeness'':
\begin{itemize}
    \item \textbf{Fisher-Rao:} Information-theoretic; measures how distinguishable distributions are from finite samples
    \item \textbf{Wasserstein:} Optimal transport; measures the ``work'' needed to transform one distribution into another
\end{itemize}

For scaffold similarity, Wasserstein is often more interpretable: $W_2(p_s, p_t)$ has units of the score (kcal/mol), while Fisher distance is unitless.
\end{rem}

%------------------------------------------------------------------
\section{Geodesics and Scaffold Interpolation}
\label{sec:ch10-geodesics}

Geodesics provide interpolation between probability distributions. An intermediate law need not be induced by any chemically valid scaffold or permitted chemical path.

\begin{defn}[Geodesic interpolation]
\label{defn:geodesic-interpolation}
For scaffolds $s, t$ with class distributions $p_s, p_t$, the \emph{geodesic interpolation} at parameter $\lambda \in [0,1]$ is the distribution $p_\lambda$ along the geodesic from $p_s$ to $p_t$.
\end{defn}

\begin{prop}[Wasserstein geodesics for 1D distributions]
\label{prop:wasserstein-geodesic}
For distributions with finite second moments, let $U$ be uniform on $(0,1)$ and define
\[
X_\lambda=(1-\lambda)F_p^{-1}(U)+\lambda F_q^{-1}(U).
\]
The law of $X_\lambda$ is a Wasserstein-2 geodesic. Equivalently, its generalized quantile function is
\[
F_{p_\lambda}^{-1}(u)=(1-\lambda)F_p^{-1}(u)+\lambda F_q^{-1}(u)
\]
almost everywhere. This formulation includes distributions with atoms.
\end{prop}

\begin{exmp}[Interpolating scaffold distributions]
\label{ex:interpolating-scaffolds}
Let $s_1$ have distribution $\mathcal{N}(-8, 1^2)$ and $s_2$ have distribution $\mathcal{N}(-6, 2^2)$, with the second parameter denoting variance. The Wasserstein geodesic at $\lambda = 0.5$ is:
\[
p_{0.5} = \mathcal{N}(-7, 1.5^2)
\]
(For Gaussians, Wasserstein geodesics are Gaussians with linearly interpolated parameters.)

This intermediate distribution represents a ``hypothetical scaffold'' with properties midway between $s_1$ and $s_2$. If a scaffold with this score law exists, its chemical accessibility and alchemical utility still require separate evaluation.
\end{exmp}

%------------------------------------------------------------------
\section{Geometry as Refinement Diagnostic}
\label{sec:ch10-refinement-diagnostic}

The geometry of $\{p_s\}$ can diagnose whether the scaffold partition is well-aligned with score structure.

\begin{defn}[Local distortion]
\label{defn:local-distortion}
For scaffold $s$ with neighbors $N(s)$ in the scaffold poset\index{poset}, define:
\[
\Delta(s) := \max_{t \in N(s)} \frac{d_{\mathcal{D}}(p_s, p_t)}{d_{\text{graph}}(s, t)}
\]
where $d_{\text{graph}}$ is the graph distance in the scaffold Hasse diagram.
\end{defn}

High local distortion indicates that moving one step in scaffold space causes a large change in the score distribution---suggesting the partition boundary passes through a region of high ``score gradient.''

\begin{rem}[Distortion as a hypothesis for refinement]
\label{prop:distortion-refinement}
A large edge difference identifies neighboring classes with different score laws. It does not prove that an intermediate valid scaffold exists or that splitting either class will reduce prediction error. The gain of a proposed refinement must be computed from its conditional expectations under the same sampling law, as in Chapter~\ref{ch:refinement}.
\end{rem}

%------------------------------------------------------------------
\section{Dependence, Coupling, and Factorization}
\label{sec:ch10-dependence}

Energy decomposition and statistical independence concern different objects. An energy residual is a function under a chosen partition; independence is a property of two variables under a chosen probability law.

\begin{defn}[Core and decoration random variables]
\label{defn:core-decor-rv}
Draw a molecule and, where relevant, a simulation configuration from a declared law. Set $S=\Scaf(M)$ and write $F=X+Y+B$ for the chosen core, decoration, and residual scores. In the reusable-core model, $X=a(S)$. In a pose-resolved model, $X$ may also depend on sampled core and receptor coordinates. These definitions must not be interchanged.
\end{defn}

\begin{thm}[Mutual information as KL divergence]
\label{thm:mi-geometry}
For random variables $X,Y$, define
\[
I(X;Y)=D_{\mathrm{KL}}(P_{XY}\|P_X\otimes P_Y).
\]
It is nonnegative and is zero exactly when $X$ and $Y$ are independent. When the appropriate finite entropies exist, this agrees with $H(X)+H(Y)-H(X,Y)$; the KL definition also handles a constant variable without using differential-entropy subtraction.
\end{thm}

\begin{cor}[Conditional information of a fixed core]
\label{cor:mi-interpretation}
\label{defn:coupling-index}
If $X=a(S)$, then
\[
\kappa(s):=I(X;Y\mid S=s)=0
\]
for each $s$ of positive probability. This remains true for arbitrarily large or variable $B$. Across-scaffold MI can instead reflect differences in the decoration sampling law across fibers.
\end{cor}
\begin{proof}
At fixed $s$, the joint distribution is $\delta_{a(s)}\otimes P_{Y\mid s}$.
\end{proof}

For example, fix $S=s$, let $Y$ be a fair binary variable, and take $B=100Y$. The residual is large while the fixed core carries no information about $Y$. Conversely, $F=X+Y$ is an additive expression even if $X=Y$ under a correlated sampling scheme. Neither implication between energy additivity and independence holds without additional assumptions.

For a pose-resolved ensemble one can define $\kappa_{\mathrm{pose}}(s)=I(X;Y\mid S=s)$ with genuinely varying $X$. This describes dependence in that ensemble. It is not automatically a bound on scaffold-only prediction error.

\begin{prop}[Gaussian conditional variance identity]
\label{prop:coupling-bounds-error}
If $(X,Y)$ is a nondegenerate jointly Gaussian pair with correlation $|\rho|<1$, then
\[
I(X;Y)=-\tfrac12\log(1-\rho^2),\qquad
\operatorname{Var}(Y\mid X)=\operatorname{Var}(Y)e^{-2I(X;Y)}.
\]
This is the squared-error Bayes risk for predicting $Y$ from $X$. It is not a lower bound for predicting $F$ from scaffold identity.
\end{prop}
\begin{proof}
The Gaussian conditional variance is $\operatorname{Var}(Y)(1-\rho^2)$. Gaussian entropy formulas give the displayed MI; eliminating $\rho$ yields the identity.
\end{proof}

\begin{rem}[MI and residual variance]
\label{rem:mi-vs-eta}
The ratio $\eta(s)=\operatorname{Var}(B\mid s)/\operatorname{Var}(F\mid s)$, when its denominator is positive, is nonnegative. It can exceed one because covariance terms can reduce total variance. MI is invariant under invertible changes to each specified variable, but changing the atom partition or adding information to a variable changes the question. Natural logarithms give MI in nats; base-two logarithms give bits. No universal relationship between $\eta$ and MI is established.
\end{rem}

\section{Environment information without a causal interpretation}
\label{sec:ch10-induced-fit}

\begin{defn}[Environment variable]
\label{defn:environment-variable}
Let $Z$ be a declared observation of the protein/solvent state sampled under a specified protocol. It may be a coordinate representation, a conformational state label, or another measured feature. State what sampling variability remains after conditioning on the scaffold.
\end{defn}

When the quantities are finite, the identity
\begin{equation}
\label{eq:coupling-decomp}
I(X;Y)=I(X;Y\mid Z)+I(X;Y;Z)
\end{equation}
defines the co-information $I(X;Y;Z)$. It can have either sign. It is not generally a nonnegative component of dependence mediated by the environment, and the identity alone supplies no causal direction.

\begin{exmp}[Conditioning can increase dependence]
\label{ex:rigid-vs-flexible}
Let $X,Y$ be independent fair bits and let $Z=X\mathbin{\oplus}Y$. Then $I(X;Y)=0$, while $I(X;Y\mid Z)=\log2$, so $I(X;Y;Z)=-\log2$. If $Z$ is constant, conditional and unconditional MI are equal. The first example disproves a universal decomposition into nonnegative intrinsic and environment-mediated contributions.
\end{exmp}

\begin{rem}[Interpreting environmental associations]
\label{prop:induced-fit-channels}
Quantities such as $I(Y;Z\mid S=s)$ can test associations between a decoration observable and environment state. Their values depend on the observation map and sampling law. To attribute the association to induced fit, compare controlled preparation and simulation protocols and account for common causes and sampling bias. Chapter~\ref{ch:response-certificates} uses declared response contrasts rather than treating MI alone as a causal certificate.
\end{rem}

%------------------------------------------------------------------
\section{Transport on Scaffold Space}
\label{sec:ch10-transport}

Statistical distances induce a weighted graph structure on scaffolds, enabling algorithmic treatment of scaffold hopping.

\begin{defn}[Scaffold graph with statistical weights]
\label{defn:ch10-scaffold-graph}
Define a weighted graph $G = (S, E, w)$ where:
\begin{itemize}
    \item Vertices: scaffolds $S$
    \item Edges: $E$ connects ``adjacent'' scaffolds (e.g., differing by one ring modification)
    \item Weights: $w(s, t) = d_{\mathcal{D}}(p_s, p_t)$ for some statistical distance $\mathcal{D}$
\end{itemize}
\end{defn}

\begin{defn}[Statistical scaffold distance]
\label{defn:scaffold-graph-distance}
The \emph{scaffold path distance} from $s$ to $t$ is:
\[
d_{\text{path}}(s, t) = \min_{\gamma: s \to t} \sum_{i} w(\gamma_i, \gamma_{i+1})
\]
the shortest weighted path through the scaffold graph.
\end{defn}

\begin{prop}[What a shortest-path calculation guarantees]
\label{prop:scaffold-hopping-geodesic}
On a finite graph with nonnegative edge weights, a shortest path minimizes their sum among admissible graph paths. If the weights equal expected additive computational costs of the corresponding transformations, it minimizes that cost model. Choosing weights from score-distribution distances instead optimizes a surrogate whose relationship to alchemical cost requires validation.
\end{prop}

Two scaffolds can have identical scalar score distributions while their coordinate ensembles have poor overlap. Thus zero distributional distance does not certify an easy FEP transformation. Exact free-energy differences telescope along every path; cycle closure residuals of estimated differences measure consistency and sampling errors, not a path-dependent physical boundary energy. See Proposition~\ref{thm:cross-scaffold-mapping} and the response rectangles in Chapter~\ref{ch:response-certificates}.

\begin{algor}[Scaffold hopping path selection]
\label{alg:scaffold-hopping-path}
\textbf{Input:} Source scaffold $s$, target scaffold $t$, statistical distance $d_{\mathcal{D}}$

\textbf{Output:} Candidate path under the supplied cost model $\gamma = (s = s_0, s_1, \ldots, s_k = t)$

\begin{enumerate}
    \item Construct scaffold graph with weights $w(s', t') = d_{\mathcal{D}}(p_{s'}, p_{t'})$
    \item Find shortest path from $s$ to $t$ (Dijkstra's algorithm)
    \item Report the path cost and whether its calibration domain covers every edge
    \item Return path $\gamma$
\end{enumerate}
\end{algor}

%------------------------------------------------------------------
\section{Experiments}
\label{sec:ch10-experiments}

\begin{exper}[Statistical versus structural similarity]
\label{exper:stat-vs-struct}
Fix a target, scoring protocol, and sampling measure within each scaffold. Estimate structural and statistical distances with uncertainty. Evaluate whether statistical distances add held-out predictive value for the declared task. Do not infer shared binding modes from similar one-dimensional score laws.
\end{exper}

\begin{exper}[Distance as a proposed FEP cost predictor]
\label{exper:fep-correlation}
Collect attempted transformations, including failures. Record computational cost, overlap diagnostics, replicate uncertainty, and cycle closure residuals where redundant edges exist. Fit any relation to score-distribution distances on a training set, then assess calibration on held-out transformations. Compare against atom-mapping and structural baselines. A positive correlation and useful path-selection benefit are hypotheses, not assumed outcomes.
\end{exper}

\begin{exper}[Dependence and energy reuse]
\label{exper:coupling-vs-eta}
Begin with the constant-core control, for which conditional MI is exactly zero regardless of boundary variance. Separately evaluate a pose-resolved ensemble with varying core contributions. Estimate dependence with an estimator suitable for the variables, record sampling uncertainty, and compare it with held-out reuse error and residual variance. Test whether it adds predictive information rather than assuming a monotone relationship.
\end{exper}

\section*{Scope of the geometric description}
\label{sec:ch10-summary}
Distributional distances compare class laws under a chosen measure. Fisher--Rao geometry describes a regular parametric model of those laws. Mutual information compares a joint law with its product marginals. These constructions answer different questions; none alone establishes physical response transfer, a causal mechanism, or alchemical sampling efficiency. Their value for those tasks must be tested against the response and cost definitions in Chapter~\ref{ch:response-certificates}.

%------------------------------------------------------------------
\section*{Exercises}

\begin{exer}[Prove MI = KL to product]
\label{exer:ch10-mi-kl}
For a finite joint distribution, derive the entropy form of Theorem~\ref{thm:mi-geometry}:
\begin{enumerate}
    \item[(a)] Write out $D_{\text{KL}}(p_{X,Y} \| p_X \otimes p_Y)$ in integral form.
    \item[(b)] Expand and simplify using the definition of entropy.
    \item[(c)] Show the result equals $H(X) + H(Y) - H(X,Y) = I(X; Y)$.
\end{enumerate}
\end{exer}

\begin{exer}[Fisher metric for Gaussian mixtures]
\label{exer:ch10-fisher-mixture}
Consider the two-component Gaussian mixture:
\[
p(x; \pi, \mu_1, \mu_2, \sigma) = \pi \mathcal{N}(x; \mu_1, \sigma^2) + (1-\pi)\mathcal{N}(x; \mu_2, \sigma^2)
\]
with shared variance $\sigma^2$.
\begin{enumerate}
    \item[(a)] Compute $\partial_\pi \log p(x)$, $\partial_{\mu_1} \log p(x)$, $\partial_{\mu_2} \log p(x)$.
    \item[(b)] Determine which off-diagonal terms can be nonzero, and identify parameter choices at which the Fisher matrix becomes singular.
    \item[(c)] For $\mu_1 = -\mu_2 = \mu$ (symmetric mixture), simplify the Fisher matrix.
\end{enumerate}
\end{exer}

\begin{exer}[Wasserstein-2 for Gaussians]
\label{exer:ch10-wasserstein-gaussian}
Prove that for 1D Gaussians $p = \mathcal{N}(\mu_1, \sigma_1^2)$, $q = \mathcal{N}(\mu_2, \sigma_2^2)$:
\[
W_2(p, q) = \sqrt{(\mu_1 - \mu_2)^2 + (\sigma_1 - \sigma_2)^2}
\]
\begin{enumerate}
    \item[(a)] Use Proposition~\ref{prop:wasserstein-1d} and the quantile function of a Gaussian.
    \item[(b)] Interpret: $W_2$ is Euclidean distance in $(\mu, \sigma)$ space.
    \item[(c)] For scaffold ranking by mean score, when does $W_2$ ordering match mean ordering?
\end{enumerate}
\end{exer}

\begin{exer}[Implement statistical distance comparison]
\label{exer:ch10-implementation}
Write code to:
\begin{enumerate}
    \item[(a)] Compute Tanimoto distance between scaffold fingerprints.
    \item[(b)] Estimate $W_2$ between scaffold score distributions from samples.
    \item[(c)] For a dataset of scaffolds, plot $d_{W_2}$ vs.\ $d_{\text{Tan}}$ and identify outliers.
    \item[(d)] Discuss: which outliers represent opportunities for scaffold hopping?
\end{enumerate}
\end{exer}

\begin{exer}[MI versus boundary variance]
\label{exer:ch10-mi-eta}
Prove the zero-information result for a scaffold-constant core. Construct a residual variance ratio exceeding one by cancellation between $Y$ and $B$. Then give a pose-resolved sampling model with nonzero conditional MI and state why this alone does not determine error in predicting $F$ from scaffold identity.
\end{exer}

%------------------------------------------------------------------
\section{Bibliographic Notes}

Information geometry was developed by Amari and colleagues; the standard reference is Amari \& Nagaoka (2000). The Fisher-Rao metric on statistical manifolds has a long history going back to Rao (1945).

Wasserstein distances and optimal transport are surveyed in Villani (2008); computational aspects are covered in Peyré \& Cuturi (2019). The use of Wasserstein distance for comparing distributions in machine learning has grown substantially.

Related topics for a literature review include:
\begin{itemize}
    \item Activity landscapes in QSAR: Maggiora et al.\ (2006)
    \item Distribution-based molecular similarity: various works on ``fuzzy'' pharmacophores
    \item Information-theoretic drug design: Sheridan \& Kearsley (2002) on mutual information in QSAR
\end{itemize}

Probabilistic graphical models distinguish statistical independence from factorization of an energy expression. Causal claims about environmental mediation need assumptions beyond the MI identities used here.

Scaffold hopping as graph search has been explored computationally (Hu \& Bajorath, 2011); statistical edge weights are a candidate to evaluate against structural and measured-cost alternatives.

%------------------------------------------------------------------
\section{Connections}

\begin{itemize}
    \item \textbf{Chapter~\ref{ch:probability}:} Supplies the class distributions $p_s$.
    \item \textbf{Chapter~\ref{ch:energy}:} Supplies energy decompositions and their residuals; this chapter distinguishes them from statistical dependence.
    \item \textbf{Chapter~\ref{ch:refinement}:} Provides conditional-expectation tests of refinement gain; distributional distortion is only a candidate diagnostic.
    \item \textbf{Chapter~\ref{ch:screening}:} Uses statistical distances for prior structure in bandit\index{bandit algorithm} algorithms.
\end{itemize}
% END SOURCE: chapters/ch10_geometry.tex


\chapter{Response Curvature and Decision Certificates}
\label{ch:response-certificates}
% BEGIN SOURCE: chapters/response_core.tex
% Shared mathematical source for main.tex and response_curvature.tex.
% Edit this file; both documents include the same theorem statements.
\section*{The decision being certified}
\addcontentsline{toc}{section}{The decision being certified}
Fix a finite molecular domain, a response endpoint, and a measurement protocol.
A representation specifies which chemical operations may be compared across
contexts. Its response model is a hypothesis about those comparisons.
A certificate proves a declared ranking or error bound for every response
consistent with that hypothesis and the stated uncertainty set. It does not
establish that an untested response model is true, nor that a computational
score measures biological activity.

The chapter separates three tasks: learning model coordinates, obtaining new
measurements that can reject the model, and deciding whether the remaining
uncertainty can change a molecular choice. The exact linear results are
conditional on a fixed model. Their statistical use additionally requires
simultaneously valid uncertainty bounds, with fresh validation evidence or
appropriate sequential inference after adaptive selection. A model fitted to
all available responses cannot use the same residuals as independent validation.

Read the complete-table identities and decision certificates first; the sparse
graph sections then extend them to missing chemical combinations and repeated
representations of one parent molecule. Path and precision designs specialize
this framework. The retrospective examples at the end test structural identities
and expose failed transfer; they do not establish prospective assay savings.

\section{The obstruction requires two coordinates}

Let $\Sset=\{s_1,\ldots,s_m\}$ denote scaffolds or other chemical contexts, and
let $\Dset=\{d_1,\ldots,d_n\}$ denote one fixed label type.  In a
reaction-feasible implementation, a label records an atom-mapped reaction
template, building block, and attachment-site correspondence.  A structural
substituent or matched-pair transformation is identified with that label only
when its equivalence across contexts is declared before outcomes are read.  The
response $Y_{sd}\in\R$ may be a homogeneous experimental activity, a calibrated
probability of satisfying a contract, or another fixed endpoint.  Smaller
responses are assumed better only in the decision results.

The construction needs a real cross-product.  A column cannot mean ``whatever
group happens to occupy attachment site 3.''  It must identify a transported
operation, such as the same reaction template and building block applied at a
mapped attachment site.  Missing or chemically invalid combinations lead to
the graph formulation in \Cref{sec:graph}.  SAR matrices and retrosynthetically
defined matched pairs already provide closely related structural
organizations\cite{wassermann2012sarm,deleon2014recapmmp}.

\begin{lem}[One-coordinate cycle sums are vacuous]
\label{lem:telescoping}
Fix $d\in\Dset$ and a closed scaffold path
$s_0,s_1,\ldots,s_k=s_0$.  Then every response function satisfies
\[
  \sum_{i=0}^{k-1}\bigl(Y_{s_{i+1}d}-Y_{s_i d}\bigr)=0.
\]
\end{lem}

\begin{proof}
The sum telescopes to $Y_{s_kd}-Y_{s_0d}=0$.
\end{proof}

Thus a sum of raw scaffold differences around a scaffold triangle cannot
measure nonadditivity.  A nontrivial contrast must vary scaffold and decoration.
Chemical double-transformation cycles use this four-state construction
already\cite{kramer2015nonadditivity}; two-way analysis of variance uses the
same interaction contrast\cite{tukey1949nonadditivity}.

\begin{defn}[Rectangle curvature]
For $s,s'\in\Sset$ and $d,d'\in\Dset$, define
\begin{equation}
\label{eq:rectangle}
  \Omega_{ss';dd'}(Y)
  :=Y_{sd}-Y_{sd'}-Y_{s'd}+Y_{s'd'}.
\end{equation}
The sign depends on orientation.  Its magnitude does not.
\end{defn}

The contrast asks whether changing $d'$ to $d$ has the same effect in contexts
$s$ and $s'$.  A nonzero value records context-dependent response.  Assay noise,
batch effects, and mixed endpoint definitions can create the same signal, so a
chemical interpretation requires homogeneous measurements and propagated
uncertainty.

\section{Complete response tables}

\begin{defn}[Additive transfer model]
The table $Y\in\R^{m\times n}$ is additively transferable if there exist
$a\in\R^m$ and $b\in\R^n$ such that
\[
  Y_{sd}=a_s+b_d
  \qquad(s\in\Sset,d\in\Dset).
\]
The parameterization has the gauge freedom $(a,b)\mapsto(a+c\one,b-c\one)$.
\end{defn}

This is the algebraic content of a scaffold main effect plus a transferable
decoration effect.  It is the response-table form of the Free--Wilson
model\cite{free1964qsar}.

\begin{thm}[Exact curvature criterion]
\label{thm:exact}
For a complete response table $Y$, the following statements are equivalent.
\begin{enumerate}
  \item $Y_{sd}=a_s+b_d$ for some $a,b$.
  \item $\Omega_{ss';dd'}(Y)=0$ for all $s,s',d,d'$.
  \item $H_mYH_n=0$, where
  \[
    H_m=I_m-\frac{1}{m}\one\one^\top,
    \qquad
    H_n=I_n-\frac{1}{n}\one\one^\top.
  \]
\end{enumerate}
\end{thm}

\begin{proof}
Additive terms cancel in \eqref{eq:rectangle}, which proves (1)$\Rightarrow$(2).
Assume (2), fix an anchor $(s_0,d_0)$, and set
\[
  a_s=Y_{s d_0},
  \qquad
  b_d=Y_{s_0d}-Y_{s_0d_0}.
\]
The identity $\Omega_{s s_0;d d_0}(Y)=0$ gives $Y_{sd}=a_s+b_d$, proving
(2)$\Rightarrow$(1).  Every matrix in the additive subspace is annihilated by
left and right centering.  Conversely, $H_mYH_n=0$ implies
\[
Y_{sd}=\bar Y_{s\cdot}+\bar Y_{\cdot d}-\bar Y_{\cdot\cdot},
\]
which is additive.  This proves (1)$\Leftrightarrow$(3).
\end{proof}

\begin{defn}[Interaction residual]
Define
\begin{equation}
\label{eq:residual}
  R:=H_mYH_n,
  \qquad
  R_{sd}=Y_{sd}-\bar Y_{s\cdot}-\bar Y_{\cdot d}+\bar Y_{\cdot\cdot}.
\end{equation}
\end{defn}

\begin{thm}[Projection and curvature-energy identity]
\label{thm:energy}
Let
\[
  \Aspace=\{a\one^\top+\one b^\top:a\in\R^m,b\in\R^n\}.
\]
Then
\begin{align}
  \min_{A\in\Aspace}\|Y-A\|_F^2&=\|R\|_F^2,\label{eq:l2best}\\
  \sum_{s,s',d,d'}\Omega_{ss';dd'}(Y)^2&=4mn\,\|R\|_F^2.\label{eq:curvenergy}
\end{align}
Moreover, if $S,S'$ are independent uniform rows and $D,D'$ are independent
uniform columns, then
\begin{equation}
\label{eq:cellaverage}
  R_{sd}=\E\bigl[\Omega_{sS';dD'}(Y)\bigr].
\end{equation}
\end{thm}

\begin{proof}
The additive subspace and the doubly centered subspace are orthogonal in the
Frobenius inner product.  The additive projection has entries
$\bar Y_{s\cdot}+\bar Y_{\cdot d}-\bar Y_{\cdot\cdot}$, so its residual is $R$.
This proves \eqref{eq:l2best}.  Rectangle curvature ignores additive terms, so
replace $Y$ with $R$ in \eqref{eq:rectangle}.  The row means, column means, and
grand mean of $R$ vanish.  Expanding the square and summing removes every cross
term, leaving four copies of $mn\|R\|_F^2$.  Taking the expectation in
\eqref{eq:cellaverage} gives
\[
Y_{sd}-\bar Y_{s\cdot}-\bar Y_{\cdot d}+\bar Y_{\cdot\cdot}=R_{sd}.
\]
\end{proof}

\begin{rem}[Product-weighted version]
Let $p_s,q_d>0$ sum to one, and center with the product measure $p\otimes q$.
For independent $S,S'\sim p$ and $D,D'\sim q$,
\[
  \E\!\left[\Omega_{SS';DD'}(Y)^2\right]
  =4\sum_{s,d}p_sq_dR_{sd}^2,
  \qquad
  R_{sd}=\E\!\left[\Omega_{sS';dD'}(Y)\right].
\]
The proof repeats the zero-conditional-mean argument above.  General
non-product sampling needs a different projection and does not satisfy this
four-copy identity without correction.
\end{rem}

\begin{cor}[Uniform curvature bounds]
Let
\[
  \delta:=\max_{s,s',d,d'}|\Omega_{ss';dd'}(Y)|.
\]
Then
\[
  \|R\|_\infty\le\delta\le4\|R\|_\infty.
\]
Under product weights $p\otimes q$, the sharper cellwise bound is
\[
  |R_{sd}|\le(1-p_s)(1-q_d)\delta.
\]
\end{cor}

\begin{proof}
The first inequality follows from the expectation identity
$R_{sd}=\E[\Omega_{sS';dD'}]$.  The second follows because a rectangle contains
four entries of $R$.  Terms with $S'=s$ or $D'=d$ vanish in the expectation,
which gives the cellwise factor.
\end{proof}

\begin{cor}[Anchor-cross design]
\label{cor:anchor}
Fix $(s_0,d_0)$ and observe the $m+n-1$ cells in row $s_0$ or column $d_0$.
The anchored additive prediction
\begin{equation}
\label{eq:anchor}
  \widetilde Y_{sd}
  :=Y_{s d_0}+Y_{s_0d}-Y_{s_0d_0}
\end{equation}
satisfies
\[
  Y_{sd}-\widetilde Y_{sd}=\Omega_{s s_0;d d_0}(Y).
\]
The anchor cross reconstructs the full table if transfer is additive.  Each of
the remaining $(m-1)(n-1)$ cells supplies one independent fundamental
curvature contrast when measured.
\end{cor}

\begin{proof}
The error identity follows by rearranging \eqref{eq:rectangle}.  The anchor
cross forms a spanning tree of the complete bipartite graph.  Adding any missing
cell closes its unique fundamental rectangle with the anchor tree.  The number
of missing cells equals $mn-(m+n-1)=(m-1)(n-1)$.
\end{proof}

The corollary sets an evidence boundary.  An anchor cross identifies an
additive model; it cannot certify arbitrary unmeasured interactions.  Each
untested cell remains a possible counterexample unless a statistical model or a
physical bound supplies additional information.
Three-neighbor prediction of a missing core--substituent cell is already used in
the SAR-matrix literature under additive core and substituent
contributions\cite{gupta2014neighborhood}.  The role of the anchor construction
here is to expose exactly which held-out cells falsify that assumption.

\section{Interaction modes}

The singular value decomposition of $R$ separates coherent response failures
from diffuse cell noise.

\begin{thm}[Optimal interaction modes]
\label{thm:svd}
Let $R=U\Sigma V^\top$ and let $R_q$ retain its $q$ largest singular values.
Among models consisting of an additive table plus an interaction matrix of rank
at most $q$, the Frobenius-optimal approximation is
\[
  \widehat Y_q
  =\bar Y_{s\cdot}+\bar Y_{\cdot d}-\bar Y_{\cdot\cdot}+R_q,
\]
and
\[
  \min_{A\in\Aspace,\ \rank(Q)\le q}\|Y-A-Q\|_F^2
  =\sum_{k>q}\sigma_k(R)^2.
\]
Also $\rank(R)\le\min(m-1,n-1)$.
\end{thm}

\begin{proof}
For any $Q$, the best additive term removes the additive projection of $Y-Q$.
The remaining error is $R-H_mQH_n$, whose second term has rank at most $q$.
The Eckart--Young theorem gives $R_q$ as the best rank-$q$ approximation.
Centering makes $\one$ a left and right null vector of $R$, which gives the rank
bound.
\end{proof}

The singular vectors define empirical context and decoration coordinates.  A
small interaction rank supports a compact extension of Free--Wilson.  Matrix
completion can estimate a low-rank residual from a subset of cells only under
conditions such as incoherence and suitable sampling\cite{recht2011matrix}.
Medicinal chemistry tables often violate random-sampling assumptions, so the
completion claim must be tested through held-out blocks and time splits.

\subsection{A generically dimension-matched low-rank assay cross}

The anchor construction turns interaction learning into a matrix cross problem.
Fix $(s_0,d_0)$ and index the remaining rows and columns by
$\Sset^\circ=\Sset\setminus\{s_0\}$ and
$\Dset^\circ=\Dset\setminus\{d_0\}$.  Define the anchored curvature matrix
\begin{equation}
\label{eq:anchored-k}
  K_{sd}:=Y_{sd}-Y_{s d_0}-Y_{s_0d}+Y_{s_0d_0},
  \qquad (s,d)\in\Sset^\circ\times\Dset^\circ.
\end{equation}

\begin{prop}[Intrinsic curvature rank]
\label{prop:curvature-rank}
The following quantities are equal:
\[
  \rank(K),
  \qquad
  \rank(R),
  \qquad
  \min\{\rank(L):Y=A+L,\ A\in\Aspace\}.
\]
Thus curvature rank is independent of the anchor and equals the minimum number
of separable interaction modes beyond additive main effects.
\end{prop}

\begin{proof}
Mixed differencing annihilates $A$, so any representation $Y=A+L$ gives
$\rank(K)\le\rank(L)$.  Conversely, let $A$ be the anchored additive table from
\eqref{eq:anchor} and embed $K$ into a matrix $L$ whose anchor row and column are
zero.  Then $Y=A+L$ and $\rank(L)=\rank(K)$.  The centered residual $R$ is
another representative of the same quotient by $\Aspace$.  Row and column
differencing are invertible between the centered interaction coordinates and
the anchored coordinates, so they preserve rank.
\end{proof}

\begin{thm}[Exact rank-$q$ assay cross]
\label{thm:rankq-cross}
Suppose $\rank(K)=q$.  Choose sets
$I\subseteq\Sset^\circ$ and $J\subseteq\Dset^\circ$ with
$|I|=|J|=q$ such that the pivot $P:=K_{I,J}$ is nonsingular.  Observe the
anchor row and column, the $q$ columns indexed by $J$, and the $q$ rows indexed
by $I$.  Then every response is recovered by
\begin{equation}
\label{eq:skeleton}
  K=K_{:,J}P^{-1}K_{I,:},
  \qquad
  Y_{sd}=Y_{s d_0}+Y_{s_0d}-Y_{s_0d_0}+K_{sd}.
\end{equation}
The number of distinct observed cells is
\begin{equation}
\label{eq:rankq-count}
  N_q=m+n-1+q(m+n-2-q).
\end{equation}
\end{thm}

\begin{proof}
The anchor row and column determine the additive term and convert every further
cell into an entry of $K$.  Since $K$ has rank $q$ and $P$ is nonsingular, its
selected columns form a column basis.  Write $K=K_{:,J}X$.  Restriction to rows
$I$ gives $K_{I,:}=PX$, hence $X=P^{-1}K_{I,:}$ and
\eqref{eq:skeleton} follows.  The anchor cross costs $m+n-1$ cells.  The selected
curvature columns and rows contain $q(m-1)+q(n-1)-q^2$ further cells, proving
\eqref{eq:rankq-count}.
\end{proof}

\begin{thm}[Schur-chart curvature certificate]
\label{thm:schur-chart}
Keep an invertible $q\times q$ pivot $P=K_{I,J}$ but do not assume
$\rank(K)=q$.  Reorder the anchored curvature matrix as
\[
  K=\begin{pmatrix}P&C\\ R_0&D\end{pmatrix}
\]
and define its Schur-curvature residual
\begin{equation}
\label{eq:schur-residual}
  S:=D-R_0P^{-1}C.
\end{equation}
Then
\begin{enumerate}
  \item $\rank(K)=q+\rank(S)$;
  \item the skeleton predictor has zero error on its observed rows and columns
  and error matrix $S$ on the complement;
  \item for each held-out cell $(s,d)$,
  \begin{equation}
  \label{eq:minor-certificate}
    S_{sd}
    =\frac{\det K_{I\cup\{s\},J\cup\{d\}}}{\det P},
  \end{equation}
  up to the declared row and column orientation;
  \item the exact uniform skeleton error is $\|S\|_\infty$.
\end{enumerate}
On the chart $\det P\ne0$, the skeleton values together with $S$ form rational
coordinates for the full response table.  The model $\rank(K)=q$ is the
coordinate plane $S=0$.
\end{thm}

\begin{proof}
Block Gaussian elimination gives
\[
  \begin{pmatrix}I&0\\-R_0P^{-1}&I\end{pmatrix}
  K
  \begin{pmatrix}I&-P^{-1}C\\0&I\end{pmatrix}
  =\begin{pmatrix}P&0\\0&S\end{pmatrix},
\]
which proves the rank identity.  The skeleton reconstruction replaces $D$ by
$R_0P^{-1}C$, so its error is exactly $S$.  Applying the block determinant
formula to the pivot augmented by row $s$ and column $d$ proves
\eqref{eq:minor-certificate}.  The remaining statements follow by reading the
skeleton and complementary coordinates from the block form.
\end{proof}

The complement contains
\[
  (m-q-1)(n-q-1)=mn-N_q
\]
independent local falsification coordinates.  This generalizes one rectangle
test per additive chord: $N_q$ cells identify a $q$-mode chart, and each
remaining cell buys one higher-curvature Schur test.  A single test uses a
$(q+2)\times(q+2)$ raw-response biclique including the anchors; holding out its
target leaves $(q+2)^2-1$ measured cells.  The equality in
\eqref{eq:minor-certificate} becomes an empirical certificate only after the
target cell is measured or its magnitude is bounded on untouched circuits.

\begin{thm}[Generic local minimality]
\label{thm:rankq-minimal}
The smooth stratum of response tables with interaction rank exactly $q$ has
dimension
\[
  m+n-1+q(m+n-2-q).
\]
Consequently, a fixed design observing fewer than $N_q$ scalar cells cannot be
locally identifying on a generic rank-$q$ table.
\end{thm}

\begin{proof}
The additive quotient has dimension $m+n-1$.  An
$(m-1)\times(n-1)$ matrix of rank $q$ has dimension
$q((m-1)+(n-1)-q)$.  These coordinates are independent because the anchored
curvature map annihilates the additive subspace.  A scalar-cell observation map
into $\R^N$ has Jacobian rank at most $N$.  If $N<N_q$, its differential has a
nontrivial kernel on the smooth rank-$q$ stratum, so the map cannot be locally
injective at a generic point.
\end{proof}

\begin{thm}[Uniform cell-query no-free-lunch]
\label{thm:rankq-nofree}
Let $q\ge1$.  Any deterministic cell-query policy that guarantees exact recovery
for every table with curvature rank at most $q$ must query all $mn$ cells.  The
same lower bound holds for a guaranteed exact optimum or top-$k$ set when the
unqueried response can have either sign.
\end{thm}

\begin{proof}
Run the policy on the zero table.  If it stops before querying every cell, choose
an unqueried cell $e$.  The zero table and $\lambda E_e$ agree on the entire
query history, so the deterministic policy follows the same path and returns
the same answer.  Yet $\lambda E_e$ has rank one and hence curvature rank at
most one.  Choosing the sign and magnitude of $\lambda$ changes the optimum and
can change top-$k$.  Therefore the policy fails on one of the two tables.
\end{proof}

The generic count $N_q$ and the uniform count $mn$ address different claims.
The smaller design works on a rank-revealing chart with a nonsingular pivot.
Uniform savings require spread or incoherence assumptions that exclude hidden
single-cell interactions.  Sparsity alone gives no savings for ordinary cell
queries because an unmeasured spike remains invisible.

\begin{prop}[Noisy Schur stability]
\label{prop:schur-stability}
For a predicted curvature entry $k=pP^{-1}c$, suppose the measured quantities
are $p+e_p$, $c+e_c$, and $P+E_P$, with
\[
  \|e_p\|\le\eta_p,
  \quad \|e_c\|\le\eta_c,
  \quad \|E_P\|\le\eta_P<\gamma:=\sigma_{\min}(P).
\]
Then
\begin{align*}
|\widehat k-k|
\le{}&
\frac{\eta_p(\|c\|+\eta_c)+\|p\|\eta_c}{\gamma-\eta_P}\\
&+\frac{\|p\|\|c\|\eta_P}{\gamma(\gamma-\eta_P)}.
\end{align*}
\end{prop}

\begin{proof}
Add and subtract terms in
$(p+e_p)(P+E_P)^{-1}(c+e_c)-pP^{-1}c$, use
\[
(P+E_P)^{-1}-P^{-1}=-(P+E_P)^{-1}E_PP^{-1},
\]
and apply
$\|(P+E_P)^{-1}\|\le(\gamma-\eta_P)^{-1}$.
\end{proof}

If every raw skeleton-cell mean has error at most $\delta$, each computed
curvature entry has error at most $4\delta$.  One may take
$\eta_p,\eta_c\le4\sqrt q\,\delta$ and $\eta_P\le4q\delta$, which requires
$4q\delta<\sigma_{\min}(P)$.  This makes pivot conditioning part of the
certificate rather than an implementation detail.  If the true rank is below
$q$, every $q\times q$ pivot is singular and one must switch to a lower-rank
chart; rank-selection inference lies on a singular model stratum.

For $q=0$, \eqref{eq:rankq-count} gives the anchor cross $m+n-1$.  If
$q=\min(m-1,n-1)$, it gives $mn$, so unrestricted interaction requires the full
table.  The identity in \eqref{eq:skeleton} is the classical skeleton or CUR
decomposition\cite{goreinov1997pseudoskeleton}.  Here it is specialized to
separate additive assay cost from interaction-rank cost; no novelty is claimed
for the skeleton identity, dimension count, or determinantal circuits.  A
compatibility mask must contain the anchor cross and the chosen pivot rows and
columns.  Reaction constraints can make that pattern unavailable.

Noise is amplified by $P^{-1}$.  A prospective design should seek a
well-conditioned or high-volume pivot and report its smallest singular value.
An adaptive search for the pivot must spend measurements to establish rank; the
count in \eqref{eq:rankq-count} assumes that suitable $I,J$ have been identified.

\section{Decision certificates}

The response residual becomes useful when it controls a decision.  The next
results apply to any approximation, including the additive projection, an
anchored predictor, or an additive-plus-mode model.

\begin{thm}[Simple-regret and threshold certificate]
\label{thm:regret}
Let $\widehat Y$ satisfy
\[
  \|Y-\widehat Y\|_\infty\le\eps
\]
on a declared finite candidate domain.  If smaller responses are better and
$\hat x\in\arg\min_x\widehat Y_x$, then
\[
  Y_{\hat x}-\min_xY_x\le2\eps.
\]
For a hit threshold $\tau$,
\begin{align*}
  \widehat Y_x\le\tau-\eps&\implies Y_x\le\tau,\\
  \widehat Y_x>\tau+\eps&\implies Y_x>\tau.
\end{align*}
\end{thm}

\begin{proof}
Let $x^*$ minimize $Y$.  Then
\[
Y_{\hat x}\le\widehat Y_{\hat x}+\eps
\le\widehat Y_{x^*}+\eps
\le Y_{x^*}+2\eps.
\]
The threshold statements follow from the same uniform bound.
\end{proof}

\begin{cor}[Top-$k$ stability]
\label{cor:topk}
Let $K$ be the indices of the $k$ smallest entries of $\widehat Y$.  If
\[
  \min_{j\notin K}\widehat Y_j
  -\max_{i\in K}\widehat Y_i>2\eps,
\]
then $K$ is the true top-$k$ set of $Y$.
\end{cor}

\begin{proof}
For $i\in K$ and $j\notin K$,
$Y_j-Y_i\ge\widehat Y_j-\widehat Y_i-2\eps>0$.
\end{proof}

Without the separation condition, the true top-$k$ set still satisfies the
one-sided inclusion
\[
  T_k(Y)\subseteq
  \{x:\widehat Y_x\le\widehat Y_{(k)}+2\eps\},
\]
where ties require a declared deterministic rule.  This identifies the smallest
candidate band that a uniform certificate alone can guarantee.

For the SVD model in \Cref{thm:svd}, the conservative bound
\[
\|Y-\widehat Y_q\|_\infty
\le\|Y-\widehat Y_q\|_F
=\left(\sum_{k>q}\sigma_k(R)^2\right)^{1/2}
\]
turns spectral tail energy into a decision certificate.  A held-out or
distributional bound will often be tighter, but it needs a declared sampling
model.

\section{Sparse response graphs}
\label{sec:graph}

Many decorations cannot be installed on every scaffold.  Let
$G=(V,E)$ be the bipartite observation graph with
$V=\Sset\sqcup\Dset$.  An edge $e=(s,d)$ exists when the combination is valid
and measured.  Orient edges from scaffold to decoration.  Let
$B\in\R^{|V|\times|E|}$ be the incidence matrix whose edge column is
$b_e=e_s-e_d$, and collect observed responses in $y\in\R^{|E|}$.

Set $\phi_s=a_s$ and $\phi_d=-b_d$.  The additive model becomes
\[
  y=B^\top\phi.
\]

Matched-series graphs, graph-mined parallel series, and SAR matrices already
organize structurally related compound series
\cite{wawer2011mmsgraph,gupta2012graphmining,wassermann2012sarm}.  The object
here places factor levels at vertices and measured responses on incidence edges
for additive-consistency analysis.

\begin{defn}[Cycle holonomy]
For an oriented simple cycle $C$, let $\chi_C\in\{-1,0,1\}^{E}$ be its signed
edge-incidence vector.  Define
\[
  \kappa_C(y):=\chi_C^\top y.
\]
On a four-cycle this is the rectangle curvature in \eqref{eq:rectangle}, up to
orientation.
\end{defn}

\begin{thm}[Graph transfer criterion]
\label{thm:graphcriterion}
The following statements are equivalent.
\begin{enumerate}
  \item There is a potential $\phi$ such that $y=B^\top\phi$.
  \item $\kappa_C(y)=0$ for every cycle $C$.
  \item $\kappa_C(y)=0$ on any cycle basis.
\end{enumerate}
\end{thm}

\begin{proof}
The cycle space is $\ker B$.  Linear algebra gives
$(\ker B)^\perp=\im B^\top$.  Cycle vectors span $\ker B$, so $y$ lies in
$\im B^\top$ exactly when it has zero inner product with a cycle basis.
\end{proof}

Four-cycles suffice on a complete bipartite graph because they generate its
cycle space.  An incomplete graph may contain chordless cycles of length six or
more.  Testing observed rectangles alone can then miss a global inconsistency.

\begin{thm}[Hodge residual and curvature degrees of freedom]
\label{thm:hodge}
Let $c$ be the number of connected components of $G$.  The least-squares
decomposition is
\[
  y=B^\top\hat\phi+r,
  \qquad
  r\in\ker B,
  \qquad
  \hat\phi=(BB^\top)^+By.
\]
It satisfies
\[
  \|r\|_2^2=\min_\phi\|y-B^\top\phi\|_2^2,
  \qquad
  \dim\ker B=|E|-|V|+c.
\]
\end{thm}

\begin{proof}
Orthogonal projection onto $\im B^\top$ gives the stated decomposition.  The
rank of the incidence matrix is $|V|-c$, so rank--nullity gives the dimension.
\end{proof}

This is the graph-Hodge decomposition used to separate global potentials from
cyclic inconsistency in ranking data\cite{jiang2011hodge}.  In the chemical
response graph, $r$ records the part of the observed table that no additive
scaffold and decoration effects can fit.

\begin{cor}[One independent curvature test per chord]
A forest has no curvature degrees of freedom, and any responses on its edges
admit an additive fit.  Adding an edge within a connected component creates one
new independent cycle and raises $|E|-|V|+c$ by one.
\end{cor}

This explains the anchor-cross design.  The first $m+n-1$ cells identify main
effects.  Each chord buys one falsification test.  A deterministic certificate
for unrestricted responses on the complete table needs all $(m-1)(n-1)$
curvature coordinates.

\subsection{The compatibility-restricted response matroid}

Freeze the chemically valid edge set $E_*$ before reading outcomes
\cite{deleon2014recapmmp,konze2019reactional}.  For a
$q$-mode model, use the local parameterization
\[
  Y_{sd}=a_s+b_d+u_s^\top v_d,
  \qquad (s,d)\in E_*,
\]
with the usual gauges.  At a generic parameter point, let $J_{E_*}$ be the
Jacobian of all valid responses and let $J_A$ contain the rows for assayed cells
$A\subseteq E_*$.

\begin{thm}[Response-matroid closure]
\label{thm:response-matroid}
The assayed set $A$ locally determines every valid response if and only if
\[
  \ker J_A=\ker J_{E_*},
\]
equivalently $\rank J_A=\rank J_{E_*}$.  A missing valid cell $e$ is locally
completable if and only if its Jacobian row $j_e$ lies in
$\operatorname{rowspan}J_A$.
\end{thm}

\begin{proof}
An infinitesimal parameter perturbation preserves the assays exactly when it
lies in $\ker J_A$.  All valid responses are fixed by those assays exactly when
every such perturbation also lies in $\ker J_{E_*}$.  Since
$\ker J_{E_*}\subseteq\ker J_A$, equality of kernels is equivalent to equality
of ranks.  The single-cell statement is the corresponding row-span criterion.
\end{proof}

This closure defines an algebraic matroid on the reaction-valid cells, adapting
determinantal-matroid completion\cite{kiraly2015algebraic}.  An assay outside the
closure learns a new model coordinate.  An assay inside the closure creates an
overidentifying circuit that can falsify the model.  Any row basis of
$J_{E_*}$ is a minimum generic local design.  Global uniqueness may retain
finite branches, and lower-rank points are singular strata.

For $q=0$, $J_{E_*}$ is the graph incidence design: closure means that the two
endpoints already lie in the same connected component.  Thus graphic bridges
and chords are the additive case of this higher-rank learn-versus-falsify rule.
For $q>0$, a reaction-valid Schur chart supplies an explicit circuit, while the
Jacobian criterion covers masks that contain no such biclique chart.

\subsection{Quotienting repeated molecular representations}

Matched-pair fragmentation can represent one measured parent molecule by many
context--decoration edges\cite{hussain2010mmp,wassermann2012sarm}.  Those edge
responses are copies of one parent assay, not independent observations.  Let
$M$ index assayable parent molecules and let
$C\in\{0,1\}^{|E|\times|M|}$ copy their latent or measured responses
$f\in\R^{|M|}$ to
decomposition edges, so $y=Cf$.

\begin{thm}[Response-visible curvature quotient]
\label{thm:response-visible}
Let $\mathcal Z=\ker B$ be the topological cycle space and define
\[
  \mathcal H_{\mathrm{resp}}:=C^\top\mathcal Z,
  \qquad
  \mathcal N:=\mathcal Z\cap\ker C^\top.
\]
Then
\[
  \mathcal H_{\mathrm{resp}}\cong\mathcal Z/\mathcal N.
\]
If the columns of $Z$ span $\mathcal Z$, then
\[
  \dim\mathcal H_{\mathrm{resp}}=\rank(C^\top Z).
\]
Moreover, a parent-response vector $f$ is additively transportable on the
decomposition graph if and only if
\[
  h^\top f=0
  \qquad\text{for all }h\in\mathcal H_{\mathrm{resp}}.
\]
\end{thm}

\begin{proof}
For $z\in\mathcal Z$, its observed holonomy is
$z^\top y=z^\top Cf=(C^\top z)^\top f$.  The first isomorphism theorem gives
$C^\top\mathcal Z\cong\mathcal Z/(\mathcal Z\cap\ker C^\top)$ and the rank
formula.  By \Cref{thm:graphcriterion}, additive transport holds exactly when
every cycle holonomy vanishes, which is the final condition.
\end{proof}

\begin{prop}[Chemical typing removes unsafe cycles]
\label{prop:typed-cycles}
Let $B_0$ be an incidence matrix whose decoration nodes ignore attachment
environment, and let $B_\sigma$ be obtained by splitting each decoration into
typed nodes $(\sigma,d)$.  If $Q$ merges the typed vertices back to their
untyped labels, then
\[
  B_0=QB_\sigma,
  \qquad
  \ker B_\sigma\subseteq\ker B_0,
  \qquad
  C^\top\ker B_\sigma\subseteq C^\top\ker B_0.
\]
Thus untyped node identification can create cycle comparisons that disappear
when the transported operations are given chemically distinct meanings.
\end{prop}

\begin{proof}
If $z\in\ker B_\sigma$, then $B_0z=QB_\sigma z=0$.  Applying the linear map
$C^\top$ gives the response-visible inclusion.
\end{proof}

\subsection{From response-visible constraints to parent assays}

Let
\[
  r_H:=\dim\mathcal H_{\mathrm{resp}}=\rank(C^\top Z),
  \qquad N:=|M|.
\]
Choose a matrix $H\in\R^{r_H\times N}$ with orthonormal rows
spanning the response-visible constraint space; one construction orthonormalizes
a row basis of $Z^\top C$.  Then the parent responses
allowed by additive transport form the linear subspace
\[
  \mathcal P:=\ker H,
  \qquad
  \dim\mathcal P=N-r_H.
\]
We use ``check matrix,'' ``information set,'' and ``syndrome'' in the algebraic
sense of systematic linear coding\cite{hamming1950codes,prange1962information}.
The response space is over $\R$, so no finite-field error-correction claim is
intended.

For an assayed parent set $A\subseteq M$, write $U=M\setminus A$ and let $H_U$
contain the corresponding columns.  Define the unresolved model ambiguity and
the algebraic model-check codimension by
\begin{align}
  \mathsf D(A)
  &:=\dim\{p\in\mathcal P:p_A=0\},\label{eq:ambiguity}\\
  \mathsf F(A)
  &:=|A|-\dim\pi_A(\mathcal P),\label{eq:falsification-df}
\end{align}
where $\pi_A$ restricts a parent-response vector to $A$.

\begin{thm}[Exact learn--versus--falsify law]
\label{thm:learn-falsify}
For every $A\subseteq M$,
\begin{align}
  \mathsf D(A)&=|U|-\rank(H_U),\label{eq:ambiguity-rank}\\
  \mathsf F(A)&=r_H-\rank(H_U),\label{eq:falsification-rank}\\
  \mathsf F(A)-\mathsf D(A)&=|A|-(N-r_H).
  \label{eq:learn-falsify-conservation}
\end{align}
Both $\mathsf D$ and $\mathsf F$ are supermodular set functions.  If
$e\in U$ is assayed next, exactly one of the following occurs:
\begin{enumerate}
  \item if removing column $e$ lowers $\rank(H_U)$, then $f_e$ was already
  determined by $f_A$, $\mathsf D$ is unchanged, and $\mathsf F$ rises by one;
  \item otherwise $\mathsf D$ falls by one and $\mathsf F$ is unchanged.
\end{enumerate}
Thus one scalar parent assay either resolves one response ambiguity or creates
one independent model check.  It cannot do both.
\end{thm}

\begin{proof}
A response in the kernel of $\pi_A|_{\mathcal P}$ is supported on $U$ and has
coordinates in $\ker H_U$.  Rank--nullity gives
\eqref{eq:ambiguity-rank}.  Since
$\dim\pi_A(\mathcal P)=N-r_H-\mathsf D(A)$,
\eqref{eq:falsification-rank} and the conservation identity follow.

The column-rank function $S\mapsto\rank(H_S)$ is a matroid rank and is
submodular.  Its nullity is supermodular; taking complements in
\eqref{eq:ambiguity-rank} and \eqref{eq:falsification-rank} proves the two
supermodularity claims.  Finally, removing one column lowers rank by either
zero or one.
Substitution into the two rank formulas gives the stated dichotomy.  When rank
drops, $H_e$ is independent of $H_{U\setminus\{e\}}$, so every vector in
$\ker H_U$ has zero $e$ coordinate; this is exactly predictability from $A$.
\end{proof}

Falsification therefore has increasing returns.  A parent that resolves an
ambiguity early can become a pure model check after other parents have been
assayed.  Maximizing immediate curvature gain from an empty design is the wrong
combinatorial objective: the design must first acquire enough model coordinates
to make held-out responses predictable.

\begin{thm}[Systematic parent assay design]
\label{thm:systematic-parent}
Choose $U\subseteq M$ with $|U|=r_H$ such that $H_U$ is nonsingular, and set
$A=M\setminus U$.  Define
\[
  Q:=H_U^{-1}H_A,
  \qquad
  \widetilde f_A:=f_A,
  \qquad
  \widetilde f_U:=-Qf_A.
\]
Then:
\begin{enumerate}
  \item every $f\in\mathcal P$ is recovered exactly from the $N-r_H$ assays in
  $A$;
  \item no fixed coordinate set with fewer than $N-r_H$ assays can recover every
  $f\in\mathcal P$;
  \item for arbitrary $f\in\R^N$, the held-out transfer residual is exactly
  \begin{equation}
  \label{eq:systematic-syndrome}
    f_U-\widetilde f_U
    =H_U^{-1}Hf.
  \end{equation}
\end{enumerate}
Equivalently, after ordering coordinates as $(A,U)$,
\[
  H_U^{-1}H=[Q\ \ I_{r_H}].
\]
The information set $A$ identifies the model and has no lack-of-fit degree;
each subsequently assayed coordinate in $U$ adds one algebraically independent
systematic transfer check.
\end{thm}

\begin{proof}
The columns of $H_U$ form a basis of the column space of $H$.  The equation
$Hf=0$ therefore has the unique completion $f_U=-H_U^{-1}H_Af_A$.
For arbitrary $f$, multiplication gives
$H_U^{-1}Hf=Qf_A+f_U=f_U-\widetilde f_U$.
Since $\mathcal P$ has dimension $N-r_H$, an injective coordinate map from
$\mathcal P$ to $\R^{|A|}$ requires $|A|\ge N-r_H$.  The displayed construction
attains that count for every response in the fixed linear model.
\end{proof}

When $r_H=0$, the model constrains no parent response and all $N$ coordinates must
be assayed.  When $r_H=N$, the model is the zero subspace and phase-one prediction
needs no assay, although evidence for that model still requires validation.
The lower bound concerns recovery of the full parent-response vector; a declared
scalar decision can require fewer coordinates.

The columns of $H$ represent a linear matroid, and $A$ is a basis of its dual
matroid\cite{whitney1935matroid}.  If parent $i$ has fixed additive assay cost
$c_i\ge0$, then minimizing the phase-one cost $\sum_{i\in A}c_i$ is equivalent
to choosing a maximum-weight column basis $U$ of $H$.  Sorting parents by
decreasing $c_i$ and greedily retaining a column exactly when it raises rank is
the classical matroid-basis algorithm\cite{edmonds1971greedy}.  This result
optimizes noiseless identification cost only.  It does not optimize precision,
power, shared route costs, batch economies, or replicate allocation.

For fixed $U$, suppose $z_A=f_A+\eta_A$ and a fresh validation panel gives
$z_U=f_U+\eta_U$, with independent Gaussian errors of covariances
$\Sigma_{AA}$ and $\Sigma_{UU}$.  Under the transport null, the validation
innovation
\[
  \nu_U:=z_U+Qz_A
\]
has covariance
\[
  V_U=\Sigma_{UU}+Q\Sigma_{AA}Q^\top.
\]
If $V_U\succ0$, then
\[
  \nu_U^\top V_U^{-1}\nu_U\sim\chi_{r_H}^2.
\]
If $V_U$ is singular, replace the inverse by $V_U^+$ and the degrees of freedom
by $\rank(V_U)$.
A cheap information set can be nearly singular and amplify noise through
$H_U^{-1}$.  Statistical design therefore needs a conditioning or
A-, D-, E-, or G-optimality criterion beyond rank and additive cost
\cite{kiefer1959design,kiefer1960equivalence}.  The basis and its calibration
must also be frozen before reading the responses, or tested on fresh data.
Measuring every held-out parent eventually returns the total assay cost to the
full panel; the construction minimizes the initial identifying phase, not the
cost of exhaustive validation.

\subsection{Decision-specific estimability and robust ranking}

Full response recovery can cost more than a declared molecular decision.  To
separate measurements from decisions, let $X\in\R^{L\times d}$ have assay rows
$x_a^\top$ and let $Z\in\R^{N_D\times d}$ have decision rows $z_i^\top$:
\[
  y=X\theta,
  \qquad
  g=Z\theta.
\]
The parent-coordinate case has $L=N_D$, $X=Z$, $g=f$, and
$\im X=\mathcal P=\ker H$.  The next result is the classical estimability
criterion for a linear contrast\cite{rao1962generalized,searle1965estimable},
specialized to molecular ranking.  Partial monitoring uses the corresponding
observation-span condition for decision differences\cite{bartok2014partial}.

\begin{thm}[Decision-contrast sufficiency]
\label{thm:decision-estimability}
Assume the parameter $\theta\in\R^d$ is unrestricted and the exact observations
$y_A=X_A\theta$ are model-consistent.  For decision candidates $i,j$, the following are
equivalent:
\begin{enumerate}
  \item $g_i-g_j$ has the same value for every model-consistent completion;
  \item $z_i-z_j\in\operatorname{rowspan}X_A$;
  \item there is a vector $\lambda$ such that
  \[
    z_i-z_j=X_A^\top\lambda,
    \qquad
    g_i-g_j=\lambda^\top y_A.
  \]
\end{enumerate}
In the parent-coordinate case, these are also equivalent to the existence of
$\lambda,\mu$ such that
\[
  e_i-e_j=E_A^\top\lambda+H^\top\mu,
  \qquad E_Af=f_A.
\]
Let smaller responses be better.  A declared set $K$, $|K|=k$, is the unique
top-$k$ set for every exact model completion if and only if every boundary
contrast $z_j-z_i$, $i\in K$, $j\notin K$, lies in
$\operatorname{rowspan}X_A$ and every determined margin satisfies
\[
  g_j-g_i>0.
\]
\end{thm}

\begin{proof}
The completion parameters form $\theta_0+\ker X_A$.  The contrast is constant on
that affine space exactly when $z_i-z_j$ annihilates $\ker X_A$, equivalently
when it lies in $(\ker X_A)^\perp=\operatorname{rowspan}X_A$.  This gives the
coefficient representation and observed contrast.  In the parent-coordinate
case, the ambiguity space is $\mathcal P\cap\ker E_A$ and
\[
  (\mathcal P\cap\ker E_A)^\perp
  =\mathcal P^\perp+\im E_A^\top
  =\im H^\top+\im E_A^\top,
\]
which is exactly the check-matrix condition.

If every cross-boundary contrast is determined and positive, all completions
have the same strict top-$k$ boundary.  If one contrast is not determined, some
$v\in\ker X_A$ has $(z_j-z_i)^\top v\ne0$.  Scaling
$\theta_0+tv$ in the two signs preserves the observations and reverses that
comparison.  This proves necessity.
\end{proof}

\pagebreak[3]
The unrestricted-parameter assumption is essential.  Physical bounds can fix
the sign of a contrast without identifying its exact value.  In the
parent-coordinate specialization, let
\[
  \mathcal F_A
  :=\{f:Hf=0,\ \underline f\le f\le\overline f,\quad
        \ell_A\le f_A\le u_A\}
\]
be a nonempty polyhedral set combining the response model, physical bounds, and
simultaneous assay intervals.  For $i\in K$, $j\notin K$, solve
\begin{equation}
\label{eq:robust-margin-lp}
  \underline\Delta_{ij}
  :=\min_{f\in\mathcal F_A}(f_j-f_i).
\end{equation}
These are $k(N-k)$ linear programs.  The exact robust certificate is
\begin{equation}
\label{eq:robust-topk}
  \min_{i\in K,\ j\notin K}\underline\Delta_{ij}>0.
\end{equation}
If a minimum is nonpositive, its optimizer is a concrete counterexample
completion.  An infeasible $\mathcal F_A$ records model--data conflict, and an
unbounded program, possible when some declared bounds are infinite, records
missing structural control; neither is a ranking certificate.  For a fixed
design $A$, a simultaneous $1-\alpha$ interval for the assayed responses
transfers its coverage to \eqref{eq:robust-topk}, conditional on the declared
model and physical bounds.  Outcome-adaptive assay selection requires a
time-uniform or independently validated interval.  This is standard robust
linear optimization\cite{soyster1973robust,bental1998robust}; ellipsoidal or
nonlinear uncertainty sets require different optimization classes.

For a proposed $K$, define the boundary-contrast space
\[
  \mathcal W_K
  :=\operatorname{span}\{z_j-z_i:i\in K,\ j\notin K\}.
\]
Selecting a minimum-cost exact decision design asks for
\[
  \min_A\sum_{a\in A}c_a
  \quad\text{subject to}\quad
  \mathcal W_K\subseteq\operatorname{rowspan}X_A.
\]
When $\mathcal W_K=\operatorname{rowspan}X$, this reduces to the full-model
matroid basis problem and greedy is exact.  For a proper target subspace and
arbitrary assay rows, it contains minimum-support real linear representation
and is NP-hard even when $\dim\mathcal W_K=1$
\cite{natarajan1995sparse}.  In the reduction, dictionary atoms become cheap
assay rows while two separate decision vectors have difference equal to
the target vector; a cheap design spans that single contrast exactly when the
chosen atoms represent the target.  Verifying a fixed design remains a
polynomial rank test, and verifying a fixed bounded decision uses the LPs above.
This proves hardness for the general transductive problem with distinct assay
and decision sets; it does not prove the same hardness when every assayed parent
is also ranked in one common candidate set.  This separates
three claims that the original screening theory conflated: full identification
is matroidal, fixed-decision verification is tractable, and optimizing the
cheapest decision-specific assay set is hard without additional chemical
structure in the transductive setting.  Contrast-specific precision is the
classical $c$-optimal-design problem\cite{elfving1952allocation}; transductive
linear bandits study its noisy sequential ranking version
\cite{fiez2019transductive,reda2021topm}.

The accompanying module
\path{linear_decision_certificate.py} in the response-curvature experiment
directory
implements the row-span test and the robust LP certificate.  It returns a
counterexample completion when a finite margin fails, distinguishes unbounded
ambiguity from model--data conflict, and treats ties through an explicit policy.

\subsection{Graph-structured decision designs: paths and Steiner trees}

The general contrast-cover problem is hard, but the additive graph potential
has an exact combinatorial specialization.  This subsection assumes one assay
action per response edge.  Repeated-parent bundles require the parent check
matrix above.

Let $D_0\subseteq E$ be decision edges, let $c\in\R^{|D_0|}$ be a declared
linear contrast, and define its vertex demand
\[
  d:=B_{D_0}c.
\]
Under $y_e=b_e^\top\phi$, the target is
$c^\top y_{D_0}=d^\top\phi$.

\begin{thm}[Component-balance estimability]
\label{thm:component-balance}
For assayed edges $A\subseteq E$, the following are equivalent:
\begin{enumerate}
  \item $c^\top y_{D_0}$ is determined by $y_A$ for every potential $\phi$;
  \item there is a flow $\lambda\in\R^{|A|}$ with
  \[
    B_A\lambda=d,
    \qquad
    c^\top y_{D_0}=\lambda^\top y_A;
  \]
  \item every connected component $C$ of the assayed subgraph $(V,A)$ is
  balanced:
  \[
    \sum_{v\in C}d_v=0.
  \]
\end{enumerate}
\end{thm}

\begin{proof}
The target is determined by $y_A=B_A^\top\phi$ exactly when $d$ annihilates
$\ker B_A^\top$, equivalently when
$d\in(\ker B_A^\top)^\perp=\im B_A$.  This is condition~(2), and the equality
$\lambda^\top y_A=d^\top\phi$ for every $\phi$ is equivalent to
$B_A\lambda=d$.  The image of an incidence matrix consists exactly of demands
whose sum is zero on every connected component: incidence columns have zero
component sum, and a balanced demand can be routed on a spanning tree of each
component\cite{whitney1935matroid,jiang2011hodge}.  Equivalently,
$\im B_A=(\ker B_A^\top)^\perp$, while $\ker B_A^\top$ is precisely the space
of vectors constant on each component.
\end{proof}

\begin{cor}[Shortest-path assay design]
\label{cor:shortest-path-design}
Let a held-out decision edge $e=(u,v)$ be excluded from the assay set.  Its
response is predictable exactly when $u$ and $v$ are connected by assayed
edges.  With nonnegative additive assay costs, the cheapest exact design is a
shortest $u$--$v$ path.

Likewise, if $e=(s,u)$ and $f=(s,v)$ are oriented consistently away from one
shared endpoint $s$, then $y_e-y_f=\phi_v-\phi_u$.  The cheapest design for that
comparison is a shortest $u$--$v$ path.
\end{cor}

\begin{proof}
The relevant demands are $e_u-e_v$ up to sign.  Component balance is therefore
equivalent to endpoint connectivity.  Every feasible subgraph contains an
endpoint path, and deleting all other edges cannot raise nonnegative cost.
Shortest-path algorithms give the optimum\cite{dijkstra1959shortest}.
\end{proof}

\begin{thm}[Same-context ranking is a Steiner design]
\label{thm:steiner-design}
Let $q\ge2$ decision edges $e_i=(s,u_i)$ share one endpoint and orientation,
and let $K\subseteq[q]$ satisfy $0<|K|=k<q$.  After deduplicating the
opposite endpoints, put $T=\{u_i:1\le i\le q\}$.  All selected--unselected
boundary contrasts are estimable if and only if all vertices in $T$ lie in one
connected component of $(V,A)$.

On the declared assay ground graph, the minimum nonnegative assay cost is
exactly the weighted Steiner-tree cost connecting $T$.  Holding the decision
edges out deletes them from that ground graph; a finite prohibitive cost
simulates deletion when it exceeds the cost of a known feasible connector.  The corresponding decision problem is
NP-complete for variable $|T|$ and nonnegative rationally encoded costs, even
on bipartite response graphs.  It is
fixed-parameter tractable in $|T|$; the Dreyfus--Wagner algorithm runs in
$3^{|T|}\operatorname{poly}(|V|,|E|)$ time, and $|T|=2$ reduces to
\Cref{cor:shortest-path-design}.
\end{thm}

\begin{proof}
For two decision edges,
$y_i-y_j=\phi_{u_j}-\phi_{u_i}$.  The cross-boundary comparison graph
$K_{|K|,q-|K|}$ is connected, so component balance for every comparison is
equivalent to placing all terminals in one assayed component.  A minimum
nonnegative connector can be pruned to a tree, and every Steiner tree is
sufficient.

For hardness, start with a connected nonnegatively weighted Steiner instance
$(G_0,T,w,W)$ with rationally encoded data and
subdivide every original edge $a$--$b$ by a new vertex $x_e$.  Give each of the
two half-edges cost $w(e)$ and double the budget.  Original vertices and
subdivision vertices form the two bipartition classes.  Pruning nonterminal
subdivision leaves and contracting every remaining degree-two $x_e$ gives a
cost-preserving correspondence up to the factor two.  Add a new decision hub
$s$ on the subdivision side and decision edges $(s,t)$ for $t\in T$.  Choose
$(s,t_0)$ as the proposed top-1 edge.  Its boundary contrasts are estimable
exactly when every terminal connects to $t_0$.  If these decision edges remain
assayable, assign each cost $1+2\sum_{e\in E_0}w(e)$, so no optimum uses one.
Connectivity and estimability are computed in the assayed subgraph; the
unassayed decision star does not connect terminals.
This proves bipartite NP-hardness; feasibility verification is polynomial.
Classical Steiner results give NP-completeness and the fixed-terminal dynamic
program\cite{garey1977steiner,muller1987bipartite,dreyfus1971steiner}.
\end{proof}

Connectivity certifies estimability, not statistical precision.  If
$z_A=B_A^\top\phi+\varepsilon_A$ with covariance $\Sigma_A$, any routing flow
has
\[
  \widehat t=\lambda^\top z_A,
  \qquad
  \Var(\widehat t)=\lambda^\top\Sigma_A\lambda.
\]
For independent path-edge noise this is the sum of path variances; overlapping
top-$k$ paths are correlated.  More generally, two routing flows satisfy
\[
  \Cov(\lambda^\top z_A,\mu^\top z_A)
  =\lambda^\top\Sigma_A\mu.
\]
For independent path noise this reduces to the signed sum of variances on
shared path edges:
\[
  \Cov(\widehat t_P,\widehat t_Q)
  =\sum_{a\in P\cap Q}s_{P,a}s_{Q,a}\sigma_a^2.
\]
Minimum assay cost, shortest-path cost, and minimum
estimator variance therefore need not agree.  When $\Sigma_A$ is positive
diagonal, the minimum-variance flow is the effective-resistance problem from
\Cref{sec:graph}; for correlated $\Sigma_A$ it is a constrained quadratic-flow
problem.  The Steiner
parameter is the total number of distinct decision terminals, not the requested
$k$.  Duplicate terminals on opposite sides of a strict ranking boundary are
algebraically tied and cannot receive a unique strict order.

The incidence criterion, shortest paths, Steiner trees, and their algorithms
are classical.  The chemical task is to construct the reaction-typed assay
graph, declare the molecular decision terminals, and validate the resulting
path or tree certificate on future measurements.

The module \path{graph_potential_decision_design.py} implements component-balance
flows and exact shortest paths.  For small terminal graphs it also performs an
exhaustive exact Steiner search with a declared edge-count limit; it raises an
error rather than silently substituting a heuristic beyond that limit.

\subsection{Decision precision as a group-flow design}

Fixed edge-opening costs lead to Steiner structure.  A per-replicate budget has
a different, convex solution.  Let $p=1,\ldots,P$ index routable vertex demands
$d_p$ with criterion weights $\omega_p>0$.  Choose unbiased routing flows
$B\lambda_p=d_p$.  Edge $e$ receives a continuous number $n_e$ of independent
replicates, each with variance $\sigma_e^2>0$ and cost $c_e>0$, under budget
\[
  \sum_ec_en_e\le\mathcal B,
  \qquad \mathcal B>0.
\]
Define the squared decision load
\[
  h_e(\Lambda):=\sum_{p=1}^P\omega_p\lambda_{ep}^2.
\]

\begin{thm}[Joint routing--replication reduction]
\label{thm:group-flow}
Assume every demand $d_p$ is routable and at least one demand is nonzero.
For fixed feasible flows, the weighted total decision variance is
\[
  \mathcal V(\Lambda,n)
  =\sum_e\frac{\sigma_e^2h_e(\Lambda)}{n_e}.
\]
Use the closed perspective convention $0/0:=0$ and $a/0:=+\infty$ for $a>0$.
Its exact continuous allocation is
\begin{equation}
\label{eq:decision-replicates}
  n_e^*
  =\frac{\mathcal B\,\sigma_e\sqrt{h_e/c_e}}
  {\sum_f\sigma_f\sqrt{h_fc_f}},
\end{equation}
on positive-load edges, with $n_e^*=0$ when $h_e=0$.  The minimum for those
flows is
\begin{equation}
\label{eq:fixed-flow-variance}
  \frac{\left(\sum_e\sigma_e\sqrt{h_ec_e}\right)^2}{\mathcal B}.
\end{equation}

Put $\ell_e:=\sigma_e\sqrt{c_e}$.  Jointly optimizing routing and replication
gives
\begin{equation}
\label{eq:group-flow-radius}
  \mathcal V_*=\frac{\rho_*^2}{\mathcal B},
  \qquad
  \rho_*:=\min_{B\lambda_p=d_p}
  \sum_e\ell_e
  \left\|(\sqrt{\omega_p}\lambda_{ep})_{p=1}^P\right\|_2.
\end{equation}
The inner problem is the second-order cone program
\[
\begin{aligned}
  \min_{\Lambda,t}\quad&\sum_e\ell_et_e\\
  \text{subject to}\quad&B\lambda_p=d_p\quad(p=1,\ldots,P),\\
  &\left\|(\sqrt{\omega_p}\lambda_{ep})_p\right\|_2\le t_e
    \quad(e\in E).
\end{aligned}
\]
If $W=\diag(\omega_p)$, $D=[d_1\ \cdots\ d_P]$, and
$\Gamma=\Lambda W^{1/2}$, then $B\Gamma=DW^{1/2}$ and the dual is
\begin{equation}
\label{eq:group-flow-dual}
  \max_{\Psi\in\R^{|V|\times P}}\ \langle\Psi,DW^{1/2}\rangle_F
  \quad\text{subject to}\quad
  \|(B^\top\Psi)_{e,:}\|_2\le\ell_e\quad(e\in E).
\end{equation}
Thus the dual assigns vector-valued potentials to vertices with an edgewise
Lipschitz bound.
\end{thm}

\begin{proof}
Cauchy--Schwarz gives
\[
  \left(\sum_e\sigma_e\sqrt{h_ec_e}\right)^2
  \le
  \left(\sum_e\frac{\sigma_e^2h_e}{n_e}\right)
  \left(\sum_ec_en_e\right).
\]
The allocation in \eqref{eq:decision-replicates} attains equality.  Eliminating
$n$ yields \eqref{eq:group-flow-radius}; its objective is a weighted sum of
row Euclidean norms, giving the displayed SOCP.  Standard conic duality gives
\eqref{eq:group-flow-dual}.  Equivalently, the original edge terms are
quadratic-over-linear perspectives, so the joint problem is convex.
\end{proof}

If every demand is zero, choose zero flows and use no replication; the optimum
variance is zero and the ratio formulas above are not invoked.

This is the square-root allocation, group-norm, and multiresponse SOCP structure
from classical optimal design\cite{neyman1934stratified,
sagnol2011multiresponse,lobo1998socp,yuan2006group}, specialized to
response-graph demands.

\begin{cor}[One-contrast precision path]
\label{cor:precision-path}
For one unit $u$--$v$ demand of weight $\omega$, the optimal routing radius and
variance are
\[
  \rho_*=\sqrt\omega\,\operatorname{dist}_\ell(u,v),
  \qquad
  \mathcal V_*=\frac{\omega\,\operatorname{dist}_\ell(u,v)^2}{\mathcal B},
  \qquad
  \ell_e=\sigma_e\sqrt{c_e}.
\]
Any shortest $\ell$-path is optimal.  On a chosen path,
\[
  n_e^*
  =\frac{\mathcal B\,\sigma_e/\sqrt{c_e}}
  {\sum_{f\in P}\sigma_f\sqrt{c_f}}.
\]
\end{cor}

\begin{proof}
For one demand, the group norm becomes weighted $\ell_1$ flow cost.  Delete
cycles and decompose a unit flow into $u$--$v$ paths.  Its cost is at least the
shortest path length, attained by one shortest-path flow.
\end{proof}

\begin{cor}[Fixed-tree top-$k$ allocation]
\label{cor:tree-topk-allocation}
Fix a tree connecting $q$ distinct decision terminals and a nontrivial selected set
$K$, $|K|=k$.  Give every selected--unselected pair common weight $\omega$.
Write $K^c=[q]\setminus K$.
Removing tree edge $e$ leaves a terminal side $S_e$.  Write
\[
  a_e:=|K\cap S_e|,
  \qquad
  b_e:=|K^c\cap S_e|,
  \qquad m:=q-k.
\]
Then the squared edge load is
\begin{equation}
\label{eq:tree-cut-load}
  h_e=\omega\{a_e(m-b_e)+(k-a_e)b_e\}.
\end{equation}
The exact tree-conditioned allocation and total boundary-variance criterion are
given by \eqref{eq:decision-replicates} and \eqref{eq:fixed-flow-variance} with
these loads.
\end{cor}

\begin{proof}
Each pair has a unique tree path.  It crosses $e$ exactly when its two terminals
lie on opposite sides of the cut.  The two terms in
\eqref{eq:tree-cut-load} count the two possible selected--unselected
orientations across that cut.  The allocation theorem then applies.
\end{proof}

The formulas also give simple guarantees.  Let
$P_+:=|\{p:d_p\ne0\}|$ and let $r_p$ be the independent minimum-$\ell_1$
routing cost of each nonzero demand.  Then
\[
  \rho_*
  \ge
  \max\left\{
    \max_p\sqrt{\omega_p}r_p,
    \frac1{\sqrt{P_+}}\sum_p\sqrt{\omega_p}r_p
  \right\}.
\]
Independently minimum-$\ell_1$-routing every demand has group objective at most
$\sum_p\sqrt{\omega_p}r_p$, hence is a $\sqrt{P_+}$-approximation in $\rho$
and a $P_+$-approximation in total variance.  On a pruned tree with unit weights
for the $R=k(q-k)$ boundary pairs, every used edge has $1\le h_e\le R$.
Thus every such tree $T$ satisfies
\[
  C_\ell(T)\le\rho(T)\le\sqrt R\,C_\ell(T),
  \qquad
  C_\ell(T):=\sum_{e\in T}\ell_e,
  \quad
  \rho(T):=\sum_{e\in T}\ell_e\sqrt{h_e}.
\]
Consequently, a minimum-$\ell$ Steiner tree is a $\sqrt R$-approximation to the
best tree-restricted group-flow radius and an $R$-approximation after squaring.
These are tree-restricted statements, not bounds against every fractional SOCP
solution.

The distinction between cost models is essential.  The convex theorem charges
only $c_en_e$ and permits split real flows.  Adding a fixed activation or
synthesis cost $F_e\mathbf1_{\{n_e>0\}}$, integer replicates, batch constraints,
or minimum-one-replicate rules creates a mixed-integer problem.  Correlated
edge errors destroy the separable load formula.  The trace objective minimizes
the weighted sum of marginal contrast variances, not maximum variance,
familywise ranking error, or simultaneous interval width.  Flow and allocation
chosen after inspecting the same noise require selective or fresh-data
inference.

On a fixed tree, the cut-load accounting is classical communication-tree
algebra\cite{hu1974communication}.  No novelty is claimed for square-root
allocation, group norms, electrical flow, or SOCP representability.  The
chemical contribution is the construction of molecular decision demands on a
reaction-typed graph and their prospective coupling to assay variances and
replicate costs.

The module \path{decision_precision_design.py} implements the fixed-flow,
one-contrast, and fixed-tree formulas, with a guarded callback interface for a
general conic solver.  No undeclared conic dependency or heuristic is silently
used.

Cycles in $\mathcal N$ vanish because the same parent was copied through several
fragmentations.  They contain no response evidence.  A topological Betti number
therefore overstates chemical curvature dimension unless $C$ is injective on the
cycle space.  If parent measurement errors have covariance $\Sigma_f$, cycle
covariance is $Z^\top C\Sigma_fC^\top Z$.

An assay action reveals one parent response and activates its full bundle of
decomposition edges.  The parent-level check matrix, rather than individual
graph edges, is therefore the correct acquisition object.  Treating copied
edges as separate assays inflates evidence.

\section{The exact uniform obstruction}

Mean-square interaction energy can hide a small number of decisive cliffs.  The
uniform approximation problem has an exact cycle formula.

The formula is a symmetric potential-fitting form of classical mean-cycle and
minimum-balance duality\cite{karp1978cyclemean,young1991minbalance}.

\begin{thm}[Cycle-mean theorem]
\label{thm:cyclemean}
For a finite observation graph, define the best uniform additive error
\[
  \rho_\infty(y):=\min_\phi\|y-B^\top\phi\|_\infty.
\]
Then
\begin{equation}
\label{eq:cyclemean}
  \rho_\infty(y)
  =\max_{C\text{ simple cycle}}\frac{|\kappa_C(y)|}{|C|},
\end{equation}
with the maximum set to zero for a forest.
\end{thm}

\begin{proof}
Distance to a subspace in the $\ell_\infty$ norm has the dual form
\[
\rho_\infty(y)
=\max\{z^\top y:z\in(\im B^\top)^\perp,\ \|z\|_1\le1\}
=\max\{z^\top y:z\in\ker B,\ \|z\|_1\le1\}.
\]
Every circulation $z\in\ker B$ decomposes into oriented simple-cycle flows
without increasing its $\ell_1$ mass.  Write
$z=\sum_C\alpha_C\chi_C$ with $\alpha_C\ge0$ and
$\sum_C\alpha_C|C|=\|z\|_1$.  Then
\[
|z^\top y|
\le\sum_C\alpha_C|\kappa_C(y)|
\le\left(\max_C\frac{|\kappa_C(y)|}{|C|}\right)\|z\|_1.
\]
The reverse inequality follows by choosing
$z=\pm\chi_C/|C|$ for a maximizing cycle.
\end{proof}

\begin{cor}[Sharp observed-domain decision bound]
Let $\hat\phi$ solve the Chebyshev problem in \Cref{thm:cyclemean}.  Choosing the
best observed edge under $B^\top\hat\phi$ incurs response regret at most
$2\rho_\infty(y)$.  The constant uses the exact maximum normalized cycle
holonomy rather than an average interaction energy.
\end{cor}

The theorem replaces the vague statement that small local curvature should
imply global factorization.  The maximum normalized holonomy equals the sharp
uniform factorization error on the observed graph.  It says nothing about an
unobserved valid combination without a completion assumption.

\section{Noise and two experimental objectives}

Assume the measured edge vector is
\[
  z=y+\eps,
  \qquad
  \eps\sim\mathcal N(0,\sigma^2W^{-1}),
\]
where $W=\diag(w_e)$ contains known relative precisions.
Write
\[
  r_G:=|E|-|V|+c=\dim\ker B.
\]

\begin{thm}[Finite-sample global additivity test]
\label{thm:chisq}
Under the additive null $y=B^\top\phi$, the weighted residual statistic
\[
  T:=\frac{1}{\sigma^2}\min_\phi
  (z-B^\top\phi)^\top W(z-B^\top\phi)
\]
has a $\chi^2_{r_G}$ distribution.
\end{thm}

\begin{proof}
Whitening turns the problem into Gaussian linear regression with design
$W^{1/2}B^\top$, whose rank is $|V|-c$.  The squared norm of the residual
projection has $|E|-(|V|-c)$ degrees of freedom.
\end{proof}

Failure to reject does not prove unrestricted additivity.  It supplies a test
on the measured graph at the declared noise model and significance level.

\subsection{Learning additive potentials}

Let $L_W=BWB^\top$.  Under the additive model, the weighted least-squares
estimator has, on a fixed gauge,
\[
  \operatorname{Cov}(\hat\phi)=\sigma^2L_W^+.
\]
For a query contrast $b_f^\top\phi$,
\begin{equation}
\label{eq:reff}
  \operatorname{Var}(b_f^\top\hat\phi)
  =\sigma^2b_f^\top L_W^+b_f.
\end{equation}
The quadratic form is effective resistance.  It supports established
potential-estimation and free-energy-network designs
\cite{ghosh2008resistance,osting2014rankings,xu2019measurement}.

If the graph is connected and a new observation on edge $e$ has relative
precision $w$, the variance of query $f$ decreases by
\begin{equation}
\label{eq:update}
  \sigma^2
  \frac{w\,(b_f^\top L_W^+b_e)^2}
       {1+w\,b_e^\top L_W^+b_e}.
\end{equation}
This is the Sherman--Morrison update on the subspace orthogonal to $\one$.

\begin{thm}[Prospective chord innovation]
\label{thm:innovation}
Fit the additive model on a graph $G$.  Let a new edge $e$ join vertices already
in the same connected component, let its independent measurement variance be
$\tau_e^2$, and define
\[
  \nu_e:=z_e-b_e^\top\hat\phi,
  \qquad
  V_e:=\sigma^2 b_e^\top L_W^+b_e.
\]
Under the additive Gaussian model,
\[
  \nu_e\sim\mathcal N(0,\tau_e^2+V_e),
  \qquad
  \frac{\nu_e^2}{\tau_e^2+V_e}\sim\chi_1^2.
\]
The global weighted lack-of-fit statistic updates by
\begin{equation}
\label{eq:sequential-update}
  T_{G+e}=T_G+\frac{\nu_e^2}{\tau_e^2+V_e}.
\end{equation}
If one starts from a spanning forest and adds chords in a fixed sequence, the
standardized innovations are independent under the null.
\end{thm}

\begin{proof}
The fitted potential is linear in the old Gaussian measurements.  Its prediction
at $e$ has variance $V_e$ and is independent of the new measurement error, which
gives the innovation variance.  Recursive least squares or a Schur-complement
calculation gives \eqref{eq:sequential-update}.  Each standardized recursive
residual is an orthogonal Gaussian projection, so the sequence is independent.
\end{proof}

The theorem creates a prospective falsification loop.  A chord closes one new
cycle and returns a calibrated one-degree-of-freedom test.  An edge joining two
components behaves differently: it fixes a relative gauge, reduces the number
of components by one, and creates no cycle degree of freedom.  Such a bridge
expands prediction reach but cannot test transfer by itself.

The independence statement requires a fixed chord sequence.  If earlier assay
outcomes choose the next chord, fixed-design $\chi^2$ calibration does not follow
from this theorem.  Fresh circuit measurements permit the conditional test below
under its Gaussian, known-covariance, and independence assumptions.

\begin{thm}[Anytime-valid adaptive circuit test]
\label{thm:eprocess}
Let $\mathcal F_{t-1}$ contain all data before round $t$.  Choose an
$\mathcal F_{t-1}$-measurable linear circuit $c_t$ that annihilates every response
under a declared linear null model.  After choosing it, obtain a fresh independent
measurement vector
\[
  z_t\sim\mathcal N(y,\Sigma_t)
\]
on its support and set
\[
  W_t:=\frac{c_t^\top z_t}{\sqrt{c_t^\top\Sigma_tc_t}}.
\]
Under the null, $W_t\mid\mathcal F_{t-1}\sim\mathcal N(0,1)$.  For any predictable
$\lambda_t$, the products
\begin{align}
  E_t^{\mathrm{signed}}
  &=\prod_{i=1}^t
    \exp\!\left(\lambda_iW_i-\frac{\lambda_i^2}{2}\right),\\
  E_t^{\mathrm{two}}
  &=\prod_{i=1}^t
    \exp\!\left(-\frac{\lambda_i^2}{2}\right)
    \cosh(\lambda_iW_i)
\end{align}
are nonnegative test martingales.  Therefore
\[
  \Pr_0\!\left(\sup_t E_t\ge\frac1\alpha\right)\le\alpha
\]
for either process, under arbitrary adaptive circuit choice and stopping.
\end{thm}

\begin{proof}
Freshness makes the selected contrast conditionally standard normal under the
null.  Its conditional moment-generating function gives
\[
\E_0[e^{\lambda_tW_t-\lambda_t^2/2}\mid\mathcal F_{t-1}]=1.
\]
Symmetrizing $\lambda_t$ gives the cosh factor with the same conditional mean.
The products are martingales, and Ville's inequality gives the crossing bound.
\end{proof}

This is an e-process construction in the sense of anytime-valid inference
\cite{howard2021timeuniform,koolen2022evalues}.  Reusing noisy anchors selected
from history does not meet the freshness condition.  One may instead reserve a
blinded validation bank or remeasure every edge in the chosen circuit.

If a training model forecasts circuit curvature $\widehat\kappa_t$ and the fresh
contrast variance is $v_t$, choosing
$\lambda_t=\widehat\kappa_t/\sqrt{v_t}$ leaves null validity intact.  Under true
curvature $\kappa_t$, the oracle expected log-evidence is
$\kappa_t^2/(2v_t)$, the Gaussian KL information.  This yields the cost-aware
acquisition score
\begin{equation}
\label{eq:curvature-information}
  \frac{\E[\kappa_t^2\mid\mathcal F_{t-1}]}
       {2v_t\,\operatorname{cost}_t},
\end{equation}
predicted curvature evidence per assay cost.  Exact Gaussian calibration applies
to linear circuits.  Noisy low-rank Schur circuits are polynomial contrasts and
need fresh-pivot uncertainty propagation or bootstrap calibration.

\begin{prop}[Sequential curvature-information lower bound]
\label{prop:information-lower}
Consider any predictable circuit policy and integrable stopping rule that
distinguishes an additive null from an alternative with type-I error at most
$\alpha$ and power at least $1-\beta>\alpha$.  If round $t$ uses a fresh Gaussian contrast with variance
$v_t$ and alternative mean $\kappa_t$, then
\begin{equation}
\label{eq:information-lower}
  \E_1\!\left[
    \sum_{t\le\tau}\frac{\kappa_t^2}{2v_t}
  \right]
  \ge
  \operatorname{kl}(1-\beta,\alpha),
\end{equation}
where $\operatorname{kl}$ is binary relative entropy.
\end{prop}

\begin{proof}
Conditionally on the past, the Gaussian alternative and null at round $t$ have
relative entropy $\kappa_t^2/(2v_t)$.  The chain rule makes the left side the
relative entropy of the stopped data laws.  Mapping the stopped data to the
binary reject-or-not decision and applying data processing gives the right side.
\end{proof}

Thus no adaptive policy can beat the cumulative information requirement in
\eqref{eq:information-lower}.  The score in
\eqref{eq:curvature-information} is the greedy oracle for information per cost,
while the e-process preserves error control when its curvature forecasts are
wrong.

\subsection{Learning curvature}

Potential accuracy and curvature accuracy require different allocations.  Let
$Z\in\R^{|E|\times r_G}$ have orthonormal columns spanning $\ker B$.  The
curvature coordinates are $h=Z^\top y$, and the cycle projector is
\[
  P_{\mathrm{cyc}}=ZZ^\top=I-B^\top(BB^\top)^+B.
\]

\begin{defn}[Cycle leverage]
Define
\[
  \ell_e:=(P_{\mathrm{cyc}})_{ee}=\|Z_{e,:}\|_2^2.
\]
For a unit-weight graph,
\[
  \ell_e=1-b_e^\top(BB^\top)^+b_e
  =1-R_{\mathrm{eff}}(e).
\]
Thus bridges have cycle leverage zero.
\end{defn}

\begin{thm}[Curvature-directed replicate allocation]
\label{thm:allocation}
Let $E_{\mathrm{cyc}}=\{e:\ell_e>0\}$.  After any required minimum coverage has
been fixed, suppose edge $e\in E_{\mathrm{cyc}}$ receives $n_e>0$ independent
replicates, each with variance $\sigma_e^2$, and one replicate costs $c_e$.  The
estimator formed from edge means satisfies
\[
  \E\|\hat h-h\|_2^2
  =\sum_e\frac{\sigma_e^2\ell_e}{n_e}.
\]
Under the continuous budget constraint
$\sum_{e\in E_{\mathrm{cyc}}}c_en_e\le B$, the curvature-MSE
minimizer is
\begin{equation}
\label{eq:allocation}
  n_e^*
  =\frac{B}{\sum_f\sigma_f\sqrt{\ell_fc_f}}
    \frac{\sigma_e\sqrt{\ell_e}}{\sqrt{c_e}},
\end{equation}
and its minimum MSE is
\begin{equation}
\label{eq:minmse}
  \frac{\left(\sum_e\sigma_e\sqrt{\ell_ec_e}\right)^2}{B}.
\end{equation}
\end{thm}

\begin{proof}
The edge-mean covariance is
$\diag(\sigma_e^2/n_e)$.  Therefore
\[
\tr\operatorname{Cov}(Z^\top\hat y)
=\tr\!\left(Z^\top\diag(\sigma_e^2/n_e)Z\right)
=\sum_e\sigma_e^2\|Z_{e,:}\|_2^2/n_e.
\]
The row norms are invariant under a change of orthonormal cycle basis.  A
Lagrange multiplier for the cost constraint gives \eqref{eq:allocation};
substitution gives \eqref{eq:minmse}.
\end{proof}

The duality is operational:
\begin{center}
\begin{tabular}{lll}
\toprule
Goal & Graph quantity & High-priority edges \\
\midrule
Estimate additive potentials & effective resistance & weakly connected contrasts \\
Estimate curvature & cycle leverage & edges carrying cycle-space mass \\
\bottomrule
\end{tabular}
\end{center}

An experimental program seeking transferable main effects should not use the
same allocation as a program seeking context dependence.  Formula
\eqref{eq:allocation} addresses repeat allocation on a fixed graph.  Choosing
new missing combinations is a separate design problem: every chord creates a
new cycle coordinate, while its expected chemical utility and synthesis cost
determine which coordinate deserves measurement.

\section{Refinement spends falsification degrees}

Let a coarse and refined representation label the same measured edge set $E$.
Assume vertex splitting supplies a coarsening map, so every coarse potential can
be lifted constantly to its refined children.  The corresponding edge-potential
spaces satisfy
\[
  \im B_0^\top\subseteq\im B_1^\top.
\]

\begin{thm}[Refinement conservation law]
\label{thm:refinement-conservation}
Write
\[
  g_i:=\rank(B_i)=|V_i|-c_i,
  \qquad
  r_i:=|E|-g_i.
\]
If $d=g_1-g_0$, then $d=r_0-r_1$.  Let $P_0,P_1$ be the $W$-orthogonal
projections onto the coarse and refined potential spaces.  For every response
vector $y$,
\begin{equation}
\label{eq:refinement-pythagoras}
  \|y-P_0y\|_W^2
  =\|(P_1-P_0)y\|_W^2+\|y-P_1y\|_W^2.
\end{equation}
\end{thm}

\begin{proof}
Rank--nullity gives $r_i=|E|-g_i$, proving the dimension identity.  Nested
orthogonal projections make $(P_1-P_0)y$ orthogonal to $(I-P_1)y$, which gives
\eqref{eq:refinement-pythagoras}.
\end{proof}

Each new refined main-effect degree removes exactly one independent curvature
degree.  A split with $d=0$ changes labels without enlarging the response model.
Splitting until every edge gets its own parameter makes $r_1=0$, fits every
response, and leaves no falsification test.  Refinement therefore cannot explain
curvature for free.

For a fixed pre-outcome refinement, suppose
$z\sim\mathcal N(y,\sigma^2W^{-1})$ under the coarse null
$y\in\im B_0^\top$.  Then
$\|(P_1-P_0)z\|_W^2/\sigma^2\sim\chi_d^2$ and is independent of the refined
residual.  An outcome-selected refinement requires untouched validation data or
multiplicity control.  A policy may choose a split on training curvature, then
accept it only when it reduces fresh circuit innovations enough to justify the
exact $d$ falsification degrees it consumes.

\section{Experimental semantics}

Two kinds of cycles must remain separate.

\begin{enumerate}
  \item An experimental response rectangle compares four molecular states.  A
  nonzero contrast can record real context dependence, assay noise, or batch
  mismatch.
  \item An exact thermodynamic free energy is a state function.  A nonzero sum
  of estimated relative free energies around a closed FEP loop diagnoses
  sampling or estimator inconsistency.  It does not by itself prove physical
  SAR nonadditivity.  Network methods enforce or analyze that closure
  separately\cite{giese2021barnet}.
\end{enumerate}

The analogy is algebraic: both use cycles.  Their physical meanings differ.

\section{Prospective reaction-feasible response graph}

A testable implementation can proceed in two stages.

\subsection*{Stage 1: define the graph before reading outcomes}
\begin{enumerate}
  \item Choose one homogeneous endpoint and assay protocol.
  \item Define context nodes through scaffolds, binding modes, or reaction
  environments.
  \item Define decoration nodes through executable reaction templates and
  mapped attachment sites.
  \item Add an edge only when the construction is valid and its product has a
  retrievable route.
  \item Freeze the graph, costs, and initial measurement set.
\end{enumerate}
Reaction-aware matched pairs and reaction-based enumeration are established
methods\cite{deleon2014recapmmp,konze2019reactional}; their role here is to
freeze the ground set on which transfer will be tested.

\subsection*{Stage 2: acquire evidence}
\begin{enumerate}
  \item Measure a connected anchor design to identify additive main effects.
  \item Add chords to create curvature tests.  Select them by expected
  contract value, model disagreement, and cycle precision per unit cost.
  Reserve independent confirmatory measurements when this selection uses prior
  responses.
  \item Allocate replicates with \Cref{thm:allocation} when curvature precision
  is the declared goal.
  \item Fit the weighted Hodge decomposition and report both the potential and
  cycle components.
  \item Fit low-rank interaction modes only after block-held-out validation.
  \item Apply \Cref{thm:regret,cor:topk} using a bound calibrated on untouched
  rounds.  Confirm selected compounds with orthogonal measurements.
\end{enumerate}
Generic molecular active learning and noisy batched optimization provide
necessary baselines\cite{graff2021molpal,bellamy2022bo}.  The distinct target
here is independent nonadditivity evidence and a transfer certificate, not
predicted potency alone.

\subsection*{Structural matched-pair opportunity audit}

The script
\texttt{experiments/response\_curvature/chembl\_mmp\_response\_complex.py}
implements a conservative retrospective graph construction.  It makes one
eligible single-bond cut at a time, assigns the larger fragment as context,
restricts the smaller decoration to at most 12 heavy atoms and one third of the
parent, checks exact isomeric reassembly, and types the decoration by the rooted
context-side attachment environment.  The filename retains ``complex'' for
compatibility; the constructed object is the response graph defined here.

On the five-target exact-IC50 extract, 8,677 source rows collapse to 8,676
target--structure observations.  Across 8,152 connected parents with accepted
decompositions, all 69,573 eligible cuts pass the serialization and reassembly
check.  After size filtering and target expansion, the graph has 51,552 response
edges.  Its curvature dimensions contract as
\[
  \begin{aligned}
  9{,}510\ \text{naive untyped cycles}
  &\longrightarrow 6{,}732\ \text{attachment-typed cycles}\\
  &\longrightarrow 3{,}771\ \text{response-visible dimensions}.
  \end{aligned}
\]
Thus only 39.65\% of the naive cycle rank survives both chemical typing and the
repeated-parent quotient in \Cref{prop:typed-cycles,thm:response-visible}.  The
graph contains 68,756 typed rectangles with four distinct parents, representing
56,737 unique molecular quartets.  These counts are structural opportunities,
not nonadditivity evidence: the input lacks assay identifiers, so the four
responses cannot be required to come from one protocol.

Applying \Cref{thm:systematic-parent} to the five target blocks splits 8,154
parent-response coordinates into 4,383 minimum-count information assays and
3,771 systematic check holdouts.  The implementation orthonormalizes the
response-visible row space before selecting columns, so the chart is invariant
to the arbitrary cycle basis; all reconstruction identities hold to numerical
precision.  The worst selected check block has condition number 80.23.  On the
already observed responses, the pooled held-out prediction RMSE is 2.001
pChEMBL units.  This large retrospective residual rejects the frozen additive
transport model on the aggregated table.  It does not validate a chemical
mechanism because assay and protocol heterogeneity remain uncontrolled.

The companion script
\path{chembl_path_certificate_audit.py} then holds out entire decision-parent
decomposition bundles and asks whether two responses sharing one context have an
alternate typed-graph path between their decoration endpoints.  Among 71,248
eligible pairs, 23,165 (32.51\%) have such a path; 23,131 paths use distinct
nondecision parents.  Of those, 18,239 are two-edge rectangle certificates and
4,892 require longer paths, up to length 12.  Removing the whole decision-context
star still leaves 22,661 observable pairs (31.81\%).  An edge-only holdout would
accept 221 paths that re-enter through another decomposition of a held-out
parent, showing why parent-bundle removal is necessary.  These are structural
path identities, not assay-valid transfer results.

\subsection*{Exploratory ChEMBL replay}

The accompanying script
\texttt{experiments/response\_curvature/chembl\_target\_molecule\_pilot.py}
applies the sparse-graph construction to the existing five-target ChEMBL
extract.  It uses targets as contexts, identical molecules as design units, and
keeps aggregated records whose standard type is exactly IC50.  Molecules need
measurements against at least two targets.  The resulting graph contains three
targets, 17 shared molecules, 34 edges, and 15 independent cycles.

\begin{center}
\begin{tabular}{lr}
\toprule
Diagnostic & Value \\
\midrule
Weighted Hodge residual RMSE & 0.546 pChEMBL \\
Best uniform additive radius & 1.218 pChEMBL \\
Online additive chord MSE & 1.515 \\
Online target-mean chord MSE & 0.564 \\
Additivity $p$-value at assumed $\sigma=0.5$ & $6.73\times10^{-5}$ \\
\bottomrule
\end{tabular}
\end{center}

The standardized chord innovations reproduce the final weighted residual sum
of squares to numerical precision, which checks the recursive identity in
\Cref{thm:innovation}.  The predictive and lack-of-fit results reject the simple
target-plus-molecule additive representation on this sample.  The replay is
small and retrospective.  Target selectivity, assay heterogeneity, and
publication selection remain confounded, so the rejection does not identify a
mechanism.

The primary comparison should include additive Free--Wilson; matched-pair and
matched-series transfer; SAR-matrix neighborhood prediction; fingerprint or
graph models; random, diversity, and generic reaction-aware active learning;
and a fixed-scaffold bandit.  Cost to a correct additivity verdict and
certificate coverage should be co-primary with confirmed contract-satisfying
compounds per total cost.

\section{Evidence and novelty boundary}

The mathematical ingredients and close chemical representations have
established antecedents:
\begin{itemize}
  \item Free--Wilson and SAR-matrix methods already provide additive
  core--substituent tables, virtual missing cells, three-corner predictions,
  and prospective synthesis
  \cite{free1964qsar,wassermann2012sarm,gupta2014neighborhood,gupta2015sarm};
  \item reagent and side-chain nonadditivity was used to guide combinatorial
  library design before and alongside double-transformation-cycle analysis
  \cite{ge2002nonadditivity,patel2008nonadditivity,
        kramer2015nonadditivity,kramer2019nonadditivity};
  \item matched-pair and matched-series graphs, cross-core SAR transfer, and
  retrosynthetic pairs predate this response graph
  \cite{wawer2011mmsgraph,gupta2012graphmining,
        zhang2012sartransfer,deleon2014recapmmp};
  \item graph Hodge theory, effective-resistance design, skeleton completion,
  algebraic matroids, and anytime-valid test martingales are established
  \cite{jiang2011hodge,xu2019measurement,goreinov1997pseudoskeleton,
        kiraly2015algebraic,howard2021timeuniform,koolen2022evalues};
  \item check matrices, information sets, syndromes, weighted matroid bases,
  and optimum linear design are established
  \cite{hamming1950codes,prange1962information,whitney1935matroid,
        edmonds1971greedy,kiefer1959design};
  \item linear contrast estimability, robust optimization, and
  decision-directed sequential linear design are established
  \cite{rao1962generalized,searle1965estimable,elfving1952allocation,
        bartok2014partial,soyster1973robust,bental1998robust,
        fiez2019transductive,reda2021topm};
  \item shortest-path and Steiner-tree design, hardness, and fixed-terminal
  algorithms are established
  \cite{dijkstra1959shortest,dreyfus1971steiner,garey1977steiner,
        muller1987bipartite};
  \item square-root allocation, multiresponse SOCP design, and tree traffic
  accounting are established
  \cite{neyman1934stratified,sagnol2011multiresponse,lobo1998socp,
        hu1974communication};
  \item reaction-based enumeration and generic chemical active learning also
  predate this chapter
  \cite{konze2019reactional,graff2021molpal,bellamy2022bo}.
\end{itemize}

Accordingly, this chapter does not claim novelty for response tables, rectangle
contrasts, chemical-series graphs, reaction-aware enumeration, systematic-code,
matroid, linear-estimability, robust-ranking, shortest-path, Steiner-tree,
conic-optimal-design, or generic active-learning methods individually.  The prospective
research target is their integration after chemical typing and parent-response
quotienting: a reaction-feasible graph with systematic transfer residuals,
decision-specific molecular terminals, contrasts, and assay precision costs,
anytime-valid
outcome-adaptive circuit testing, and certificates evaluated on untouched
future assays.
Whether that integration is novel remains to be established by a
literature-complete comparison and prospective evidence.
% END SOURCE: chapters/response_core.tex


%--------Part IV: Algorithms and Complexity--------
\part{Algorithms under Explicit Assumptions}
\label{part:algorithms}

\chapter{Screening Complexity}
\label{ch:complexity}
% BEGIN SOURCE: chapters/ch11_complexity.tex
% ======================================================================
% Chapter 11 content file: chapters/ch11_complexity.tex
% Screening Complexity
% (Do NOT include \chapter{...} here; main.tex already does that.)
% ======================================================================

\begin{quote}
\textit{A guarantee is just a bound you agreed to pay.}
\end{quote}

\section{The Two Faces of Screening\index{screening} Hardness}
\label{sec:ch11-intro}

How hard is it to find the best molecules in chemical space\index{chemical space}? The answer depends on what ``hard'' means and what guarantees we require.

The preceding chapters developed two distinct frameworks for understanding scaffold\index{scaffold}-based screening:
\begin{itemize}
    \item \textbf{Combinatorial:} The Galois closure (Chapter~\ref{ch:galois}) determines the minimum output size for any scaffold\index{scaffold}-sound algorithm. The overhead $\ScafOverhead(T_k)=|c(T_k)|/|T_k|$ measures multiplicative output inflation.

    \item \textbf{Statistical:} The variance spectrum\index{variance spectrum} (Chapter~\ref{ch:variance}) and class distributions (Chapter~\ref{ch:probability}) determine how much information scaffold identity provides. The irreducible error\index{irreducible error} $\Err(\Scaf)$ bounds prediction accuracy.
\end{itemize}

This chapter distinguishes three resources in molecular screening\index{screening}:
\begin{enumerate}
    \item \textbf{Output-size bounds:} Scaffold-sound inclusion forces whole classes into the output. Oracle-query bounds require a separate observation model.
    \item \textbf{Statistical lower bounds:} In a specified Gaussian observation model, identification needs enough queries to distinguish instances with different correct answers.
    \item \textbf{Computational costs:} Assigning supplied interaction terms to a fixed partition is linear in the term list. Optimizing the partition is a different problem; the nonnegative hypergraph cut considered here is polynomial.
\end{enumerate}

The synthesis is a \emph{hardness taxonomy}: a two-dimensional characterization of screening targets by their combinatorial overhead and statistical irreducibility. The resulting descriptors motivate hypotheses about screening performance; they do not determine query complexity.

\begin{framed}
\noindent\textbf{Why complexity matters for screening.}

\medskip
\noindent
A useful screening guarantee specifies the output, the observations, and the loss.
Returning the whole library guarantees inclusion without scoring anything.
Identifying its best molecule is a different task. Likewise, a low average
prediction error need not preserve rare top hits. The response-specific
certificates in Chapter~\ref{ch:response-certificates} make the missing
comparison assumptions explicit.

\end{framed}

%------------------------------------------------------------------
\section{Formal Screening Models}
\label{sec:ch11-screening-models}

We begin by formalizing what screening algorithms do and what they try to achieve.

\subsection{Screening Objectives}

\begin{defn}[Top-$k$ recovery]
\label{defn:ch11-top-k}
Work on a finite library $\Mol$ of size $N$ and fix a deterministic rule for
breaking score ties. Let $T_k$ contain the $k$ smallest scores, with $1\leq k<N$.
\emph{Inclusion recovery} returns $\hat T\supseteq T_k$; \emph{exact recovery}
returns $\hat T=T_k$. Inclusion recovery needs an output-size or false-positive
constraint to exclude the zero-query solution $\hat T=\Mol$.
\end{defn}

\begin{defn}[Threshold hits]
\label{defn:ch11-threshold-hits}
For threshold $\tau \in \R$, the \emph{hit set} is:
\[
H_\tau := \{m \in \Mol : F(m) \leq \tau\}
\]
The \emph{threshold hit recovery problem} is to output a set $\hat{H} \supseteq H_\tau$.
\end{defn}

\begin{defn}[Quantile objective]
\label{defn:quantile-objective}
For quantile level $\alpha \in (0,1)$, the \emph{$\alpha$-quantile score} of scaffold $s$ is:
\[
q_s(\alpha) := \inf\{x : \Pr(F \leq x \mid \Scaf = s) \geq \alpha\}
\]
The \emph{scaffold quantile ranking problem} is to identify scaffolds with smallest $q_s(\alpha)$.
\end{defn}

These objectives differ in their tolerance for errors:
\begin{itemize}
    \item Top-$k$: Exact recovery fixes the answer size; inclusion recovery also needs an output constraint
    \item Threshold hits: Must find \emph{all} molecules below threshold
    \item Quantile ranking: Tolerates individual molecule\index{molecule} errors if scaffold ranking is correct
\end{itemize}

\subsection{Query Models}

\begin{defn}[Sequential query model]
\label{defn:query-model}
An algorithm operates as follows:
\begin{enumerate}
    \item Initialize with prior information (scaffold structure, historical data)
    \item Repeatedly: select molecule\index{molecule} $m_t$, observe score $F(m_t)$ (or noisy version)
    \item After budget $B$ queries, output candidate set $\hat{T}$
\end{enumerate}
\end{defn}

\begin{defn}[Multi-fidelity query model]
\label{defn:multi-fidelity-query}
A request specifies a molecule, a scoring protocol $f$, and a replicate or
randomization identity. Protocol $f$ returns an observation $Y^{(f)}(m)$ at cost
$c_f$. Its mean, bias relative to a chosen endpoint, and noise distribution are
separate quantities. Docking, MM-GBSA, and FEP may supply different protocols;
their names alone imply neither a common estimand nor decreasing error.
The algorithm chooses both the molecule and the protocol.

\end{defn}

\subsection{Algorithm Classes}

\begin{defn}[Scaffold-sound algorithm]
\label{defn:scaffold-sound}
An algorithm is \emph{scaffold-sound} if its decision to include molecule $m$ in the output depends only on $\Scaf(m)$:
\[
m \in \hat{T} \iff \Scaf(m) \in \hat{S}
\]
for some scaffold set $\hat{S} \subseteq S$.

Equivalently: the output is $\gamma(\hat{S})$ for some $\hat{S}$.
\end{defn}

\begin{defn}[Scaffold-aware algorithm]
\label{defn:scaffold-aware}
An algorithm is \emph{scaffold-aware} if it uses scaffold identity for stratification or prior structure, but can use additional features (decorations, fingerprints) for final selection.
\end{defn}

Scaffold-sound algorithms are a strict subset of scaffold-aware algorithms. The closure bounds of Chapter~\ref{ch:galois} apply to scaffold-sound algorithms; scaffold-aware algorithms can potentially do better by exploiting within-scaffold information.

%------------------------------------------------------------------
\section{Deterministic Lower Bounds from Closure}
\label{sec:ch11-deterministic-bounds}

We now develop the complexity implications of the Galois closure.

\subsection{The Closure Lower Bound}

\begin{thm}[Closure lower bound for top-$k$]
\label{thm:closure-lower-bound}
Any scaffold-sound output containing $T_k$ contains $c(T_k)$ and therefore has size at least $|c(T_k)|$. This is an output-cardinality bound.
\end{thm}

\begin{proof}
By definition of scaffold-sound (Definition~\ref{defn:scaffold-sound}), the output has the form $\gamma(\hat{S})$ for some scaffold set $\hat{S}$.

For the output to contain $T_k$, we need $T_k \subseteq \gamma(\hat{S})$. Applying $\alpha$ to both sides:
\[
\alpha(T_k) \subseteq \alpha(\gamma(\hat{S})) = \hat{S}
\]
Here $S$ is the realized image of the scaffold map, so each selected class is
nonempty and $\alpha(\gamma(\hat S))=\hat S$. Equivalently, the implication
follows directly from the adjunction.

Therefore $\hat{S} \supseteq \alpha(T_k)$, which implies:
\[
\gamma(\hat{S}) \supseteq \gamma(\alpha(T_k)) = c(T_k)
\]

Thus $|\text{output}| = |\gamma(\hat{S})| \geq |c(T_k)|$.

The proof uses no score-query assumption and therefore yields no oracle-query lower bound.
\end{proof}

\begin{defn}[Overhead hardness index]
\label{defn:overhead-hardness}
The \emph{overhead hardness index} at level $k$ is:
\[
\mathfrak{h}_{\text{closure}}(k) := \ScafOverhead(T_k) = \frac{|c(T_k)|}{k}
\]
This measures the multiplicative overhead imposed by scaffold-sound guarantees.
\end{defn}

\begin{exmp}[Overhead in practice]
\label{ex:overhead-practice}
For an illustrative finite library, suppose the counts are:
\begin{center}
\begin{tabular}{rrrr}
\toprule
$k$ & $|T_k|$ & $|c(T_k)|$ & $\mathfrak{h}_{\text{closure}}(k)$ \\
\midrule
10 & 10 & 847 & 84.7 \\
100 & 100 & 3,214 & 32.1 \\
1000 & 1000 & 12,891 & 12.9 \\
\bottomrule
\end{tabular}
\end{center}

These hypothetical counts illustrate decreasing overhead. Neither the values nor monotonicity are empirical claims.
\end{exmp}

\begin{prop}[Closure grows; its ratio need not decrease]
\label{prop:overhead-monotonicity}
For $k_1<k_2$, $c(T_{k_1})\subseteq c(T_{k_2})$, but
$|c(T_k)|/k$ need not be monotone.
\end{prop}
\begin{proof}
The inclusion follows from $T_{k_1}\subseteq T_{k_2}$ and monotonicity of closure.
If the first hit occupies a singleton class and the second belongs to a new
class of size $100$, the overhead increases from $1$ to $101/2$. If the first
two hits instead occupy the same class of size $100$, it decreases from $100$
to $50$.
\end{proof}

\begin{prop}[An unstructured exact-query lower bound]
\label{prop:unstructured-query-bound}
For arbitrary distinct real scores on a finite library, any deterministic
algorithm guaranteed to recover $T_k$ exactly from noiseless score queries,
without prior score restrictions, must query all $N$ molecules in the worst case.
\end{prop}
\begin{proof}
If it stops with an unqueried molecule $x$, the observed transcript is compatible
both with $x$ scoring below every other molecule and with $x$ scoring above
every other molecule. Complete all other unobserved scores identically and
without ties. The two top-$k$ sets differ, yet the algorithm returns the same
answer. Hence such a stopping rule cannot be correct on every score assignment.
\end{proof}

A scaffold or surrogate assumption can exclude these adversarial completions.
The saving then comes from that assumption and its validation, not from the
closure cardinality alone.

\subsection{Order-Saturated Variants}

Some algorithms impose additional structure: they must select \emph{downsets} in the scaffold poset\index{poset} (all predecessors of selected scaffolds).

\begin{defn}[Downset-constrained algorithm]
\label{defn:downset-constrained}
\index{downset}
An algorithm is \emph{downset-constrained} if its scaffold selection $\hat{S}$ must be a downset: $t \preceq s \in \hat{S} \Rightarrow t \in \hat{S}$.
\end{defn}

This constraint arises naturally in hierarchical search: if you include a complex scaffold, you implicitly include all its fragments.

\begin{thm}[Downset closure lower bound]
\label{thm:downset-lower-bound}
Any downset-constrained scaffold-sound algorithm outputting a superset of $T_k$ must output at least $|\gamma(\downarrow \alpha(T_k))|$ molecules, where $\downarrow B = \{s : \exists t \in B, s \preceq t\}$ is the downward closure.
\end{thm}

\begin{proof}
The output is $\gamma(\hat{S})$ for a downset $\hat{S}$. To contain $T_k$, we need $\alpha(T_k) \subseteq \hat{S}$. Since $\hat{S}$ is a downset:
\[
\downarrow \alpha(T_k) \subseteq \hat{S}
\]
Therefore $|\gamma(\hat{S})| \geq |\gamma(\downarrow \alpha(T_k))|$.
\end{proof}

\begin{cor}[Downset overhead ratio]
\label{cor:downset-ratio}
The downset constraint increases overhead by a factor:
\[
\frac{|\gamma(\downarrow \alpha(T_k))|}{|c(T_k)|} = \frac{|\gamma(\downarrow \alpha(T_k))|}{|\gamma(\alpha(T_k))|}
\]
This ratio measures how ``hierarchically spread'' the hit scaffolds are.
\end{cor}

%------------------------------------------------------------------
\section{Statistical Lower Bounds}
\label{sec:ch11-statistical-bounds}

Deterministic bounds address worst-case guarantees. Real algorithms accept probabilistic guarantees: ``find top-$k$ with probability $\geq 1 - \delta$.'' This section develops information-theoretic lower bounds.

\subsection{The Gap Parameter}

\begin{defn}[Scaffold gap]
\label{defn:scaffold-gap}
For scaffold $s$ with mean score $\mu_s = \E[F \mid \Scaf = s]$, the \emph{gap} to the best scaffold is:
\[
\Delta_s := \mu_s - \min_{t \in S} \mu_t
\]
The best scaffold $s^* = \arg\min_s \mu_s$ has gap $\Delta_{s^*} = 0$.
\end{defn}

The gap measures how suboptimal scaffold $s$ is on average. Small gaps make scaffolds hard to distinguish; large gaps make identification easier.

\begin{defn}[Hardness parameter]
\label{defn:hardness-parameter}
The \emph{hardness parameter} for scaffold identification is:
\[
H := \sum_{s \neq s^*} \frac{\sigma_s^2}{\Delta_s^2}
\]
where $\sigma_s^2$ denotes the observation variance for independent sampling within scaffold $s$. It is not automatically the variance of repeated measurements on one fixed molecule.
\end{defn}

Intuitively, $H$ sums the ``difficulty'' of distinguishing each suboptimal scaffold from the best: high variance and small gap both increase difficulty.

\subsection{Best Scaffold Identification}

\begin{thm}[Gaussian change-of-measure lower bound]
\label{thm:minimax-scaffold}
Suppose each query to scaffold $s$ returns an independent observation from
$\mathcal N(\mu_s,v_s)$, with known $v_s>0$ and unknown unrestricted means.
Assume a unique best mean. An algorithm is $\delta$-correct if, on every such
instance, it stops almost surely and identifies that scaffold with probability
at least $1-\delta$, where $0<\delta<1/2$. Then its stopping time $T$ satisfies
\[
\E[T]\ \geq\ 2\,\operatorname{kl}(1-\delta,\delta)
       \sum_{s\ne s^*}\frac{v_s}{\Delta_s^2},
\]
where $\operatorname{kl}(p,q)=p\log(p/q)+(1-p)\log((1-p)/(1-q))$.
\end{thm}
\begin{proof}
Fix $s\ne s^*$ and replace only its mean by $\mu_{s^*}-\eta$, $\eta>0$.
Let $P,Q$ be the laws of the stopped transcript under the original and altered
instances. The likelihood-ratio chain rule for adaptive sampling gives
\[
\operatorname{KL}(P\Vert Q)
 =\E_P[N_s(T)]\frac{(\Delta_s+\eta)^2}{2v_s}.
\]
For the event that the algorithm recommends the original best scaffold,
$P(E)\geq1-\delta$ and $Q(E)\leq\delta$. Data processing for relative entropy
therefore bounds this transcript divergence below by
$\operatorname{kl}(1-\delta,\delta)$. Let $\eta\downarrow0$ and sum the resulting
bounds over the suboptimal scaffolds. If the expected stopping time is
infinite, the stated inequality already holds.
\end{proof}

\begin{rem}[What the distribution assumption supplies]
\label{rem:gap-dependent}
The Gaussian assumption supplies the explicit KL divergence used in the proof.
A sub-Gaussian tail bound alone supplies an upper concentration bound, not this
KL identity for every distribution. The theorem is a valid hard-instance
lower bound for classes that contain these Gaussian instances. With equal
variances and equal suboptimal gaps it grows linearly in $|S|-1$.
A bound using a known positive minimum gap remains gap-dependent. For exact
identification there is no finite uniform sample bound as gaps approach zero.
\end{rem}

\subsection{Top-$k$ Scaffold Set Identification}
\begin{defn}[Top-$k$ scaffolds]
\label{defn:top-k-scaffolds}
Order means increasingly as $\mu_{(1)}\leq\cdots\leq\mu_{(|S|)}$ and suppose
$\mu_{(k)}<\mu_{(k+1)}$. Let $S_k^*$ contain the $k$ smallest means. This
objective concerns class means, not the $k$ best individual molecules.
\end{defn}

\begin{thm}[Boundary-gap lower bound for Gaussian top-$k$ identification]
\label{thm:fano-top-k}
Under the Gaussian observation model above, define
\[
d_s=\begin{cases}
\mu_{(k+1)}-\mu_s,&s\in S_k^*,\\
\mu_s-\mu_{(k)},&s\notin S_k^*.
\end{cases}
\]
Every algorithm that is $\delta$-correct on all instances with a strict
boundary gap satisfies
\[
\E[T]\geq2\,\operatorname{kl}(1-\delta,\delta)
                   \sum_{s\in S}\frac{v_s}{d_s^2}.
\]
\end{thm}
\begin{proof}
For an outside scaffold, lower only its mean to just below $\mu_{(k)}$; for an
inside scaffold, raise only its mean to just above $\mu_{(k+1)}$. In each case
the correct top-$k$ set changes. Apply the same transcript argument, let the
strict crossing margin tend to zero, and sum the per-scaffold query bounds.
\end{proof}

\begin{cor}[Scaffold-count dependence]
\label{cor:scaling-scaffold-count}
In the Gaussian model, if all variances equal $v$ and all boundary gaps equal
$d>0$, top-$k$ identification needs at least
$2|S|v d^{-2}\operatorname{kl}(1-\delta,\delta)$ expected queries.
Counting possible answers by $\log\binom{|S|}{k}$ alone does not justify a
logarithmic-in-$|S|$ adaptive query guarantee.
\end{cor}

\subsection{The Hardness Taxonomy}

We now synthesize the combinatorial and statistical bounds into a unified framework.

\begin{defn}[Two-axis hardness descriptor]
\label{defn:hardness-descriptor}
For a screening target with scoring function\index{scoring function} $F$, define:
\begin{enumerate}
    \item \textbf{Combinatorial hardness (x-axis):}
    \[
    \mathfrak{h}_{\text{comb}} := \mathfrak{h}_{\text{closure}}(k) = \frac{|c(T_k)|}{k}
    \]
    the scaffold overhead for top-$k$ recovery.

    \item \textbf{Statistical hardness (y-axis):}
    \[
    \mathfrak{h}_{\text{stat}} := \frac{\Err(\Scaf)}{\Var(F)}
    \]
    the fraction of score variance that is irreducible given scaffold identity, defined when $\Var(F)>0$. Neither axis encodes the mean gaps or rare-event probabilities required for query bounds.
\end{enumerate}
\end{defn}

\begin{figure}[h]
\centering
\begin{tikzpicture}[scale=1.2]
    % Axes
    \draw[->] (0,0) -- (5,0) node[right] {$\mathfrak{h}_{\text{comb}}$ (overhead)};
    \draw[->] (0,0) -- (0,4) node[above] {$\mathfrak{h}_{\text{stat}}$ (irreducibility)};

    % Quadrant labels
    \node at (1.2, 3.2) {\footnotesize Low overhead};
    \node at (1.2, 2.9) {\footnotesize High irreducibility};
    \node at (1.2, 2.5) {\footnotesize \textbf{Use features}};

    \node at (3.8, 3.2) {\footnotesize High overhead};
    \node at (3.8, 2.9) {\footnotesize High irreducibility};
    \node at (3.8, 2.5) {\footnotesize \textbf{Hard target}};

    \node at (1.2, 1.2) {\footnotesize Low overhead};
    \node at (1.2, 0.9) {\footnotesize Low irreducibility};
    \node at (1.2, 0.5) {\footnotesize \textbf{Scaffold wins}};

    \node at (3.8, 1.2) {\footnotesize High overhead};
    \node at (3.8, 0.9) {\footnotesize Low irreducibility};
    \node at (3.8, 0.5) {\footnotesize \textbf{Refine partition}};

    % Dividing lines
    \draw[dashed] (2.5, 0) -- (2.5, 4);
    \draw[dashed] (0, 2) -- (5, 2);

    % Example points
    \fill[blue] (0.8, 0.6) circle (2pt) node[right, font=\tiny] {Kinase A};
    \fill[red] (4.2, 3.5) circle (2pt) node[left, font=\tiny] {GPCR X};
    \fill[green!60!black] (3.8, 0.8) circle (2pt) node[left, font=\tiny] {Protease B};
    \fill[orange] (1.0, 3.2) circle (2pt) node[right, font=\tiny] {Target C};
\end{tikzpicture}
\caption{The hardness plane for scaffold-based screening. The plotted target names are schematic. The axes describe output inflation and mean-square residual error; placement does not certify query savings or top-hit recall.}
\label{fig:hardness-plane}
\end{figure}

\begin{rem}[Interpreting the two descriptors]
\label{prop:hardness-quadrants}
Low overhead limits output inflation; low irreducibility supports prediction
from scaffold identity under the chosen sampling law. High overhead motivates
finer output distinctions, and high irreducibility motivates decoration-level
features. None of these statements ranks extreme molecules. A rare severe
activity cliff can have little effect on average squared error while changing
the best molecule. Likewise, even one hit scaffold can have enormous overhead
if its fiber is large. The descriptors should be reported beside gap, tail,
and observed-recall diagnostics.
\end{rem}

%------------------------------------------------------------------

\section{Complexity of Factorization}
\label{sec:ch11-factorization-complexity}

We now address computational complexity: how hard is it to compute the energy decomposition of Chapter~\ref{ch:energy}?

\subsection{Tractability for Pairwise Potentials}

\begin{thm}[Term assignment for a supplied partition]
\label{thm:linear-time-decomposition}
Suppose a core/decoration atom assignment and the numerical interaction terms
are supplied. If $T$ terms involve at most $k$ atoms each, assigning each term
to the core, decoration, or boundary sum costs $O(kT+|m|+|P|)$ operations.
For ligand--protein pair terms with bounded evaluation cost, $T\leq|m||P|$,
so this is $O(|m||P|)$ when both atom sets are nonempty.
\end{thm}
\begin{proof}
Store the supplied membership of each atom, inspect the incident atoms of
each term, and add its value to the designated sum. Each inspection uses at
most $k$ membership tests. This proves the stated term-list bound.
\end{proof}

\begin{cor}[Scope of the operation count]
\label{cor:practical-efficiency}
The bound concerns assignment of terms at a fixed geometry. Scaffold matching,
pose search, relaxation, and computation of expensive interaction values are
additional costs. Higher-order terms do not by themselves make assignment
hard when their list and the partition are given.
\end{cor}

\subsection{Optimizing a Boundary Cut}
\begin{defn}[Interaction hypergraph]
\label{defn:interaction-hypergraph}
An interaction hypergraph $H=(V,E)$ has one vertex per atom and one hyperedge
per interaction term. For the cut problem below, each edge has a supplied
nonnegative rational penalty $w_e$; this need not equal a signed physical
interaction energy.
\end{defn}

\begin{defn}[Boundary hyperedge cut]
\label{defn:boundary-cut}
For a partition $V=C\sqcup D$, define
$\operatorname{Cut}(C,D)=\{e:e\cap C\ne\varnothing,\ e\cap D\ne\varnothing\}$.
The cut cost is $\sum_{e\in\operatorname{Cut}(C,D)}w_e$.
\end{defn}

\begin{thm}[Nonnegative terminal-constrained hypergraph cut]
\label{thm:np-hard-decomposition}
With disjoint sets of atoms fixed to $C$ and $D$, minimizing the nonnegative
boundary cut over all remaining assignments is polynomial-time solvable,
including for hyperedges of size three or larger.
\end{thm}
\begin{proof}
Use a directed minimum-cut network. For each hyperedge $e$, introduce auxiliary
vertices $a_e,b_e$ and an arc $a_e\to b_e$ of capacity $w_e$. For every $v\in e$,
add arcs $v\to a_e$ and $b_e\to v$ of capacity
$M>\sum_{e\in E}w_e$. Connect fixed core vertices from the source and fixed
decoration vertices to the sink with capacity $M$.
For any finite cut below $M$, a hyperedge touching both sides forces $a_e$ to
the source side and $b_e$ to the sink side, paying $w_e$. An edge wholly on
one side admits both auxiliaries on that side at zero cost. Thus minimizing
over auxiliaries reproduces exactly the hypergraph cut objective. The network
has $|V|+2|E|+2$ vertices and $O(|V|+|E|+\sum_e|e|)$ arcs, so ordinary
polynomial-time minimum cut solves the problem~\cite{lawler1973hypergraph}.
\end{proof}

\begin{rem}[Different constraints define different problems]
\label{rem:structure-matters}
A balance constraint, a chemically valid-core constraint, or signed edge
weights changes this optimization problem. Its complexity must be established
for those constraints. It cannot be inferred from hyperedge arity. Nor does
an efficiently found atom cut establish a transferable thermodynamic
factorization across relaxed molecules.
\end{rem}

\begin{cor}[Exact solution of the stated cut problem]
\label{cor:approximation-hardness}
No approximation barrier follows for the cut problem just stated: the
network construction computes its optimum exactly. The numerical precision
and bit complexity of its rational capacities are part of the input size.
\end{cor}

%------------------------------------------------------------------

\section{Combined Complexity Analysis}
\label{sec:ch11-combined}

Output size, oracle calls, and arithmetic work are distinct resources.

\begin{prop}[Separate accounting for a composite campaign]
\label{thm:end-to-end-complexity}
Suppose a campaign has two explicitly required outputs: a scaffold-sound
superset $\hat T\supseteq T_k$ from a finite library, and identification of the
best scaffold in the independent Gaussian query model of
Theorem~\ref{thm:minimax-scaffold}. Charge $c_{\rm out}$ for materializing each
output item and $c_{\rm query}$ per observation. If $T$ counts observations,
\[
\E[\mathrm{Cost}]=c_{\rm out}\E[|\hat T|]+c_{\rm query}\E[T]
\geq c_{\rm out}|c(T_k)|+
2c_{\rm query}\operatorname{kl}(1-\delta,\delta)
          \sum_{s\ne s^*}\frac{v_s}{\Delta_s^2}.
\]
\end{prop}
\begin{proof}
Apply the output-cardinality bound to the first deliverable and the Gaussian
query bound to the second, then add their separately charged costs.
\end{proof}

This is an accounting result for the stated composite task. Exact molecular
top-$k$ recovery need not require identifying class means, and identifying
those means need not recover the molecular top-$k$. There is no general
product or maximum of closure size and MLMC variance that bridges the two.

\begin{defn}[A descriptive screening index]
\label{defn:effective-hardness}
One may report the dimensionless descriptive index
\[
\mathfrak H(k)=\max\{\mathfrak h_{\rm comb}(k),
\mathfrak h_{\rm stat}\log|S|\}.
\]
It is not a lower or upper bound on oracle cost. Report its two components,
the sampling measure, and the objective so the index remains interpretable.
\end{defn}

%------------------------------------------------------------------

\section{Experiments}
\label{sec:ch11-experiments}

These are proposed diagnostics, not completed empirical results.

\begin{exper}[Hardness plane across targets]
\label{exper:hardness-plane}
On fully scored finite libraries, compute closure overhead at prespecified
values of $k$ and estimate irreducibility under a stated sampling measure.
Then compare these descriptors with held-out recall at matched total query
cost. Keep scaffold sizes, rare-hit frequencies, and baseline policies in the
analysis. The hypothesis is that the descriptors help explain performance;
target-class rankings and numerical gains are not assumed.
\end{exper}

\begin{exper}[Checking the closure bound]
\label{exper:closure-tightness}
For a known $T_k$, construct $c(T_k)$ directly. It attains the smallest
scaffold-sound superset exactly. For each algorithm record output cardinality
and oracle calls separately. A scaffold-aware output can be smaller than the
closure because it solves a less constrained output problem; it does not
violate the bound. Returning the entire library with no queries is the
required zero-information control.
\end{exper}

\begin{exper}[Statistical bound validation]
\label{exper:statistical-validation}
Use Gaussian scaffold arms with specified means and variances, vary the gaps,
and measure expected stopping time and empirical error of a fixed-confidence
procedure over independent repetitions. Compare with the explicit
change-of-measure lower bound. A comparison using real molecule scores must
separately justify the sampling law and concentration model; matching a fitted
$H$ does not verify those assumptions.
\end{exper}

\section{Summary}
\label{sec:ch11-summary}
Scaffold closure constrains output size. Statistical identification depends on
an observation family and gaps between competing answers. Fixed-partition
term assignment depends on the interaction-list size. Keeping these quantities
separate prevents an output-size argument from becoming an unsupported claim
about oracle efficiency. Chapter~\ref{ch:screening} uses explicit confidence
and molecule-level exclusion certificates to connect observations to outputs.

%------------------------------------------------------------------

\section{Exercises}

\begin{exer}[Prove closure lower bound]
\label{exer:ch11-closure}
Prove Theorem~\ref{thm:closure-lower-bound} directly from the Galois connection\index{Galois connection} properties of Chapter~\ref{ch:galois}:
\begin{enumerate}
    \item[(a)] Show that $\alpha \circ \gamma \circ \alpha = \alpha$ (closure property).
    \item[(b)] Use this to prove $T_k \subseteq \gamma(\hat{S}) \Rightarrow \alpha(T_k) \subseteq \hat{S}$.
    \item[(c)] Conclude the lower bound on $|\gamma(\hat{S})|$.
\end{enumerate}
\end{exer}

\begin{exer}[Two-scaffold Le Cam bound]
\label{exer:ch11-lecam}
Consider the simplest case: two scaffolds $s_1, s_2$ with distributions $\mathcal{N}(\mu_1, \sigma^2)$ and $\mathcal{N}(\mu_2, \sigma^2)$.
\begin{enumerate}
    \item[(a)] Write the log-likelihood ratio for $n$ observations from scaffold $s_i$.
    \item[(b)] Using Le Cam's two-point method, show that identifying the better scaffold requires:
    \[
    n \geq \frac{c\sigma^2}{(\mu_1 - \mu_2)^2} \log\frac{1}{\delta}
    \]
    \item[(c)] Interpret: why does small gap $|\mu_1 - \mu_2|$ make identification hard?
\end{enumerate}
\end{exer}

\begin{exer}[Downset constraint impact]
\label{exer:ch11-downset}
For a scaffold poset\index{poset} with depth $d$ (maximum chain length):
\begin{enumerate}
    \item[(a)] Give a depth-two poset whose greatest element has arbitrarily many predecessors, disproving a bound by $d\,|\alpha(T_k)|$ based on depth alone.
    \item[(b)] Prove the bound $|\downarrow A|\leq d|A|$ when every principal downset is a chain with at most $d$ elements. Give an example attaining order $d$ overhead.
    \item[(c)] When is the downset constraint ``free'' (no additional overhead)?
\end{enumerate}
\end{exer}

\begin{exer}[Implement hardness visualization]
\label{exer:ch11-implementation}
Write code to:
\begin{enumerate}
    \item[(a)] Compute $\mathfrak{h}_{\text{comb}}(k)$ from docking scores and scaffold assignments.
    \item[(b)] Compute $\mathfrak{h}_{\text{stat}}$ from the variance spectrum\index{variance spectrum}.
    \item[(c)] Generate a hardness plane plot for a multi-target dataset.
    \item[(d)] Color points by target class and identify patterns.
\end{enumerate}
\end{exer}

\begin{exer}[Hypergraph cut reduction]
\label{exer:ch11-hypergraph}
Verify the network construction in Theorem~\ref{thm:np-hard-decomposition}:
\begin{enumerate}
    \item[(a)] Enumerate all partitions of a three-vertex hyperedge.
    \item[(b)] Minimize the auxiliary-vertex cut cost for each partition.
    \item[(c)] Check equality with the all-or-nothing hyperedge penalty.
    \item[(d)] Explain why a signed interaction energy is not a nonnegative cut capacity.
\end{enumerate}

\end{exer}

%------------------------------------------------------------------
\section{Bibliographic Notes}

The change-of-measure argument is standard in best-arm identification;
see~\cite{lattimore2020bandits}. It is given explicitly here for Gaussian
observations so that its distribution assumptions and query counts can be
checked. General entropy bounds require a specified family of hypotheses;
counting possible answers is insufficient~\cite{cover2006information}.

The hypergraph network construction is the nonnegative cut reduction of
Lawler~\cite{lawler1973hypergraph}. It establishes polynomial solvability,
not NP-hardness. Closure cardinality follows directly from the abstraction
adjunction; any claim that it predicts practical screening savings requires
separate empirical evidence.

%------------------------------------------------------------------

\section{Connections}

\begin{itemize}
    \item \textbf{Chapter~\ref{ch:galois}:} Provides the Galois closure operator\index{closure operator} underlying combinatorial bounds.
    \item \textbf{Chapter~\ref{ch:variance}:} Provides the variance spectrum for the statistical irreducibility axis.
    \item \textbf{Chapter~\ref{ch:energy}:} Factorization complexity results unified here.
    \item \textbf{Chapter~\ref{ch:mlmc}:} Uses hardness to guide multilevel allocation.
    \item \textbf{Chapter~\ref{ch:screening}:} Designs algorithms to match the hardness taxonomy.
\end{itemize}
% END SOURCE: chapters/ch11_complexity.tex


\chapter{Multilevel Monte Carlo on Scaffold Chains}
\label{ch:mlmc}
% BEGIN SOURCE: chapters/ch12_mlmc.tex
% ======================================================================
% Chapter 12 content file: chapters/ch12_mlmc.tex
% Multilevel Monte Carlo on Scaffold Chains
% (Do NOT include \chapter{...} here; main.tex already does that.)
% ======================================================================

\begin{quote}
\textit{If corrections get smaller with scale, sample the coarse scales often and the fine scales rarely.}
\end{quote}

\section{The Multilevel Opportunity}
\label{sec:ch12-intro}

The preceding chapters established a hierarchical view of chemical space\index{chemical space}: molecules live in scaffold classes, which themselves form refinement towers $\mathcal{S}_0 \subset \mathcal{S}_1 \subset \cdots \subset \mathcal{S}_L$. The variance spectrum\index{variance spectrum} (Chapter~\ref{ch:variance}) quantifies how much information each refinement level contributes.

Multilevel Monte Carlo\index{multilevel Monte Carlo} (MLMC) can exploit a
hierarchy when we can actually sample its corrections cheaply and with low
variance. A hierarchy of exact conditional means is a useful mathematical
decomposition, but it is not yet such a sampling algorithm. This chapter
separates those two objects and states the conditions for valid allocation.

\begin{framed}
\noindent\textbf{The MLMC\index{multilevel Monte Carlo} principle.}

\medskip
\noindent
If you want to estimate a quantity at fine resolution:
\begin{enumerate}
    \item Specify an evaluable coarse approximation and the intended finest-level estimand
    \item Add corrections for each refinement level
    \item Allocate according to actual correction variance and full sampling cost
    \item Optimal allocation minimizes total cost for fixed accuracy
\end{enumerate}

Allocation depends on the variances of the correction samples actually
computed and on their full costs. The population spectrum
$\{\|D_\ell\|^2\}$ can suggest useful refinements, but it does not replace
measurement of those correction variances. A speedup is a comparison of two
costs at the same bias and mean-square error.
\end{framed}

This chapter develops:
\begin{enumerate}
    \item \textbf{Multilevel estimators} for scaffold\index{scaffold} statistics (means, hit probabilities, quantiles)
    \item \textbf{Optimal allocation} across levels, driven by sampled correction variances and costs
    \item \textbf{Parent-child correlation} as a diagnostic for refinement quality
    \item \textbf{Multi-fidelity extension} combining scaffold hierarchy with scoring method hierarchy
\end{enumerate}

%------------------------------------------------------------------
\section{What We Estimate}
\label{sec:ch12-quantities}

Before building multilevel estimators, we specify the target quantities.

\subsection{Scaffold-Level Statistics}

\begin{defn}[Scaffold mean]
\label{defn:scaffold-mean}
The mean score for scaffold $s$ at refinement level $\ell$ is:
\[
\mu_s := \E[F \mid \Scaf_\ell = s]
\]
where $\Scaf_\ell: \Mol \to S_\ell$ is the scaffold map at level $\ell$.
\end{defn}

\begin{defn}[Scaffold hit probability]
\label{defn:scaffold-hit-prob}
For threshold $\tau$, the hit probability is:
\[
p_s(\tau) := \Pr(F \leq \tau \mid \Scaf_\ell = s)
\]
\end{defn}

\begin{defn}[Scaffold quantile]
\label{defn:scaffold-quantile}
The $\alpha$-quantile of scaffold $s$ is:
\[
q_s(\alpha) := \inf\{x : \Pr(F \leq x \mid \Scaf_\ell = s) \geq \alpha\}
\]
\end{defn}

These statistics serve different screening\index{screening} objectives:
\begin{itemize}
    \item $\mu_s$: Rank scaffolds by average quality (expected value optimization)
    \item $p_s(\tau)$: Estimate enrichment at a threshold (hit rate optimization)
    \item $q_s(\alpha)$: Rank a specified lower quantile, which need not rank the single best molecule
\end{itemize}

\subsection{Level-Indexed Statistics}

\begin{defn}[Level-$\ell$ class statistic]
\label{defn:level-statistic}
For a functional $\Phi$ (mean, hit probability, quantile), define:
\[
\theta_\ell(u) := \Phi(\text{Law}(F \mid \Scaf_\ell = u))
\]
where $u \in S_\ell$ is a class at level $\ell$.
\end{defn}

\begin{exmp}[Means across levels]
\label{ex:means-across-levels}
Consider a three-level tower:
\begin{itemize}
    \item $\ell = 0$: A fixed coarse graph descriptor
    \item $\ell = 1$: That descriptor together with the Murcko scaffold
    \item $\ell = 2$: The preceding labels together with mapped attachment features
\end{itemize}

Keeping the preceding labels makes nesting explicit. Names such as
``generic scaffold'' do not establish a refinement relation: deleting atom or
bond types usually coarsens a scaffold. Each level map must determine its
parent map for the tower identities to apply.
\end{exmp}

%------------------------------------------------------------------
\section{Telescoping Along Scaffold Chains}
\label{sec:ch12-telescoping}

The key to MLMC is the telescoping identity: a fine-level statistic equals the coarsest approximation plus cumulative corrections.

\subsection{The Telescoping Sum}

\begin{prop}[Telescoping decomposition]
\label{prop:telescoping}
For a scaffold chain $u_0 \to u_1 \to \cdots \to u_L$ where $u_\ell$ is the ancestor of $u_L$ at level $\ell$:
\[
\theta_L(u_L) = \theta_0(u_0) + \sum_{\ell=1}^{L} \bigl[\theta_\ell(u_\ell) - \theta_{\ell-1}(u_{\ell-1})\bigr]
\]
\end{prop}

\begin{proof}
Immediate from telescoping:
\begin{align}
\theta_L(u_L) &= \theta_0(u_0) + (\theta_1(u_1) - \theta_0(u_0)) + (\theta_2(u_2) - \theta_1(u_1)) + \cdots + (\theta_L(u_L) - \theta_{L-1}(u_{L-1})) \\
&= \theta_0(u_0) + \sum_{\ell=1}^{L} (\theta_\ell(u_\ell) - \theta_{\ell-1}(u_{\ell-1}))
\end{align}
\end{proof}

\begin{defn}[Correction random field]
\label{defn:correction-field}
The \emph{level-$\ell$ correction} at scaffold $u_\ell$ is:
\[
\Delta_\ell(u_\ell) := \theta_\ell(u_\ell) - \theta_{\ell-1}(\text{parent}(u_\ell))
\]
where $\text{parent}(u_\ell) \in S_{\ell-1}$ is the coarser scaffold containing $u_\ell$.
\end{defn}

The correction $\Delta_\ell$ measures how much the statistic changes when refining from level $\ell-1$ to level $\ell$.

\subsection{MLMC Estimators}

The telescope is an identity among numbers. To turn it into an estimator,
we must supply random variables with those expectations.

\begin{defn}[Implementable multilevel estimator]
\label{defn:mlmc-estimator}
For a fixed chain, suppose $Q_\ell$ is an integrable, evaluable random variable
with $\E[Q_\ell]=\theta_\ell(u_\ell)$. Set $Y_0=Q_0$. For $\ell\geq1$, specify a
joint law for $(Q_\ell,Q_{\ell-1})$ with the required marginals and set
$Y_\ell=Q_\ell-Q_{\ell-1}$. Draw independent batches across levels and i.i.d.
pairs within each batch. Then
\[
\widehat\theta_L=\sum_{\ell=0}^L\frac1{n_\ell}
                  \sum_{i=1}^{n_\ell}Y_{\ell,i},\qquad
\E[\widehat\theta_L]=\theta_L(u_L).
\]
Write $v_\ell=\Var(Y_\ell)$ and let $c_\ell>0$ include both members of a paired
sample and its preparation cost.
\end{defn}

For class means, a valid pair is
$Q_\ell=F(M_\ell)$ and $Q_{\ell-1}=F(M_{\ell-1})$ with
$M_j\sim\mu(\cdot\mid\Scaf_j=u_j)$. Independent draws have the right marginals,
but may have large difference variance and require two full score evaluations.
An efficient coupling is additional structure. Merely assigning a molecule to
a coarser class does not produce a cheaper score observation.

A useful alternative holds the sampling law fixed, for example
$\nu=\mu(\cdot\mid\Scaf_L=u_L)$, and evaluates successive scoring
approximations $Q_\ell$ on that law. The telescope then estimates
$\E_\nu[Q_L]$, provided each correction preserves its stated marginals.

\begin{rem}[Estimated parents, reuse, and bias]
\label{rem:delta-estimation}
An unknown parent mean cannot be inserted as though it were an exact control
variate. An independent unbiased parent estimate preserves the correction
mean but adds its own variance. Reusing parent estimates across levels creates
covariance; in that case
\[
\Var(\widehat\theta_L)=\sum_\ell\Var(\widehat Y_\ell)
 +2\sum_{j<\ell}\Cov(\widehat Y_j,\widehat Y_\ell),
\]
and the independent-batch allocation below does not apply unchanged.
A fitted regression may also introduce bias. Train it on separate data,
condition on that fit, and correct its residual under the target sampling law;
account for training, correction variance, and any remaining bias. Shared-parent
covariance is especially relevant to the contrasts in
Chapter~\ref{ch:response-certificates}.
\end{rem}

\begin{rem}[Hit probabilities and quantiles]
For a hit probability, use paired indicators
$\mathbf 1\{Q_\ell\leq\tau\}-\mathbf 1\{Q_{\ell-1}\leq\tau\}$ under the
specified laws; their variance must be measured separately. Quantiles are
nonlinear functionals. A difference of raw scores is not an unbiased quantile
correction. One route estimates a CDF and inverts a confidence band. Translating
a CDF error $\eta$ into a quantile error $\eta/a$ requires a density bounded
below by $a>0$ near that quantile, with the band remaining in that region.
\end{rem}

%------------------------------------------------------------------

\section{The Variance Spectrum as the MLMC Driver}
\label{sec:ch12-variance-driver}

The population variance spectrum describes exact conditional expectations. The variances $v_\ell$ in an implemented estimator describe random correction observations. They coincide only under additional sampling and evaluation assumptions.

\subsection{Corrections as Martingale\index{martingale} Differences}

\begin{prop}[Mean case = martingale\index{martingale} differences]
\label{prop:martingale-differences}
When $\theta_\ell(u) = \E[F \mid \Scaf_\ell = u]$ (the conditional mean), the correction at a random molecule $m$ satisfies:
\[
\Delta_\ell(\Scaf_\ell(m)) = D_\ell(m)
\]
where $D_\ell$ is the martingale difference from Chapter~\ref{ch:variance}.

Consequently:
\[
\Var(\Delta_\ell) = \|D_\ell\|^2
\]
is literally the level-$\ell$ component of the variance spectrum.
\end{prop}

\begin{proof}
Recall from Chapter~\ref{ch:variance}:
\[
D_\ell(m) := A_\ell(m) - A_{\ell-1}(m) = \E[F \mid \Scaf_\ell](m) - \E[F \mid \Scaf_{\ell-1}](m)
\]

For a random molecule $m$ with $\Scaf_\ell(m) = u_\ell$ and $\Scaf_{\ell-1}(m) = u_{\ell-1} = \text{parent}(u_\ell)$:
\begin{align}
D_\ell(m) &= \E[F \mid \Scaf_\ell = u_\ell] - \E[F \mid \Scaf_{\ell-1} = u_{\ell-1}] \\
&= \theta_\ell(u_\ell) - \theta_{\ell-1}(u_{\ell-1}) \\
&= \Delta_\ell(u_\ell)
\end{align}
\end{proof}

This identity is for $M\sim\mu$ and exact conditional means. In particular,
$\E[D_\ell]=0$ for $\ell\geq1$, and $\E[A_\ell]=\E[F]$ at every level. If the
maps $A_\ell$ were already known and evaluable, their differences would be
available; learning the maps is an additional statistical task. Conditioning
on one chosen finest class changes the law, so the global spectrum does not
give that class estimator's correction variance.

For example, let every child class contain equally likely scores $-1$ and $1$.
All conditional means vanish, so $\|D_\ell\|^2=0$. Nevertheless, independent
child and parent score draws have difference variance $2$. Treating the
zero population increment as zero sampling cost would be incorrect.

\subsection{The MLMC Cost Functional}

\begin{defn}[MLMC cost functional]
\label{defn:mlmc-variance-functional}
For implementable correction variances $v=(v_0,\ldots,v_L)$ and costs $c$, set
\[
V(v;c)=\left(\sum_{\ell=0}^L\sqrt{v_\ell c_\ell}\right)^2.
\]
This notation replaces the use of population projection energies as if they
were observable correction variances.
\end{defn}

\begin{thm}[Optimal independent-batch allocation]
\label{thm:optimal-allocation}
Under Definition~\ref{defn:mlmc-estimator}, assume $v_\ell>0$ and $c_\ell>0$.
Let $b=\theta_L-\theta_\star$ be the bias relative to the intended target and
let $q=\epsilon^2-b^2>0$. Over positive real sample sizes, the least cost giving
MSE at most $\epsilon^2$ is
\[
n_\ell^*=\frac1q\sqrt{\frac{v_\ell}{c_\ell}}
                \sum_{j=0}^L\sqrt{v_jc_j},\qquad
\mathrm{Cost}^*=\frac{V(v;c)}q.
\]
Taking $n_\ell=\lceil n_\ell^*\rceil$ preserves the MSE bound and adds at most
$\sum_\ell c_\ell$ cost. Pilot and model-training costs are additional.
\end{thm}
\begin{proof}
Independence gives
$\operatorname{MSE}=b^2+\sum_\ell v_\ell/n_\ell$. Cauchy--Schwarz yields
\[
\left(\sum_\ell\sqrt{v_\ell c_\ell}\right)^2
\leq\left(\sum_\ell v_\ell/n_\ell\right)
       \left(\sum_\ell n_\ell c_\ell\right).
\]
Thus feasible cost is at least $V(v;c)/q$. The stated allocation has sampling
variance $q$ and attains equality. Rounding up decreases each variance term.
\end{proof}

If a correction has zero variance and its constant value is known, it needs
no sampling. If that constant is unknown, at least one evaluation is needed;
a zero-variance formula does not reveal its value. Estimated pilot variances
also need safeguards: a small pilot can report zero despite nonzero true
variance.

\begin{cor}[Comparison with direct estimation]
\label{cor:mlmc-vs-naive}
Let a direct observation of $Q_L$ cost $c_{\rm direct}$ and have variance
$v_{\rm direct}$. At the same bias and variance budget $q$, the continuous
cost ratio is
\[
\frac{\mathrm{Cost}_{\rm direct}}{\mathrm{Cost}_{\rm MLMC}}
=\frac{v_{\rm direct}c_{\rm direct}}{V(v;c)}.
\]
MLMC saves sampling cost exactly when this ratio exceeds one. Integer, pilot,
and training costs must be included in a practical comparison.
\end{cor}

\begin{exmp}[Geometric correction variances]
\label{ex:exponential-decay}
Suppose actual correction variances satisfy $v_\ell=v_0\rho^\ell$ and costs
$c_\ell=c_0\gamma^\ell$, with $0<\rho<1<\gamma$. Then
\[
V(v;c)=v_0c_0\left(\sum_{\ell=0}^L(\rho\gamma)^{\ell/2}\right)^2.
\]
As $L$ grows, this is bounded if $\rho\gamma<1$, grows as $(L+1)^2$ if
$\rho\gamma=1$, and grows as $(\rho\gamma)^L$ if $\rho\gamma>1$.
The last case can still beat a direct cost growing as $\gamma^L$; it does not
by itself imply failure. To turn these statements into asymptotic complexity
in $\epsilon$, specify how the bias decreases with $L$~\cite{giles2015mlmc}.
\end{exmp}

%------------------------------------------------------------------

\section{Parent-Child Correlation Diagnostics}
\label{sec:ch12-correlation}

The variance spectrum can be measured via parent-child correlations---a quantity directly observable from data.

\subsection{Correlation-Spectrum Identity}

\begin{defn}[Parent-child correlation]
\label{defn:parent-child-correlation}
At level $\ell$, the parent-child correlation is:
\[
\rho_\ell := \text{Corr}(A_\ell, A_{\ell-1})
\]
where $A_\ell = \E[F \mid \Scaf_\ell]$ is the level-$\ell$ conditional expectation\index{conditional expectation}.
\end{defn}

\begin{prop}[Correlation-spectrum identity]
\label{prop:correlation-spectrum}
Assume $\Var(A_{\ell-1})>0$ and $\Var(A_\ell)>0$. Because $A_{\ell-1} = \E[A_\ell \mid \mathcal{S}_{\ell-1}]$ (tower property):
\[
\text{Cov}(A_\ell, A_{\ell-1}) = \Var(A_{\ell-1})
\]

Therefore:
\[
\rho_\ell^2 = \frac{\Var(A_{\ell-1})}{\Var(A_\ell)}
\]

And the variance spectrum component satisfies:
\[
\|D_\ell\|^2 = \Var(A_\ell) - \Var(A_{\ell-1}) = \Var(A_\ell)(1 - \rho_\ell^2)
\]
\end{prop}

\begin{proof}
By the tower property:
\[
\text{Cov}(A_\ell, A_{\ell-1}) = \E[A_\ell \cdot A_{\ell-1}] - \E[A_\ell]\E[A_{\ell-1}]
\]

Since $\E[A_\ell \mid \mathcal{S}_{\ell-1}] = A_{\ell-1}$:
\[
\E[A_\ell \cdot A_{\ell-1}] = \E[\E[A_\ell \mid \mathcal{S}_{\ell-1}] \cdot A_{\ell-1}] = \E[A_{\ell-1}^2]
\]

Therefore:
\[
\text{Cov}(A_\ell, A_{\ell-1}) = \E[A_{\ell-1}^2] - \E[A_\ell]\E[A_{\ell-1}] = \E[A_{\ell-1}^2] - \mu^2 = \Var(A_{\ell-1})
\]

(using $\E[A_\ell] = \E[A_{\ell-1}] = \E[F] = \mu$).

The correlation is:
\[
\rho_\ell = \frac{\text{Cov}(A_\ell, A_{\ell-1})}{\sqrt{\Var(A_\ell)\Var(A_{\ell-1})}} = \frac{\Var(A_{\ell-1})}{\sqrt{\Var(A_\ell)\Var(A_{\ell-1})}} = \sqrt{\frac{\Var(A_{\ell-1})}{\Var(A_\ell)}}
\]

So $\rho_\ell^2 = \Var(A_{\ell-1})/\Var(A_\ell)$, and:
\[
\|D_\ell\|^2 = \Var(A_\ell) - \Var(A_{\ell-1}) = \Var(A_\ell) - \rho_\ell^2 \Var(A_\ell) = \Var(A_\ell)(1 - \rho_\ell^2)
\]
\end{proof}

\subsection{Interpretation and Diagnostics}

\begin{prop}[High correlation = small correction]
\label{prop:high-correlation}
~
\begin{enumerate}
    \item If $\rho_\ell^2\approx1$: The projection increment is small relative to $\Var(A_\ell)$; its absolute size also depends on that variance.
    \item If $\rho_\ell^2\approx0$: Most of $\Var(A_\ell)$ is newly resolved at this level. Small correlation does not imply independence.
\end{enumerate}
\end{prop}

\begin{exmp}[Correlation decay pattern]
\label{ex:correlation-decay}
Illustrative values, not an observed or required pattern:
\begin{center}
\begin{tabular}{ccc}
\toprule
Level $\ell$ & $\rho_\ell$ & Interpretation \\
\midrule
1 & 0.95 & First nested feature: small relative increment \\
2 & 0.85 & Second nested feature: larger relative increment \\
3 & 0.70 & Third nested feature: larger relative increment \\
4 & 0.50 & Fourth nested feature: large relative increment \\
\bottomrule
\end{tabular}
\end{center}

The decreasing correlations describe relative projection increments in this illustration. They imply neither decreasing absolute increment variance nor a good cost--accuracy tradeoff. Empirical fitted means need held-out assessment; a zero-variance parent has undefined correlation.
\end{exmp}

\begin{rem}[Outliers indicate bad refinements]
\label{rem:bad-refinements}
An unexpectedly large projection increment can indicate an informative refinement, a sparse-cell estimation artifact, or a changed sampling law. Distinguish these explanations with held-out data and class occupancy diagnostics.
\end{rem}

%------------------------------------------------------------------
\section{Cluster Expansion and Cumulant Decay}
\label{sec:ch12-cluster}

Locality can motivate a low-variance approximation, but interaction range,
interaction order, and refinement level are distinct indices. A bound on one
coefficient does not bound the sum of all unresolved terms; the number of
terms and their dependence also matter.

\subsection{A Sufficient Approximation Condition}
\begin{prop}[An approximation bound for projection increments]
\label{prop:cluster-multilevel}
Let $\mathcal S_{\ell-1}\subseteq\mathcal S_\ell$, and suppose there exists an
$\mathcal S_{\ell-1}$-measurable approximation $G_{\ell-1}$ with
\[
\|F-G_{\ell-1}\|_2^2\leq C\rho^\ell,\qquad 0<\rho<1.
\]
Then the exact projection increment satisfies $\|D_\ell\|_2^2\leq C\rho^\ell$.
\end{prop}
\begin{proof}
Conditional expectation is the best $L^2$ approximation on its sigma-algebra,
so $\|F-A_{\ell-1}\|_2\leq\|F-G_{\ell-1}\|_2$.
Also $D_\ell=\E[F-A_{\ell-1}\mid\mathcal S_\ell]$; contraction of conditional
expectation proves the bound.
\end{proof}

For a convergent cumulant or interaction expansion, one sufficient route is
to show that all retained terms are measurable at level $\ell-1$ and that the
sum of the $L^2$ norms of the omitted terms is at most
$\sqrt C\rho^{\ell/2}$. The triangle inequality then supplies the hypothesis.
A locality conjecture, or a decay bound on individual terms without controlling
their multiplicity, is not enough.

\begin{rem}[From physical hypothesis to an estimator]
\label{rem:physical-intuition}
A chemically useful truncation should be tested for both approximation error
and evaluation cost. Even after establishing the projection bound, the MLMC
allocation requires the variance of the implemented coupled corrections.
Relaxation, shared measurement error, and fitted-parent uncertainty can change
that variance. These are properties to measure rather than consequences of
the scaffold notation.
\end{rem}

%------------------------------------------------------------------

\section{Multi-Fidelity MLMC}
\label{sec:ch12-multifidelity}

Scoring protocols provide another possible hierarchy. Their common target,
units, coupling, and bias must be specified before differences are meaningful.

\begin{defn}[Fidelity hierarchy]
\label{defn:fidelity-hierarchy}
Choose evaluable $F^{(0)},\ldots,F^{(K)}$ under one molecular sampling law $\nu$.
Each returns a quantity on a specified common scale; the intended computational
endpoint is $\E_\nu[F^{(K)}]$. Costs and correction variances are measured.
Calling docking, MM-GBSA, and FEP successive levels does not establish a common
physical observable or a monotone error ordering. Any mapping between their
outputs is part of the protocol and contributes uncertainty.
\end{defn}

For integrable outputs on this common law,
\[
\E_\nu[F^{(K)}]=\E_\nu[F^{(0)}]
 +\sum_{f=1}^K\E_\nu[F^{(f)}-F^{(f-1)}].
\]
Evaluating both members on the same molecule is a natural coupling. The
variance is
$\Var(F^{(f)})+\Var(F^{(f-1)})-2\Cov(F^{(f)},F^{(f-1)})$;
high correlation alone is insufficient if their scales differ greatly.

\subsection{Combining Two Hierarchies}
\begin{defn}[Mixed corrections and multi-fidelity cost]
\label{defn:ch12-mf-mlmc}
Let $Q_{\ell,f}$ be coupled, integrable approximations on a rectangular index
set, with $Q_{-1,f}=Q_{\ell,-1}=0$. Define
\[
Z_{\ell,f}=Q_{\ell,f}-Q_{\ell-1,f}-Q_{\ell,f-1}+Q_{\ell-1,f-1}.
\]
The double telescope gives
$\sum_{\ell=0}^L\sum_{f=0}^K\E[Z_{\ell,f}]=\E[Q_{L,K}]$.
Set $v_{\ell,f}=\Var(Z_{\ell,f})$ and let $c_{\ell,f}$ be the cost of the full
coupled correction. Then
\[
V_{\rm MF}(v;c)=\left(\sum_{\ell,f}\sqrt{v_{\ell,f}c_{\ell,f}}\right)^2.
\]
\end{defn}

\begin{prop}[Independent mixed-correction allocation]
\label{prop:mf-allocation}
For independent batches of the specified mixed corrections, positive
variances, and sampling-variance budget $q$, the continuous optimum is
\[
n_{\ell,f}=q^{-1}\sqrt{v_{\ell,f}/c_{\ell,f}}
                   \sum_{j,g}\sqrt{v_{j,g}c_{j,g}}.
\]
The proof is Theorem~\ref{thm:optimal-allocation} with a two-dimensional index.
Summing arbitrary single-axis corrections would double-count terms and would
not estimate the stated endpoint.
\end{prop}

\begin{exmp}[Randomized calibration versus shortlist promotion]
\label{ex:mf-screening}
A campaign can evaluate a cheap protocol broadly and draw randomized paired
higher-cost evaluations within each scaffold to estimate corrections. Every
part of the target population must retain sampling support; known unequal
inclusion probabilities require appropriate weighting. Sending only the top
10\% of cheap scores to the next protocol estimates a selected population's
correction, not the full scaffold correction. Such a shortlist is a screening
heuristic and needs separate recall validation.
\end{exmp}

%------------------------------------------------------------------

\section{Experiments}
\label{sec:ch12-experiments}

The following studies separate representation diagnostics from estimator
validation. No savings or decay pattern is presumed.

\begin{exper}[Parent-child projection diagnostics]
\label{exper:correlation-plot}
Construct genuinely nested partitions and compute exact finite-population
conditional means on a fully scored reference library. Verify the
correlation-spectrum identity where its variances are nonzero. Separately
fit the means on training data and assess prediction error on held-out data.
Report how the fitted spectrum differs from the reference spectrum, including
rare and empty classes.
\end{exper}

\begin{exper}[MLMC versus direct estimation]
\label{exper:mlmc-cost}
Specify the estimand, each correction sampler, and its cost. Use pilot data
to estimate $v_\ell$, freeze a conservative allocation, and evaluate bias,
variance, MSE, and total cost over independent repeated runs. Compare with
direct estimation at the same endpoint and target MSE. Include a control with
zero population projection increments but noisy raw score corrections to
check that the implementation does not confuse those two variances.
\end{exper}

\begin{exper}[Multi-fidelity validation]
\label{exper:multifidelity}
On a tractable common reference population, compare direct high-level mean
estimation, randomized coupled corrections, and deterministic shortlist
promotion. The first two target the same population mean; the third generally
does not. Measure correction variance and selection bias separately. For a
ranking objective, record ranking error as an additional outcome rather than
inferring it from mean-estimation MSE.
\end{exper}

\section*{Summary}
\label{sec:ch12-summary}
Exact projection increments describe information in a refinement tower.
Implemented correction observations determine Monte Carlo variance and cost.
For unbiased, independent correction batches the allocation is
$n_\ell\propto\sqrt{v_\ell/c_\ell}$. Unknown parents, shared observations,
truncation bias, and selective high-fidelity evaluation each require explicit
accounting. These conditions are what make the telescope an estimator.

%------------------------------------------------------------------

\section*{Exercises}

\begin{exer}[Derive optimal allocation]
\label{exer:ch12-allocation}
Prove Theorem~\ref{thm:optimal-allocation} via Lagrange multipliers:
\begin{enumerate}
    \item[(a)] Write the Lagrangian for minimizing cost subject to variance constraint.
    \item[(b)] Derive the first-order conditions $\partial \mathcal{L}/\partial n_\ell = 0$.
    \item[(c)] Solve for $n_\ell^*$ and verify the optimal cost formula.
    \item[(d)] Show that if all $c_\ell$ are equal, the allocation reduces to $n_\ell \propto \sqrt{v_\ell}$.
\end{enumerate}
\end{exer}

\begin{exer}[Prove correlation-spectrum identity]
\label{exer:ch12-correlation}
Prove Proposition~\ref{prop:correlation-spectrum} carefully:
\begin{enumerate}
    \item[(a)] Show that $\E[A_\ell \mid \mathcal{S}_{\ell-1}] = A_{\ell-1}$ (tower property).
    \item[(b)] Use this to prove $\text{Cov}(A_\ell, A_{\ell-1}) = \Var(A_{\ell-1})$.
    \item[(c)] Derive the formula $\rho_\ell^2 = \Var(A_{\ell-1})/\Var(A_\ell)$.
    \item[(d)] Verify that $\sum_{\ell=1}^L \|D_\ell\|^2 = \Var(A_L) - \Var(A_0)$.
\end{enumerate}
\end{exer}

\begin{exer}[Implement MLMC estimator]
\label{exer:ch12-implementation}
Write code to:
\begin{enumerate}
    \item[(a)] Given a specified coupled sampler, estimate its correction variances $v_\ell$; compute the population projection spectrum separately.
    \item[(b)] Compute optimal allocation $\{n_\ell^*\}$ for a given cost structure $\{c_\ell\}$ and target MSE $\epsilon^2$.
    \item[(c)] Implement the MLMC estimator for scaffold means.
    \item[(d)] Validate that achieved MSE matches the theoretical prediction.
\end{enumerate}
\end{exer}

\begin{exer}[Optimal allocation comparison]
\label{exer:ch12-comparison}
For actual correction variances $v_\ell=\sigma_0^2\,0.5^\ell$, costs $c_\ell=c_0\,2^\ell$, zero bias, and a direct estimator with variance $\sum_{\ell=0}^L v_\ell$ and cost $c_0 2^L$:
\begin{enumerate}
    \item[(a)] Compute the optimal allocation $\{n_\ell^*\}$ for levels $\ell=0,\ldots,5$.
    \item[(b)] Compute the MLMC cost for target MSE $\epsilon^2 = 0.01$.
    \item[(c)] Compute the naive (direct level-$L$ estimation) cost for the same MSE.
    \item[(d)] What is the MLMC speedup factor? How does it depend on $L$?
\end{enumerate}
\end{exer}

%------------------------------------------------------------------
\section{Bibliographic Notes}

The estimator and allocation argument follow standard MLMC
analysis~\cite{heinrich2001mlmc,giles2008mlmc,giles2015mlmc}. The essential
quantity is the variance of a coupled correction that can actually be sampled.
Cost comparison also requires control of approximation bias. Multi-fidelity
methods with fitted surrogates require additional care about training and
sampling laws~\cite{peherstorfer2018multifidelity}. The scaffold projection
identities provide diagnostics; their use does not by itself establish a new
Monte Carlo complexity theorem.

%------------------------------------------------------------------

\section{Connections}

\begin{itemize}
    \item \textbf{Chapter~\ref{ch:variance}:} Defines exact projection increments, distinguished here from sampled corrections.
    \item \textbf{Chapter~\ref{ch:energy}:} Supplies hypotheses for constructing and checking physically motivated approximations.
    \item \textbf{Chapter~\ref{ch:complexity}:} Separates estimation cost from output size and exact identification.
    \item \textbf{Chapter~\ref{ch:refinement}:} Optimal refinement selection determines which tower to use.
    \item \textbf{Chapter~\ref{ch:screening}:} Uses valid estimation subroutines while requiring separate molecule-level recovery certificates.
\end{itemize}
% END SOURCE: chapters/ch12_mlmc.tex


\chapter{Refinement Selection}
\label{ch:refinement}
% BEGIN SOURCE: chapters/ch13_refinement.tex
% ======================================================================
% Chapter 13 content file: chapters/ch13_refinement.tex
% Optimal Refinement Selection
% (Do NOT include \chapter{...} here; main.tex already does that.)
% ======================================================================

\begin{quote}
\textit{A refinement that does not reduce error is just a larger taxonomy.}
\end{quote}

\section{From Hierarchy to Optimization}
\label{sec:ch13-intro}

Refinement chooses which molecular distinctions a representation retains.
The objective depends on the response and the sampling population: a feature
that predicts one assay may add no useful information for another. This
chapter studies reductions in population squared error, then asks what can
be estimated from finite data and at what complexity cost.

A locally best split need not be part of a globally useful hierarchy.
Complementary features can have no individual predictive gain and a large
joint gain. We therefore separate oracle limits, feasible chemical features,
and the assumptions needed for any greedy guarantee.

We develop three complementary perspectives:
\begin{itemize}
    \item \textbf{Oracle refinements} (§\ref{sec:ch13-oracle}): What is theoretically achievable if we could cluster by score?
    \item \textbf{Actionable refinements} (§\ref{sec:ch13-actionable}): What is achievable with chemically meaningful splits?
    \item \textbf{Physics-guided refinements} (§\ref{sec:ch13-physics}): How can energy factorization and scaffold cumulants (Chapter~\ref{ch:energy}) guide split selection?
\end{itemize}

\begin{framed}
\noindent\textbf{For the medicinal chemist:} This chapter addresses a practical question: when you're building a SAR table, which structural distinctions actually matter for your target? Not all refinements are created equal---distinguishing ortho from para substitution matters enormously for some targets and not at all for others. The projection increment measures the population squared-error gain of a specified refinement. Estimating that gain and choosing among competing refinements requires held-out data and attention to joint effects.
\end{framed}

%------------------------------------------------------------------
\section{Refinement as an Optimization Problem}
\label{sec:ch13-optimization}

\subsection{The Refinement Order}

We begin by formalizing what it means to refine a scaffold abstraction.

\begin{defn}[Refinement order on $\sigma$-algebras]
\label{defn:refinement-order}
Let $\mathcal{S}$ and $\mathcal{S}'$ be $\sigma$-algebras on $\Mol$ generated by scaffold maps. We say $\mathcal{S}'$ \emph{refines} $\mathcal{S}$, written $\mathcal{S} \preceq \mathcal{S}'$, if $\mathcal{S} \subseteq \mathcal{S}'$.
\end{defn}

Equivalently, $\mathcal{S}'$ refines $\mathcal{S}$ if knowing which $\mathcal{S}'$-class a molecule\index{molecule} belongs to determines its $\mathcal{S}$-class. Every class of $\mathcal{S}$ is a union of classes of $\mathcal{S}'$.

\begin{exmp}[Bemis-Murcko\index{Bemis--Murcko scaffold} refinements]
\label{ex:bm-refinements}
Consider the hierarchy:
\begin{itemize}
    \item $\mathcal{S}_0$: Ring system\index{ring system} count (monocyclic, bicyclic, tricyclic, ...)
    \item $\mathcal{S}_1$: Ring connectivity (fused vs.\ spiro vs.\ bridged)
    \item $\mathcal{S}_2$: Ring sizes (6-6, 6-5, 5-5, ...)
    \item $\mathcal{S}_3$: Heteroatom positions (quinoline vs.\ isoquinoline vs.\ quinazoline)
    \item $\mathcal{S}_4$: Full Bemis-Murcko scaffold
\end{itemize}
Define each level to retain all preceding labels; the final label contains
the full scaffold together with any earlier descriptors it does not determine.
Then $\mathcal S_0\subseteq\cdots\subseteq\mathcal S_4$ by construction.
Names of chemical descriptors alone do not prove nesting.
\end{exmp}

\begin{prop}[Refinement is a partial order\index{partial order}]
\label{prop:refinement-partial-order}
The refinement relation $\preceq$ is a partial order on $\sigma$-algebras: it is reflexive, antisymmetric, and transitive.
\end{prop}

\begin{proof}
Reflexivity: $\mathcal{S} \subseteq \mathcal{S}$. Antisymmetry: $\mathcal{S} \subseteq \mathcal{S}'$ and $\mathcal{S}' \subseteq \mathcal{S}$ implies $\mathcal{S} = \mathcal{S}'$. Transitivity: $\mathcal{S} \subseteq \mathcal{S}'$ and $\mathcal{S}' \subseteq \mathcal{S}''$ implies $\mathcal{S} \subseteq \mathcal{S}''$.
\end{proof}

\subsection{The Refinement Objective}

The central quantity we seek to minimize is the irreducible error.

\begin{defn}[Refinement objective]
\label{defn:refinement-objective}
For a scoring function $F \in L^2(\mu)$ and $\sigma$-algebra\index{$\sigma$-algebra@$\sigma$-algebra} $\mathcal{S}$, define the \emph{refinement objective}:
\[
J(\mathcal{S}) := \mathrm{Err}(\mathcal{S}) = \E[\Var(F \mid \mathcal{S})]
\]
This is the expected within-class variance---the irreducible prediction error from Chapter~\ref{ch:variance}.
\end{defn}

From Theorem~\ref{thm:monotonicity} (monotonicity under refinement), we know:

\begin{prop}[Refinement reduces error]
\label{prop:refinement-reduces-error}
If $\mathcal{S} \preceq \mathcal{S}'$, then $J(\mathcal{S}') \leq J(\mathcal{S})$.
\end{prop}

Finer refinements always reduce (or maintain) the irreducible error. The challenge is that finer refinements also increase complexity.

\subsection{Complexity Constraints}

Refinement cannot be pursued without bound. Finer abstractions require more samples to estimate reliably and may lose interpretability.

\begin{defn}[Complexity measures]
\label{defn:complexity-measures}
Common complexity measures for a scaffold $\sigma$-algebra $\mathcal{S}$ include:
\begin{enumerate}
    \item \textbf{Class count:} $K(\mathcal{S}) = |\mathrm{Im}(\mathrm{Scaf})|$, the number of distinct scaffold classes
    \item \textbf{Entropy:} $H(\mathcal{S}) = -\sum_s \mu(\mathrm{Fib}(s)) \log \mu(\mathrm{Fib}(s))$
    \item \textbf{Sample complexity:} The number of samples required to estimate class statistics to precision $\epsilon$ with confidence $1-\delta$
\end{enumerate}
\end{defn}

\begin{rem}[Sample complexity scaling]
\label{rem:sample-complexity}
For independent observations with a known sub-Gaussian variance proxy $v_s$,
$n_s\geq2v_s\epsilon^{-2}\log(2K/\delta)$ is sufficient for simultaneous
fixed-sample mean error at most $\epsilon$ across $K$ fixed classes. Adaptive
partition selection, repeated looks, and unknown variance need additional
control. Finer classes also create rare or empty cells; population projection
error does not account for the error of learning their means.

\end{rem}

\subsection{The Constrained Optimization Problem}

Combining the objective and constraints yields the central optimization problem of this chapter.

\begin{defn}[Optimal refinement problem]
\label{defn:optimal-refinement}
Given a base $\sigma$-algebra $\mathcal{S}_0$ and a complexity budget $B$, the \emph{optimal refinement problem} is:
\[
\min_{\mathcal{S} : \mathcal{S}_0 \preceq \mathcal{S}} J(\mathcal{S}) \quad \text{subject to} \quad K(\mathcal{S}) \leq B
\]
\end{defn}

\begin{rem}[Why this is hard]
\label{rem:why-hard}
The optimal refinement problem is computationally challenging for several reasons:
\begin{enumerate}
    \item The space of refinements is enormous (all partitions refining $\mathcal{S}_0$)
    \item The objective $J(\mathcal{S})$ depends on the unknown scoring function\index{scoring function} $F$
    \item We want refinements that are chemically interpretable, not arbitrary partitions
\end{enumerate}
The remainder of this chapter develops tractable approaches to this problem.
\end{rem}

%------------------------------------------------------------------
\section{Oracle Refinements and Quantization Limits}
\label{sec:ch13-oracle}

Before considering chemically interpretable refinements, we establish theoretical limits: what is the best possible error achievable by \emph{any} partition into $B$ classes?

\subsection{The Oracle Partition}

\begin{defn}[Oracle $B$-partition]
\label{defn:oracle-partition}
An \emph{oracle $B$-partition} of $\Mol$ is a partition into at most $B$ classes that minimizes $J(\mathcal{S})$ among all such partitions, without any constraint on chemical interpretability.
\end{defn}

This is the ``cheating'' baseline: we are allowed to look at the scores $F(m)$ when choosing how to partition.

\begin{thm}[Oracle optimality via quantization]
\label{thm:oracle-quantization}
Among all partitions of $\Mol$ into at most $B$ classes, the partition minimizing $J(\mathcal{S}) = \E[\Var(F \mid \mathcal{S})]$ solves a \emph{quantization problem}: it minimizes
\[
\min_{\text{partitions } \pi} \sum_{C \in \pi} \mu(C) \cdot \Var(F \mid m \in C)
\]
Equivalently, if we let $\hat{F}_\pi(m) = \E[F \mid m \in C]$ for the class $C \ni m$, the oracle partition minimizes
\[
\E\left[(F - \hat{F}_\pi)^2\right]
\]
over all $B$-class partitions.
\end{thm}

\begin{proof}
By the law of total variance, $\Var(F) = \E[\Var(F \mid \mathcal{S})] + \Var(\E[F \mid \mathcal{S}])$. Since $\Var(F)$ is fixed, minimizing $\E[\Var(F \mid \mathcal{S})]$ is equivalent to maximizing $\Var(\E[F \mid \mathcal{S}])$---the between-class variance.

For any partition, retain its class centroids and reassign each score to a
nearest centroid. This cannot increase pointwise squared error. Replacing the
centroids by the resulting class means can only improve it further. Thus an
optimal partition can be chosen as a scalar quantizer of $F$; the population
projection objective is the same distortion criterion.
\end{proof}

\begin{cor}[Connection to $k$-means]
\label{cor:kmeans}
When $F$ is one-dimensional, the oracle $B$-partition corresponds to optimal scalar quantization. Lloyd-Max iterations seek a stationary quantizer, without guaranteeing the global optimum:
\begin{enumerate}
    \item Given centroids $\{c_1, \ldots, c_B\}$, assign each molecule to the nearest centroid in $F$-value
    \item Update centroids to class means
\end{enumerate}
This is $k$-means clustering on the score $F$.
\end{cor}

\begin{rem}[The oracle is unimplementable]
\label{rem:oracle-unimplementable}
The oracle partition requires knowing $F(m)$ for all molecules---precisely the information we are trying to predict. The oracle serves as a \textbf{lower bound} on achievable error: no scaffold-based method can do better than the oracle at the same complexity level.
\end{rem}

\subsection{The Diminishing Returns Curve}

As $B$ increases, the optimal achievable error decreases. The rate of decrease characterizes the intrinsic ``refinement difficulty'' of the target.

\begin{defn}[Oracle error curve]
\label{defn:oracle-error-curve}
The \emph{oracle error curve} is the function
\[
\mathrm{Err}_{\text{oracle}}(B) := \min_{\text{$B$-class partitions}} J(\mathcal{S})
\]
\end{defn}

\begin{prop}[Properties of the oracle error curve]
\label{prop:oracle-curve-properties}
The oracle error curve satisfies:
\begin{enumerate}
    \item $\mathrm{Err}_{\text{oracle}}(1) = \Var(F)$ (single class = no refinement)
    \item On a finite library, $\mathrm{Err}_{\text{oracle}}(|\Mol|)=0$ (one class per molecule)
    \item $\mathrm{Err}_{\text{oracle}}(B)$ is non-increasing in $B$
    \item For $F\in L^2$, the error is nonnegative and at most $\Var(F)$
\end{enumerate}
\end{prop}

\begin{proof}
(1) and the finite-library statement (2) are immediate. For (3), the feasible
set of partitions into at most $B$ classes grows with $B$. Nonnegativity and
comparison with the single-class partition give (4). These facts do not
establish a greedy algorithm, a convex error curve, or diminishing marginal
gains along a chosen hierarchy.

\end{proof}

\begin{defn}[Intrinsic refinement rate]
\label{defn:intrinsic-rate}
The \emph{intrinsic refinement rate} is characterized by how fast $\mathrm{Err}_{\text{oracle}}(B)$ decays. Define
\[
\rho^*(B) := \frac{\mathrm{Err}_{\text{oracle}}(B)}{\Var(F)}
\]
the fraction of variance remaining after optimal $B$-class quantization, defined when $\Var(F)>0$.
\end{defn}

\begin{exmp}[Fast vs.\ slow decay]
\label{ex:fast-slow-decay}
\begin{itemize}
    \item \textbf{Fast decay:} If $F$ is well-approximated by a few discrete levels (e.g., ``active'', ``moderate'', ``inactive''), then $\rho^*(B)$ drops quickly for small $B$.
    \item \textbf{Slow decay:} If $F$ is smoothly distributed with no natural clusters, $\rho^*(B)$ decreases slowly, requiring many classes to capture variance.
\end{itemize}
Fast decay means a few score-dependent classes can represent the response. It does not show that an available chemical feature identifies those classes without observing the scores.
\end{exmp}

\subsection{The Oracle Gap}

The gap between oracle and achievable error quantifies the cost of chemical interpretability.

\begin{defn}[Oracle gap]
\label{defn:oracle-gap}
For a restricted family $\mathfrak{R}$ of chemically interpretable refinements, the \emph{oracle gap} at complexity $B$ is:
\[
\mathrm{Gap}(B) := \min_{\mathcal{S} \in \mathfrak{R}, K(\mathcal{S}) \leq B} J(\mathcal{S}) - \mathrm{Err}_{\text{oracle}}(B)
\]
\end{defn}

\begin{prop}[Gap bounds actionable methods]
\label{prop:gap-bounds}
For any feasible $\mathcal S\in\mathfrak R$,
\[
\mathrm{Err}_{\rm oracle}(B)+\mathrm{Gap}(B)
 =\min_{\mathcal T\in\mathfrak R,\ K(\mathcal T)\leq B}J(\mathcal T)
 \leq J(\mathcal S)\leq\Var(F).
\]
The first inequality is attained by an optimal feasible refinement when a
minimizer exists. For a heuristic, its excess above the restricted optimum is
an additional optimization error, not part of the oracle gap.

\end{prop}

\begin{framed}
\noindent\textbf{Interpretation:} The oracle gap measures how much predictive power is sacrificed by requiring chemically meaningful distinctions. A small gap indicates that chemistry-based refinements capture most of the score-relevant structure; a large gap indicates that score variation does not align with chemical intuition.
\end{framed}

%------------------------------------------------------------------
\section{Actionable Refinements and Greedy Splits}
\label{sec:ch13-actionable}

We now restrict attention to refinements that are chemically interpretable and algorithmically tractable.

\subsection{The Actionable Refinement Family}

\begin{defn}[Actionable refinement family]
\label{defn:actionable-family}
An \emph{actionable refinement family} $\mathfrak{R}$ is a collection of scaffold $\sigma$-algebras satisfying:
\begin{enumerate}
    \item \textbf{Chemical interpretability:} Each $\mathcal{S} \in \mathfrak{R}$ corresponds to chemically meaningful structural distinctions
    \item \textbf{Closure under refinement:} If $\mathcal{S} \in \mathfrak{R}$ and $\mathcal{S}'$ refines $\mathcal{S}$ via an interpretable split, then $\mathcal{S}' \in \mathfrak{R}$
    \item \textbf{Computability:} Class membership can be determined algorithmically from molecular structure
\end{enumerate}
\end{defn}

\begin{exmp}[Common actionable refinements]
\label{ex:actionable-refinements}
\begin{itemize}
    \item \textbf{Ring-based:} Fused vs.\ spiro; ring size; aromaticity
    \item \textbf{Heteroatom-based:} N vs.\ O vs.\ S in ring; position of heteroatoms
    \item \textbf{Substituent-based:} Ortho vs.\ meta vs.\ para; alkyl vs.\ aryl vs.\ halogen
    \item \textbf{Attachment-point-based:} Number of attachment points; attachment positions
    \item \textbf{Stereochemistry:} R vs.\ S; cis vs.\ trans
\end{itemize}
\end{exmp}

\subsection{Variance Reduction as the Split Criterion}

Given a current refinement $\mathcal{S}$, how do we choose the next split?

\begin{defn}[Split]
\label{defn:split}
A \emph{split} of $\mathcal{S}$ is a refinement $\mathcal{S}'$ obtained by partitioning exactly one class of $\mathcal{S}$ into two or more subclasses, leaving other classes unchanged.
\end{defn}

\begin{prop}[Variance reduction identity]
\label{prop:variance-reduction}
For a split $\mathcal{S} \to \mathcal{S}'$:
\[
J(\mathcal{S}) - J(\mathcal{S}') = \|\E[F \mid \mathcal{S}'] - \E[F \mid \mathcal{S}]\|_{L^2}^2
\]
The variance reduction equals the squared $L^2$ norm of the new martingale\index{martingale} increment.
\end{prop}

\begin{proof}
By the law of total variance applied to the tower $\mathcal{S} \subseteq \mathcal{S}'$:
\begin{align}
\Var(F) &= \E[\Var(F \mid \mathcal{S}')] + \Var(\E[F \mid \mathcal{S}']) \\
&= \E[\Var(F \mid \mathcal{S}')] + \E[\Var(\E[F \mid \mathcal{S}'] \mid \mathcal{S})] + \Var(\E[F \mid \mathcal{S}])
\end{align}
Subtracting the analogous expansion for $\mathcal{S}$:
\[
J(\mathcal{S}) - J(\mathcal{S}') = \E[\Var(\E[F \mid \mathcal{S}'] \mid \mathcal{S})] = \|\E[F \mid \mathcal{S}'] - \E[F \mid \mathcal{S}]\|^2
\]
\end{proof}

\begin{rem}[Connection to decision trees]
\label{rem:decision-trees}
Proposition~\ref{prop:variance-reduction} is the \emph{impurity reduction} criterion used in regression tree learning (CART, etc.). Greedy variance reduction builds the scaffold hierarchy exactly as a regression tree would partition feature space.
\end{rem}

\subsection{The Greedy Algorithm}

\begin{algorithm}
\caption{Greedy refinement selection}
\label{alg:greedy-refinement}
\mbox{}\\
\textbf{Input:} Base refinement $\mathcal{S}_0$, complexity budget $B$, actionable split family $\mathfrak{R}$

\textbf{Output:} Refinement tower $\mathcal{S}_0 \preceq \mathcal{S}_1 \preceq \cdots \preceq \mathcal{S}_L$

\begin{enumerate}
    \item Initialize $\ell \leftarrow 0$
    \item While $K(\mathcal{S}_\ell) < B$:
    \begin{enumerate}
        \item For each candidate split $\mathcal{S}_\ell\to\mathcal{S}'$ in $\mathfrak R$ with $K(\mathcal S')\leq B$:
        \begin{itemize}
            \item Estimate variance reduction $\Delta J = J(\mathcal{S}_\ell) - J(\mathcal{S}')$
        \end{itemize}
        \item Select split $\mathcal{S}^* = \arg\max_{\mathcal{S}'} \Delta J$
        \item If no feasible split remains, or estimated gain is below $\epsilon_{\min}$: \textbf{break}. This is a local stopping rule, not a certificate that joint refinements have small gain.
        \item Set $\mathcal{S}_{\ell+1} \leftarrow \mathcal{S}^*$, $\ell \leftarrow \ell + 1$
    \end{enumerate}
    \item Return $\mathcal{S}_0, \mathcal{S}_1, \ldots, \mathcal{S}_\ell$
\end{enumerate}
\end{algorithm}

\begin{exmp}[Complementary features defeat one-step greedy search]
\label{ex:refinement-xor}
Let $X,Y,Z$ be independent fair bits and let
$F=\mathbf1\{X\ne Y\}+\varepsilon Z$, where $0<\varepsilon<1$.
With a constant base label, observing $X$ alone or $Y$ alone explains no
variance, whereas observing both explains $1/4$. Thus explained variance is
not submodular: the marginal gain of $Y$ rises from zero to $1/4$ after $X$
is observed. The score is bounded and the feature family is finite.

With a budget of two features, greedy first selects $Z$, whose gain is
$\varepsilon^2/4>0$. Adding either remaining feature still gives zero gain,
leaving error $1/4$. Selecting $X,Y$ instead leaves error $\varepsilon^2/4$.
The ratio of greedy error to optimal error is $\varepsilon^{-2}$ and can be
arbitrarily large. More data cannot repair this population-level failure.
\end{exmp}

\begin{thm}[Greedy guarantee under explicit diminishing returns]
\label{thm:greedy-guarantee}
Let $\mathcal E$ be a fixed finite set of features and
$\mathcal S_A=\sigma(\mathcal S_0,X_e:e\in A)$. Define
$g(A)=J(\mathcal S_0)-J(\mathcal S_A)$. Suppose $g$ is submodular:
\[
g(A\cup\{e\})-g(A)\geq g(B\cup\{e\})-g(B)
\quad(A\subseteq B,\ e\notin B).
\]
Under a cardinality budget $|A|\leq b$, with $1\leq b\leq|\mathcal E|$,
exact-gain greedy addition for $b$ steps returns $A_g$ satisfying
\[
g(A_g)\geq\bigl[1-(1-1/b)^b\bigr]g(A^*)\geq(1-e^{-1})g(A^*),
\]
where $A^*$ maximizes $g$ under that budget. Equivalently,
\[
J(\mathcal S_{A_g})\leq J(\mathcal S_{A^*})
       +e^{-1}\bigl[J(\mathcal S_0)-J(\mathcal S_{A^*})\bigr].
\]
\end{thm}
\begin{proof}
Projection monotonicity gives $g\geq0$ and monotonicity. Submodularity and
monotonicity imply
\[
g(A^*)-g(A_i)\leq\sum_{e\in A^*\setminus A_i}
                [g(A_i\cup\{e\})-g(A_i)]
\leq b[g(A_{i+1})-g(A_i)].
\]
Hence the remaining gap contracts by at most $1-1/b$ at each step. Iterate
from $g(\varnothing)=0$ and use $(1-1/b)^b\leq e^{-1}$.
\end{proof}

Submodularity is an additional hypothesis, disproved in general by
Example~\ref{ex:refinement-xor}. A cardinality budget on globally available
features also differs from Algorithm~\ref{alg:greedy-refinement}'s budget on
leaf classes and state-dependent splits. The theorem does not supply a
constant-factor residual-error guarantee for that algorithm. Joint-feature
candidates, bounded lookahead, and held-out comparison are appropriate when
response curvature suggests complementary effects
(Chapter~\ref{ch:response-certificates}).

\begin{rem}[Practical considerations]
\label{rem:practical-greedy}
In practice, the greedy algorithm requires:
\begin{enumerate}
    \item Sufficient samples to estimate variance reduction reliably
    \item A tractable enumeration of candidate splits
    \item A stopping criterion (minimum variance reduction, maximum depth, or cross-validation)
\end{enumerate}
The sample requirement can be substantial for rare scaffolds. Multilevel estimation helps only when the implemented correction costs and variances satisfy the conditions of Chapter~\ref{ch:mlmc}.
\end{rem}

%------------------------------------------------------------------
\section{Physics-Guided Refinement from Factorization}
\label{sec:ch13-physics}

The greedy algorithm treats refinement selection as a purely statistical problem. This section shows how the energy factorization and incidence-algebra cumulants from Chapter~\ref{ch:energy} provide \emph{physical guidance} for choosing splits.

\subsection{The Key Insight}

Recall from Chapter~\ref{ch:energy} that scoring functions decompose as:
\[
F = F_{\text{core}} + F_{\text{decor}} + F_{\text{boundary\index{boundary}}}
\]
If $F_{\rm core}$ is measurable with respect to $\mathcal S$, residual
variation arises from $F_{\rm decor}+F_{\rm boundary}$. For relaxed poses,
core geometry may vary with decoration; that case requires a further core
residual term. The boundary term is an aggregate quantity: after choosing an attachment-site poset, M\"obius inversion further decomposes it into scaffold cumulants $\kappa_s(I)$ indexed by sets of sites. Refinement should therefore ask which structural distinctions reveal large unresolved cumulants.

\begin{prop}[Variance decomposition by energy components]
\label{prop:energy-variance}
If $F_{\rm core}$ is constant within the conditioned class and all terms are square integrable, then
\[
\Var(F \mid \mathrm{Scaf} = s) = \Var(F_{\text{decor}} \mid s) + \Var(F_{\text{boundary}} \mid s) + 2\Cov(F_{\text{decor}}, F_{\text{boundary}} \mid s)
\]
\end{prop}

Boundary terms and response contrasts can propose features to test.
The total response, including covariance and relaxation, determines whether
those features reduce prediction error.

\subsection{Attachment-Local Features}

\begin{defn}[Attachment-local feature map]
\label{defn:attachment-local}
For a scaffold $s$ with attachment sites $a_1, \ldots, a_k$, define local decoration\index{decoration} descriptors:
\[
\phi_i(d) = (\text{substituent class}, \text{size}, \text{charge}, \text{HBD/HBA}, \text{lipophilicity})
\]
capturing the local chemistry at each attachment point\index{attachment point}.
\end{defn}

\begin{exmp}[Substituent classes]
\label{ex:substituent-classes}
A simple classification:
\begin{itemize}
    \item H (hydrogen)
    \item Alkyl (methyl, ethyl, isopropyl, ...)
    \item Aryl (phenyl, naphthyl, ...)
    \item Halogen (F, Cl, Br, I)
    \item Polar (OH, NH$_2$, COOH, ...)
    \item Other
\end{itemize}
\end{exmp}

\subsection{Boundary-Aware Split Score}

\begin{prop}[Component-wise refinement gain]
\label{prop:boundary-aware}
Assume $F=F_{\rm core}+U+V$ with $F_{\rm core}$ measurable on $\mathcal S$,
$U=F_{\rm decor}$ and $V=F_{\rm boundary}$. For $\mathcal S\subseteq\mathcal S'$
write $g_U=\E[U\mid\mathcal S']-\E[U\mid\mathcal S]$ and similarly $g_V$.
Then
\[
J(\mathcal S)-J(\mathcal S')=
 \|g_U\|_2^2+\|g_V\|_2^2+2\langle g_U,g_V\rangle.
\]
In particular, if $g_U=0$, the reduction equals $\|g_V\|_2^2$.
\end{prop}
\begin{proof}
Linearity of conditional expectation makes the total new projection increment
$g_U+g_V$. Expand its squared norm using
Proposition~\ref{prop:variance-reduction}.
\end{proof}

A boundary-only explained-variance score therefore equals the actual split
gain only under conditions such as $g_U=0$. If $g_U=-g_V$, a large apparent
boundary gain is canceled entirely. Estimates of boundary variance or its
correlation with a feature cannot establish gain in the total response.

\begin{defn}[Cumulant-aware split score]
\label{defn:cumulant-split-score}
For a scaffold $s$ with attachment-site cumulants $\kappa_s(I)$, define the squared coefficient diagnostic
\[
U_{\geq 2}(s) := \sum_{|I|\geq 2} \kappa_s(I)^2.
\]
For a candidate refinement that distinguishes a block of sites or features $B$, define its cumulant score
\[
\Gamma(B;s) := \sum_{\substack{|I|\geq 2\\ I\cap B\neq\emptyset}} w_I \kappa_s(I)^2,
\]
where $w_I\geq0$ can downweight poorly estimated, high-order, or spatially
diffuse coefficients. In an anchored M\"obius basis these squared coefficients
are not generally orthogonal variance contributions. Their values depend on
the anchor and response scale. Measurement uncertainty must be reported; large
squared estimates can reflect noise rather than interaction.
\end{defn}

\begin{rem}[Why cumulants improve the split criterion]
\label{rem:cumulant-split}
The boundary ratio $\eta(s)$ is a model-specific diagnostic; it does not distinguish a one-site effect from a two-site cooperative effect. Correlations with $\phi_i$ isolate the one-site case; the score $\Gamma(B;s)$ isolates coupled-site cases. A large two- or three-site cumulant suggests refining a coupled block of sites together. This is the algebraic version of a medicinal chemist's observation that some SAR tables cannot be explained one column at a time.
\end{rem}

\begin{framed}
\noindent\textbf{A testable proposal:} A mapped steric interaction at a binding
site can motivate a joint feature split. Its value is the held-out reduction
in total response error at an acceptable measurement cost, not the magnitude
of a component term alone.

\end{framed}

\subsection{The Physics-Guided Algorithm}

\begin{algorithm}
\caption{Physics-guided refinement selection}
\label{alg:physics-guided}
\mbox{}\\
\textbf{Input:} Scaffold set $S$, scoring function $F$, energy decomposition oracle, optional SAR-grid cumulant estimates

\textbf{Output:} Candidate refinement tower, assessed on held-out total-response error

\begin{enumerate}
    \item \textbf{Compute boundary variance:} For each scaffold $s$, estimate
    \[
    \eta(s) := \frac{\Var(F_{\text{boundary}} \mid s)}{\Var(F \mid s)}
    \]
    when $\Var(F\mid s)>0$. This ratio can exceed one because component covariances may cancel; it is not a variance fraction.

    \item \textbf{Estimate cumulant energy when possible:} On scaffolds with local SAR grids, compute
    \[
    \eta_\kappa(s) := \frac{\sum_{|I|\geq 2}\kappa_s(I)^2}{\sum_I \kappa_s(I)^2}.
    \]
    This is a coefficient diagnostic when its denominator is positive, not a
variance decomposition. Report anchors and uncertainty before interpreting it
as evidence of non-additivity.

    \item \textbf{Prioritize high-$\eta$ or high-$\eta_\kappa$ scaffolds:} Use these diagnostics to propose candidate splits, then assess gain on the total response. They do not order achievable gains by themselves.

    \item \textbf{Identify predictive features:} For scaffolds with high $\eta(s)$, compute correlations
    \[
    \rho_i(s) := \mathrm{Corr}(\phi_i(d), F_{\text{boundary}}(s, d))
    \]
    for each attachment-local feature $\phi_i$. When cumulants are available, also compute $\Gamma(B;s)$ for single-site and coupled-site blocks $B$.

    \item \textbf{Refine by top features or blocks:} Split scaffolds using features $\phi_i$ with largest $|\rho_i(s)|$, or coupled blocks $B$ with largest $\Gamma(B;s)$.

    \item \textbf{Validate and iterate:} Retain splits supported by held-out total-response error, respecting class-count and measurement budgets; consider joint blocks when one-feature gains vanish.
\end{enumerate}
\end{algorithm}

\begin{rem}[Computational requirements]
\label{rem:computational-physics}
The physics-guided algorithm requires:
\begin{enumerate}
    \item Energy decomposition for sampled molecules (Chapter~\ref{ch:energy} methods)
    \item Sufficient sampling within each scaffold to estimate correlations
    \item Feature engineering for attachment-local descriptors
\end{enumerate}
The overhead is justified when the resulting refinements significantly outperform purely statistical approaches. Cumulant estimation should be used cautiously on sparse tables; low-order or regularized cumulants are preferable unless the design is close to factorial.
\end{rem}

\subsection{Connection to SAR Tables}

\begin{rem}[SAR tables as refinements]
\label{rem:SAR-tables}
A structure-activity relationship (SAR) table is implicitly a refinement: rows correspond to scaffolds, columns to substitution patterns, and entries to activities. The physics-guided algorithm provides a \emph{principled reason} why certain SAR tables are effective: they are partial observations of the scaffold cumulant expansion. Mapped, protocol-compatible measurements can estimate selected response contrasts. Which orders dominate is an empirical question, and missing cells can prevent identification.
\end{rem}

\begin{exmp}[Free-Wilson analysis]
\label{ex:free-wilson}
The classical Free-Wilson model assumes additive contributions:
\[
\text{Activity} = \mu + \sum_i \alpha_i \cdot \mathbf{1}[\text{substituent $i$ present}]
\]
By Theorem~\ref{thm:sar-cumulant-truncation}, this is exact when all higher-order scaffold cumulants vanish, $\kappa_s(I)=0$ for $|I|\geq 2$. Thus additivity can hold even with nonzero one-site boundary effects; what breaks Free-Wilson is not boundary energy itself but irreducible multi-site coupling. The physics-guided refinement extends Free-Wilson by identifying when and where those cumulants matter.
\end{exmp}

%------------------------------------------------------------------
\section{Intrinsic Depth as a Stopping Rule}
\label{sec:ch13-stopping}

How do we know when to stop refining?

\subsection{The Diminishing Returns Curve}

\begin{defn}[Empirical error curve]
\label{defn:empirical-error-curve}
For a refinement tower $\mathcal{S}_0 \preceq \mathcal{S}_1 \preceq \cdots \preceq \mathcal{S}_L$, the \emph{population error curve} is:
\[
\mathrm{Err}(\ell) := J(\mathcal{S}_\ell) = \E[\Var(F \mid \mathcal{S}_\ell)]
\]
as a function of level $\ell$.
\end{defn}

From Chapter~\ref{ch:variance}, we know $\mathrm{Err}(\ell)$ is non-increasing. A good refinement tower makes $\mathrm{Err}(\ell)$ decrease rapidly.

\begin{defn}[Intrinsic scaffold depth (refined)]
\label{defn:intrinsic-depth-refined}
For tolerance $\epsilon > 0$, the \emph{intrinsic scaffold depth} is:
\[
L^*(\epsilon) := \min\left\{\ell : \frac{\mathrm{Err}(\ell)}{\Var(F)} \leq \epsilon\right\}
\]
The smallest such level, with $\Var(F)>0$; set $L^*(\epsilon)=\infty$ if none exists. The depth is relative to the chosen tower and sampling law.
\end{defn}

\begin{prop}[Stopping rule]
\label{prop:stopping-rule}
If $L^*(\epsilon)<\infty$, then for every later refinement
$\mathcal T\supseteq\mathcal S_{L^*}$,
\[
0\leq J(\mathcal S_{L^*})-J(\mathcal T)
   \leq J(\mathcal S_{L^*})\leq\epsilon\Var(F).
\]
This bounds the total remaining population gain. It establishes neither
diminishing individual marginal gains nor a cost-optimal stopping point.

\end{prop}

\subsection{Adaptive Stopping via Cross-Validation}

In practice, we estimate $\mathrm{Err}(\ell)$ from data and must guard against overfitting.

\begin{algorithm}
\caption{Cross-validated stopping}
\label{alg:cv-stopping}
\mbox{}\\
\begin{enumerate}
    \item Split data by the intended unit of transfer into training, validation, and final test sets; keep shared-parent observations and assay replicates in the appropriate groups
    \item Build refinement tower on training data using greedy or physics-guided algorithm
    \item For each level $\ell$, estimate validation error $\hat{J}_{\text{val}}(\ell)$
    \item Select a level using a prespecified validation rule; report final performance once on the untouched test set. Repeated validation-based selection needs its own error control
\end{enumerate}
\end{algorithm}

%------------------------------------------------------------------
\section{Experiments}
\label{sec:ch13-experiments}

\begin{exper}[Oracle vs.\ actionable gap]
\label{exp:oracle-gap}
\mbox{}\\
\textbf{Setup:}
\begin{itemize}
    \item Dataset: Docking scores for 100K molecules against 5 targets
    \item Oracle: exact scalar quantization on the finite score sample for $B=10,50,100,500$ classes; ordinary $k$-means supplies a local solution, not a certified oracle optimum
    \item Actionable: Bemis-Murcko hierarchy, substituent-based refinements
\end{itemize}

\textbf{Measure:}
\begin{itemize}
    \item $\mathrm{Err}_{\text{oracle}}(B)$ vs.\ $\mathrm{Err}_{\text{actionable}}(B)$ curves
    \item Gap = $\mathrm{Err}_{\text{actionable}} - \mathrm{Err}_{\text{oracle}}$
\end{itemize}

\textbf{Evaluation:} Report the restricted optimum or heuristic excess separately from the oracle distortion. Do not prescribe a numerical gap in advance.
\end{exper}

\begin{exper}[Intrinsic depth across targets]
\label{exp:intrinsic-depth}
\mbox{}\\
\textbf{Setup:}
\begin{itemize}
    \item Multiple targets with diverse binding sites
    \item Build candidate towers with greedy and joint-feature alternatives
    \item Plot $\mathrm{Err}(\ell)/\Var(F)$ vs.\ $\ell$
\end{itemize}

\textbf{Measure:}
\begin{itemize}
    \item Intrinsic depth $L^*(0.1)$ for each target
    \item Held-out error, class occupancy, and marginal gain per measurement cost
\end{itemize}

\textbf{Evaluation:} Estimate tower-specific depth with uncertainty; no typical range or curve shape is assumed.
\end{exper}

\begin{exper}[Physics-guided vs.\ naive refinement]
\label{exp:physics-guided}
\mbox{}\\
\textbf{Setup:}
\begin{itemize}
    \item Single target with energy decomposition available
    \item Compare three strategies:
    \begin{enumerate}
        \item Random splitting
        \item Chemistry-based splitting (Bemis-Murcko hierarchy)
        \item Physics-guided splitting (boundary-variance-aware)
    \end{enumerate}
\end{itemize}

\textbf{Measure:}
\begin{itemize}
    \item Variance reduction per split
    \item Total error at fixed complexity budget
\end{itemize}

\textbf{Evaluation:} Compare strategies at the same total feature, fitting, and oracle costs, including targets where the physical proposal performs worse.
\end{exper}

%------------------------------------------------------------------
\section*{Summary}

This chapter established that refinement selection is fundamentally a statistical optimization problem:
\begin{enumerate}
    \item \textbf{Oracle refinements} (§\ref{sec:ch13-oracle}) set theoretical limits via quantization theory
    \item \textbf{Actionable refinements} (§\ref{sec:ch13-actionable}) restrict to chemically interpretable splits, with a greedy explained-variance guarantee only under explicitly verified submodularity and a feature-cardinality budget
    \item \textbf{Physics-guided refinements} (§\ref{sec:ch13-physics}) use energy components and response contrasts to propose splits, whose gain must be checked on the total response
\end{enumerate}

The key formula connecting refinement to the variance spectrum:
\[
J(\mathcal{S}) - J(\mathcal{S}') = \|D_{\text{new}}\|^2
\]
where $D_{\text{new}}$ is the martingale increment from the new split.

Chapter~\ref{ch:screening}\index{screening} uses these refined towers in end-to-end screening algorithms, where the quality of refinement directly impacts budget efficiency.

%------------------------------------------------------------------
\section*{Exercises}

\begin{exer}[Variance reduction identity]
\label{exer:ch13-variance-reduction}
Prove Proposition~\ref{prop:variance-reduction} in full detail, starting from the law of total variance.
\end{exer}

\begin{exer}[Oracle quantization]
\label{exer:ch13-quantization}
For a uniformly distributed score $F \sim \mathrm{Uniform}[0, 1]$:
\begin{enumerate}
    \item Compute $\mathrm{Err}_{\text{oracle}}(B)$ explicitly for equal-width quantization
    \item Show that equal-width is optimal for uniform distributions
    \item Derive the decay rate: $\mathrm{Err}_{\text{oracle}}(B) = O(1/B^2)$
\end{enumerate}
\end{exer}

\begin{exer}[Greedy vs.\ optimal]
\label{exer:ch13-greedy}
Construct an example where the greedy algorithm produces a refinement with $J(\mathcal{S}_{\text{greedy}}) > 1.1 \cdot J(\mathcal{S}_{\text{opt}})$, demonstrating that greedy can be suboptimal.
\end{exer}

\begin{exer}[Physics-guided refinement implementation]
\label{exer:ch13-implementation}
For a scaffold family of your choice:
\begin{enumerate}
    \item Compute energy decomposition for 50 molecules (docking scores with atom contributions)
    \item Estimate $\eta(s) = \Var(F_{\text{boundary}} \mid s) / \Var(F \mid s)$
    \item Identify the attachment-local feature most correlated with $F_{\text{boundary}}$
    \item Propose a refinement based on this feature and estimate the variance reduction
\end{enumerate}
\end{exer}

\begin{exer}[Stopping rule design]
\label{exer:ch13-stopping}
Design a cross-validation-based stopping rule for greedy refinement that:
\begin{enumerate}
    \item Accounts for estimation noise in $\hat{J}(\mathcal{S})$
    \item Balances bias (stopping too early) vs.\ variance (stopping too late)
    \item Provides a confidence interval for the optimal stopping level
\end{enumerate}
\end{exer}
% END SOURCE: chapters/ch13_refinement.tex


\chapter{Scaffold-Stratified Screening}
\label{ch:screening}
% BEGIN SOURCE: chapters/ch14_screening.tex
% ======================================================================
% Chapter 14 content file: chapters/ch14_screening.tex
% Scaffold-Stratified Screening with Guarantees
% (Do NOT include \chapter{...} here; main.tex already does that.)
% ======================================================================

\begin{quote}
\textit{Don't sample molecules; allocate risk.}
\end{quote}

\section{From Estimation to Screening Decisions}
\label{sec:ch14-intro}

The preceding chapters developed a mathematical infrastructure for reasoning about chemical space\index{chemical space}:
\begin{itemize}
    \item Chapter~\ref{ch:galois}: The Galois connection\index{Galois connection} quantifies the \emph{combinatorial overhead} of scaffold\index{scaffold}-based selection
    \item Chapter~\ref{ch:variance}: The variance spectrum\index{variance spectrum} quantifies how much information lives at each refinement level
    \item Chapter~\ref{ch:energy}: Energy and response decompositions state conditions under which structural information transfers
    \item Chapter~\ref{ch:mlmc}: Implementable coupled corrections support multilevel mean estimation under explicit bias and variance conditions
    \item Chapter~\ref{ch:refinement}: Optimal refinement selection aligns the hierarchy with score information
\end{itemize}

This chapter separates two decisions: identifying a class with a favorable
mean and recovering the best individual molecules. Scaffold classes can be
bandit arms for the first task. The second requires valid bounds on molecular
scores or on class extremes. A mean-estimation guarantee cannot supply that
missing bound.

\begin{framed}
\noindent\textbf{The central question:} Given a computational budget, how should we allocate evaluations across scaffolds to maximize the probability of finding the top molecules?
\end{framed}

The answer involves three nested levels of optimization:
\begin{enumerate}
    \item \textbf{Scaffold selection} (§\ref{sec:ch14-bandit}): Which scaffolds deserve more evaluation budget?
    \item \textbf{Within-scaffold search} (§\ref{sec:ch14-within}): Given a scaffold, how do we find the best decorations?
    \item \textbf{Multi-fidelity allocation} (§\ref{sec:ch14-multifidelity}): How do we combine cheap (docking) and expensive (FEP) evaluations?
\end{enumerate}

Theorem~\ref{thm:budget-scaling} gives an explicit recovery guarantee after a
valid selection certificate. Its cost adds the work of selection and final
verification. Closure size describes retained output; correction variances
describe an estimator. Neither is a substitute for that certificate.

%------------------------------------------------------------------

\section{Screening Objectives and Pipeline Overview}
\label{sec:ch14-objectives}

\subsection{What Are We Trying to Find?}

Different screening campaigns have different objectives. We formalize the three most common.

\begin{defn}[Top-$k$ recovery]
\label{defn:ch14-top-k}
For a finite library of $N$ molecules and distinct endpoint scores $F:\Mol\to\R$ (lower is better), the \emph{top-$k$ set}, $1\leq k<N$, is:
\[
T_k := \{m \in \Mol : |\{m' : F(m') < F(m)\}| < k\}
\]
The top-$k$ recovery problem is to identify $T_k$ with high probability using limited evaluations.
\end{defn}

\begin{defn}[Threshold hit recovery]
\label{defn:ch14-threshold-hits}
For threshold $\tau$, the \emph{hit set} is:
\[
H_\tau := \{m \in \Mol : F(m) \leq \tau\}
\]
The threshold hit recovery problem is to identify all (or a specified fraction of) $H_\tau$.
\end{defn}

\begin{defn}[Top scaffold identification]
\label{defn:top-scaffolds}
Let $\mu_s := \E[F \mid \mathrm{Scaf} = s]$ be the mean score for scaffold $s$. The \emph{top-$r$ scaffolds} are:
\[
S_r := \{s \in S : |\{t : \mu_t < \mu_s\}| < r\}
\]
Identifying top scaffolds is often a precursor to within-scaffold optimization.
\end{defn}

\begin{rem}[Relationship between objectives]
\label{rem:objective-relationship}
A best-mean class need not contain the best molecule. Let class $A$ have
scores $-10,10$ and class $B$ have scores $-1,-1$. The better mean is in $B$,
but the best molecule is in $A$. Even exact class means therefore do not
justify pruning $A$ for molecular top-1 recovery. A small mean-square
prediction error can also hide a rare extreme molecule. A uniform score
bound or a tail statement matched to the requested recall is needed.

\end{rem}

\subsection{The Query Model}

\begin{defn}[Evaluation oracle]
\label{defn:evaluation-oracle}
An evaluation request specifies a molecule $m$, protocol $f$, and replicate
identity. It returns $Y^{(f)}(m)$ at cost $c_f(m)$. Fix the endpoint $F$ to be
ranked and state how the observations relate to it. For example, a noisy
endpoint oracle may satisfy $\E[Y^{(f)}(m)]=F(m)$ with a known sub-Gaussian
variance proxy. A docking score used as a surrogate for another endpoint
requires a separate, validated bias or error model.

\end{defn}

\begin{defn}[Budget constraint]
\label{defn:budget}
A \emph{total budget} $B$ limits the sum of evaluation costs:
\[
C_{\rm setup}+\sum_j c_{f_j}(m_j)\leq B
\]
where $j$ indexes requests and $C_{\rm setup}$ includes preparation, pilot, and training costs not already counted as requests.
\end{defn}

\subsection{The End-to-End Pipeline}

Algorithm~\ref{alg:full-pipeline} gives the high-level structure. Subsequent sections develop each component.

\begin{algorithm}
\caption{Scaffold-stratified screening pipeline}
\label{alg:full-pipeline}
\mbox{}\\
\textbf{Input:} Finite library $\Mol$, endpoint $F$, cost budget $B$, target
error probability $\delta$, and a specified top-$k$ or threshold objective.

\textbf{Output:} A scored shortlist and an explicit certificate status.
A fixed budget can expire before exact recovery is certified.

\begin{enumerate}
    \item Build a nested representation and fit any surrogate on designated
    training data; record preparation, training, and pilot costs.
    \item Maintain score intervals or valid class-wide envelopes. Spend a
    portion of $\delta$ on their simultaneous coverage.
    \item Use scaffold means, uncertainty, and response contrasts to propose
    the next molecule or protocol. Mean-based acquisition is a heuristic for
    molecular top-$k$ recovery unless it also supplies a valid exclusion bound.
    \item Prune only when the molecule-level certificate below permits it;
    retain unbounded candidates as unresolved.
    \item Verify retained candidates at the specified endpoint. Stop with a
    certified answer when the intervals separate, or return the observed
    shortlist and unresolved set when the budget ends.
\end{enumerate}

\end{algorithm}

%------------------------------------------------------------------
\section{Scaffold-Level Exploration and Exploitation}
\label{sec:ch14-bandit}

The scaffold selection problem is a \emph{multi-armed bandit} problem: each scaffold is an arm, and pulling an arm means evaluating a random molecule from that scaffold.

\subsection{The Bandit Formulation}

\begin{defn}[Scaffold as bandit arm]
\label{defn:scaffold-arm}
Each scaffold $s \in S$ defines an arm with:
\begin{itemize}
    \item \textbf{Reward distribution:} $p_s = \mathrm{Law}(-F \mid \mathrm{Scaf} = s)$ (negated so higher is better)
    \item \textbf{Mean reward:} $\mu_s = -\E[F \mid \mathrm{Scaf} = s]$
    \item \textbf{Variance:} $v_s=\Var(F\mid\mathrm{Scaf}=s)$; the known sub-Gaussian proxy $\sigma_s^2$ used below may be larger than $v_s$
\end{itemize}
\end{defn}

\begin{defn}[Gap]
\label{defn:gap}
The \emph{gap} of scaffold $s$ relative to the best scaffold is:
\[
\Delta_s := \mu_{s^*} - \mu_s = \E[F \mid s] - \min_t \E[F \mid t]
\]
where $s^* = \arg\min_t \E[F \mid t]$ is the best scaffold.
\end{defn}

Scaffolds with small gaps are hard to distinguish from the best; scaffolds with large gaps can be eliminated quickly.

\subsection{Confidence Intervals for Scaffold Means}

The following bounds concern reward means, so larger $\mu_s$ is better throughout this section. Elsewhere in the chapter $F$ remains a score to minimize.

\begin{thm}[Sub-Gaussian confidence interval]
\label{thm:confidence-interval}
Let reward observations be independent draws from
$-F\mid\mathrm{Scaf}=s$, with
$\E[\exp(\lambda(X-\mu_s))]\leq\exp(\lambda^2\sigma_s^2/2)$ for every
$\lambda\in\R$, where $\sigma_s>0$ is a known variance-proxy bound. After $n_s$ i.i.d.\ samples from scaffold $s$, the sample mean $\hat{\mu}_s$ satisfies:
\[
\Pr\left(|\hat{\mu}_s - \mu_s| \geq \sigma_s \sqrt{\frac{2\log(2/\delta)}{n_s}}\right) \leq \delta
\]
\end{thm}

\begin{proof}
Standard Hoeffding/sub-Gaussian concentration. See Chapter~\ref{ch:probability} for details.
\end{proof}

\begin{defn}[Confidence bounds valid at every sample count]
\label{defn:UCB-LCB}
For $K=|S|$ fixed arms and $n\geq1$, set
\[
r_s(n)=\sigma_s\sqrt{\frac{2\log(4Kn^2/\delta)}n},\qquad
\mathrm{UCB}_s(n)=\widehat\mu_s(n)+r_s(n),\qquad
\mathrm{LCB}_s(n)=\widehat\mu_s(n)-r_s(n).
\]
At $n=0$ use the uninformative interval $[-\infty,\infty]$.
\end{defn}

At a fixed arm and sample count the failure probability is at most
$\delta/(2Kn^2)$. Summing over all arms and counts gives at most
$(\delta/2)\sum_{n\geq1}n^{-2}<\delta$. Thus every interval covers its reward
mean simultaneously with probability at least $1-\delta$, including when the
algorithm adaptively chooses the next arm or stopping time. A sample standard
deviation cannot simply replace $\sigma_s$ in this proof. Unknown scale
requires an appropriate self-normalized bound or a known bound on rewards.

\subsection{Confidence-Based Scaffold Elimination}
\begin{algorithm}
\caption{Round-robin elimination for the best reward mean}
\label{alg:UCB-elimination}
\mbox{}\\
\textbf{Input:} Fixed scaffold arms, known sub-Gaussian bounds, confidence
$\delta$, and a maximum cost budget.
\begin{enumerate}
    \item Initialize all scaffolds as active and their sample counts at zero.
    \item In each round, take one new independent reward observation from
    each active scaffold, if the complete round fits the remaining budget.
    \item Compute the bounds in Definition~\ref{defn:UCB-LCB}; let
    $b=\max_{t\ \mathrm{active}}\mathrm{LCB}_t$.
    \item Eliminate $s$ only if $\mathrm{UCB}_s<b$.
    \item Stop when one scaffold remains, or when the budget cannot fund
    another round. In the latter case return all survivors as unresolved.
\end{enumerate}
\end{algorithm}

\begin{prop}[Correctness of mean elimination]
\label{prop:elimination-correct}
On the simultaneous coverage event, the true best-mean scaffold is never
eliminated. If only one remains, it is the true best. This statement does not
claim that eliminated scaffolds contain no top-scoring molecules.
\end{prop}
\begin{proof}
For the best arm $s^*$, $\mathrm{UCB}_{s^*}\geq\mu_{s^*}$, whereas
$\mathrm{LCB}_t\leq\mu_t\leq\mu_{s^*}$ for every active $t$. Therefore
$\mathrm{UCB}_{s^*}<\max_t\mathrm{LCB}_t$ is impossible.
\end{proof}

\begin{thm}[A finite stopping bound with explicit gaps]
\label{thm:ucb-sample-complexity}
Assume a unique best arm and sufficient budget to continue the rounds.
For each $s\ne s^*$ let $n_s^\star$ be any integer for which
\[
2\bigl[r_s(n_s^\star)+r_{s^*}(n_s^\star)\bigr]<\Delta_s.
\]
On the simultaneous coverage event, arm $s$ is eliminated no later than
round $n_s^\star$. Consequently the total number of observations before
identification is at most
\[
\sum_{s\ne s^*}n_s^\star+\max_{s\ne s^*}n_s^\star.
\]
Such finite integers exist because $r_s(n)\to0$.
\end{thm}
\begin{proof}
Every surviving arm has one observation per completed round. On coverage,
$\mathrm{UCB}_s\leq\mu_s+2r_s$ and
$\mathrm{LCB}_{s^*}\geq\mu_{s^*}-2r_{s^*}$. The displayed inequality separates
these bounds. Count the observations of each suboptimal arm and then of the
best arm, which survives until the last elimination.
\end{proof}

\begin{rem}[The role of the gap]
\label{rem:gap-bounds}
The bound includes uncertainty in both arms and the repeated-look logarithm
through $r_s(n)$. Exact identification has no finite gap-independent
sample bound over Gaussian means with arbitrarily small positive gaps.
An $\epsilon$-optimal mean is a different objective with an explicit tolerance.
For finite molecular libraries sampled without replacement, use finite
population confidence bounds; the i.i.d. model here refers to independent
sampling under a fixed class law. An adaptive decoration policy changes that
law and cannot automatically reuse this arm-mean analysis.
\end{rem}

\subsection{Thompson Sampling Alternative}

\begin{algorithm}
\caption{Thompson sampling for scaffold selection}
\label{alg:thompson}
\mbox{}\\
\textbf{Input:} Prior distributions $p_0(\mu_s, \sigma_s^2)$ for each scaffold

\textbf{While} budget remains:
\begin{enumerate}
    \item For each scaffold $s$, sample $(\tilde{\mu}_s, \tilde{\sigma}_s^2) \sim p_t(\mu_s, \sigma_s^2 \mid \text{data})$
    \item Select scaffold $s_t = \arg\max_s \tilde{\mu}_s$
    \item Evaluate a molecule from $s_t$; observe score
    \item Update the posterior using the reward $-F$ and the specified likelihood
\end{enumerate}
\end{algorithm}

\begin{rem}[Thompson vs.\ UCB]
\label{rem:thompson-vs-ucb}
Thompson sampling is an alternative acquisition rule under a specified likelihood and prior. Its posterior ranking does not inherit the simultaneous confidence guarantee above; compare its recall and calibration empirically and use a separately justified stopping certificate.
\end{rem}

%------------------------------------------------------------------
\section{Within-Scaffold Search}
\label{sec:ch14-within}

Once promising scaffolds are identified, we must find the best molecules within each scaffold---the \emph{decoration optimization} problem.

\subsection{The Core Caching Advantage}

Chapter~\ref{ch:energy} established that scores factor as:
\[
F(s, d) = F_{\text{core}}(s) + F_{\text{decor}}(d) + F_{\text{boundary\index{boundary}}}(s, d)
\]
This representation permits caching only if the core geometry, atom mapping,
receptor state, and scoring protocol are shared, so that $F_{\rm core}(s)$ is
actually constant across the compared decorations. Relaxing each molecule can
invalidate that condition.

\begin{prop}[Core caching]
\label{prop:core-caching}
Under the fixed-core conditions just stated, $F_{\text{core}}(s)$ can be:
\begin{enumerate}
    \item Computed once and reused for all decorations
    \item Ignored for ranking (it's a constant offset)
\end{enumerate}
The effective search problem is to minimize $F_{\text{decor}}(d) + F_{\text{boundary}}(s, d)$ over decorations $d$.
\end{prop}

\begin{cor}[Speedup from factorization]
\label{cor:factorization-speedup}
If one core evaluation costs $C$, each decoration plus boundary evaluation
costs $D$, and no further recomputation is needed, scoring $n$ decorations
costs $n(C+D)$ without reuse and $C+nD$ with reuse. The speedup is
$n(C+D)/(C+nD)$. Atom counts alone do not determine these costs or establish
a speedup for a docking pipeline with pose search.

\end{cor}

\subsection{Decoration Space Enumeration vs.\ Generation}

Two paradigms for within-scaffold search:

\begin{defn}[Library-restricted search]
\label{defn:library-restricted}
Given a finite library $\mathrm{Fib}(s) = \{m : \mathrm{Scaf}(m) = s\}$, enumerate all decorations in the library.
\end{defn}

\begin{defn}[Generative search]
\label{defn:generative-search}
Generate decorations from a model (VAE\index{VAE (variational autoencoder)}, diffusion, SMILES-based) conditioned on scaffold $s$.
\end{defn}

\begin{rem}[Trade-offs]
\label{rem:enum-vs-gen}
\begin{itemize}
    \item Library-restricted: Finite candidate set; chemical validity and preparation quality still require checks
    \item Generative: Explores novel decorations; may generate invalid/unsynthesizable molecules
\end{itemize}
The choice depends on campaign goals. This chapter focuses on library-restricted search; Chapter~\ref{ch:generative} addresses generative models.
\end{rem}

\subsection{Stratified Sampling Within Scaffolds}

Even within a single scaffold, decorations can be stratified by attachment-point features (Chapter~\ref{ch:refinement}).

\begin{algorithm}
\caption{Stratified within-scaffold search}
\label{alg:stratified-within}
\mbox{}\\
\textbf{Input:} Scaffold $s$, decoration set $\mathrm{Fib}(s)$, attachment-point features $\phi$

\begin{enumerate}
    \item \textbf{Stratify:} Partition $\mathrm{Fib}(s)$ by attachment-point features into strata $\{D_1, \ldots, D_K\}$

    \item \textbf{Allocate for mean estimation:} With weights $w_j$ and standard deviations $\sigma_j$, use $n_j\propto w_j\sigma_j$ at equal cost, or $n_j\propto w_j\sigma_j/\sqrt{c_j}$ under a cost budget

    \item \textbf{Identify:} Within each stratum, track best observed molecule

    \item \textbf{Return:} Top molecules across all strata
\end{enumerate}
\end{algorithm}

\begin{prop}[Variance reduction from stratification]
\label{prop:stratification-variance}
For fixed strata, independent within-stratum samples with replacement, and
proportional allocation $n_j=nw_j$, the weighted mean estimator has variance
\[
\Var(\widehat\mu_{\rm strat})=n^{-1}\sum_j w_j\sigma_j^2,
\qquad
\frac{\Var(\widehat\mu_{\rm simple})}{\Var(\widehat\mu_{\rm strat})}
 =\frac{\Var(F\mid s)}{\sum_jw_j\sigma_j^2}\geq1,
\]
when the denominator is positive and the allocations are integral. At equal
cost, optimal Neyman allocation instead gives variance
$n^{-1}(\sum_jw_j\sigma_j)^2$. Both are statements about mean estimation,
not guarantees of recovering extreme molecules.

\end{prop}

%------------------------------------------------------------------
\section{Multi-Fidelity Evaluation}
\label{sec:ch14-multifidelity}

Real screening campaigns have access to multiple evaluation methods at different costs and accuracies. The key question is how to combine them.

\subsection{The Multi-Fidelity Hierarchy}

\begin{defn}[Fidelity levels]
\label{defn:fidelity-levels}
A \emph{fidelity hierarchy} is a sequence of scoring functions:
\[
F^{(0)}, F^{(1)}, \ldots, F^{(K)}
\]
under a specified common sampling law and on a justified common scale. Costs and errors relative to the chosen endpoint are measured, rather than ordered by method name.
\end{defn}

\begin{exmp}[A protocol comparison to measure]
\label{ex:fidelity-hierarchy}
Docking, MM-GBSA, and relative free-energy calculations can be compared as
candidate evaluation protocols. Record the actual endpoint, receptor and
ligand preparation, applicability domain, paired residuals, replicate
variability, and measured CPU/GPU cost. Neither their numerical scales nor
their error ordering is universal; a nominally expensive method is not a
reference truth merely because it is expensive.
\end{exmp}



\subsection{Telescoping Estimators}

From Chapter~\ref{ch:mlmc}, the MLMC identity:
\[
\E[F^{(K)}] = \E[F^{(0)}] + \sum_{k=1}^{K} \E[F^{(k)} - F^{(k-1)}]
\]
allows us to estimate the high-fidelity expectation using mostly low-fidelity evaluations.

\begin{defn}[Multi-fidelity MLMC estimator]
\label{defn:ch14-mf-mlmc}
Given sample allocations $\{n_k\}_{k=0}^K$, the multi-fidelity estimator is:
\[
\hat{\mu}^{\text{MF}} = \frac{1}{n_0} \sum_{i=1}^{n_0} F^{(0)}(m_i) + \sum_{k=1}^{K} \frac{1}{n_k} \sum_{i=1}^{n_k} \left[F^{(k)}(m_i^{(k)}) - F^{(k-1)}(m_i^{(k)})\right]
\]
where each pair has the required marginal law and the batches are independent
across levels. Choosing only high-scoring molecules for correction evaluation
changes the estimand unless the design and weighting correct that selection.
\end{defn}

\begin{thm}[Optimal multi-fidelity allocation]
\label{thm:mf-allocation}
Let $\Delta_0=F^{(0)}$ and $\Delta_k=F^{(k)}-F^{(k-1)}$ for $k\geq1$.
For independent correction batches unbiased for the displayed telescope,
let $v_k=\Var(\Delta_k)$ and $d_k$ be the full paired-evaluation cost.
With zero bias relative to $\E[F^{(K)}]$, the continuous optimal allocation is
\[
n_k=\epsilon^{-2}\sqrt{v_k/d_k}\sum_j\sqrt{v_jd_j},\qquad
\mathrm{Cost}=\epsilon^{-2}\left(\sum_k\sqrt{v_kd_k}\right)^2.
\]
Rounding, pilot costs, zero-variance cases, and bias relative to a different
endpoint are handled as in Theorem~\ref{thm:optimal-allocation}.

\end{thm}

\begin{proof}
Lagrangian optimization over sample sizes subject to fixed variance constraint. See Chapter~\ref{ch:mlmc}, Theorem~\ref{thm:optimal-allocation}.
\end{proof}

\begin{rem}[When multi-fidelity wins]
\label{rem:mf-wins}
Useful empirical conditions to assess are:
\begin{enumerate}
    \item Level correlations are high ($\Var(\Delta_k) \ll \Var(F^{(k)})$)
    \item Cost ratios are large ($c_k / c_{k-1} \gg 1$)
\end{enumerate}
These conditions must be measured on the intended population. High correlation alone does not imply a small difference variance if the outputs have different scales. The exact cost comparison is Corollary~\ref{cor:mlmc-vs-naive}.
\end{rem}

\subsection{Two-Stage Within-Scaffold Search}

Combining factorization and multi-fidelity:

\begin{algorithm}
\caption{Three-stage shortlist heuristic}
\label{alg:two-stage-within}
\mbox{}\\
\textbf{Input:} Scaffold $s$, decoration set $\mathrm{Fib}(s)$, budget $B_s$

\begin{enumerate}
    \item \textbf{Stage 1 (Docking sweep):}
    \begin{itemize}
        \item Evaluate all decorations with $F^{(0)}$ (docking)
        \item Rank by $F^{(0)}$; select top $N_1$ candidates
    \end{itemize}

    \item \textbf{Stage 2 (MM-GBSA refinement):}
    \begin{itemize}
        \item Evaluate top $N_1$ candidates with $F^{(1)}$ (MM-GBSA)
        \item Re-rank; select top $N_2$ candidates
    \end{itemize}

    \item \textbf{Stage 3 (FEP validation):}
    \begin{itemize}
        \item Evaluate top $N_2$ candidates with $F^{(2)}$ (FEP)
        \item Return final ranking
    \end{itemize}
\end{enumerate}

\textbf{Budget accounting:} For fixed per-candidate costs, choose affordable $N_1,N_2$:
\[
c_0 \cdot |\mathrm{Fib}(s)| + c_1 \cdot N_1 + c_2 \cdot N_2 \leq B_s
\]
\end{algorithm}

%------------------------------------------------------------------
\section{The Combined Guarantee}
\label{sec:ch14-guarantee}

A useful guarantee must certify every discarded molecule, not merely the
average of its scaffold.

\subsection{Score Intervals and Safe Pruning}
\begin{prop}[A top-$k$ inclusion and separation certificate]
\label{prop:molecular-topk-certificate}
Suppose intervals $[L_m,U_m]$ cover every endpoint score $F(m)$ simultaneously
at every decision time with probability at least $1-\delta$. Let $u_{(k)}$
be the $k$-th smallest upper bound over the finite library and set
\[
R=\{m:L_m\leq u_{(k)}\}.
\]
Then $T_k\subseteq R$ on the coverage event. Moreover, if a set $Q$ of size $k$
satisfies
\[
\max_{m\in Q}U_m<\min_{m\notin Q}L_m,
\]
then $Q=T_k$ on that event.
\end{prop}
\begin{proof}
Coordinatewise $U_m\geq F(m)$ implies $u_{(k)}\geq F_{(k)}$.
For $m\in T_k$, $L_m\leq F(m)\leq F_{(k)}\leq u_{(k)}$, proving inclusion.
The strict separation condition puts every score in $Q$ below every score
outside it, proving exact recovery.
\end{proof}

Unqueried molecules have $[-\infty,\infty]$ unless a validated structural or
surrogate argument provides a narrower interval. To prune an entire scaffold,
one needs a lower bound on \emph{every} member. For example, a justified
uniform residual envelope $F(m)\geq\mu_s-a_s$ for all $m\in\mathrm{Fib}(s)$,
together with a confidence lower bound on the score mean $\mu_s$, supplies
$L_m=\mathrm{LCB}^{\rm score}_s-a_s$. A variance estimate alone is not such an
envelope. Chapter~\ref{ch:response-certificates} studies related certificates
for controlled response comparisons.

\subsection{What the MLMC Functional Measures}
\begin{defn}[Estimator cost functional]
\label{defn:spectrum-complexity}
For implemented correction variances $v_\ell$ and full correction costs $c_\ell$,
\[
V(v;c)=\left(\sum_{\ell=0}^L\sqrt{v_\ell c_\ell}\right)^2.
\]
This includes the level-zero term and has the meaning established in
Chapter~\ref{ch:mlmc}: cost times sampling variance for an independent
multilevel estimator.
\end{defn}
\begin{rem}[Scope of the functional]
\label{rem:V-interpretation}
A small $V(v;c)$ makes the specified expectation cheaper to estimate.
To use that estimate in screening, derive a confidence interval with the
required simultaneous coverage and connect its estimand to the exclusion
condition. Substituting $V$ for the price of an exact molecular score does
neither of these things.
\end{rem}

\subsection{Selection Plus Verification}
\begin{thm}[Conditional budget for exact top-$k$ recovery]
\label{thm:budget-scaling}
Suppose a selection stage costs $B_{\rm sel}$ and returns a finite set $R$ with $|R|\geq k$
with $\Pr(T_k\subseteq R)\geq1-\delta_{\rm sel}$. Then:
\begin{enumerate}
    \item With an exact endpoint oracle costing $c_m$ per molecule, evaluating
    all of $R$ and returning its best $k$ costs
    $B_{\rm sel}+\sum_{m\in R}c_m$ and is correct with probability at least
    $1-\delta_{\rm sel}$.
    \item Suppose instead that fresh verification observations are independent,
    unbiased for $F(m)$, and have known sub-Gaussian proxies $v_m>0$.
    If a known $\Delta>0$ bounds the global gap
    $F_{(k+1)}-F_{(k)}$ from below, use
    \[
    n_m=\max\left\{1,\left\lceil\frac{32v_m}{\Delta^2}
                      \log\frac{2|R|}{\delta_{\rm ver}}\right\rceil\right\}.
    \]
    Ranking these sample means is correct with probability at least
    $1-\delta_{\rm sel}-\delta_{\rm ver}$, at total cost
    $B_{\rm sel}+\sum_{m\in R}n_mc_m$.
\end{enumerate}
\end{thm}
\begin{proof}
Condition on the selected set and the event $T_k\subseteq R$. Exact scores
settle the first claim. For the second, concentration and a union bound give
$|\widehat F(m)-F(m)|\leq\Delta/4$ simultaneously for all $m\in R$, except
with conditional probability $\delta_{\rm ver}$. A true top-$k$ molecule and
a non-top molecule then remain separated by at least $\Delta/2$ in estimated
score. Add the selection failure probability. Fresh verification noise makes
the conditional concentration argument valid after selection.
\end{proof}

\begin{cor}[Output closure enters only when imposed]
\label{cor:factor-interpretation}
If the selected set is scaffold-sound, $R=\gamma(\widehat S)$, then successful
inclusion forces $|R|\geq|c(T_k)|$. This is a lower bound on the number of
items in that retained set. It does not show that the selection stage can
find the true hit scaffolds, or that its cost is logarithmic in $|S|$.
The theorem's verification cost uses the actual $R$, not an unknown oracle
closure.
\end{cor}

\begin{rem}[When early pruning saves work]
\label{rem:when-scaffold-excels}
Savings require a small retained set, valid exclusion bounds, and a selection
stage whose cost is below the work avoided. Compare total training, pilot,
selection, and verification costs. If the evidence cannot bound unobserved
molecules, retaining the full library is the correct certificate outcome.
\end{rem}

\subsection{Exact Recovery and Approximate Means}
\begin{prop}[No finite gap-free exact-recovery budget]
\label{cor:gap-independent}
For noisy Gaussian endpoint observations, no finite budget uniformly guarantees
exact top-$k$ recovery with error below $1/2$ over all positive boundary gaps.
\end{prop}
\begin{proof}
The two-candidate case reduces to the Gaussian identification lower bound in
Chapter~\ref{ch:complexity}, which diverges as its positive mean gap tends to
zero. A uniform bound for the larger problem would also solve this case.
\end{proof}

An $\epsilon$-optimal mean is a distinct, gap-free objective. With $K$ arms
sharing sub-Gaussian proxy $v$, sampling each arm
$n\geq8v\epsilon^{-2}\log(2K/\delta)$ times ensures every reward-mean estimate
is within $\epsilon/2$ simultaneously. Its empirical maximizer is then within
$\epsilon$ of the best mean. This does not identify the exact winner or the
best molecule.

%------------------------------------------------------------------

\section{Algorithm Variants}
\label{sec:ch14-variants}

\subsection{Successive Halving}

\begin{algorithm}
\caption{Successive halving for scaffold selection}
\label{alg:successive-halving}
\mbox{}\\
\textbf{Input:} Scaffolds $S$, total budget $B$, rounds $R = \lceil \log_2 |S| \rceil$

\begin{enumerate}
    \item Initialize $S_0 \leftarrow S$, budget per round $B_r = B / R$

    \item For $r = 0, 1, \ldots, R-1$:
    \begin{enumerate}
        \item Allocate $n_r = B_r / |S_r|$ samples per surviving scaffold
        \item Evaluate; compute $\hat{\mu}_s$ for each $s \in S_r$
        \item Keep the $\lceil|S_r|/2\rceil$ largest estimated reward means; integer sample counts must fit the round budget
    \end{enumerate}

    \item Return $S_R$ (single scaffold) or top scaffolds from final round
\end{enumerate}
\end{algorithm}

\begin{prop}[Successive halving guarantee]
\label{prop:sh-guarantee}
Assume independent rewards with common sub-Gaussian proxy $v$, a unique best
mean, and a known lower bound $\Delta>0$ on its gap to every other mean.
Let $R=\lceil\log_2K\rceil$ for $K=|S|\geq2$ and use fresh samples at each
round. It is sufficient that each surviving arm receive at least
\[
n_{\min}=\left\lceil32v\Delta^{-2}\log(2KR/\delta)\right\rceil
\]
samples in every round. For the equal-round-budget version, taking at least
$K R n_{\min}$ total samples suffices. On simultaneous round-wise errors at
most $\Delta/4$, the best empirical mean is the true best and survives every
halving. A union bound over at most $KR$ arm--round estimates proves the
$1-\delta$ statement. This is a conservative sufficient bound.

\end{prop}

\begin{rem}[Successive halving vs.\ UCB]
\label{rem:sh-vs-ucb}
Successive halving is simpler and parallelizes well (all scaffolds in a round can be evaluated simultaneously). The round-robin elimination procedure also admits parallel batches but uses confidence separation rather than deleting a fixed fraction. These rules need separate error analyses.
\end{rem}

\subsection{Best-Arm Identification for Top-$r$ Scaffolds}

\begin{algorithm}
\caption{Top-$r$ scaffold identification}
\label{alg:top-r}
\mbox{}\\
\textbf{Goal:} Identify the $r$ largest reward means.
\begin{enumerate}
    \item Maintain simultaneous confidence bounds for every arm, including
    already provisionally selected arms.
    \item Form $Q$, the $r$ arms with largest empirical reward means.
    \item Certify only when
    $\min_{s\in Q}\mathrm{LCB}_s>\max_{t\notin Q}\mathrm{UCB}_t$.
    \item Otherwise sample unresolved boundary comparisons until the budget
    ends; return an unresolved status if separation has not occurred.
\end{enumerate}
\end{algorithm}

\begin{prop}[A sufficient top-$r$ mean budget]
\label{prop:top-r-complexity}
Suppose the reward means have a known boundary gap of at least $d>0$ and
all observations share proxy $v$. Uniformly sampling every arm
$\lceil32v d^{-2}\log(2K/\delta)\rceil$ times is sufficient for correct
empirical top-$r$ selection with probability $1-\delta$.
\end{prop}
\begin{proof}
All sample means are within $d/4$ simultaneously by concentration and a union
bound. No inside--outside ordering can then reverse. Both sides of the
boundary need observations; summing only over the outside arms omits part
of the estimation problem.
\end{proof}

%------------------------------------------------------------------

\section{Experiments}
\label{sec:ch14-experiments}

These are proposed validation studies. They do not report observed speedups.

\begin{exper}[Molecular recall at matched total cost]
\label{exp:recall-at-k}
Use a fully scored finite reference library with a prespecified endpoint and
$k$ values. Compare uniform sampling, fixed scaffold allocation, mean-based
bandits, molecular-feature active learning, and any certificate-driven policy.
Report observed top-$k$ recall, unresolved candidates, certificate coverage,
and total costs, including training and pilot labels. Include an oracle
mean-ranking policy as a diagnostic: even exact means need not recover rare
extremes. Use held-out contexts and repeated seeds; do not presume a positive
ranking of methods.
\end{exper}

\begin{exper}[Factorization and response-comparison ablation]
\label{exp:factorization-ablation}
Compare validated fixed-context response bounds with heuristic energy-component
scores and a no-factorization baseline. Record total-response residuals,
shared-parent uncertainty, coverage violations, and recall. A low component
variance ratio is not itself a guarantee; test whether it predicts any
practical saving after costs and covariance are included.
\end{exper}

\begin{exper}[Multi-fidelity versus single-fidelity]
\label{exp:multi-fidelity}
Fix a common endpoint and sampling population. Compare randomized paired
correction estimation, deterministic shortlist promotion, and direct endpoint
evaluation at the same total compute cost. Report mean-estimation error and
top-$k$ recall separately. Score correlation and nominal protocol price are
explanatory measurements, not assumed error guarantees.
\end{exper}

\begin{exper}[Certificate and budget validation]
\label{exp:budget-scaling}
On synthetic libraries, supply known valid and deliberately misspecified
score envelopes. Verify top-$k$ inclusion on every valid-envelope case and
measure how an invalid envelope can discard a true hit. For the noisy
verification theorem, specify the gap and variance proxies in advance,
measure failure rates across independent repetitions, and audit the additive
selection-plus-verification cost. This tests the actual assumptions of
Theorem~\ref{thm:budget-scaling} rather than fitting an unsupported product law.
\end{exper}

\section*{Summary}
Mean-based bandits rank class averages; molecular recovery needs bounds on
individual scores. Time-uniform confidence controls adaptive elimination,
implemented correction variances control MLMC allocation, and safe pruning
connects those observations to a recoverable molecule set. If a fixed budget
ends before separation, the output should identify unresolved candidates
rather than assert exact recovery.

%------------------------------------------------------------------

\section*{Exercises}

\begin{exer}[UCB confidence intervals]
\label{exer:ch14-ucb}
Derive the UCB confidence interval (Definition~\ref{defn:UCB-LCB}) from first principles:
\begin{enumerate}
    \item State the sub-Gaussian assumption precisely
    \item Derive the single-scaffold concentration bound
    \item Apply union bound to get simultaneous validity
    \item Discuss how to handle unknown $\sigma_s$
\end{enumerate}
\end{exer}

\begin{exer}[Thompson sampling implementation]
\label{exer:ch14-thompson}
Implement Thompson sampling for scaffold selection:
\begin{enumerate}
    \item Choose an appropriate prior (e.g., Normal-Gamma for unknown mean and variance)
    \item Derive the posterior update after observing scores
    \item Implement the sampling and selection loop
    \item Compare to UCB on a synthetic bandit problem with known ground truth
\end{enumerate}
\end{exer}

\begin{exer}[Multi-fidelity allocation]
\label{exer:ch14-mf}
For a two-level hierarchy (docking at cost 1, MM-GBSA at cost 100):
\begin{enumerate}
    \item Assume $\Var(F^{(0)})=4$, $\Var(F^{(1)}-F^{(0)})=1$, and a paired correction costs $101$; assume zero bias and independent batches
    \item Compute the optimal allocation ratio $n_0 / n_1$
    \item Calculate the cost to achieve MSE $\leq 0.01$ for the high-fidelity mean
    \item Specify $\Var(F^{(1)})$ before comparing with direct high-level estimation; it is not determined by the two variances above alone
\end{enumerate}
\end{exer}

\begin{exer}[Budget formula verification]
\label{exer:ch14-budget}
For a finite synthetic screening campaign:
\begin{enumerate}
    \item Construct simultaneous valid score intervals and the retained set
    $R$ of Proposition~\ref{prop:molecular-topk-certificate}.
    \item Verify that $R$ contains the exact top-$k$ set; compare $|R|$ with
    closure size when the selection is constrained to whole scaffold classes.
    \item Choose a known boundary gap and variance proxies, and compute the
    verification allocation in Theorem~\ref{thm:budget-scaling}.
    \item Audit total cost and failure probability over repeated independent
    observations. Explain why class means alone would not certify $R$.
\end{enumerate}

\end{exer}

\begin{exer}[Successive halving analysis]
\label{exer:ch14-sh}
Prove Proposition~\ref{prop:sh-guarantee}:
\begin{enumerate}
    \item Show that the best scaffold survives each round with high probability
    \item Bound the total number of samples across all rounds
    \item Compare to UCB sample complexity; identify when each method is preferable
\end{enumerate}
\end{exer}
% END SOURCE: chapters/ch14_screening.tex


% Exploratory source retained outside the revised core; see RESEARCH_AUDIT.md.

%--------Part V: Implications and Extensions--------
\part{Transfer, Generation, and Research Design}
\label{part:extensions}

\chapter{Compatibility and Response Transfer}
\label{ch:transfer}
% BEGIN SOURCE: chapters/ch16_transfer.tex
% Local compatibility, calibration, and transfer. Chapter heading is in main.tex.

\section{What can be transferred?}

A predictor trained on one chemical series may disagree with another predictor on molecules both can evaluate. There are three different questions: whether the predictions define one function on their combined domains; whether that function belongs to a chosen model class; and whether either predictor is accurate on unseen molecules. Compatibility answers the first question. It does not answer the other two.

Exact scaffold fibers form a partition. If $s\ne t$, then $\Fib(s)\cap\Fib(t)=\varnothing$. Local functions on distinct fibers therefore combine without any overlap constraint. To study overlapping models, we must specify larger domains: substructure families, unions of fibers, or explicit sets of molecules on which several predictors are evaluated. Sharing a fragment or a binding hypothesis does not by itself make two exact fibers overlap.

Chapter~\ref{ch:response-certificates} studies a more specific transfer question: when a response to a chemical edit in one context can be reused in another. Here we establish the compatibility mathematics needed to distinguish that question from ordinary function gluing.

\section{Domains and restrictions}

\begin{defn}[Alexandrov topology on a scaffold order]
\label{def:alexandrov}
For a scaffold poset $(S,\preceq)$, a set $U\subseteq S$ is open in the upper Alexandrov topology if $s\in U$ and $s\preceq t$ imply $t\in U$. Write
\[
\gamma(U)=\{m\in\Mol:\Scaf(m)\in U\}.
\]
The principal upset $\uparrow s$ represents molecules whose scaffolds extend $s$ under the specified order.
\end{defn}

This topology organizes domains. It imposes no monotonicity on activity, binding affinity, or prediction error. A property observed for $s$ need not persist in its extensions.

\begin{defn}[Presheaf of predictors]
\label{def:presheaf-predictors}
A presheaf assigns a space $\mathcal F(U)$ to each open domain and maps $\rho_{UV}:\mathcal F(U)\to\mathcal F(V)$ for $V\subseteq U$, with
\[
\rho_{UU}=\mathrm{id},\qquad \rho_{VW}\rho_{UV}=\rho_{UW}.
\]
For unrestricted real-valued predictors, take $\mathcal F(U)=\mathbb R^{\gamma(U)}$ with ordinary restriction. A restricted model family defines such a presheaf only if its declared spaces and restriction maps satisfy these requirements.
\end{defn}

\begin{defn}[Sheaf condition]
\label{def:sheaf-condition}
A presheaf is a sheaf if every family $f_i\in\mathcal F(U_i)$ agreeing on all overlaps has a unique section on $U=\bigcup_i U_i$ that restricts to each $f_i$.
\end{defn}

\begin{thm}[Gluing compatible functions]
\label{thm:gluing-obstruction}
Let $X_i=\gamma(U_i)$ cover $X$. For arbitrary functions $f_i:X_i\to\mathbb R$, there exists a function $f:X\to\mathbb R$ satisfying $f|_{X_i}=f_i$ if and only if
\[
f_i(m)=f_j(m)\quad\text{for every }m\in X_i\cap X_j.
\]
When it exists, this function is unique.
\end{thm}
\begin{proof}
Necessity follows by restricting $f$. For sufficiency, set $f(m)=f_i(m)$ for any domain containing $m$. Agreement makes this definition independent of the choice of $i$. The covering property gives existence and uniqueness at every point.
\end{proof}

There is no additional $H^1$ test for this theorem. In particular, exactly compatible real-valued functions cannot become inconsistent only on a triple overlap.

\begin{exmp}[A restriction on the global model]
On $X_1=\{0,1\}$ prescribe values $(0,0)$, and on $X_2=\{1,2\}$ prescribe $(0,1)$. They agree at $1$ and glue to a function on $\{0,1,2\}$. Each local function is the restriction of an affine function, but their union is not affine. The failure concerns the specified affine model class; it is not a failure to glue functions.
\end{exmp}

\section{The degree of the cohomological objects}

For a cover and a presheaf of vector spaces, the \v{C}ech cochains are
\begin{defn}[\v{C}ech complex]
\label{def:cech-complex}
\[
C^q(\mathcal U,\mathcal F)
=\prod_{i_0<\cdots<i_q}\mathcal F(U_{i_0}\cap\cdots\cap U_{i_q}),
\]
with alternating restriction maps. In degrees zero and one,
\begin{align*}
(\delta^0 f)_{ij}&=f_j|_{ij}-f_i|_{ij},\\
(\delta^1 g)_{ijk}&=g_{jk}|_{ijk}-g_{ik}|_{ijk}+g_{ij}|_{ijk}.
\end{align*}
The notation $|_{ij}$ denotes restriction to $U_i\cap U_j$, and likewise for triple intersections.
\end{defn}

\begin{prop}[Differences of local functions have zero second difference]
\label{prop:cech-square-zero}
For every $f\in C^0$, whether or not its components agree,
\[
\delta^1\delta^0 f=0.
\]
\end{prop}
\begin{proof}
On a triple intersection the expression is
$(f_k-f_j)-(f_k-f_i)+(f_j-f_i)=0$.
Restriction compatibility ensures that every occurrence of a function is evaluated on the same domain.
\end{proof}

\begin{defn}[\v{C}ech cohomology]
\label{def:cech-cohomology}
\[
H^0=\ker\delta^0,\qquad
H^1=\ker\delta^1/\operatorname{im}\delta^0.
\]
A family of local predictors is a degree-zero cochain. Compatible families lie in $H^0$ and are global sections when $\mathcal F$ is a sheaf. Its family of pairwise differences is an exact degree-one cochain and therefore has zero class in $H^1$.
\end{defn}

Taking norms of differences destroys their signed linear structure. A triangle residual formed from $\|f_i-f_j\|$ is not a \v{C}ech coboundary and cannot be used to declare $H^1\ne0$. Independently fitted edge corrections are different objects: they need not be differences of any common set of node values.

\section{Edge calibration and path consistency}

Suppose an oriented finite graph describes comparisons between local models. Let $g_{ij}=-g_{ji}$ be a separately estimated additive correction along edge $i\to j$. The calibration question is whether there are constants $b_i$ such that $g_{ij}=b_j-b_i$.

\begin{prop}[Graph calibration criterion]
\label{prop:contractible-gluing}
\label{prop:acyclic-gluing}
On a connected graph, edge corrections have the form $g_{ij}=b_j-b_i$ if and only if their signed sum along every cycle is zero. The $b_i$ are then unique up to one common additive constant. In particular, every edge assignment on a tree admits such constants.
\end{prop}
\begin{proof}
Differences telescope around cycles. Conversely, choose a root and define $b_i$ as the sum along a path from the root to $i$. Zero cycle sums make this independent of the path.
\end{proof}

\begin{cor}[What a tree removes]
\label{cor:hierarchical-covers}
A tree removes cycle constraints for this scalar calibration problem. It provides no test of whether the fitted corrections predict held-out data, and no guarantee that arbitrary local predictors agree on overlaps.
\end{cor}

\begin{exmp}[Three independently fitted corrections]
For edges $1\to2$, $2\to3$, and $1\to3$, take corrections $1$, $1$, and $3$. The cycle residual is $1+1-3=-1$, so no constants $b_i$ reproduce all three. If the triangle is filled as a two-simplex, this cochain is not closed and therefore does not represent an $H^1$ class. If the complex contains only its boundary edges, every one-cochain is closed, and this assignment represents a nonzero class. The chosen cells matter.
\end{exmp}

Appendix~\ref{app:sheaf} gives the matrix calculation and a counterexample to the claim that a contractible nerve forces vanishing cohomology for every coefficient system. A tree result for scalar identity transports must not be generalized to arbitrary presheaves.

\begin{rem}[Response transfer and thermodynamic cycles]
\label{rem:path-transfer}
Estimated response-transfer maps can also disagree across paths. Their domains, composition law, calibration data, and loss must be specified before a residual has meaning. Chapter~\ref{ch:response-certificates} supplies the response setting. An SAR rectangle compares the effect of an edit across contexts; its mixed difference can be nonzero for an exact scalar response. A thermodynamic loop instead sums exact differences of a state function around a closed path and is zero. These two calculations should not share an interpretation merely because both use four molecular states.
\end{rem}

\section{Approximate agreement and a usable global predictor}

\begin{defn}[Weighted patching objective]
\label{def:ls-global}
On a finite evaluation set $X$, assign nonnegative weights $a_i(m)$ to models defined at $m$. Assume $A(m)=\sum_i a_i(m)>0$ and minimize over all functions $f:X\to\mathbb R$:
\[
L(f)=\sum_{m\in X}\sum_{i:m\in X_i}a_i(m)(f(m)-f_i(m))^2.
\]
\end{defn}

\begin{prop}[Optimal patching and its residual]
\label{prop:optimal-patching}
Writing $w_i(m)=a_i(m)/A(m)$, the unique minimizer and pointwise minimum are
\begin{align*}
f^*(m)&=\sum_i w_i(m)f_i(m),\\
\min_x\sum_i a_i(m)(x-f_i(m))^2
&=\frac{1}{A(m)}\sum_{i<j}a_i(m)a_j(m)(f_i(m)-f_j(m))^2.
\end{align*}
\end{prop}
\begin{proof}
Complete the square about the weighted mean. The remaining weighted variance equals the stated pairwise expression.
\end{proof}

\begin{prop}[A compatibility error bound]
\label{prop:obstruction-transfer-bound}
If two local models disagree by $d$ at the same molecule, every scalar global prediction differs from at least one of them by at least $|d|/2$. If all active local values at a molecule lie in an interval of width $\epsilon$, any convex combination differs from each by at most $\epsilon$.
\end{prop}
\begin{proof}
The first claim is the triangle inequality. The second follows because a convex combination remains in the interval.
\end{proof}

\begin{rem}[Disagreement is not irreducible biological uncertainty]
\label{rem:uncertainty-obstruction}
These bounds measure error relative to the supplied models. They do not lower-bound error relative to the unknown response: one model may be correct and the others biased, or all may share the same error. Predictive intervals need held-out calibration under the intended sampling and assay protocol. Weight selection also needs validation; an inverse-variance rule assumes a specified error covariance model.
\end{rem}

\begin{algor}[Patching with recorded support]
\label{alg:optimal-patching}
For a query molecule, identify the models whose declared domains contain it. If none do, report unsupported transfer. Otherwise return the validated weighted average, the active model identities, and their disagreement. Apply an uncertainty calibration fitted on held-out observations. Keep model disagreement and calibrated predictive uncertainty as separate outputs.
\end{algor}

\section{An empirical transfer audit}

\begin{exper}[Compatibility and calibration audit]
\label{exp:h1-detection}
\label{exp:cohomology-computation}
Fix molecular identity rules, assay context, model domains, and a held-out evaluation set. Evaluate each pair of models on the same molecules in their domain intersection; empty intersections provide no agreement evidence. Record signed residuals, their uncertainty, and support size. If independent edge corrections are fitted, test them on a separate calibration split and compare cycle sums to their sampling uncertainty. Report the coefficient system and complex if computing a cohomology group. Neither a nonzero norm triangle nor model disagreement alone estimates that group.
\end{exper}

\begin{exper}[Does a transfer diagnostic predict error?]
\label{exp:transfer-vs-obstruction}
\label{exp:meta-learning}
Compare local-only, pooled, weighted, and adapted predictors under scaffold-held-out splits. Test whether calibration residuals or the response-curvature diagnostic of Chapter~\ref{ch:response-certificates} improve prediction of transfer error beyond training-set size, structural distance, and assay shift. Freeze diagnostic thresholds before evaluating the test split. No correlation magnitude or guaranteed advantage is assumed.
\end{exper}

The variance spectrum measures which response variation is lost by a chosen abstraction. It does not determine model compatibility. A constant response can have arbitrarily disagreeing fitted predictors; a highly variable response can be represented by perfectly compatible local restrictions of the same function. Experiments must distinguish response variation from estimation error and model restrictions.

\section*{Exercises}

\begin{exer}[Compatible functions on a cover]
\label{exer:ch16-cech}
Let $X=\{a,b,c\}$ and $X_1=\{a,b\}$, $X_2=\{b,c\}$, $X_3=\{a,c\}$. Give local restrictions of one function and verify they glue. Change one shared value and prove that exact gluing fails. State which values are constrained when the domains are disjoint exact scaffold fibers.
\end{exer}

\begin{exer}[Signs and degrees]
\label{exer:ch16-cocycle}
Prove $\delta^1\delta^0=0$. Explain why $f\in C^0$ does not define a class $[f]\in H^1$, and why replacing signed differences by norms invalidates the calculation.
\end{exer}

\begin{exer}[Tree calibration]
\label{exer:ch16-hierarchy}
Construct node offsets from arbitrary scalar edge corrections on a tree. Give two predictors on overlapping domains that disagree despite a two-node tree comparison graph.
\end{exer}

\begin{exer}[A triangle with and without its face]
\label{exer:ch16-cycles}
\label{exer:ch16-minimal}
For constant real coefficients, compute $H^1$ of the triangle boundary and of the filled triangle. Classify the correction assignment $(g_{12},g_{23},g_{13})=(1,1,3)$ in each complex.
\end{exer}

\begin{exer}[Variance does not determine compatibility]
\label{exer:ch16-variance-obstruction}
Construct both counterexamples in the last paragraph of the chapter. What additional assumptions would connect a response approximation error to disagreement between fitted models?
\end{exer}

\begin{exer}[Weighted patching]
\label{exer:ch16-patching}
Derive Proposition~\ref{prop:optimal-patching}. For predictions $1$ and $3$ with equal weights, compute the optimum and its residual. Give a true response for which only the first model is correct.
\end{exer}

\begin{exer}[Held-out transfer comparison]
\label{exer:ch16-meta}
Specify a train/calibration/test design for comparing pooled and adapted predictors. Identify which quantities can be used for weight selection, diagnostic thresholds, and final accuracy reporting without reusing test responses.
\end{exer}
% END SOURCE: chapters/ch16_transfer.tex


\chapter{Generative Proposals and Structural Validation}
\label{ch:generative}
% BEGIN SOURCE: chapters/ch17_generative.tex
% Generative proposals and structural validation. Chapter heading is in main.tex.

\section{Continuous probabilities and discrete molecules}

A molecular generator proposes discrete structures. Its internal computations may use continuous vectors, categorical probabilities, graph states, or combinations of these. The relevant questions are whether its outputs are valid under the chosen molecular representation, whether they satisfy the requested scaffold relation, and whether they improve the response under a fixed evaluation budget.

Connected latent spaces impose a simple restriction on deterministic discrete labels. That restriction does not imply that VAEs, diffusion models, or autoregressive models must generate invalid molecules. We first state the restriction precisely, then show why common probabilistic decoders do not satisfy its continuity hypothesis.

\section{A narrow topological statement}

\begin{thm}[Continuous discrete labels on a connected domain]
\label{thm:scaffold-impossibility}
Let $\mathcal Z$ be connected and give the scaffold set $S$ the discrete topology. If $h:\mathcal Z\to S$ is continuous, then $h$ is constant. Consequently a deterministic decoder $D$ attaining at least two scaffold labels cannot have a continuous scaffold-label map $h=\Scaf\circ D$.
\end{thm}
\begin{proof}
If $h$ were nonconstant, choose a value $s$ with a nonempty proper preimage. Both $h^{-1}(\{s\})$ and its complement would be nonempty open sets, giving a separation of $\mathcal Z$. This contradicts connectedness.
\end{proof}

The proof uses connectedness, not path connectedness. It also requires only two distinct scaffold labels; incomparability in the scaffold order plays no role. Changing the topology on $S$ changes the statement.

\begin{cor}[Scope of the discontinuity]
\label{cor:discontinuity}
A deterministic decoder attaining multiple scaffold classes must have a discontinuity in its discrete scaffold labels. This permits valid outputs on both sides of a decision boundary. It gives no lower bound on invalid-output probability, boundary probability, or prediction error.
\end{cor}

\begin{exmp}[A valid decoder with a single label transition]
Let $m_0,m_1$ be two valid molecules with different scaffolds and let $\mathcal Z=[-1,1]$. Define
\[
D(z)=\begin{cases}m_0,&z<0,\\m_1,&z\geq0.\end{cases}
\]
Every output is valid and both scaffolds are represented. Only the label map is discontinuous. Interpolating from $-1$ to $1$ changes the decoded scaffold once; it does not pass through an invalid molecule.
\end{exmp}

\begin{defn}[Stochastic decoder]
\label{def:scaffold-consistent}
A stochastic decoder is a probability kernel $K(z,\cdot)$ on molecular outputs, including a separate invalid-output symbol if the decoder can fail. For a finite molecular set, continuity of the decoder can mean continuity of each probability $K(z,\{m\})$. The induced scaffold probabilities are
\[
p_s(z)=\sum_{m:\Scaf(m)=s}K(z,\{m\}).
\]
\end{defn}

\begin{exmp}[Smooth probabilities across scaffold classes]
On $[0,1]$, set $K(z,\{m_0\})=1-z$ and $K(z,\{m_1\})=z$. This kernel is continuous in total variation, generates only valid molecules, and assigns positive probability to both scaffold classes in the interior. Such a mixture violates no scaffold-order axiom. Sampling a discrete output is not a continuous map into discrete labels, so Theorem~\ref{thm:scaffold-impossibility} does not apply to that sampling step.
\end{exmp}

Neural logits followed by categorical sampling or an argmax have this same distinction between continuous probabilities and discrete decisions. Autoregressive models need not have a single global latent space at all. Continuous relaxations used during training also need not be valid intermediate molecules. Their final outputs must be checked under the chosen decoding and validation procedure.

\section{Boundary probability needs assumptions on the prior}

\begin{defn}[A decoder's scaffold boundary neighborhood]
\label{def:boundary-measure}
For a deterministic label map $h$ on a metric latent space, define
\[
B_\epsilon(h)=\{z:\exists z'\in\mathcal Z,\ d(z,z')<\epsilon,
\ h(z)\ne h(z')\},\qquad
\mu_\partial(\epsilon)=P_Z(B_\epsilon(h)).
\]
This definition concerns label changes, not molecular invalidity. Incomparability can be recorded as a separate event.
\end{defn}

\begin{prop}[No positive universal boundary-mass bound]
\label{thm:boundary-lower}
For any fixed $0<\epsilon<1$, a connected latent support and positive probabilities for two decoded scaffolds do not imply a positive lower bound on $\mu_\partial(\epsilon)$ independent of the prior.
\end{prop}
\begin{proof}
Use the threshold decoder on $[-1,1]$. Let $Q$ assign equal mass to small intervals near $-1$ and $1$, each at distance greater than $\epsilon$ from zero. Set
\[
P_Z=(1-\eta)Q+\eta\,\mathrm{Unif}[-1,1],\qquad 0<\eta<1.
\]
The support is all of $[-1,1]$, each scaffold has probability $1/2$, and the $Q$ component assigns zero mass to $B_\epsilon$. Thus $\mu_\partial(\epsilon)=\eta\epsilon$, up to endpoints of measure zero, and tends to zero with $\eta$.
\end{proof}

Concentration claims require assumptions on the prior, the geometry and regularity of the decoding regions, and the neighborhood radius. Dimension and the number of scaffold labels alone do not supply such assumptions. Sampling random perturbation pairs estimates the probability of a label change under that perturbation kernel; it does not directly estimate the existential event in Definition~\ref{def:boundary-measure}.

\section{What scaffold constraints actually guarantee}

A scaffold constraint can specify exact membership $\Scaf(m)=s$, containment of a fragment, or a predecessor/successor relation. These are different predicates. Exact scaffold preservation is stronger than retaining a supplied fragment, since new ring construction may change the scaffold.

\begin{defn}[Scaffold-first generation]
\label{def:scaffold-first}
Choose a scaffold $s\sim p_S$ and sample a molecular proposal from a conditional distribution $K_s$. If every accepted output satisfies $\Scaf(m)=s$, the accepted distribution is supported on $\Fib(s)$.
\end{defn}

\begin{defn}[Fiber-restricted accepted generator]
\label{def:fiber-restricted}
A proposal mechanism together with a sound membership validator is fiber-restricted when it returns either an accepted molecule in $\Fib(s)$ or an explicit failure. Soundness means no molecule outside that fiber is accepted. It does not imply completeness, diversity, fast acceptance, or synthetic feasibility.
\end{defn}

\begin{prop}[Conditional sampling by rejection]
\label{prop:fiber-avoids-obstruction}
If a proposal law $Q_s$ has $q_s=Q_s(\Fib(s))>0$, independent proposals accepted exactly when they lie in $\Fib(s)$ have law $Q_s(\cdot\mid\Fib(s))$. The expected proposals per accepted molecule are $1/q_s$. With at most $B$ attempts, the failure probability is $(1-q_s)^B$.
\end{prop}
\begin{proof}
For an event $A$, summing over the first accepted attempt gives
$\sum_{n\geq1}(1-q_s)^{n-1}Q_s(A\cap\Fib(s))=Q_s(A\mid\Fib(s))$.
The waiting time is geometric.
\end{proof}

Training only on one scaffold does not establish this support guarantee. It can be supplied by a construction grammar with proved invariants or by post-generation validation. A single conditional model can cover many scaffolds; separate models per fiber are optional.

For the operators in this book, retain the requested operator, source identity, proposed target, and result of the exact structural predicate. Invalid proposals and valid molecules with the wrong relation are separate failures. A valid relation certificate establishes that the chemical move is admissible; response transfer requires the additional evidence developed in Chapter~\ref{ch:response-certificates}.

\section{Diagnostics with specified sampling rules}

\begin{defn}[Local scaffold agreement]
\label{def:scaffold-coherence}
Choose a latent prior $P_Z$, a perturbation kernel $Q_\epsilon(dz'\mid z)$, and a decoding rule. For deterministic decoding, define
\[
C_\epsilon=P[\Scaf(D(Z))=\Scaf(D(Z'))\mid
D(Z),D(Z')\text{ valid}].
\]
Report the probability of the conditioning event as well. For stochastic decoding, state whether decoder random seeds are shared or independent: agreement then reflects both latent perturbations and sampling variability.
\end{defn}

\begin{defn}[Held-out scaffold recoverability]
\label{def:scaffold-recoverability}
For models with an encoder $E$, train a classifier of $\Scaf(m)$ from $E(m)$ on training data and evaluate on a held-out split. Report class-balanced performance and compare with a scaffold-frequency baseline. This measures recoverability under the selected data distribution, not validity of generation.
\end{defn}

\begin{defn}[Order-respecting sampled path]
\label{def:scaffold-monotone}
For decoded labels $s_0,\ldots,s_T$ at a specified finite sequence of latent locations, the path is order-preserving if $s_0\preceq s_1\preceq\cdots\preceq s_T$. A chain only requires pairwise comparability and may move in both directions. Neither condition is necessary for molecular validity or useful scaffold hopping.
\end{defn}

\begin{defn}[Comparability fraction]
\label{def:monotonicity-score}
For a fully valid sampled path of $T+1\geq2$ labels, define
\[
\mathrm{Comp}(s_{0:T})=
\frac{1}{\binom{T+1}{2}}\sum_{i<j}
\mathbf1\{s_i\preceq s_j\text{ or }s_j\preceq s_i\}.
\]
Report invalid steps separately instead of deleting them without disclosure. This finite statistic measures comparability, not a direction of travel.
\end{defn}

\begin{defn}[Scaffold distribution divergence]
\label{def:scaffold-divergence}
For distributions $P,Q$ on valid molecules, compare their pushforwards $P_S=\Scaf_*P$ and $Q_S=\Scaf_*Q$. With consistent supports, KL divergence obeys
\[
D_{\mathrm{KL}}(P_S\|Q_S)\leq D_{\mathrm{KL}}(P\|Q).
\]
On a finite space this follows from the decomposition
\[
D_{\mathrm{KL}}(P\|Q)=D_{\mathrm{KL}}(P_S\|Q_S)
+\sum_s P_S(s)D_{\mathrm{KL}}(P(\cdot\mid s)\|Q(\cdot\mid s)).
\]
The identity uses the usual extended-value conventions. An empirical zero count need not imply a true zero probability.
\end{defn}

Thus small molecule-level KL forces small scaffold-level KL. The converse fails because distributions may agree on scaffold frequencies while differing within fibers. Report both scaffold coverage and within-fiber diversity; neither replaces the other.

\section{A comparison protocol}

\begin{exper}[Interpolation and output validation]
\label{exp:interpolation-validity}
\label{exp:monotonicity-audit}
For models with a declared latent interpolation rule, choose endpoint pairs and a finite evaluation grid before testing. Record representation validity, requested scaffold relation, scaffold transitions, and response quality at every decoded step. For stochastic decoders, repeat decoding and record variability. A jump in labels is not an invalid molecule. Valid endpoints do not guarantee valid intermediate outputs, but the failure rate is an empirical property of the model.
\end{exper}

\begin{exper}[Conditional generation under a shared budget]
\label{exp:scaffold-coherence}
\label{exp:scaffold-recovery}
\label{exp:scaffold-regularized}
\label{exp:fiber-restricted}
Compare an unconditional generator with filtering, a scaffold-conditional generator, and a constrained construction method on the same requested scaffolds. Match training-data access and total proposal/evaluation budgets. Report acceptance rate, unique accepted structures, relation correctness, held-out novelty, and downstream response. Include preprocessing and validator costs. Treat any latent separation regularizer as an ablation; it supplies no support guarantee on its own.
\end{exper}

\section{Generation and screening are competing proposal mechanisms}

Screening selects from a declared library; generation can propose outside it. Neither has an unconditional efficiency advantage. A generator concentrated on a productive region may find a hit without covering a fiber, while a poor generator may repeatedly propose the same molecule. Complete fiber coverage and hit discovery are different tasks.

\begin{prop}[A simple cost comparison]
\label{thm:regime-characterization}
Under independent proposals with fixed hit probability $p>0$ and fixed cost $c$ per proposal, the expected cost to the first hit is $c/p$. Under a fixed budget of $n$ proposals, the probability of at least one hit is $1-(1-p)^n$.
\end{prop}
\begin{proof}
The first-hit time is geometric, with mean $1/p$, and the event of no hit has probability $(1-p)^n$.
\end{proof}

For adaptive, correlated, or finite-library proposals these formulas need replacement by the appropriate sampling model. In practice, compare methods at equal total cost with duplicate handling, constraints, and evaluation fidelity fixed. A threshold on fiber size alone cannot choose the better method.

\begin{defn}[Scaffold-guided proposal pipeline]
\label{def:hybrid-pipeline}
Use observed data to choose promising scaffold regions; propose validated chemical moves within or between those regions; evaluate accepted molecules under a fixed response protocol; update the search policy. Retain failed proposals, computational costs, and provenance. The response-transfer tests of Chapter~\ref{ch:response-certificates} determine where previous response information can be reused.
\end{defn}

\section*{Exercises}

\begin{exer}[Connectedness and discrete labels]
\label{exer:ch17-impossibility}
Prove Theorem~\ref{thm:scaffold-impossibility} without assuming path connectedness. Explain why a continuous kernel on two valid molecules is not a counterexample to the theorem, but is a counterexample to a claimed universal invalid-generation consequence.
\end{exer}

\begin{exer}[Boundary mass]
\label{exer:ch17-boundary}
Compute $\mu_\partial(\epsilon)$ for the threshold decoder under a uniform prior on $[-1,1]$. Then construct the mixture prior in Proposition~\ref{thm:boundary-lower} and make this probability smaller than any prescribed positive number while preserving connected support.
\end{exer}

\begin{exer}[Finite path audit]
\label{exer:ch17-transitions}
\label{exer:ch17-monotone-subspace}
Give a chain-valued sampled path that is not order-preserving. Audit a trained model on a declared grid and report validity, transition counts, and comparability separately. State how endpoint sampling affects each statistic.
\end{exer}

\begin{exer}[Conditional generation and acceptance]
\label{exer:ch17-scaffold-first}
\label{exer:ch17-separation-loss}
Compare two proposal models whose scaffold acceptance probabilities are $0.2$ and $0.8$. Compute expected proposal counts for one accepted molecule and failure probabilities after ten attempts. Explain why those quantities alone do not rank diversity or response quality.
\end{exer}

\begin{exer}[Latent and diffusion diagnostics]
\label{exer:ch17-vae-topology}
\label{exer:ch17-diffusion}
Choose a model with a documented decoding rule. Identify which intermediate states can be decoded to valid molecules, and report scaffold labels only where the representation supports that operation. Distinguish a two-dimensional visualization from a proof of latent topology.
\end{exer}

\begin{exer}[Generation versus screening]
\label{exer:ch17-breakeven}
Under the independent-proposal model, compare a screening method with cost $c_s$ and hit probability $p_s$ to a generator with $c_g,p_g$. State the breakeven inequality and explain which assumptions fail for adaptive search without replacement.
\end{exer}
% END SOURCE: chapters/ch17_generative.tex


% Exploratory source retained outside the revised core; see RESEARCH_AUDIT.md.

\chapter{A Research Program for Molecular Decisions}
\label{ch:open}
% BEGIN SOURCE: chapters/ch19_open.tex
% Research questions are separated from proved finite statements.
\section{The first research question}
\label{sec:first-research-question}

Fix a finite candidate family $\mathcal M$, one response endpoint, a measurement
protocol, and a set size $k$. Let $f\in\mathbb R^{|\mathcal M|}$ be the unknown
response vector, with smaller values preferred. A representation $\alpha$
defines chemical contexts, comparable operations, and a candidate model class
$\mathcal P_\alpha$. A policy selects observations and decides when to stop.

\begin{quest}[Decision cost under a testable representation]
\label{prob:optimal-abstraction}
Can a policy that tests and, when necessary, refines $\mathcal P_\alpha$ reduce
the total cost of a correct top-$k$ decision relative to strong acquisition
baselines, while maintaining a declared error probability $\delta$ and an
explicit rate of returning no certificate?
\end{quest}

The question has an algebraic part and an empirical part. For a fixed model and
observation set, Chapter~\ref{ch:response-certificates} characterizes estimable
contrasts and robust decision margins. Whether a chemical representation makes
those conditions economical is an experimental question. Rejecting additivity
can be a useful result even when no alternative representation is yet adequate.

\section{A contract for the first study}

Before inspecting held-out outcomes, declare:
\begin{enumerate}
 \item the finite candidate domain, canonical parent identities, and allowed
 context--decoration constructions;
 \item the endpoint, its direction and units, measurement or simulation protocol,
 and rules for missing, censored, failed, or repeated observations;
 \item the target decision, tie handling, tolerance, confidence level, and
 maximum budget;
 \item the initial observations, candidate representations, permitted refinements,
 and independent confirmation procedure;
 \item all policy comparisons and the complete cost accounting.
\end{enumerate}

A computational first study can fix a docking or simulation endpoint under one
protocol. Its certificate then concerns that endpoint. Experimental activity
requires assay-specific data and separate validation. The present aggregated
ChEMBL response tables cannot supply that validation: they lack the assay
identifiers needed to establish a common protocol across their comparisons.
JAK2 is a candidate for continuity with RLMM2, subject to this data-feasibility
check; it is not a selected homogeneous assay panel.

\subsection{What counts as success}

Report the cost to a correct decision, the frequency of false certificates,
the frequency of abstention, and the cost to reject transfer violations above a
declared tolerance. Report discovery or top-hit recovery separately from these
calibration outcomes. A method that always abstains can be well calibrated but
useless; a method that always returns a ranking can be useful but uncertified.
Both rates are needed.

An accepted run ends with one of three outcomes: a decision certified within
the stated evidence model, a rejected representation, or exhausted budget
without a certificate. If refinement follows rejection, validation for the new
representation must account for the adaptive choice. The absence of a rejection
is not a proof of the model on unmeasured molecules.

\subsection{Comparisons and controls}

Start with random feasible selection, diversity selection, greedy predicted
response, and molecular uncertainty-based acquisition. Use an actual learned
molecular predictor for the latter; a hash feature control does not establish
performance against molecular active learning. Compare the proposed policy
with the same policy without transfer tests and with a fixed representation.
If a learned controller is added later, report the observations used to train
it as well as its evaluation budget.

Synthetic additive and interacting response tables test known failure modes.
A frozen retrospective benchmark then tests held-out predictions under explicit
selection limitations. A prospective or independently acquired homogeneous
panel is needed to test chemical transfer and measurement savings. Candidate
selection, calibration, and final confirmation must not silently reuse the
same responses for different evidential roles.

\section{How RLMM enters the experiment}

Legacy RLMM connects molecular changes to pose transfer, system preparation,
and simulation. RLMM2 has reusable oracle adapters, budgets, caches, and
experiment records. A measurement-design implementation would use these
components while adding explicit observation and protocol identities. The
choice to evaluate a new molecule differs from the choice to repeat a noisy
measurement on an old one.

The ordinary RLMM2 ChEMBL path currently provides candidate records without
parent lineage, whereas its scaffold graph builder creates edges from that
lineage. Saved graph topology and action traces must therefore establish which
runs navigate chemical relationships and which rely on global restarts. A
lineage edge also does not, by itself, certify an executable reaction.

The local June 2026 RLMM2 records preserve a positive ChEMBL1024 retrospective
result and a ChEMBL2048 scale boundary: the latter does not support superiority
over scaffold-greedy. These are historical computational comparisons under
their original baselines and budgets. Their primary remote artifacts have not
been reverified in this manuscript revision. A successor method must have a
new protocol and its own evidence; it cannot turn an earlier negative result
into a positive one by changing the comparison after observing outcomes.

\section{Open mathematical work}

\begin{prob}[Certificates under bounded model error]
\label{prob:non-additive}
For a fixed check matrix $H$, replace exact transfer $Hf=0$ by a declared
uncertainty set for $Hf$. Determine decision margins under that set and derive
procedures that calibrate its size using independent evidence. Separate
uncertainty about measurement noise from error in the transfer model.
\end{prob}

The exact-model linear programs do not establish this calibration. Allowing an
arbitrary discrepancy at every unmeasured molecule can remove all transfer
savings; useful bounds require assumptions that the study can test or justify.

\begin{prob}[Adaptive representation selection]
\label{prob:hierarchy-learning}
Choose among a fixed family of representations using earlier observations,
then construct valid decision uncertainty after that choice. Quantify the
cost of fresh validation, multiplicity correction, or a sequential procedure
whose guarantees cover the selection rule.
\end{prob}

The population variance gain from refinement is nonnegative but is not
generally submodular. Synergistic features can reveal information only when
used together. A greedy rule is therefore a baseline unless stronger structure
has been proved for the admissible refinement family.

\begin{prob}[Executable fidelity couplings]
\label{prob:flexibility}
Given two specified molecular evaluation protocols, construct paired
estimators of their difference and estimate the bias, variance, covariance, and
cost needed for a multilevel allocation. Determine whether the resulting
allocation improves a declared decision or only estimates a population mean.
\end{prob}

A scaffold hierarchy is a hierarchy of information. A fidelity hierarchy is a
family of executable evaluators. They may be related, but one does not provide
the other. Receptor adaptation additionally changes the evaluation context and
needs a common final protocol for fair comparison.

\begin{prob}[Generalization beyond the finite domain]
\label{quest:universality}
Which assumptions on a construction language and its response function allow
evidence collected on a finite family to control decisions about newly
constructed molecules? State the distribution, approximation error, or proof
obligation that connects the observed and unobserved domains.
\end{prob}

This is where the broader problem of construction and verification enters.
The present finite certificates are a starting point. They do not resolve
unrestricted molecular design, identify a universally best abstraction, or
prove that a scaffold representation is necessary for discovery.

Appendix~\ref{app:counterexample-design} explores a related extension: use
counterexamples to a proposed decision to design measurements and revise the
rules that permit transfer from one molecule to another. Its proposed
controller and representation-learning mechanisms are research hypotheses;
the finite results above do not depend on them.
% END SOURCE: chapters/ch19_open.tex


%--------Appendices--------
\appendix
\part*{Appendices}
\addcontentsline{toc}{part}{Appendices}

\chapter{Order Theory Primer}
\label{app:order}
% BEGIN SOURCE: appendices/app_order.tex
\section{Orders, embeddings, and bounds}

A preorder is reflexive and transitive. A partial order is also antisymmetric.
For finite typed graphs, non-induced subgraph embeddings give a preorder on labeled representatives and a partial order on isomorphism classes.
Mutual embeddings force equality of vertex and edge counts; a vertex-bijective embedding alone can still omit host edges.

A meet is a greatest lower bound; a join is a least upper bound.
Maximal lower bounds and minimal upper bounds need not have these properties.
A lattice has meets and joins for pairs; a complete lattice has them for every
family, including the empty family.

\section{Downsets and principal generators}

For $A\subseteq P$, write
\[
\downarrow A=\{x\in P:\text{there is }a\in A\text{ with }x\preceq a\}.
\]
A downset $I$ satisfies $I=\downarrow I$.
Every downset is the union of its principal downsets:
$I=\bigcup_{x\in I}\downarrow x$.
In a finite downset, every element lies below a maximal element, so
$I=\downarrow\max(I)$.
For an infinite downset this need not hold: $\mathbb{N}$ in its usual order has no maximal elements.

All downsets form a complete distributive lattice under inclusion, with union as join and intersection as meet.
Its atoms are singletons of minimal elements of $P$.
Even a finite downset lattice need not be atomistic: for $a<b$, its only atom is $\{a\}$, which cannot generate $\{a,b\}$ by union.

The word ideal can mean a downset or, more restrictively, a nonempty directed
downset. This book uses all downsets; arbitrary unions need not preserve
directedness.

\section{Compactness and finite representation}

An element $c$ of a complete lattice is compact if, whenever $c\leq\bigvee D$ for a nonempty directed family $D$, some $d\in D$ satisfies $c\leq d$.
The compact elements of the downset lattice are exactly the finitely generated downsets $\downarrow A$ with $A$ finite.
To prove the converse, express a compact $I$ as the directed union of $\downarrow A$ over all finite $A\subseteq I$ and apply compactness to that directed family.
The principal downsets alone need not form a directed family.

These are representation identities. They do not bound the cost of graph matching, enumerating generators, or measuring molecular responses.
% END SOURCE: appendices/app_order.tex


\chapter{Measure-Theoretic Background}
\label{app:measure}
% BEGIN SOURCE: appendices/app_measure.tex
\section{Conditional expectation as projection}
Let $(\Omega,\mathcal F,\mu)$ be a probability space, $F\in L^2(\mu)$,
and $\mathcal S\subseteq\mathcal F$ a sub-sigma-algebra. The conditional
expectation $P_{\mathcal S}F=\E[F\mid\mathcal S]$ is the orthogonal projection
of $F$ onto the closed subspace of $\mathcal S$-measurable square-integrable
functions. Hence, for every such $g$,
\[
 \|F-g\|_2^2=\|F-P_{\mathcal S}F\|_2^2+
 \|P_{\mathcal S}F-g\|_2^2.
\]
The minimizer is unique up to equality almost surely. For a finite molecular
library with positive weights, it is the weighted mean in each abstraction
fiber. Null-probability fibers have no identified conditional mean under this
measure.

\section{Nested information}
If $\mathcal S\subseteq\mathcal T$, the tower property gives
$P_{\mathcal S}P_{\mathcal T}=P_{\mathcal S}$ and orthogonality gives
\[
 \|F-P_{\mathcal S}F\|_2^2-
 \|F-P_{\mathcal T}F\|_2^2
 =\|P_{\mathcal T}F-P_{\mathcal S}F\|_2^2.
\]
The information sets must be nested. A list of independently chosen molecular
representations need not form a refinement tower. Population projections also
differ from predictors fitted on a finite, adaptively selected sample.

\section{Fixed-time and simultaneous uncertainty}
An interval with coverage $1-\delta$ at a fixed sample count is not automatically
valid when checked repeatedly until a favorable result appears. For a finite
planned collection of $J$ checks, intervals with failure probabilities
$\delta_j$ satisfying $\sum_j\delta_j\leq\delta$ provide simultaneous coverage
by a union bound. A confidence sequence or the explicitly defined fresh-data
procedure in Chapter~\ref{ch:response-certificates} supplies a different route
for sequential work. Reusing a cached response does not produce a fresh sample.
% END SOURCE: appendices/app_measure.tex


\chapter{Sheaf Cohomology Primer}
\label{app:sheaf}
% BEGIN SOURCE: appendices/app_sheaf.tex
\section{The objects and their degrees}

For this book, sheaf language describes local data, its restrictions, and exact compatibility. It does not by itself establish a claim about learnability or molecular response.

For a cover $\mathcal U=\{U_i\}$ and a presheaf of real vector spaces, set
\[
C^q=\prod_{i_0<\cdots<i_q}\mathcal F(U_{i_0}\cap\cdots\cap U_{i_q}).
\]
Restriction maps must compose consistently. The alternating differential then satisfies $\delta^{q+1}\delta^q=0$ and
\[
H^q=\ker\delta^q/\operatorname{im}\delta^{q-1}.
\]
In degree zero the denominator is zero. For a sheaf, $H^0$ is the space of global sections on the covered domain. For an arbitrary presheaf, $H^0$ consists of compatible local families; a separate gluing property is required to identify these with global sections.

A family $(f_i)$ of local predictors is a zero-cochain. Its pairwise differences $g_{ij}=f_j-f_i$ form an exact one-cochain. On every triple intersection,
\[
g_{jk}-g_{ik}+g_{ij}=0.
\]
Thus $\delta^1\delta^0 f$ cannot diagnose an obstruction. If this expression is nonzero numerically, check domains, sign conventions, restriction maps, and whether the edge quantities were independently estimated rather than obtained from the same local functions.

\section{Constant coefficients on a finite complex}

Let a finite simplicial complex have $n_0$ vertices, $n_1$ edges, and $n_2$ triangles. With one scalar per simplex and identity restriction maps, write
\[
B:C^0\to C^1,\qquad C:C^1\to C^2.
\]
For an oriented edge $i\to j$, $(Bb)_{ij}=b_j-b_i$. For a triangle $i<j<k$, $(Cg)_{ijk}=g_{jk}-g_{ik}+g_{ij}$. Then $CB=0$ and
\[
\dim H^1=n_1-\operatorname{rank}C-\operatorname{rank}B.
\]
This dimension depends on the complex and coefficient maps, not on measured prediction disagreements. Large scaffold covers need sparse algorithms; a cubic complexity claim in the number of domains alone ignores the potentially much larger number of overlaps.

For the triangle boundary, $n_0=n_1=3$, $C=0$, and $\operatorname{rank}B=2$, so $\dim H^1=1$. Filling its face adds one independent row to $C$ and gives $\dim H^1=0$. The assignment $(g_{12},g_{23},g_{13})=(1,1,3)$ has a nonzero cycle sum. On the boundary it is a nontrivial class; on the filled triangle it is not a cocycle.

\section{Why the coefficient system matters}

A contractible nerve has no positive-degree cohomology for its simplicial cochain complex with constant coefficients. This assertion does not extend to every presheaf.

For a finite explicit example, use the poset with two minimal points $a,b$ and one point $c$ above both. Its principal opens $U=\{a,c\}$ and $V=\{b,c\}$ cover it and have intersection $\{c\}$. Define a sheaf by giving zero-dimensional stalk spaces at $a,b$, the space $\mathbb R$ at $c$, and zero maps from the first two stalks to the third. Its sections satisfy
\[
\mathcal F(U)=0,\quad \mathcal F(V)=0,\quad
\mathcal F(U\cap V)=\mathbb R.
\]
For this cover, $C^0=0$, $C^1=\mathbb R$, $C^2=0$, so its \v{C}ech $H^1$ is $\mathbb R$, although the nerve is an interval. The example does not contradict the sheaf axiom: the only compatible sections on $U$ and $V$ are both zero, and they glue uniquely.

For the unrestricted function sheaf on a finite molecular set, compatible sections always glue pointwise. Indeed its positive-degree \v{C}ech cohomology vanishes for any finite cover: for each molecule $m$, the indices of domains containing $m$ form a full simplex, whose positive-degree cohomology is zero. The full cochain complex is a finite product of these pointwise complexes.

\section{What an obstruction statement must specify}

An $H^1$ obstruction may arise in a defined lifting problem or in independently supplied transition data, with a specified coefficient sheaf. One must give the one-cocycle, prove it is closed, and state what solving its exactness equation would accomplish. Nonzero $H^1$ only states that some nonexact cocycles exist; it says nothing about the class of a particular datum until that datum is supplied.

Higher degrees require the same care. A triangulated sphere can have $H^2\ne0$ without any three-simplices: $H^2=\ker\delta^2/\operatorname{im}\delta^1$ can be nonzero when $C^3=0$. No empirical frequency of such groups in QSAR is asserted here.

For molecular transfer, begin with the response and edit maps in Chapter~\ref{ch:response-certificates}. A scalar additive graph calibration problem, a nonlinear learned transport map, and a molecular response rectangle are different constructions. Their diagnostics should retain those distinctions.
% END SOURCE: appendices/app_sheaf.tex


\chapter{Code and Data}
\label{app:code}
% BEGIN SOURCE: appendices/app_code.tex
\section{One mathematical source, two documents}
The book is assembled by \texttt{main.tex}. The standalone note
\texttt{response\_curvature.tex} and the book both include
\texttt{chapters/response\_core.tex}; that file contains the shared response
identities, assumptions, proofs, and evidence limits. The command
\texttt{python} followed by \path{build_helper_scripts/verify_revision.py} runs the bounded
verification suite and writes its result manifest. LaTeX compilation is
recorded separately in the revision output.

\section{What the verifiers establish}
The directory \texttt{experiments/response\_curvature/} contains checks of
finite algebra, graph paths, robust linear decisions, precision designs, and
counterexamples used in this revision. The verifiers test deterministic cases,
small exhaustive problems, or simulations with recorded seeds. Passing these
checks does not replace a proof, establish novelty, calibrate chemical assay
noise, or demonstrate discovery benefit.

\section{Stored chemical evidence}
The response-graph outputs record molecular reassembly, graph dimensions,
parent-response dependencies, and retrospective transfer diagnostics. Their
source aggregates target-level activities without assay identifiers. Those
outputs establish structural opportunities; they cannot establish homogeneous
experimental nonadditivity. Chapter~\ref{ch:response-certificates} reports the
recorded rejection of the simple additive model.

The claim ledger \texttt{RESEARCH\_AUDIT.md} records the relation to legacy
RLMM and RLMM2. The local June 2026 RLMM2 summaries are historical computational
evidence. No remote run, new molecular simulation, or RLMM implementation
change is part of this revision.

\section{Revision provenance}
The original source files and their SHA-256 hashes are preserved in a dated
snapshot:
\begin{quote}
\path{tmp/revision-baselines/20260915-212126/}
\end{quote}
The revision output contains source comparisons, evidence hashes, command logs,
and verification results. The snapshot supplies a baseline because this
repository had no commits before the revision.
% END SOURCE: appendices/app_code.tex


\chapter{Verified Examples and Counterexamples}
\label{app:solutions}
% BEGIN SOURCE: appendices/app_verified_examples.tex
This appendix contains selected checked examples supporting the revised core.
The earlier full solution collection is retained in the project sources but is
not included here: its arguments have not all been reconciled with the new
chapter statements.

% BEGIN SOURCE: appendices/app_foundation_examples.tex
\section{Checked examples for the graph and order foundations}
\label{app:foundation-examples}

The examples below use non-induced embeddings of uniformly typed simple graphs.
Their carbon-skeleton interpretation treats hydrogen counts as implicit, not as preserved atom labels.
They accompany the finite verifier
\texttt{experiments/response\_curvature/verify\_foundation\_repairs.py}.
The verifier uses the Python standard library and does not claim empirical chemical validation.

\subsection{Molecular meet and join failure}

For Exercise~\ref{exer:ch2-chem}, the skeletons are
\[
E(G)=\{12,13,14,45\},\qquad
E(H)=\{12,23,13,14\}.
\]
Here $G$ is 2-methylbutane and $H$ is methylcyclopropane.
The star $K_{1,3}$ and the path $P_4$ embed in both.
Every common subgraph has at most four vertices, because $H$ has four vertices, and is acyclic, because $G$ is a tree.
A four-vertex forest containing either four-vertex tree has exactly its three edges.
It cannot contain the other tree, whose degree sequence differs.
Thus $G,H$ have no meet; the star and path have no join because such a join would have to embed in both $G,H$.

For comparison, replacing methylcyclopropane by methylcyclobutane would invalidate the argument: the entire five-vertex tree $G$ embeds in that five-vertex ring-and-pendant graph.
Equal vertex counts do not imply isomorphism under a one-way non-induced embedding.
Exercise~\ref{exer:ch2-math} instead uses mutual embeddings, which force equality of edge counts as well as vertex counts.

\subsection{The 2-core map}

For Exercise~\ref{exer:ch3-math}, surjectivity onto the image is automatic.
Injectivity depends on the domain: a library containing toluene and ethylbenzene has a collision, whereas a one-molecule library does not.
The graph 2-core is idempotent when extended to all finite typed graphs, including the empty graph.
This extension does not require core-only graphs to satisfy the original chemical-validity rules.

For an embedding $G\hookrightarrow H$, the image of $G$'s 2-core has internal minimum degree at least two.
Those vertices survive the pruning of $H$, so the original map restricts to an embedding of the two cores.
The verifier checks this containment for every vertex/edge-subgraph inclusion with hosts on at most five vertices, including empty and disconnected cases.
Changing atom types during normalization is outside that claim.

\subsection{Scaffold cones and order sizes}

For Exercise~\ref{exer:cone-cardinality}, the lower cone of the diamond graph, in the full connected 2-core domain with the empty scaffold, is
\[
\{\bot,C_3,C_4,D\}.
\]
The four-cycle $C_4$ is obtained by removing the diagonal from $D$.
It cannot be found by examining vertex-induced subgraphs alone.
Atom count stays at four while cycle rank rises from one to two, answering the first part of Exercise~\ref{exer:grading-comparison}.

For nonempty connected 2-cores, proper containment strictly increases cycle rank by Proposition~\ref{prop:ring-grading}.
Equal atom counts can have different cycle ranks ($C_4$ and $D$), while equal cycle ranks can have different atom counts ($C_3$ and $C_4$).
Neither descriptor's level partition therefore refines the other.

For Exercise~\ref{exer:fragment-merging}, let $Q$ be a triangle and a square sharing one vertex.
Both $D$ and $Q$ contain $C_3,C_4$.
Any upper bound of those two cycles lying below $D$ must be $D$: four vertices and a square require four edges, and a triangle requires the fifth.
The only triangle and square in $Q$ together use all of $Q$'s edges, so it too is a minimal common upper bound.
The graph $D$ has two triangles and cannot embed in $Q$; $Q$ has six vertices and cannot embed in $D$.
They are incomparable minimal upper bounds, so there is no join.

For Exercise~\ref{exer:complexity}, enumerate the $n!/(n-k)!$ injective maps from a $k$-vertex pattern to an $n$-vertex host, checking vertex and edge types for each.
With adjacency lookup, each check takes $O(k+|E_{\mathrm{pattern}}|)$ time.
The resulting bound is polynomial in $n$ for fixed $k$, but the exponent depends on $k$.
A path pattern has treewidth one; if its size equals the host size, injective matching asks for a Hamiltonian path.
Pattern treewidth alone therefore does not justify the unrestricted polynomial bound previously suggested by a local-consistency argument.

\subsection{Finite downsets and closure}

For Exercise~\ref{exer:ch5-math}, arbitrary unions and intersections preserve downward closure.
When $P=\{a<b\}$, the downsets are
\[
\emptyset,\quad \{a\},\quad \{a,b\}.
\]
The unique atom is $\{a\}$; this finite lattice is atomic but not atomistic.
It is nevertheless generated by principal downsets, since both nonempty downsets are principal.

For Exercise~\ref{exer:ch5-ml}, the smallest admissible downset containing all active scaffolds is $I^*=\downarrow\alpha(A)$.
Provided $\mu(A)<1$, its false-positive rate is
\[
\frac{\mu(\gamma(I^*))-\mu(A)}{1-\mu(A)}.
\]
The fiber closure $\gamma(\alpha(A))$ can be smaller than $\gamma(I^*)$ because it does not add sub-scaffolds.
These sets quantify classification output, not how many evaluations are needed to learn the active set.
% END SOURCE: appendices/app_foundation_examples.tex


\section{An ordinary-mean scaffold can contain the best molecule}
Consider two equally weighted fibers with scores $(-10,10)$ and $(-2,-1)$,
where smaller is better. Their means are $0$ and $-1.5$. Selecting the better
mean discards the molecule with score $-10$. Thus even exact knowledge of both
means cannot justify a claim about the best individual molecule without an
additional within-fiber condition.

\section{Refinement gains can be complementary}
Let $X_1,X_2$ be independent fair bits and let $F=X_1\mathbin{\mathrm{xor}}X_2$.
Then $\E[F\mid X_1]=\E[F\mid X_2]=1/2$ but
$\E[F\mid X_1,X_2]=F$. Either single split explains zero variance; both together
explain $\Var(F)=1/4$. The marginal gain from the second split increases after
the first. This violates diminishing returns and disproves generic
submodularity of explained variance over feature splits.

\section{A response rectangle and a telescoping loop differ}
For a $2\times2$ response table with entries $0,0,0,1$, the mixed contrast is
one. A loop formed from differences of a single state potential always sums to
zero. The rectangle tests whether two coordinate effects are additive. An
inconsistent loop of estimated free-energy differences tests whether those
estimates agree with a state-potential model. They have different experimental
meanings.

\section{Two elementary probability checks}
For a Rademacher response $X\in\{-B,B\}$, $\Var(X)=B^2$. Any MGF variance
proxy must be at least this variance, so $B/\sqrt2$ cannot be an admissible
sub-Gaussian scale. For a Gaussian $X\sim N(0,\sigma^2)$,
\[
 \E\exp(X^2/t^2)=(1-2\sigma^2/t^2)^{-1/2}
 \quad(t^2>2\sigma^2).
\]
Setting this expression equal to two gives the Orlicz norm
$t=\sqrt{8/3}\,\sigma$ for the definition used in Chapter~\ref{ch:probability}.
Neither calculation supplies a population tail bound from a finite sample.
% END SOURCE: appendices/app_verified_examples.tex


\chapter{Chemical Theory, Regret, and the Geometry of Transfer}
\label{app:counterexample-design}
% BEGIN SOURCE: appendices/app_counterexample_design.tex
% Research extension: the main book does not depend on this appendix.
% Source conversation and revision history: research/thread_idea_review.md.
\section{Begin with a theory worth using}
\label{sec:ced-question}

Chemical theory makes molecular search possible by allowing us to leave things
out. A scaffold representation groups structures into families. A pairwise
score replaces a coupled physical system with a collection of tractable terms.
A conformer-generation protocol keeps a manageable account of the configurations
a molecule can occupy. Each supplies a useful starting description on which to
make a choice.

This appendix takes that starting description seriously. Assume that the
representation, scoring terms, and preparation protocol have been accepted for
the campaign. Use them to predict, compare, and choose. The question is what to
do next when the chosen molecule could be wrong for a reason that the working
theory suppresses.

The proposed answer is to let regret guide the acquisition of missing detail.
We seek an omitted interaction, state, or distinction that could make another
choice materially better. We then construct an experiment that discriminates
that possibility from the working account. A useful correction restores a
shared term in the theory and changes what can be inferred about molecules
that have not yet been measured.

\begin{quote}
Chemical theory supplies the initial compression. Regret determines where
additional resolution is worth acquiring. Curvature records the context
dependence that a transferable correction must explain.
\end{quote}

The ambition is to learn which distinctions a molecular decision requires.
Chemistry can contain far more structure than a particular choice needs.
A successful campaign may leave substantial response error unresolved while
learning a small interaction that determines which lead to pursue. Conversely,
a model can fit most measured responses well and omit precisely the term that
reverses its recommendation.

That difference changes the purpose of refinement. Its target is an adequate
account of the consequences of choosing, with enough retained structure to
make the next choice useful. The initial theory remains valuable throughout:
it supplies candidate molecules, proposes plausible omissions, and gives each
challenger a margin it must overcome.

The exact results below describe finite response vectors and declared sets of
possible corrections. The proposal is to use those results to organize a
chemical campaign. Whether the resulting procedure reduces experimental cost
is a research question. Its value is not established by assuming that the
starting theory is useful; that assumption gives the procedure a place to begin.

\section{What a representation permits us to ignore}
\label{sec:ced-compression}

A representation carries a modeling commitment when a response rule is built
on it. If a predictor depends only on scaffold identity, it treats decoration
differences within a fiber as irrelevant to that prediction. An additive
context--decoration model allows decorations to matter but shares their effects
across contexts. A fixed-pose score can resolve many local interactions while
leaving changes in conformational populations outside the calculation.
These compressions omit different things, and an experiment should challenge
the particular omission on which the decision relies.

Here ``smooth'' means that the working theory shares effects, ties parameters,
or shrinks certain interactions toward zero. It need not mean differentiability
in a continuous molecular coordinate. A scaffold hypergraph supplies a chemical
organization over which those assumptions can be made; the response model
specifies how the organization is used. For this appendix, both are adopted as
the working theory rather than held in suspension pending a proof of all their
physical premises.

Let its current predictions on $N$ candidates be $g\in\mathbb R^N$.
For a response to be minimized, choose
\[
 i\in\operatorname*{arg\,min}_{1\leq j\leq N}g_j.
\]
Write the actual response as $f=g+r$. The residual $r$ can include parameter
error, omitted physics, or the discrepancy between a computational endpoint
and the specified reference endpoint. The identity does not require that we
learn every coordinate of $r$.

The distinction that matters for a challenger $j$ is
\[
 f_i-f_j=(r_i-r_j)-(g_j-g_i).
\]
The working theory grants $i$ a margin $g_j-g_i$. A correction matters to this
comparison only through its difference between $i$ and $j$, and only to the
extent that it consumes that margin. A large common offset disappears. A
large error at an uncompetitive molecule may leave it uncompetitive. A small
error across a narrow margin can change the decision.

We can therefore retain some approximation error deliberately. What may be
ignored depends on the endpoint, candidate set, tolerance, and decision horizon.
An offset irrelevant to ranking can matter to a potency threshold; an interaction
irrelevant to today's four candidates can matter to a molecule we plan to make
next. The appropriate resolution belongs to the task as well as to the molecule.

\section{Regret specifies the resolution the decision needs}
\label{sec:ced-witness}

Let $\mathcal D_t\subseteq\mathbb R^N$ be a nonempty compact set of admissible
corrections to $g$, given the current evidence and declared challenge assumptions.
It may be a finite collection of alternatives, a bounded polyhedron, or a
parameterized uncertainty set. It records which departures we are prepared to
investigate. It is not defined merely by setting every omitted term to zero.

The accepted approximations suggest how to construct this set. Release a
suppressed conformer or protonation state, allow a particular receptor response,
or introduce a candidate cooperative term. A tractable auxiliary calculation
can propose the direction and scale of its effect. Chemical feasibility and
existing evidence narrow the challengers to investigate. These calculations
supply challenge hypotheses; a claimed statistical envelope needs its own
calibration. The campaign can begin with a finite set of such alternatives
rather than an unconstrained search over all possible errors.

Put $\mathcal F_t=g+\mathcal D_t$ and define
\begin{equation}
 R_t(i)=\sup_{f\in\mathcal F_t}\left[f_i-\min_j f_j\right].
 \label{eq:ced-regret}
\end{equation}
This is the largest loss from choosing $i$ among the admitted response vectors.
It is an analysis of a choice under a supplied family of challenges.

\begin{prop}[An omitted effect must overcome a margin]
\label{prop:ced-margin}
Let $\Delta_j=g_j-g_i\geq0$ and define the support function
$\sigma_{\mathcal D_t}(c)=\max_{r\in\mathcal D_t}c^\top r$.
For the working minimizer $i$,
\begin{equation}
 R_t(i)=\max_j\left\{
 \sigma_{\mathcal D_t}(e_i-e_j)-\Delta_j\right\}.
 \label{eq:ced-support}
\end{equation}
The term $j=i$ is zero. A candidate $j$ can improve on $i$ by more than
$\epsilon$ exactly when some admitted correction satisfies
$r_i-r_j>\Delta_j+\epsilon$.
\end{prop}
\begin{proof}
For each $r$, $g_i+r_i-\min_j(g_j+r_j)
=\max_j\{r_i-r_j-\Delta_j\}$. The maximum over finitely many $j$ commutes
with the maximum over the compact correction set. The last statement is the
strict version of the same comparison.
\end{proof}

The formula measures an omission in the direction of a decision contrast.
It supplies an algebraic meaning to buying resolution: acquire information
about corrections whose contrast exceeds, or could exceed, a relevant margin.
The largest residual norm and the most consequential residual direction need
not identify the same experiment.

\begin{prop}[A region of errors that the choice can tolerate]
\label{prop:ced-decision-region}
For $\epsilon\geq0$, define
\[
 \mathcal C_i(\epsilon)=
 \bigcap_j\{r:(e_i-e_j)^\top r\leq\Delta_j+\epsilon\}.
\]
Then $R_t(i)\leq\epsilon$ if and only if
$\mathcal D_t\subseteq\mathcal C_i(\epsilon)$.
If the true residual belongs to $\mathcal D_t$, this condition makes $i$
$\epsilon$-optimal for the true response.
\end{prop}
\begin{proof}
The set inclusion states that every admitted correction satisfies every
competitor inequality. Taking maxima gives exactly
$\max_r\max_j(r_i-r_j-\Delta_j)\leq\epsilon$.
\end{proof}

The region $\mathcal C_i(\epsilon)$ is a polyhedron, generally not a cone:
multiplying an admissible error by a large scalar can cross a fixed margin.
It may contain response vectors with large absolute errors. It represents
errors tolerable for this decision rather than accurate reconstructions of
the whole response surface. At $\epsilon=0$ the condition permits ties;
strict competitor inequalities establish a unique winner.

\begin{figure}[htbp]
\centering
\begin{tikzpicture}[scale=1.65]
  \fill[blue!5] (-1.5,-1.5) rectangle (1.5,1.5);
  \fill[orange!28] (.5,1.5)--(1.5,.5)--(1.5,1.5)--cycle;
  \draw[blue!65!black,thick] (-1.5,-1.5) rectangle (1.5,1.5);
  \draw[->] (-1.75,0)--(1.9,0) node[right] {$b_1$};
  \draw[->] (0,-1.75)--(0,1.9) node[above] {$b_2$};
  \draw[orange!80!black,thick] (.22,1.78)--(1.78,.22);
  \node[rotate=-45,fill=white,inner sep=1pt,font=\scriptsize]
    at (1.06,1.06) {$b_1+b_2=2$};
  \fill (0,0) circle (.025) node[below right,font=\scriptsize] {working theory};
  \fill[blue!65!black] (-1.25,-1.25) circle (.025);
  \node[align=center,font=\scriptsize] at (-.68,-.8)
    {large correction,\\same preferred candidate};
  \node[align=center,font=\scriptsize] at (.25,-1.95)
    {admitted corrections: $|b_1|,|b_2|\leq1.5$};
  \node[align=left,font=\scriptsize,anchor=west] at (1.85,1.4)
    {a challenger\\can win};
  \draw[->,thin] (1.84,1.25)--(1.36,1.3);
\end{tikzpicture}
\caption{One comparison in correction coordinates. The working margin is two
and the residual contrast is $b_1+b_2$. Only the shaded corner of the admitted
box permits a reversal; corrections elsewhere can remain without defeating
this choice. Additional candidates supply additional comparison boundaries.}
\label{fig:ced-decision-region}
\end{figure}

For a simple bound, suppose $N\geq2$ and every admitted correction has
$\|r\|_\infty\leq\eta$. With
$\Delta_{\min}=\min_{j\ne i}\Delta_j$,
\[
 R_t(i)\leq\max\{0,2\eta-\Delta_{\min}\}.
\]
Zero gaps from tied working minima must be included. A large margin can make
a coarse model sufficient; a near tie makes small unresolved terms expensive.
No universal numerical accuracy target is appropriate for both cases.

\begin{prop}[An attained counterexample to the decision]
\label{prop:ced-witness}
For a nonempty compact $\mathcal F_t$, each fixed recommendation has
\[
 R_t(i)=\max_j\max_{f\in\mathcal F_t}(f_i-f_j).
\]
Both maxima are attained. If $R_t(i)>\epsilon$, an attaining pair
$(f^\star,j^\star)$ exhibits an admitted response under which the recommendation
loses by more than $\epsilon$.
\end{prop}
\begin{proof}
Each contrast is continuous on a compact set and there are finitely many
contrasts. The equality follows by expanding the minimum in the regret.
\end{proof}

Such a witness identifies an unresolved scientific possibility. Experimental
design must still turn it into an available measurement. Finding a witness is
not the same as discovering a physical mechanism. Failing to find one with an
approximate optimizer is not an upper bound on regret. Empty response sets
mean model--data conflict; unbounded sets may leave regret unbounded. The
finite certificates of Chapter~\ref{ch:response-certificates} provide suitable
solver behavior once those distinctions are explicit.

\section{A worked example: restore the missing interaction}
\label{sec:ced-example}

Consider a deliberately simplified lead-design problem. There are two chemical
families, one preorganized and one flexible, and two decorations, small and
bulky. A supplied fixed-pose, pairwise scoring model predicts the response in
Table~\ref{tab:ced-working-theory}. Scores are illustrative units of a fixed
endpoint, not measured affinity data. Lower is better. Pairwise atomic scoring
does not generally imply additive context--decoration responses: geometry,
parameterization, and pose selection can all produce context dependence.
The equal substitution gains below are a feature of this particular working
prediction, chosen to make its challenged simplification explicit.

\begin{table}[htbp]
\centering
\begin{tabular}{lrr}
\toprule
Working theory & Small decoration & Bulky decoration\\
\midrule
Preorganized family $A$ & $-8.0$ & $-10.0$\\
Flexible family $B$ & $-8.4$ & $-10.4$\\
\bottomrule
\end{tabular}
\caption{An accepted working theory shares the two-unit benefit of the bulky
substitution across both families. It recommends $B$-bulky.}
\label{tab:ced-working-theory}
\end{table}

We act on this theory and evaluate $B$-bulky. Suppose its response is $-9.2$.
The observation differs from the prior by $+1.2$. That discrepancy alone does
not identify which term was missing. Three simple working challengers explain
it exactly:
\begin{enumerate}
 \item $H_1$: a $+1.2$ penalty applies to every bulky decoration;
 \item $H_2$: a $+1.2$ offset applies to every member of family $B$;
 \item $H_3$: a $+1.2$ interaction applies only to the flexible--bulky combination.
\end{enumerate}

These hypotheses restore different distinctions. The first recalibrates a
shared decoration effect. The second recalibrates a family. The third says that
a previously shared effect must depend on context. One possible interpretation
of $H_3$ is an omitted cost associated with the population of bound-compatible
conformers. That interpretation motivates a testable correction; the toy scores
do not establish that mechanism.

The decision guides the next experiment. Under $H_1$, $A$-bulky is predicted
at $-8.8$ and the observed $B$-bulky remains best. Under $H_2$ or $H_3$,
$A$-bulky stays at $-10.0$ and becomes the preferred candidate. Measuring
$A$-bulky directly resolves this disagreement. Suppose it is $-10.0$.
We retain $H_2$ and $H_3$ and choose $A$-bulky among the original four.

At this point, a measurement of $B$-small would distinguish the remaining
hypotheses: $H_2$ predicts $-7.2$, whereas $H_3$ predicts $-8.4$.
But both hypotheses already select $A$-bulky. For the exact four-candidate
problem, with these as the only remaining alternatives, its worst-case regret
is zero. Further explanation is not required to make that particular choice.
We can leave the distinction unresolved if the campaign ends there.

Now extend the declared decision to a proposed conformationally locked analogue
$C$. It belongs to family $B$ and has the bulky decoration, but the predeclared
flexibility descriptor is zero. The working score is $-10.3$.
The hypotheses make different predictions:
\[
 H_2(C)=-9.1,\qquad H_3(C)=-10.3.
\]
Under $H_2$, $A$-bulky remains preferred; under $H_3$, the proposed analogue is
better. The previously uncompetitive $B$-small compound is now an instrument
for deciding whether $C$ deserves a more costly synthesis or evaluation.
Its information value arises from this new decision horizon, not from an
unexplained claim that every diagnostic molecule improves current regret.

If $B$-small evaluates at $-8.4$, the retained toy hypothesis is $H_3$.
The bulky change then improves family $A$ by two units and flexible family $B$
by only $0.8$. The correction can be written
\[
 f(s,d)=a_s+b_d+
 \beta\,\chi_{\rm flexible}(s)\chi_{\rm bulky}(d),
 \qquad\beta=1.2.
\]
The theory has acquired one context-dependent interaction rather than an
independent offset for every molecule. Its scope includes predeclared new
combinations, including the locked analogue, where it makes predictions that
can fail. A fresh evaluation of $C$ tests the proposed extension.
In a real noisy experiment, overlapping intervals would replace exact
elimination of hypotheses, and additional explanations could remain.

This is the family-level consequence of a measurement. A compound can be a
poor lead and still identify a coefficient that changes an entire set of
unmeasured choices. The resulting knowledge can also create a design action:
remove the feature predicted to activate the penalty. A descriptor invented
only after observing $C$ would not supply that prospective test.

\section{What is learned when the theory changes}
\label{sec:ced-learning}

Three operations should remain distinct. Parameter fitting estimates a
coefficient for an effect already represented. Model selection chooses among
a supplied set of possible corrections. Representation learning introduces a
new reusable distinction that the old model could not express. The worked
example is model selection if the flexible--bulky interaction was listed in
advance. It becomes a representation-learning proposal if the failure leads
to a new, independently computable descriptor whose predictions are then tested.

An unrestricted per-molecule intercept can fit every observed exception. That
fit does not explain why an effect transfers or tell us how to alter a new
molecule. A useful additional term has a domain, a construction or evaluation
rule, and consequences beyond the point that prompted it. Its value comes from
what those consequences allow the campaign to do.

This supplies a stronger interpretation of ``learning what to ignore.'' We
are learning a task-appropriate effective theory: a working description that
retains the distinctions necessary for a specified class of decisions and
continues to suppress others. The word effective describes its declared scope
and resolution. It does not assert a universal physical reduction.

The approximation may remain wrong in many places. Suppose a large residual
occurs on a poor candidate that remains far from the leading set under every
admitted correction. Improving that prediction can be scientifically useful,
but it has little immediate value for the current selection. Meanwhile a
small systematic interaction across a narrow lead margin may justify an
expensive evaluation. Regret organizes the restoration of detail by its
consequences, while chemical theory organizes which restorations are plausible.

A failed challenge has value too. If a discriminating observation rules out a
specific decision-reversing correction, the campaign may continue using the
coarser account for that comparison. It has learned a reason that this
particular distinction can remain unresolved. This is a scoped conclusion
about the tested alternatives, not a proof that no unimagined effect exists.

\section{Curvature separates calibration from context-dependent error}
\label{sec:ced-prior}

Freeze a prior $f_0$ on a consistently typed context--decoration panel and
write $Y=f_0+r$. For four valid cells, define
\[
 \Omega_{ss';dd'}(Y)=Y_{sd}-Y_{sd'}-Y_{s'd}+Y_{s'd'}.
\]
The same decoration change must have a common chemical meaning across the two
contexts. On that ground set the contrast records whether its effect transfers.

\begin{prop}[Prior-relative curvature]
\label{prop:ced-prior-curvature}
For arbitrary row offsets $a_s$, column offsets $b_d$, and a fixed prior,
\begin{align}
 \Omega(Y+a_s+b_d)&=\Omega(Y),\label{eq:ced-calibration}\\
 \Omega(Y-f_0)&=\Omega(Y)-\Omega(f_0).\label{eq:ced-residual}
\end{align}
\end{prop}
\begin{proof}
Each row and column offset appears once with each sign. The second identity
is linearity of the four-corner contrast.
\end{proof}

In the example, $H_1$ and $H_2$ have zero residual curvature. $H_3$ has residual
curvature $+1.2$ for the displayed row and column ordering. It restores the
context dependence that the additive working model omitted. Regret identifies
why the failed transfer might matter; curvature distinguishes that failure
from an additive correction.

These roles are related but do not coincide. A context offset can reverse the
best molecule without creating curvature. For instance,
\[
 f_0=\begin{pmatrix}0&1\\2&3\end{pmatrix},
 \qquad
 Y=\begin{pmatrix}4&5\\2&3\end{pmatrix}
\]
have the same rectangle curvature but different preferred cells. A constant
offset preserves rankings and can still change a threshold decision.
Thus calibration remains part of choosing correctly. What survives additive
recalibration is the interaction contrast, not every decision.

The sparse-graph form is equally direct. With the parent-response check matrix
$H$ from Chapter~\ref{ch:response-certificates}, an additive residual model is
\[
 H(Y-f_0)=0,\qquad Y\in f_0+\ker H.
\]
A representation of a measurement is not another measurement: different
fragmentations of one parent retain their shared response and uncertainty.
The prior can be updated between rounds, but a declared fresh check must not
be reclassified as validation of a prior refitted to that same outcome.

For paired, aligned evaluation levels,
\begin{equation}
 \Omega(Y_H)=\Omega(Y_L)+\Omega(Y_H-Y_L).
 \label{eq:ced-fidelity}
\end{equation}
This lets the campaign ask whether a cheaper evaluator transfers the wrong
effect in a comparison that matters. The two endpoints need declared alignment
and a discrepancy model; the identity does not itself give a cheap, unbiased
estimator. Existing multifidelity methods provide a framework for measuring
accuracy and cost while exploiting cheaper evaluations~\cite{peherstorfer2018multifidelity}.

Nor should a response rectangle be confused with a thermodynamic loop.
A nonzero rectangle can represent real context dependence in four responses.
Exact free-energy differences around a state cycle telescope to zero;
nonclosure of their estimates concerns sampling or estimation consistency.
These diagnostics call for different corrections.

\section{A budget of omitted interactions}
\label{sec:ced-geometry}

One way to make the proposal operational is to maintain a chemical dictionary
of possible correction terms. For selected terms $A$, write the current model
as
\[
 g_A=f_0+\Phi_A\widehat\theta_A.
\]
The columns might describe a shared attachment effect, a conformational state,
a protonation distinction, or a typed interaction between them. Keeping a term
can require data, computation, and a more elaborate preparation protocol.
Omitting it is a modeling choice whose decision effect can be challenged.

Suppose the remaining discrepancy is represented by
$r=\Phi_B b$ with $b$ in a supplied bounded set. The coefficient uncertainty
of retained terms can also be included in this discrepancy. Choose $i_A\in\operatorname*{arg\,min}_i g_{A,i}$ and use its updated margins
$\Delta_j^{(A)}=g_{A,j}-g_{A,i_A}$.
The effect of a correction on competitor $j$ is
\[
 r_{i_A}-r_j=v_j^\top b,
 \qquad v_j=\Phi_B^\top(e_{i_A}-e_j).
\]
Only these projections enter the current regret. An interaction constant
across the compared molecules disappears from the comparison even if its
absolute contribution is large.

For an independent centered box $|b_\ell|\leq\tau_\ell$ with $\tau_\ell\geq0$,
\begin{equation}
 \max_b v_j^\top b=\sum_\ell\tau_\ell|v_{j\ell}|.
 \label{eq:ced-box}
\end{equation}
Each coefficient takes the sign maximizing its term. For the ellipsoid
$b^\top Q^{-1}b\leq\rho^2$, with $Q$ positive definite and $\rho\geq0$,
\begin{equation}
 \max_b v_j^\top b=\rho\sqrt{v_j^\top Qv_j}.
 \label{eq:ced-ellipse}
\end{equation}
This follows by writing $b=Q^{1/2}u$ and applying Cauchy--Schwarz on
$\|u\|_2\leq\rho$; the bound is attained in the direction of
$Q^{1/2}v_j$. Inserting either support into \eqref{eq:ced-support} gives a
finite regret bound under that discrepancy model.

The shapes encode scientific assumptions. A box permits every extreme sign
combination; if it encloses a correlated set, it may give a conservative bound.
The ellipsoid records correlated directions. Data must justify how these sets
are calibrated for any claimed coverage. An unlisted class of omitted physics
is not controlled just because the listed coefficients have small intervals.
The theory can still be used provisionally while those additional challenges
are investigated.

This formulation clarifies the intuition of changing a basis. An invertible
coordinate change within the same response span changes its description, not
its expressive power. If $\Phi'=\Phi T$ and $b'=T^{-1}b$, transform the entire
coefficient set as well. The admitted response vectors, support values, and
regret remain unchanged. Transforming the basis while leaving an independent
coefficient box fixed generally changes the scientific assumptions.

Restoring a missing interaction does more. Adding a column outside the current
response span admits predictions that were previously impossible. Splitting a
context can likewise allow a formerly shared effect to vary. This is the
substantive sense in which the campaign changes the geometry: the set of
permitted response relationships changes. The molecules themselves have not
been shown to inhabit a curved manifold by this algebra.

A decision-directed representation objective can therefore be posed as finding
a low-cost retained set $A$ for which the admitted omitted effects cannot make
$i_A$ lose by more than $\epsilon$. Cost includes learning the retained terms
and obtaining evidence about the terms left out. This is a proposed optimization
problem; there is no general greedy solution. Complementary features and panels
can be necessary, as Chapter~\ref{ch:refinement} demonstrates.

The ambitious hypothesis is that this sufficient representation can be much
smaller than one needed to predict every response accurately. It becomes a
measurable claim only after specifying the candidate class, available correction
language, reference endpoint, and cost of assessing its omissions.

\section{Learning within a theory and revising the theory}
\label{sec:ced-two-updates}

The campaign needs two kinds of update. Within a fixed response family,
compatible observations remove possibilities:
\[
 \mathcal F_{t+1}\subseteq\mathcal F_t.
\]
Worst-case regret for any fixed action can only decrease. This is ordinary
learning under a model. It improves estimates and resolves contenders the
model already knows how to express.

A structural challenge can reveal that the family itself excludes an important
possibility. Enlarging or changing it may increase honest regret. Start with
only the response table
\[
 g=\begin{pmatrix}0&1\\2&3\end{pmatrix}.
\]
Its preferred cell has zero regret within that singleton family. Admit also
an interaction of $-5$ confined to the bottom-right cell. The alternative table
has best response $-2$, so the original choice at zero now has worst-case
regret two. No new molecule has become worse during this change of mind.
The representation has stopped ruling out a way the choice could be wrong.

An acquisition reward based solely on decreasing the controller's own interval
width would treat this correction of overconfidence as a setback. A useful
objective must evaluate the eventual decision and the declared validity of its
evidence, while allowing a model failure to increase uncertainty temporarily.
Otherwise a controller can appear efficient by refusing distinctions that
would expose its mistakes.

The two directions also explain why changing a theory differs from merely
conditioning a fixed prior. Conditioning cannot assign positive probability
to a response excluded by the prior's support. A model that already includes
all proposed interactions can learn their weights; adding a new descriptor or
family of effects changes the model under which learning occurs. In practice
both operations can be useful, and their costs should be recorded separately.

There is an exact accounting rule in the nested linear setting. On the same
$N$ parent coordinates, if
$\mathcal P_0\subseteq\mathcal P_1$ and
$\dim\mathcal P_1-\dim\mathcal P_0=d$, then
\[
 (N-\dim\mathcal P_0)-(N-\dim\mathcal P_1)=d.
\]
Each additional response degree removes one independent linear constraint.
Fitting every parent separately leaves no transfer constraint to test. This is
why a correction should predict fresh combinations, not merely explain its
training exception. Arbitrary nonlinear or changing-domain models need their
own accounting; they do not inherit this dimension identity unchanged.

\section{Construct experiments that expose the missing term}
\label{sec:ced-bridges}

The most useful next observation need not be the molecule with the largest
predicted reward or uncertainty. In the worked example, $B$-small tests which
correction should govern the proposed analogue $C$. It is useful because of
the inference it enables, not because it is itself a promising lead.

Panels can have complementary value. If a decision requires $u-v$ and the
available exact observations reveal $u$ and $v$ separately, neither alone
identifies the contrast on an unrestricted two-parameter domain. Together
they do. In the row-span criterion of Chapter~\ref{ch:response-certificates},
$(1,-1)$ belongs to the span of both measurement vectors and neither singleton
span. A rule that stops whenever immediate estimability gain is zero misses
this panel.

A chemical bridge panel can connect two series through comparable transformations
or supply the reference measurements needed to isolate an interaction.
Connectivity and rank identify a possibility for inference; noise, conditioning,
parent dependence, and preparation cost determine whether it is useful to buy.
Generating such a panel gives molecular generation a scientific-instrument role.
The output is a set of feasible compounds that makes an important distinction
observable.

For a fixed nonempty compact response family, let
$q(\mathcal F)=\min_i R_{\mathcal F}(i)$. One possible exact-observation panel
criterion is
\begin{equation}
 q(\mathcal F_t)-
 \sup_{y:\,\mathcal F_t(P,y)\ne\emptyset}
       q\bigl(\mathcal F_t(P,y)\bigr),
 \label{eq:ced-panel}
\end{equation}
where $\mathcal F_t(P,y)$ is the family after panel outcome $y$.
This is guaranteed reduction in the best remaining worst-case regret under a
fixed model. Compactness keeps it finite; the unbounded $u-v$ example is an
estimability illustration, not an instance in which to subtract infinite
regrets. Expected-value planning instead requires a declared distribution over
outcomes. Neither criterion by itself values an unanticipated model failure.

A fuller controller must therefore retain a way to challenge the current
representation, with independent checks or exploration that can expose omitted
alternatives. Refusing all experiments that appear uninformative under the
accepted model would make the model increasingly hard to contradict. The
challenge policy should be targeted, but its definition cannot depend only on
possibilities already forced to have zero effect.

\section{RLMM as a controller of scientific resolution}
\label{sec:ced-controller}

The resulting role for RLMM is to coordinate molecular actions with actions
that improve the working theory. A state would include the current molecular
and physical context, evidence, retained correction terms, unresolved
challengers, available constructions, and remaining resources. Its actions
could propose a lead, evaluate a reference compound, collect another replicate,
change fidelity, construct a bridge panel, test a new descriptor, or stop.

Some of these actions change chemistry; others change what can be concluded
about chemistry. Their value can be delayed. A reference measurement can select
a correction, the correction can suggest an intervention, and that intervention
can produce a better molecule. The worked example makes this sequence explicit
without assigning value to explanation independently of a declared future task.
Sequential control is worth investigating when such dependencies matter.
A short-horizon planner remains an essential comparator.

Frontier models can propose the missing terms and the experiments that could
separate them. A proposed term should become an executable hypothesis: a
pre-outcome descriptor or mapping, a domain of use, a predicted signed contrast,
and a feasible intervention or measurement. A mechanistic story acquires
scientific force through those consequences. In the example, the proposal to
lock a conformer matters because it changes a prediction before the analogue
is evaluated.

The retained physical state also matters. If an earlier molecular edit changes
the receptor ensemble used for later scores, a molecule label alone may not
specify the decision state. A controlled reset protocol or a representation of
the relevant history is needed; a belief over hidden states is another possible
implementation. The abstraction must preserve the information required by its
control problem, as well as the contrasts required by its response model.

Existing RLMM components provide possible construction and evaluation operations.
The proposed controller should be compared with simpler acquisition and panel
planning policies using the same accepted theory, initial evidence, and budget.

\section{Related ideas and the proposed contribution}
\label{sec:ced-antecedents}

Goal-oriented error estimation already asks which local approximation errors
matter to a chosen quantity of interest. Dual-weighted residual methods provide
a developed version in finite-element analysis~\cite{becker2001dwr}.
Goal-oriented modeling-error analysis compares tractable physical descriptions
with richer reference models~\cite{oden2002modelingerror}. These are close
antecedents to spending effort on omitted terms that affect a molecular
contrast. Their error estimators require specified mathematical models; they
are not automatic bounds for docking or scaffold predictors.

Counterexample-guided abstraction refinement begins with an abstraction, checks
a property, and uses a counterexample to identify a necessary
refinement~\cite{clarke2000cegar}. Its classical setting uses sound
overapproximations, often removing spurious abstract behavior. A chemical
working theory can instead exclude real interactions, so adding a missing term
may enlarge its prediction family. The common principle is targeted refinement;
the direction and validity of the set update must be established separately.

Transductive experimental design distinguishes the vectors one can measure from
the candidates one must select~\cite{fiez2019transductive}. It supplies a precise
fixed-model basis for informative reference compounds and decision contrasts.
Molecular pool-based active learning provides a practical acquisition
baseline~\cite{graff2021molpal}. The additional proposal here is to let
consequential failures change the retained chemical response terms and, when
needed, the representation that defines those terms.

The contribution to test is therefore a chemical program of selective theory
refinement: use accepted structure to make decisions, challenge its consequential
omissions, and carry supported corrections into new predictions and designs.
The established components provide the comparison against which this chemical
integration should be assessed.

\section{An experiment that tests the deeper claim}
\label{sec:ced-pilot}

The first experiment should ask whether decision-directed restoration buys a
better cost--decision tradeoff than globally improving the response model.
Hold the molecular domain and evaluator fixed initially. Supply every policy
the same working prior and initial measurements. This isolates the effect of
how it spends on missing detail from the effect of having a better initial
predictor.

A possible controlled panel has 64 unique parents with compatible response
measurements, a typed construction graph fixed in advance, and known costs.
Reveal the same eight initial parents to each policy, with a budget equivalent
to 24 further parent measurements. Replication and fidelity changes consume
that budget at their recorded cost. Acquiring the complete reference table
still costs 64 parent measurements plus replication; masked replay does not
make that reference evidence free.

Synthetic controls should include a decision-changing interaction, a much
larger but decision-irrelevant interaction, additive bias with zero curvature,
a pair of complementary corrections, and a model that is confidently wrong
because an important term was excluded. Include a correctly adequate coarse
model: a method should also be able to retain a useful simplification instead
of always adding complexity.

Compare a fixed prior, ordinary molecular uncertainty acquisition, global
residual-error acquisition, and regret-directed acquisition using the same
permitted corrections. Then compare a fixed correction dictionary with a
procedure allowed to propose new descriptors, reserving fresh combinations to
assess those proposals. A separate extension can introduce new candidates and
bridge panels, with the candidate expansion and decision horizon recorded as
in the worked example. Compare short panel lookahead against learned control
before attributing an advantage to RL.

Report cost to an $\epsilon$-optimal decision, actual regret on the reference
panel, false certificates, abstention, and the number and cost of retained
interaction terms. Also report residual prediction error outside the leading
set. A decision-directed method may succeed with worse global prediction error;
that is part of the hypothesis, provided its decision performance and uncertainty
are supported. Fresh data must test whether the retained correction transfers
beyond the cases that selected it. Include a predeclared intervention intended
to activate or suppress the proposed term, such as the conformational lock in
$C$, and predict its signed response contrast before evaluation. A matched
comparison should assess other changes introduced by that intervention.
Agreement supports the proposed transferable correction without uniquely
identifying a mechanism from the contrast alone.

The main falsifiers are concrete. If the policy restores as much detail as
full-model learning at equal cost, the proposed compression advantage has not
appeared. If omitted interactions repeatedly defeat its uncertainty envelope,
the challenge language is inadequate. If its new descriptors only fit their
triggering exceptions, it has not learned reusable chemical structure. If a
simple planner does as well, the theory-refinement architecture can still be
useful without a more elaborate controller.

A run may end with a decision certified under its stated assumptions, a
rejected representation, or exhausted budget without a certificate. Future
assays are needed to move from computational or masked-replay evidence to
chemical discovery claims. The guiding scientific question remains precise:
how much of a useful working theory must be corrected to make a particular
molecular choice reliable, and can carefully chosen counterexamples teach us
that correction at substantially lower cost than resolving the full chemistry?
% END SOURCE: appendices/app_counterexample_design.tex


%--------Back Matter--------
\backmatter

% Bibliography
\addcontentsline{toc}{chapter}{Bibliography}
% Bibliography is inlined below; BibTeX is not required.
% BEGIN FROZEN BIBLIOGRAPHY
\newcommand{\etalchar}[1]{$^{#1}$}
\begin{thebibliography}{GOWWB12}

\bibitem[BAROK22]{bellamy2022bo}
Hugo Bellamy, Abbi Abdel~Rehim, Oghenejokpeme~I. Orhobor, and Ross King.
\newblock Batched bayesian optimization for drug design in noisy environments.
\newblock {\em Journal of Chemical Information and Modeling},
  62(17):3970--3981, 2022.

\bibitem[BFP{\etalchar{+}}14]{bartok2014partial}
G.~Bart{\'o}k, D.~P. Foster, D.~P{\'a}l, A.~Rakhlin, and C.~Szepesv{\'a}ri.
\newblock Partial monitoring---classification, regret bounds, and algorithms.
\newblock {\em Mathematics of Operations Research}, 39(4):967--997, 2014.

\bibitem[BM96]{bemis1996scaffold}
Guy~W. Bemis and Mark~A. Murcko.
\newblock The properties of known drugs. 1. molecular frameworks.
\newblock {\em Journal of Medicinal Chemistry}, 39(15):2887--2893, 1996.

\bibitem[BR01]{becker2001dwr}
Roland Becker and Rolf Rannacher.
\newblock An optimal control approach to a posteriori error estimation in
  finite element methods.
\newblock {\em Acta Numerica}, 10:1--102, 2001.

\bibitem[BTN98]{bental1998robust}
Aharon Ben-Tal and Arkadi Nemirovski.
\newblock Robust convex optimization.
\newblock {\em Mathematics of Operations Research}, 23(4):769--805, 1998.

\bibitem[Cay57]{cayley1857trees}
Arthur Cayley.
\newblock On the theory of the analytical forms called trees.
\newblock {\em The London, Edinburgh, and Dublin Philosophical Magazine and
  Journal of Science}, 13(85):172--176, 1857.

\bibitem[Cay74]{cayley1874isomers}
Arthur Cayley.
\newblock On the mathematical theory of isomers.
\newblock {\em The London, Edinburgh, and Dublin Philosophical Magazine and
  Journal of Science}, 47:444--446, 1874.

\bibitem[CGJ{\etalchar{+}}00]{clarke2000cegar}
Edmund~M. Clarke, Orna Grumberg, Somesh Jha, Yuan Lu, and Helmut Veith.
\newblock Counterexample-guided abstraction refinement.
\newblock In {\em Computer Aided Verification}, volume 1855 of {\em Lecture
  Notes in Computer Science}, pages 154--169, 2000.

\bibitem[CT06]{cover2006information}
Thomas~M. Cover and Joy~A. Thomas.
\newblock {\em Elements of Information Theory}.
\newblock Wiley-Interscience, 2nd edition, 2006.

\bibitem[Dij59]{dijkstra1959shortest}
Edsger~W. Dijkstra.
\newblock A note on two problems in connexion with graphs.
\newblock {\em Numerische Mathematik}, 1(1):269--271, 1959.

\bibitem[dlVdLB14]{deleon2014recapmmp}
Antonio de~la Vega~de Le{\'o}n and J{\"u}rgen Bajorath.
\newblock Matched molecular pairs derived by retrosynthetic fragmentation.
\newblock {\em MedChemComm}, 5(1):64--67, 2014.

\bibitem[DW71]{dreyfus1971steiner}
Stuart~E. Dreyfus and Robert~A. Wagner.
\newblock The steiner problem in graphs.
\newblock {\em Networks}, 1(3):195--207, 1971.

\bibitem[Edm71]{edmonds1971greedy}
Jack Edmonds.
\newblock Matroids and the greedy algorithm.
\newblock {\em Mathematical Programming}, 1(1):127--136, 1971.

\bibitem[Elf52]{elfving1952allocation}
Gustav Elfving.
\newblock Optimum allocation in linear regression theory.
\newblock {\em The Annals of Mathematical Statistics}, 23(2):255--262, 1952.

\bibitem[FJJR19]{fiez2019transductive}
Tanner Fiez, Lalit Jain, Kevin~G. Jamieson, and Lillian Ratliff.
\newblock Sequential experimental design for transductive linear bandits.
\newblock In {\em Advances in Neural Information Processing Systems},
  volume~32, 2019.

\bibitem[FW64]{free1964qsar}
S.~M. Free and J.~W. Wilson.
\newblock A mathematical contribution to structure-activity studies.
\newblock {\em Journal of Medicinal Chemistry}, 7(4):395--399, 1964.

\bibitem[GBS08]{ghosh2008resistance}
Arpita Ghosh, Stephen Boyd, and Amin Saberi.
\newblock Minimizing effective resistance of a graph.
\newblock {\em SIAM Review}, 50(1):37--66, 2008.

\bibitem[GCH{\etalchar{+}}02]{ge2002nonadditivity}
Nanxiang Ge, Sung~Jin Cho, Mark Hermsmeier, Michael Poss, and C.~Frank Shen.
\newblock Testing non-additivity of biological activity in a combinatorial
  library.
\newblock {\em Combinatorial Chemistry \& High Throughput Screening},
  5(2):147--154, 2002.

\bibitem[GGJ77]{garey1977steiner}
Michael~R. Garey, Ronald~L. Graham, and David~S. Johnson.
\newblock The complexity of computing steiner minimal trees.
\newblock {\em SIAM Journal on Applied Mathematics}, 32(4):835--859, 1977.

\bibitem[Gil08]{giles2008mlmc}
Michael~B. Giles.
\newblock Multilevel {Monte Carlo} path simulation.
\newblock {\em Operations Research}, 56(3):607--617, 2008.

\bibitem[Gil15]{giles2015mlmc}
Michael~B. Giles.
\newblock Multilevel {Monte Carlo} methods.
\newblock {\em Acta Numerica}, 24:259--328, 2015.

\bibitem[GOHO{\etalchar{+}}15]{gupta2015sarm}
Disha Gupta-Ostermann, Yoichiro Hirose, Takenao Odagami, Hiroyuki Kouji, and
  J{\"u}rgen Bajorath.
\newblock Follow-up: Prospective compound design using the `{SAR} matrix'
  method and matrix-derived conditional probabilities of activity.
\newblock {\em F1000Research}, 4:75, 2015.

\bibitem[GOSB14]{gupta2014neighborhood}
Disha Gupta-Ostermann, Veerabahu Shanmugasundaram, and J{\"u}rgen Bajorath.
\newblock Neighborhood-based prediction of novel active compounds from {SAR}
  matrices.
\newblock {\em Journal of Chemical Information and Modeling}, 54(3):801--809,
  2014.

\bibitem[GOWWB12]{gupta2012graphmining}
Disha Gupta-Ostermann, Mathias Wawer, Anne~Mai Wassermann, and J{\"u}rgen
  Bajorath.
\newblock Graph mining for {SAR} transfer series.
\newblock {\em Journal of Chemical Information and Modeling}, 52(4):935--942,
  2012.

\bibitem[GSC21]{graff2021molpal}
David~E. Graff, Eugene~I. Shakhnovich, and Connor~W. Coley.
\newblock Accelerating high-throughput virtual screening through molecular
  pool-based active learning.
\newblock {\em Chemical Science}, 12(22):7866--7881, 2021.

\bibitem[GTZ97]{goreinov1997pseudoskeleton}
Sergei~A. Goreinov, Eugene~E. Tyrtyshnikov, and Nikolai~L. Zamarashkin.
\newblock A theory of pseudoskeleton approximations.
\newblock {\em Linear Algebra and its Applications}, 261:1--21, 1997.

\bibitem[GY21]{giese2021barnet}
Timothy~J. Giese and Darrin~M. York.
\newblock Variational method for networkwide analysis of relative ligand
  binding free energies with loop closure and experimental constraints.
\newblock {\em Journal of Chemical Theory and Computation}, 17(3):1326--1336,
  2021.

\bibitem[Ham50]{hamming1950codes}
Richard~W. Hamming.
\newblock Error detecting and error correcting codes.
\newblock {\em Bell System Technical Journal}, 29(2):147--160, 1950.

\bibitem[Hei01]{heinrich2001mlmc}
Stefan Heinrich.
\newblock Multilevel {Monte Carlo} methods.
\newblock 2179:58--67, 2001.

\bibitem[HR10]{hussain2010mmp}
Jameed Hussain and Ceara Rea.
\newblock Computationally efficient algorithm to identify matched molecular
  pairs ({MMPs}) in large data sets.
\newblock {\em Journal of Chemical Information and Modeling}, 50(3):339--348,
  2010.

\bibitem[HRMS21]{howard2021timeuniform}
Steven~R. Howard, Aaditya Ramdas, Jon McAuliffe, and Jasjeet Sekhon.
\newblock Time-uniform, nonparametric, nonasymptotic confidence sequences.
\newblock {\em The Annals of Statistics}, 49(2):1055--1080, 2021.

\bibitem[Hu74]{hu1974communication}
T.~C. Hu.
\newblock Optimum communication spanning trees.
\newblock {\em SIAM Journal on Computing}, 3(3):188--195, 1974.

\bibitem[JLYY11]{jiang2011hodge}
Xiaoye Jiang, Lek-Heng Lim, Yuan Yao, and Yinyu Ye.
\newblock Statistical ranking and combinatorial {Hodge} theory.
\newblock {\em Mathematical Programming}, 127(1):203--244, 2011.

\bibitem[Kar78]{karp1978cyclemean}
Richard~M. Karp.
\newblock A characterization of the minimum cycle mean in a digraph.
\newblock {\em Discrete Mathematics}, 23(3):309--311, 1978.

\bibitem[KBD{\etalchar{+}}19]{konze2019reactional}
Kyle~D. Konze, Pieter~H. Bos, Markus~K. Dahlgren, Karl Leswing, Ivan
  Tubert-Brohman, Andrea Bortolato, Braxton Robbason, Robert Abel, and Sathesh
  Bhat.
\newblock Reaction-based enumeration, active learning, and free energy
  calculations to rapidly explore synthetically tractable chemical space and
  optimize potency of cyclin-dependent kinase 2 inhibitors.
\newblock {\em Journal of Chemical Information and Modeling}, 59(9):3782--3793,
  2019.

\bibitem[KFL15]{kramer2015nonadditivity}
Christian Kramer, Julian~E. Fuchs, and Klaus~R. Liedl.
\newblock Strong nonadditivity as a key structure--activity relationship
  feature: Distinguishing structural changes from assay artifacts.
\newblock {\em Journal of Chemical Information and Modeling}, 55(3):483--494,
  2015.

\bibitem[KG22]{koolen2022evalues}
Wouter~M. Koolen and Peter Gr{\"u}nwald.
\newblock Log-optimal anytime-valid e-values.
\newblock {\em International Journal of Approximate Reasoning}, 141:69--82,
  2022.

\bibitem[Kie59]{kiefer1959design}
Jack Kiefer.
\newblock Optimum experimental designs.
\newblock {\em Journal of the Royal Statistical Society, Series B},
  21(2):272--304, 1959.

\bibitem[Kra19]{kramer2019nonadditivity}
Christian Kramer.
\newblock Nonadditivity analysis.
\newblock {\em Journal of Chemical Information and Modeling}, 59(9):4034--4042,
  2019.

\bibitem[KTT15]{kiraly2015algebraic}
Franz~J. Kir{\'a}ly, Louis Theran, and Ryota Tomioka.
\newblock The algebraic combinatorial approach for low-rank matrix completion.
\newblock {\em Journal of Machine Learning Research}, 16(41):1391--1436, 2015.

\bibitem[KW60]{kiefer1960equivalence}
Jack Kiefer and Jacob Wolfowitz.
\newblock The equivalence of two extremum problems.
\newblock {\em Canadian Journal of Mathematics}, 12:363--366, 1960.

\bibitem[Law73]{lawler1973hypergraph}
Eugene~L. Lawler.
\newblock Cutsets and partitions of hypergraphs.
\newblock {\em Networks}, 3(3):275--285, 1973.

\bibitem[LS20]{lattimore2020bandits}
Tor Lattimore and Csaba Szepesv{\'a}ri.
\newblock {\em Bandit Algorithms}.
\newblock Cambridge University Press, 2020.

\bibitem[LVBL98]{lobo1998socp}
Miguel~Sousa Lobo, Lieven Vandenberghe, Stephen Boyd, and Herv{\'e} Lebret.
\newblock Applications of second-order cone programming.
\newblock {\em Linear Algebra and its Applications}, 284(1--3):193--228, 1998.

\bibitem[MB87]{muller1987bipartite}
Haiko M{\"u}ller and Andreas Brandst{\"a}dt.
\newblock The {NP}-completeness of steiner tree and dominating set for chordal
  bipartite graphs.
\newblock {\em Theoretical Computer Science}, 53(2--3):257--265, 1987.

\bibitem[MP09]{maurer2009bernstein}
Andreas Maurer and Massimiliano Pontil.
\newblock Empirical {Bernstein} bounds and sample variance penalization.
\newblock In {\em Proceedings of the 22nd Conference on Learning Theory
  (COLT)}, pages 115--124, 2009.

\bibitem[Nat95]{natarajan1995sparse}
Balas~K. Natarajan.
\newblock Sparse approximate solutions to linear systems.
\newblock {\em SIAM Journal on Computing}, 24(2):227--234, 1995.

\bibitem[Ney34]{neyman1934stratified}
Jerzy Neyman.
\newblock On the two different aspects of the representative method: The method
  of stratified sampling and the method of purposive selection.
\newblock {\em Journal of the Royal Statistical Society}, 97(4):558--625, 1934.

\bibitem[OBO14]{osting2014rankings}
Braxton Osting, Christoph Brune, and Stanley~J. Osher.
\newblock Optimal data collection for informative rankings expose
  well-connected graphs.
\newblock {\em Journal of Machine Learning Research}, 15(85):2981--3012, 2014.

\bibitem[OP02]{oden2002modelingerror}
J.~Tinsley Oden and Serge Prudhomme.
\newblock Estimation of modeling error in computational mechanics.
\newblock {\em Journal of Computational Physics}, 182(2):496--515, 2002.

\bibitem[PGH{\etalchar{+}}08]{patel2008nonadditivity}
Yogendra Patel, Valerie~J. Gillet, Trevor Howe, Joaquin Pastor, Julen
  Oyarzabal, and Peter Willett.
\newblock Assessment of additive/nonadditive effects in structure--activity
  relationships: Implications for iterative drug design.
\newblock {\em Journal of Medicinal Chemistry}, 51(23):7552--7562, 2008.

\bibitem[Pra62]{prange1962information}
Eugene Prange.
\newblock The use of information sets in decoding cyclic codes.
\newblock {\em IEEE Transactions on Information Theory}, 8(5):5--9, 1962.

\bibitem[PWG18]{peherstorfer2018multifidelity}
Benjamin Peherstorfer, Karen Willcox, and Max Gunzburger.
\newblock Survey of multifidelity methods in uncertainty propagation,
  inference, and optimization.
\newblock {\em SIAM Review}, 60(3):550--591, 2018.

\bibitem[Rao62]{rao1962generalized}
C.~Radhakrishna Rao.
\newblock A note on a generalized inverse of a matrix with applications to
  problems in mathematical statistics.
\newblock {\em Journal of the Royal Statistical Society, Series B},
  24(1):152--158, 1962.

\bibitem[Rec11]{recht2011matrix}
Benjamin Recht.
\newblock A simpler approach to matrix completion.
\newblock {\em Journal of Machine Learning Research}, 12:3413--3430, 2011.

\bibitem[RKDD21]{reda2021topm}
Cl{\'e}mence R{\'e}da, Emilie Kaufmann, and Andr{\'e}e Delahaye-Duriez.
\newblock Top-$m$ identification for linear bandits.
\newblock In {\em Proceedings of the 24th International Conference on
  Artificial Intelligence and Statistics}, volume 130 of {\em Proceedings of
  Machine Learning Research}, pages 1108--1116, 2021.

\bibitem[Rot64]{rota1964mobius}
Gian-Carlo Rota.
\newblock On the foundations of combinatorial theory {I}. theory of
  {M{\"o}bius} functions.
\newblock {\em Zeitschrift f{\"u}r Wahrscheinlichkeitstheorie und Verwandte
  Gebiete}, 2:340--368, 1964.

\bibitem[Sag11]{sagnol2011multiresponse}
Guillaume Sagnol.
\newblock Computing optimal designs of multiresponse experiments reduces to
  second-order cone programming.
\newblock {\em Journal of Statistical Planning and Inference},
  141(5):1684--1708, 2011.

\bibitem[Sea65]{searle1965estimable}
Shayle~R. Searle.
\newblock Additional results concerning estimable functions and generalized
  inverse matrices.
\newblock {\em Journal of the Royal Statistical Society, Series B},
  27(3):486--490, 1965.

\bibitem[Soy73]{soyster1973robust}
Allen~L. Soyster.
\newblock Convex programming with set-inclusive constraints and applications to
  inexact linear programming.
\newblock {\em Operations Research}, 21(5):1154--1157, 1973.

\bibitem[Sta11]{stanley2011enumerative}
Richard~P. Stanley.
\newblock {\em Enumerative Combinatorics, Volume 1}.
\newblock Cambridge University Press, 2nd edition, 2011.

\bibitem[Tuk49]{tukey1949nonadditivity}
John~W. Tukey.
\newblock One degree of freedom for non-additivity.
\newblock {\em Biometrics}, 5(3):232--242, 1949.

\bibitem[Ver18]{vershynin2018highdim}
Roman Vershynin.
\newblock {\em High-Dimensional Probability: An Introduction with Applications
  in Data Science}.
\newblock Cambridge University Press, 2018.

\bibitem[WB11]{wawer2011mmsgraph}
Mathias Wawer and J{\"u}rgen Bajorath.
\newblock Local structural changes, global data views: Graphical
  substructure--activity relationship trailing.
\newblock {\em Journal of Medicinal Chemistry}, 54(8):2944--2951, 2011.

\bibitem[Whi35]{whitney1935matroid}
Hassler Whitney.
\newblock On the abstract properties of linear dependence.
\newblock {\em American Journal of Mathematics}, 57(3):509--533, 1935.

\bibitem[WHWB12]{wassermann2012sarm}
Anne~Mai Wassermann, Peter Haebel, Nils Weskamp, and J{\"u}rgen Bajorath.
\newblock {SAR} matrices: Automated extraction of information-rich {SAR} tables
  from large compound data sets.
\newblock {\em Journal of Chemical Information and Modeling}, 52(7):1769--1776,
  2012.

\bibitem[Xu19]{xu2019measurement}
Huafeng Xu.
\newblock Optimal measurement network of pairwise differences.
\newblock {\em Journal of Chemical Information and Modeling},
  59(11):4720--4728, 2019.

\bibitem[YL06]{yuan2006group}
Ming Yuan and Yi~Lin.
\newblock Model selection and estimation in regression with grouped variables.
\newblock {\em Journal of the Royal Statistical Society, Series B},
  68(1):49--67, 2006.

\bibitem[YTO91]{young1991minbalance}
Neal~E. Young, Robert~E. Tarjan, and James~B. Orlin.
\newblock Faster parametric shortest path and minimum-balance algorithms.
\newblock {\em Networks}, 21(2):205--221, 1991.

\bibitem[ZWVB12]{zhang2012sartransfer}
Bijun Zhang, Anne~Mai Wassermann, Martin Vogt, and J{\"u}rgen Bajorath.
\newblock Systematic assessment of compound series with {SAR} transfer
  potential.
\newblock {\em Journal of Chemical Information and Modeling},
  52(12):3138--3143, 2012.

\end{thebibliography}
% END FROZEN BIBLIOGRAPHY

% BEGIN FROZEN INDEX
\begin{theindex}

  \item $\sigma $-algebra, \hyperpage{55}
  \item $\sigma$-algebra, \hyperpage{10}, \hyperpage{12}, 
		\hyperpage{54}, \hyperpage{154}
  \item 2-core, \hyperpage{15}

  \indexspace

  \item attachment point, \hyperpage{160}

  \indexspace

  \item bandit algorithm, \hyperpage{78}, \hyperpage{81}, 
		\hyperpage{106}
  \item Bemis--Murcko scaffold, \hyperpage{14}, \hyperpage{55}, 
		\hyperpage{154}
  \item binding affinity, \hyperpage{54, 55}
  \item boundary, \hyperpage{18}, \hyperpage{80}, \hyperpage{160}, 
		\hyperpage{170}

  \indexspace

  \item chemical space, \hyperpage{14}, \hyperpage{45}, \hyperpage{54}, 
		\hyperpage{133}, \hyperpage{143}, \hyperpage{166}
  \item closure operator, \hyperpage{23}, \hyperpage{45}, 
		\hyperpage{48}, \hyperpage{142}
  \item conditional expectation, \hyperpage{5}, \hyperpage{23}, 
		\hyperpage{54}, \hyperpage{94}, \hyperpage{148}

  \indexspace

  \item decoration, \hyperpage{18}, \hyperpage{56}, \hyperpage{76, 77}, 
		\hyperpage{160}
  \item docking score, \hyperpage{55}, \hyperpage{57}, \hyperpage{84}
  \item downset, \hyperpage{5}, \hyperpage{34}, \hyperpage{51}, 
		\hyperpage{136}

  \indexspace

  \item energy factorization, \hyperpage{18}, \hyperpage{80}

  \indexspace

  \item fiber, \hyperpage{16}, \hyperpage{48, 49}, \hyperpage{71}
  \item Fisher information, \hyperpage{99}

  \indexspace

  \item Galois connection, \hyperpage{45}, \hyperpage{54}, 
		\hyperpage{141}, \hyperpage{166}

  \indexspace

  \item information geometry, \hyperpage{94, 95}
  \item irreducible error, \hyperpage{12}, \hyperpage{22}, 
		\hyperpage{58}, \hyperpage{69}, \hyperpage{71}, 
		\hyperpage{133}
  \item isomer, \hyperpage{43}

  \indexspace

  \item KL divergence, \hyperpage{96--98}

  \indexspace

  \item lattice, \hyperpage{23}, \hyperpage{45}, \hyperpage{47}, 
		\hyperpage{54}
  \item lower cone, \hyperpage{35}

  \indexspace

  \item martingale, \hyperpage{60}, \hyperpage{146}, \hyperpage{158}
  \item molecule, \hyperpage{6}, \hyperpage{15}, \hyperpage{47, 48}, 
		\hyperpage{54, 55}, \hyperpage{74}, \hyperpage{79}, 
		\hyperpage{134}, \hyperpage{154}
  \item multilevel Monte Carlo, \hyperpage{67}, \hyperpage{81}, 
		\hyperpage{143}

  \indexspace

  \item partial order, \hyperpage{12}, \hyperpage{14}, \hyperpage{19}, 
		\hyperpage{154}
  \item poset, \hyperpage{8}, \hyperpage{19}, \hyperpage{34}, 
		\hyperpage{46}, \hyperpage{51}, \hyperpage{101}, 
		\hyperpage{136}, \hyperpage{141}
  \item pushforward measure, \hyperpage{58}

  \indexspace

  \item ring system, \hyperpage{55}, \hyperpage{59}, \hyperpage{154}

  \indexspace

  \item scaffold, \hyperpage{14}, \hyperpage{45}, \hyperpage{54}, 
		\hyperpage{69}, \hyperpage{94}, \hyperpage{133}, 
		\hyperpage{143}, \hyperpage{166}
  \item scaffold hopping, \hyperpage{94, 95}
  \item scoring function, \hyperpage{10}, \hyperpage{51}, 
		\hyperpage{55, 56}, \hyperpage{70}, \hyperpage{73}, 
		\hyperpage{95}, \hyperpage{138}, \hyperpage{155}
  \item screening, \hyperpage{21}, \hyperpage{45}, \hyperpage{55}, 
		\hyperpage{69}, \hyperpage{97}, \hyperpage{133}, 
		\hyperpage{144}, \hyperpage{165}

  \indexspace

  \item VAE (variational autoencoder), \hyperpage{171}
  \item valence, \hyperpage{6}, \hyperpage{14}
  \item variance spectrum, \hyperpage{20}, \hyperpage{23}, 
		\hyperpage{54}, \hyperpage{61}, \hyperpage{69}, 
		\hyperpage{94}, \hyperpage{133}, \hyperpage{142, 143}, 
		\hyperpage{166}

\end{theindex}
% END FROZEN INDEX

\end{document}
% END SOURCE: main.tex
